\documentclass[11pt,a4paper,twoside]{book}
\usepackage[margin=2.5cm,top=3cm,bottom=3cm,headheight=22pt]{geometry}
\usepackage{xcolor}
\usepackage{fontspec}
\usepackage{titlesec}
\usepackage{fancyhdr}
\usepackage{hyperref}
\usepackage{parskip}
\usepackage{microtype}
\usepackage{tikz}
\usepackage{bookmark}
\usepackage{tcolorbox}
\usepackage{graphicx}
\tcbuselibrary{skins,breakable}

\definecolor{bg}{HTML}{06060C}
\definecolor{neoncyan}{HTML}{00F0FF}
\definecolor{neonpink}{HTML}{FF2CF0}
\definecolor{neongreen}{HTML}{39FF14}
\definecolor{neonyellow}{HTML}{FFE500}
\definecolor{neonred}{HTML}{FF003C}
\definecolor{neonorange}{HTML}{FF6B00}
\definecolor{dimtext}{HTML}{8888AA}
\definecolor{bodytext}{HTML}{C8C8E0}
\definecolor{codebg}{HTML}{0C0C18}
\definecolor{ruleglow}{HTML}{00AACC}
\definecolor{chapterbg}{HTML}{0A0A16}

\pagecolor{bg}
\color{bodytext}

\setmonofont{Menlo}[Scale=0.82]
\setmainfont{Helvetica Neue}[UprightFont=*,BoldFont=* Bold,ItalicFont=* Italic]

\linespread{1.3}
\setlength{\parskip}{8pt}
\setlength{\parindent}{0pt}
\sloppy

\titleformat{\chapter}[display]
    {\normalfont\Huge\bfseries\color{neoncyan}}
    {}
    {0pt}
    {
\begin{tikzpicture}[remember picture, overlay]
        \fill[chapterbg] (current page.north west) rectangle ([yshift=-6cm]current page.north east);
        \draw[neoncyan, line width=1.5pt] ([yshift=-6cm]current page.north west) -- ([yshift=-6cm]current page.north east);
        \draw[neonpink, line width=0.5pt, opacity=0.4] ([yshift=-5.9cm]current page.north west) -- ([yshift=-5.9cm]current page.north east);
    \end{tikzpicture}\vspace{-20pt}}
    [\vspace{8pt}{\color{dimtext}\small\itshape // ZPWR ENCYCLOPEDIA //}]
\titlespacing*{\chapter}{0pt}{50pt}{40pt}

\titleformat{\section}
    {\normalfont\Large\bfseries\color{neonpink}}
    {}
    {0pt}
    {}
    [\vspace{2pt}{\color{ruleglow}\rule{\textwidth}{0.5pt}}\vspace{1pt}{\color{neonpink!40}\rule{\textwidth}{0.2pt}}]
\titlespacing*{\section}{0pt}{24pt}{12pt}

\pagestyle{fancy}
\fancyhf{}
\fancyhead[LE]{\color{dimtext}\small\textsc{\leftmark}}
\fancyhead[RO]{\color{dimtext}\small\textit{\rightmark}}
\fancyfoot[C]{{\color{neoncyan}\rule[4pt]{2cm}{0.3pt}}\quad{\color{dimtext}\small\thepage}\quad{\color{neoncyan}\rule[4pt]{2cm}{0.3pt}}}
\renewcommand{\headrule}{{\color{ruleglow}\hrule height 0.3pt}}
\renewcommand{\chaptermark}[1]{\markboth{#1}{}}
\renewcommand{\sectionmark}[1]{\markright{#1}}
\fancypagestyle{plain}{
    \fancyhf{}
    \fancyfoot[C]{{\color{neoncyan}\rule[4pt]{2cm}{0.3pt}}\quad{\color{dimtext}\small\thepage}\quad{\color{neoncyan}\rule[4pt]{2cm}{0.3pt}}}
    \renewcommand{\headrule}{}
}

\newtcolorbox{codebox}{
    enhanced,colback=codebg,colframe=ruleglow,boxrule=0.5pt,arc=3pt,
    left=10pt,right=10pt,top=6pt,bottom=6pt,
    fontupper=\small\ttfamily\color{neongreen},breakable,
    shadow={1mm}{-1mm}{0mm}{neoncyan!15!bg},
    borderline west={2pt}{0pt}{neonpink!60}
}

% Glowing info box for tips and notes
\newtcolorbox{glowbox}{
    enhanced,colback=bg,colframe=neonpink,boxrule=0.4pt,arc=4pt,
    left=10pt,right=10pt,top=6pt,bottom=6pt,
    fontupper=\color{bodytext},breakable,
    shadow={2mm}{-2mm}{0mm}{neonpink!10!bg},
    borderline north={2pt}{0pt}{neoncyan!40},
    borderline south={2pt}{0pt}{neoncyan!40}
}

% Page border on every page
\usepackage{background}
\backgroundsetup{
    scale=1,angle=0,opacity=1,
    contents={
\begin{tikzpicture}[remember picture, overlay]
        % Corner decorations
        \draw[neoncyan, line width=0.3pt, opacity=0.15]
            ([xshift=8mm,yshift=-8mm]current page.north west) --
            ([xshift=8mm,yshift=-25mm]current page.north west);
        \draw[neoncyan, line width=0.3pt, opacity=0.15]
            ([xshift=8mm,yshift=-8mm]current page.north west) --
            ([xshift=25mm,yshift=-8mm]current page.north west);
        \draw[neonpink, line width=0.3pt, opacity=0.15]
            ([xshift=-8mm,yshift=-8mm]current page.north east) --
            ([xshift=-8mm,yshift=-25mm]current page.north east);
        \draw[neonpink, line width=0.3pt, opacity=0.15]
            ([xshift=-8mm,yshift=-8mm]current page.north east) --
            ([xshift=-25mm,yshift=-8mm]current page.north east);
        \draw[neoncyan, line width=0.3pt, opacity=0.15]
            ([xshift=8mm,yshift=8mm]current page.south west) --
            ([xshift=8mm,yshift=25mm]current page.south west);
        \draw[neoncyan, line width=0.3pt, opacity=0.15]
            ([xshift=8mm,yshift=8mm]current page.south west) --
            ([xshift=25mm,yshift=8mm]current page.south west);
        \draw[neonpink, line width=0.3pt, opacity=0.15]
            ([xshift=-8mm,yshift=8mm]current page.south east) --
            ([xshift=-8mm,yshift=25mm]current page.south east);
        \draw[neonpink, line width=0.3pt, opacity=0.15]
            ([xshift=-8mm,yshift=8mm]current page.south east) --
            ([xshift=-25mm,yshift=8mm]current page.south east);
        % Subtle scanlines across entire page
        \foreach \y in {0,0.4,...,29.7} {
            \draw[white, line width=0.08pt, opacity=0.012]
                ([yshift=-\y cm]current page.north west) -- ([yshift=-\y cm]current page.north east);
        }
    \end{tikzpicture}}
}

% Enhanced chapter page with circuit board pattern
\titleformat{\chapter}[display]
    {\normalfont\Huge\bfseries\color{neoncyan}}
    {}
    {0pt}
    {
\begin{tikzpicture}[remember picture, overlay]
        \fill[chapterbg] (current page.north west) rectangle ([yshift=-7cm]current page.north east);
        % Neon glow lines
        \draw[neoncyan, line width=2pt] ([yshift=-7cm]current page.north west) -- ([yshift=-7cm]current page.north east);
        \draw[neonpink, line width=0.8pt, opacity=0.5] ([yshift=-6.9cm]current page.north west) -- ([yshift=-6.9cm]current page.north east);
        \draw[neoncyan, line width=0.3pt, opacity=0.2] ([yshift=-6.7cm]current page.north west) -- ([yshift=-6.7cm]current page.north east);
        % Circuit traces
        \foreach \x in {2,4,...,18} {
            \draw[neoncyan, line width=0.2pt, opacity=0.08] ([xshift=\x cm, yshift=-1cm]current page.north west) -- ([xshift=\x cm, yshift=-7cm]current page.north west);
        }
        \foreach \y in {-2,-3,...,-6} {
            \draw[neonpink, line width=0.2pt, opacity=0.06] ([yshift=\y cm]current page.north west) -- ([yshift=\y cm]current page.north east);
        }
        % Dot grid with glow
        \foreach \x in {2,4,...,18} {
            \foreach \y in {-1.5,-3,...,-6.5} {
                \fill[neoncyan, opacity=0.2] ([xshift=\x cm, yshift=\y cm]current page.north west) circle (1.2pt);
                \fill[neoncyan, opacity=0.05] ([xshift=\x cm, yshift=\y cm]current page.north west) circle (3pt);
            }
        }
        % Glowing orb
        \fill[neonpink, opacity=0.08] ([xshift=-3cm, yshift=-3.5cm]current page.north east) circle (2.5cm);
        \fill[neoncyan, opacity=0.06] ([xshift=2cm, yshift=-4cm]current page.north west) circle (2cm);
    \end{tikzpicture}\vspace{-20pt}}
    [\vspace{8pt}{\color{dimtext}\small\itshape // ZPWR ENCYCLOPEDIA // CYBERPUNK EDITION //}]

\hypersetup{colorlinks=true,linkcolor=neoncyan,urlcolor=neongreen,
    pdftitle={The ZPWR Encyclopedia},pdfauthor={MenkeTechnologies}}

\begin{document}

\begin{titlepage}

\begin{tikzpicture}[remember picture, overlay]
    % Full page gradient wash
    \fill[bg] (current page.south west) rectangle (current page.north east);
    \shade[top color=neoncyan!8!bg, bottom color=bg]
        (current page.north west) rectangle ([yshift=-10cm]current page.north east);
    \shade[bottom color=neonpink!6!bg, top color=bg]
        ([yshift=8cm]current page.south west) rectangle (current page.south east);
    % Hexagonal grid background
    \foreach \x in {0,1.5,...,21} {
        \foreach \y in {0,1.732,...,29.7} {
            \pgfmathsetmacro{\xshift}{mod(round(\y/1.732),2)*0.75}
            \draw[neoncyan, line width=0.15pt, opacity=0.04]
                ([xshift=\x cm + \xshift cm, yshift=-\y cm]current page.north west)
                ++(30:0.5) -- ++(330:0.5) -- ++(270:0.5) -- ++(210:0.5) -- ++(150:0.5) -- ++(90:0.5) -- cycle;
        }
    }
    % Scanlines
    \foreach \y in {0,0.25,...,29.7} {
        \draw[white, line width=0.1pt, opacity=0.015]
            ([yshift=-\y cm]current page.north west) -- ([yshift=-\y cm]current page.north east);
    }
    % Neon glow bars
    \draw[neoncyan, line width=3pt] ([yshift=-3cm, xshift=1.5cm]current page.north west) -- ([yshift=-3cm, xshift=-1.5cm]current page.north east);
    \draw[neonpink, line width=1.5pt, opacity=0.7] ([yshift=-3.15cm, xshift=1.5cm]current page.north west) -- ([yshift=-3.15cm, xshift=-1.5cm]current page.north east);
    \draw[neoncyan, line width=0.5pt, opacity=0.3] ([yshift=-3.35cm, xshift=1.5cm]current page.north west) -- ([yshift=-3.35cm, xshift=-1.5cm]current page.north east);
    \draw[neoncyan, line width=3pt] ([yshift=3cm, xshift=1.5cm]current page.south west) -- ([yshift=3cm, xshift=-1.5cm]current page.south east);
    \draw[neonpink, line width=1.5pt, opacity=0.7] ([yshift=2.85cm, xshift=1.5cm]current page.south west) -- ([yshift=2.85cm, xshift=-1.5cm]current page.south east);
    \draw[neoncyan, line width=0.5pt, opacity=0.3] ([yshift=2.65cm, xshift=1.5cm]current page.south west) -- ([yshift=2.65cm, xshift=-1.5cm]current page.south east);
    % Glowing orbs
    \fill[neoncyan, opacity=0.06] ([xshift=3cm, yshift=-5cm]current page.north west) circle (2cm);
    \fill[neonpink, opacity=0.04] ([xshift=-4cm, yshift=6cm]current page.south east) circle (3cm);
    \fill[neongreen, opacity=0.03] ([xshift=8cm, yshift=-15cm]current page.north west) circle (1.5cm);
    % Data streams (diagonal lines)
    \foreach \i in {1,...,8} {
        \pgfmathsetmacro{\xoff}{\i * 2.5}
        \draw[neoncyan, line width=0.2pt, opacity=0.06]
            ([xshift=\xoff cm]current page.south west) -- ([xshift=\xoff cm + 3cm]current page.north west);
    }
\end{tikzpicture}
\vspace*{5cm}
\begin{center}
{\fontsize{42}{50}\selectfont\bfseries\color{neoncyan} THE ZPWR}\\[0.5cm]
{\fontsize{42}{50}\selectfont\bfseries\color{neonpink} ENCYCLOPEDIA}\\[2cm]
{\Large\color{bodytext} A Hacker's Field Guide to the Terminal}\\[2cm]
{\large\color{neonyellow}45 Chapters\quad\color{dimtext}//\quad\color{neonorange}510+ Verbs}\\[2cm]
{\color{bodytext}\large MenkeTechnologies}\\[0.3cm]
{\color{dimtext}\small \url{https://github.com/MenkeTechnologies/zpwr}}\\[3cm]
{\color{dimtext}\small CYBERPUNK EDITION\quad//\quad Generated from \texttt{zpwr wizard}}\\[1cm]
\includegraphics[width=0.5\textwidth]{/Users/wizard/.zpwr/docs/screenshots/figlet_zpwr.png}
\end{center}
\end{titlepage}

{\hypersetup{linkcolor=neoncyan}\tableofcontents}


\chapter{THE RED PILL}


\begin{center}
\includegraphics[width=0.85\textwidth]{/Users/wizard/.zpwr/docs/screenshots/art_jack_in.png}
\end{center}

\section{Welcome to the Machine}


\vspace{10pt}
{\large\bfseries\color{neoncyan} WELCOME TO THE MACHINE}\vspace{6pt}

You found it. The rabbit hole. \textbf{ZPWR} is not a dotfiles repo. It is not a "framework." It is a fully integrated shell operating system that happens to run inside your terminal.


\vspace{12pt}
{\bfseries\color{neonyellow} WHAT IS \textbf{ZPWR}?}\vspace{4pt}

A \$\{\#\textbf{ZPWR\_VERBS}\}-verb command dispatcher built on top of zsh, tmux, vim/nvim, fzf, and git. Every verb is a surgical tool. Every keybinding is a weapon. Every alias is a shortcut through the matrix.


\vspace{12pt}
{\bfseries\color{neonyellow} THE PHILOSOPHY}\vspace{4pt}

1. Verbs, not scripts. Everything is zpwr <verb>. 2. Discoverability.   Tab completion on everything. 3. Symmetry.          vim* and emacs* verbs mirror each other. 4. Speed.             Compiled zwc files. Benchmarked startup. 5. No magic.          Read the source. It is all zsh.


\vspace{12pt}
{\bfseries\color{neonyellow} THE DIRECTORY TREE}\vspace{4pt}

\$\textbf{ZPWR} ........................ root of the kingdom \$\textbf{ZPWR\_ENV} .................... env vars and shell init \$\textbf{ZPWR\_SCRIPTS} ................ executable scripts \$\textbf{ZPWR\_AUTOLOAD} ............... autoloaded functions \$\textbf{ZPWR\_LOCAL} .................. local config (gitignored) \$\textbf{ZPWR\_INSTALL} ................ installer scripts \$\textbf{ZPWR\_TMUX} ................... tmux configuration \$\textbf{ZPWR\_TEST} ................... test suite


\vspace{12pt}
{\bfseries\color{neonyellow} YOUR FIRST COMMAND}\vspace{4pt}

\begin{codebox}
zpwr about    \# behold the banner\\
zpwr help     \# list every verb in color\\
\end{codebox}

Pro tip: zpwr is aliased to many short forms. Type zpwr <tab> and let the completions wash over you.

WARNING: Side effects may include mass alias adoption, involuntary tab-completing in other people's terminals, and an irrational hatred of vanilla bash.


\vspace{16pt}

\section{Your First 5 Minutes}


\vspace{10pt}
{\large\bfseries\color{neoncyan} YOUR FIRST 5 MINUTES}\vspace{6pt}

You just installed zpwr. The terminal blinks. What now? Here is your survival kit for the first 300 seconds.


\vspace{12pt}
{\bfseries\color{neonyellow} STEP 1: SEE THE BANNER}\vspace{4pt}

\begin{codebox}
zpwr about              \# the ASCII art welcome screen\\
zpwr banner             \# same thing, different alias\\
zpwr bannercounts       \# banner + environment statistics\\
\end{codebox}


\vspace{12pt}
{\bfseries\color{neonyellow} STEP 2: COUNT YOUR ARSENAL}\vspace{4pt}

\begin{codebox}
zpwr environmentcounts  \# how many aliases, functions, etc.\\
\end{codebox}

This prints a breakdown of your loaded shell environment: aliases, functions, builtins, parameters, completions. If the numbers are big, congratulations. You are armed.


\vspace{12pt}
{\bfseries\color{neonyellow} STEP 3: DISCOVER THE VERBS}\vspace{4pt}

\begin{codebox}
zpwr help               \# colorized list of ALL verbs\\
zpwr verbs     \# fzf fuzzy search through verbs\\
\end{codebox}

The verbsfzf verb is your panic button (alias: subcommandsfzf). Forgot what a command is called? Fuzzy search it. This pipes every verb and its description into fzf. Type a few characters. Watch the list collapse. Hit enter. Execute.


\vspace{12pt}
{\bfseries\color{neonyellow} STEP 4: VERIFY YOUR ENVIRONMENT}\vspace{4pt}

\begin{codebox}
zpwr doctor             \# run environment health check\\
zpwr envvars            \# list all zpwr env vars\\
zpwr colortest          \# verify your terminal colors\\
zpwr colortest256       \# the full 256-color palette\\
\end{codebox}


\vspace{12pt}
{\bfseries\color{neonyellow} STEP 5: TAB COMPLETE EVERYTHING}\vspace{4pt}

Type zpwr and press <tab>. You will see \$\{\#\textbf{ZPWR\_VERBS}\} verbs. Type zpwr git and press <tab>. Filtered to git verbs. Type zpwr vim and press <tab>. Filtered to vim verbs. The completion system knows every verb. Trust it.

If you only remember one command from this page: zpwr verbs It is the cheat code to everything else.


\vspace{16pt}

\section{The Anatomy of a Verb}


\vspace{10pt}
{\large\bfseries\color{neoncyan} THE ANATOMY OF A VERB}\vspace{6pt}

Every zpwr command follows one sacred pattern: zpwr <verb> [args...]


\vspace{12pt}
{\bfseries\color{neonyellow} HOW DISPATCH WORKS}\vspace{4pt}

When you type zpwr filesearch, here is what happens: 1. The zpwr function receives "filesearch" as \$1 2. It looks up \textbf{ZPWR\_VERBS}[filesearch] in the verb table 3. The table entry maps to a function or alias 4. That function is invoked with any remaining args.


\vspace{12pt}
{\bfseries\color{neonyellow} THE VERB TABLE}\vspace{4pt}

The associative array \textbf{ZPWR\_VERBS} is the brain. Each entry has the format:

\begin{codebox}
ZPWR\_VERBS[verb]='function\_or\_alias=description'\\
\end{codebox}


\vspace{12pt}
{\bfseries\color{neonyellow} MANY NAMES, ONE FUNCTION}\vspace{4pt}

Verbs often have aliases for muscle memory:

\begin{codebox}
zpwr forward  = zpwr f   = zpwr fwd\\
zpwr reverse  = zpwr r   = zpwr rvs = zpwr cdup\\
zpwr github   = zpwr gh  = zpwr mygithub\\
zpwr repo     = zpwr zp  = zpwr zpwr\\
\end{codebox}


\vspace{12pt}
{\bfseries\color{neonyellow} TAB COMPLETION IN ACTION}\vspace{4pt}

zpwr git<tab>     \# shows: gitcommit, gitconfig, gitsearch... zpwr vim<tab>     \# shows: vimall, vimrecent, vimctags... zpwr home<tab>    \# shows: home, homeenv, homelocal... zpwr forgit<tab>  \# shows: forgitadd, forgitlog, forgitdiff...


\vspace{12pt}
{\bfseries\color{neonyellow} INSPECT THE TABLE YOURSELF}\vspace{4pt}

\begin{codebox}
zpwr subcommands        \# dump all verbs to stdout\\
zpwr verbs     \# fuzzy search verbs with fzf\\
zpwr verbscount          \# how many verbs exist\\
\end{codebox}


\vspace{12pt}
{\bfseries\color{neonyellow} THE EDIT SUFFIX PATTERN}\vspace{4pt}

Many verbs come in pairs: filesearch      -> open file after selection filesearchedit  -> type into command line for editing The "edit" variant puts the result on your command line instead of executing immediately. Power and precision.


\vspace{16pt}

\section{Configuring Your Rig}


\vspace{10pt}
{\large\bfseries\color{neoncyan} CONFIGURING YOUR RIG}\vspace{6pt}

Your terminal is a cockpit. These are the dials.


\vspace{12pt}
{\bfseries\color{neonyellow} THE FOUR CONFIG SESSIONS}\vspace{4pt}

Each of these opens a curated vim session with the relevant config files already loaded in buffers:

\begin{codebox}
zpwr brc    \# shell aliases file vim session\\
\end{codebox}

Opens \textbf{\$ZPWR\_LOCAL/.tokens.sh} and related files this is where YOUR personal aliases live.

\begin{codebox}
zpwr zrc    \# zshrc vim session\\
\end{codebox}

Opens the zshrc and related zsh configs the source of all shell initialization.

\begin{codebox}
zpwr vrc    \# vimrc/init.vim vim session\\
\end{codebox}

Opens vim/nvim config files your editor's soul laid bare.

\begin{codebox}
zpwr trc    \# tmux.conf vim session\\
\end{codebox}

Opens tmux configuration pane and window management settings.


\vspace{12pt}
{\bfseries\color{neonyellow} THE TOKENS FILE}\vspace{4pt}

\begin{codebox}
zpwr tokens       \# vim the .tokens.sh file\\
zpwr vimtokens    \# same thing, explicit vim\\
zpwr emacstokens  \# same thing, emacs edition\\
\end{codebox}

Located at \textbf{\$ZPWR\_LOCAL/.tokens.sh} This file is gitignored. Put secrets and custom aliases here. Never commit this file. It is YOUR local override.


\vspace{12pt}
{\bfseries\color{neonyellow} \textbf{EDITOR} COMMANDS}\vspace{4pt}

\begin{codebox}
zpwr edit      \# launch EDITOR\\
zpwr editor    \# same thing\\
zpwr vimconfig \# edit all zpwr configs in vim\\
\end{codebox}


\vspace{12pt}
{\bfseries\color{neonyellow} THE FULL CONFIG STACK}\vspace{4pt}

\textbf{\$ZPWR\_LOCAL/.tokens.sh} ... your aliases and secrets \textbf{\$ZPWR\_ENV/.zpwr\_env.sh} .. env vars and paths \textbf{\textasciitilde{}/.zshrc} .................. shell entry point \textbf{\textasciitilde{}/.vimrc} / init.vim ....... editor config \textbf{\textasciitilde{}/.tmux.conf} .............. tmux config

Think of it as: brc for shell, zrc for zsh, vrc for vim, trc for tmux. Four letters, four kingdoms.


\vspace{16pt}


\chapter{NAVIGATION PROTOCOL}


\begin{center}
\includegraphics[width=0.85\textwidth]{/Users/wizard/.zpwr/docs/screenshots/art_the_grid.png}
\end{center}

\section{Directory Warfare}


\vspace{10pt}
{\large\bfseries\color{neoncyan} DIRECTORY WARFARE}\vspace{6pt}

In the filesystem, speed is survival. Typing full paths is for people who have given up on life.


\vspace{12pt}
{\bfseries\color{neonyellow} FORWARD AND REVERSE}\vspace{4pt}

\begin{codebox}
zpwr forward <dir>  \# cd into a child directory\\
zpwr f <dir>        \# same thing, shorter\\
zpwr fwd <dir>      \# same thing, medium length\\
\end{codebox}

\begin{codebox}
zpwr reverse <dir>  \# cd UP the tree to a parent matching <dir>\\
zpwr r <dir>        \# same thing, one letter\\
zpwr rvs <dir>      \# same thing, three letters\\
zpwr cdup <dir>     \# same thing, cd-up style\\
\end{codebox}

Example: you are in \textbf{\textasciitilde{}/code/myproject/src/lib/utils}

\begin{codebox}
zpwr r myproject    \# jumps UP to \textasciitilde{}/code/myproject\\
\end{codebox}

No more ../../../../ chains. Navigate like a human.


\vspace{12pt}
{\bfseries\color{neonyellow} BASIC CD}\vspace{4pt}

\begin{codebox}
zpwr cd <path>      \# cd to directory arg\\
\end{codebox}

Seems redundant but integrates with zpwr's dispatch.


\vspace{12pt}
{\bfseries\color{neonyellow} Z: FRECENCY JUMPING}\vspace{4pt}

\begin{codebox}
zpwr z <query>      \# jump to frecency-ranked directory\\
\end{codebox}

Z tracks where you go. The more you visit a directory, the higher it ranks. Type a fragment, land at the top match.

\begin{codebox}
zpwr zcd            \# fzf list of z-ranked dirs, pick one\\
\end{codebox}

When you cannot remember the exact fragment, zcd shows you the full ranked list in fzf.


\vspace{12pt}
{\bfseries\color{neonyellow} PRACTICAL COMBOS}\vspace{4pt}

zpwr z proj         \# land in your most-visited "proj" dir zpwr r src          \# climb back up to the "src" parent zpwr f tests        \# dive into the "tests" child.

The pattern is: z for teleportation, f for descent, r for ascent. Three directions. Total coverage.

PROTIP: If you are still typing cd ../../../, this is your intervention.


\vspace{16pt}

\section{Home Base Operations}


\vspace{10pt}
{\large\bfseries\color{neoncyan} HOME BASE OPERATIONS}\vspace{6pt}

Every zpwr directory has a one-word teleporter. No paths. No thinking. Just say where.


\vspace{12pt}
{\bfseries\color{neonyellow} THE \textbf{HOME} VERB FAMILY}\vspace{4pt}

\begin{codebox}
zpwr home                \# cd ZPWR (the root)\\
zpwr homelocal           \# cd ZPWR\_LOCAL\\
zpwr homescripts         \# cd ZPWR\_SCRIPTS\\
zpwr homeenv             \# cd ZPWR\_ENV\\
zpwr homeinstall         \# cd ZPWR\_INSTALL\\
zpwr homeautoload        \# cd ZPWR\_AUTOLOAD\\
zpwr homeautoloadcommon   \# cd ZPWR\_AUTOLOAD\_COMMON\\
zpwr homecomps           \# cd ZPWR\_COMPS\\
zpwr hometest            \# cd ZPWR\_TEST\\
zpwr hometmux            \# cd ZPWR\_TMUX\\
zpwr homeexercism        \# cd HOME/Exercism\\
\end{codebox}


\vspace{12pt}
{\bfseries\color{neonyellow} WHAT LIVES WHERE}\vspace{4pt}

home .............. the zpwr repo root, start of everything homelocal ........ gitignored personal configs and tokens homescripts ...... executable shell scripts homeenv .......... environment variables and init files homeinstall ...... installer scripts for dependencies homeautoload ..... autoloaded functions (the engine room) homeautoloadcommon  platform-independent autoloads homecomps ........ zsh completion definitions hometest ......... test suite files hometmux ......... tmux configuration files


\vspace{12pt}
{\bfseries\color{neonyellow} EXAMPLE WORKFLOW}\vspace{4pt}

You want to edit an autoload function:

\begin{codebox}
zpwr homeautoloadcommon\\
\end{codebox}

\# Now you are in the autoload directory \# ls, grep, edit to your heart's content.

You want to check your local tokens:

\begin{codebox}
zpwr homelocal\\
\end{codebox}

\# Lands you right next to .tokens.sh.

Every home* verb is a bookmark you never need to set. The bookmarks are built into the system.


\vspace{16pt}

\section{Named Directories}


\vspace{10pt}
{\large\bfseries\color{neoncyan} NAMED DIRECTORIES}\vspace{6pt}

Zsh has a secret weapon most people never learn: named directories. They turn env vars into tilde paths.


\vspace{12pt}
{\bfseries\color{neonyellow} HOW IT WORKS}\vspace{4pt}

When you do hash -d \textbf{ZPWR}=\$\textbf{ZPWR}, zsh registers a named directory. Now \textasciitilde{}\textbf{ZPWR} expands to \$\textbf{ZPWR} everywhere: in cd, in tab completion, in prompt, in arguments.


\vspace{12pt}
{\bfseries\color{neonyellow} \textbf{ZPWR}'S NAMED DIRECTORIES}\vspace{4pt}

\begin{codebox}
\textasciitilde{}ZPWR ................. ZPWR root\\
\textasciitilde{}ZPWR\_SCRIPTS ........ ZPWR\_SCRIPTS\\
\textasciitilde{}ZPWR\_AUTOLOAD ....... ZPWR\_AUTOLOAD\\
\textasciitilde{}ZPWR\_ENV ............ ZPWR\_ENV\\
\textasciitilde{}ZPWR\_LOCAL .......... ZPWR\_LOCAL\\
\textasciitilde{}ZPWR\_INSTALL ........ ZPWR\_INSTALL\\
\textasciitilde{}ZPWR\_TMUX ........... ZPWR\_TMUX\\
\textasciitilde{}ZPWR\_TEST ........... ZPWR\_TEST\\
\textasciitilde{}ZPWR\_COMPS .......... ZPWR\_COMPS\\
\end{codebox}


\vspace{12pt}
{\bfseries\color{neonyellow} WHAT THIS MEANS IN PRACTICE}\vspace{4pt}

Your prompt might show: \textasciitilde{}\textbf{ZPWR\_AUTOLOAD}/common  instead of the full path You can type:

\begin{codebox}
cd \textasciitilde{}ZPWR\_SCRIPTS  and it just works\\
vim \textasciitilde{}ZPWR\_ENV/.zpwr\_env.sh  tilde paths everywhere\\
\end{codebox}


\vspace{12pt}
{\bfseries\color{neonyellow} TAB COMPLETION WITH TILDE}\vspace{4pt}

Type cd \textasciitilde{}\textbf{ZPWR}<tab> and watch zsh expand the options: \textasciitilde{}\textbf{ZPWR}  \textasciitilde{}\textbf{ZPWR\_SCRIPTS}  \textasciitilde{}\textbf{ZPWR\_ENV}  \textasciitilde{}\textbf{ZPWR\_LOCAL} ... This works in any command, not just cd.


\vspace{12pt}
{\bfseries\color{neonyellow} THE MECHANISM}\vspace{4pt}

Zsh option AUTO\_NAME\_DIRS turns any exported scalar parameter into a named directory automatically. \textbf{ZPWR} exports its directory variables, so they become named directories for free.

Think of named directories as human-readable symlinks that live in your shell's brain instead of the filesystem. No files created. No cleanup needed. Pure memory.


\vspace{16pt}

\section{Quick Jumps}


\vspace{10pt}
{\large\bfseries\color{neoncyan} QUICK JUMPS}\vspace{6pt}

The home* verbs land you in zpwr directories. These verbs land you in the ecosystem around zpwr.


\vspace{12pt}
{\bfseries\color{neonyellow} \textbf{ZPWR} INTERNALS}\vspace{4pt}

\begin{codebox}
zpwr scripts       \# cd to scripts directory\\
\end{codebox}

Shortcut alias: zs.

\begin{codebox}
zpwr autoload      \# cd to autoload directory\\
\end{codebox}

Shortcut alias: zal.

\begin{codebox}
zpwr completions   \# cd to completion directory\\
zpwr comps         \# same thing\\
\end{codebox}

Shortcut alias: zco.


\vspace{12pt}
{\bfseries\color{neonyellow} PLUGIN ECOSYSTEM}\vspace{4pt}

\begin{codebox}
zpwr plugins       \# cd to ZSH\_CUSTOM/plugins\\
\end{codebox}

Shortcut alias: zpl This is where oh-my-zsh custom plugins live. Your third-party zsh plugins are neighbors here.

\begin{codebox}
zpwr forked        \# cd to ZPWR\_FORKED\_DIR\\
zpwr fp            \# same thing\\
\end{codebox}

Forked repos that zpwr depends on live here. Useful when you need to hack on a dependency.


\vspace{12pt}
{\bfseries\color{neonyellow} SPECIAL DIRECTORIES}\vspace{4pt}

\begin{codebox}
zpwr hidden        \# cd to ZPWR\_HIDDEN\_DIR\\
\end{codebox}

The hidden directory. Dot-prefixed internal state.

\begin{codebox}
zpwr repo          \# cd to ZPWR\_REPO\_NAME\\
zpwr zp            \# same thing\\
zpwr zpwr          \# same thing (yes, zpwr zpwr)\\
\end{codebox}


\vspace{12pt}
{\bfseries\color{neonyellow} THE NAVIGATION CHEAT SHEET}\vspace{4pt}

Use case                      Verb --------                      ---- Edit a script .............. zpwr scripts Fix a completion ........... zpwr comps Add an autoload function .. zpwr autoload Debug a plugin ............ zpwr plugins Patch a fork .............. zpwr forked

Each of these is one verb, zero path typing, and instant teleportation to where the work is.


\vspace{16pt}


\chapter{THE GIT DIMENSION}


\begin{center}

\includegraphics[width=0.85\textwidth]{/Users/wizard/.zpwr/docs/screenshots/art_data_stream.png}
\end{center}

\section{Git Basics}


\vspace{10pt}
{\large\bfseries\color{neoncyan} GIT BASICS}\vspace{6pt}

Git is already a language. \textbf{ZPWR} makes it a dialect you can actually speak without consulting the man page.


\vspace{12pt}
{\bfseries\color{neonyellow} THE COMMIT PIPELINE}\vspace{4pt}

\begin{codebox}
zpwr gitcommit "your message here"\\
\end{codebox}

This is not just git commit. It does: git add -> git commit -> git push All in one verb. The full gacp (git add commit push). Because who commits without pushing? Cowards.


\vspace{12pt}
{\bfseries\color{neonyellow} CONFIG AND IGNORES}\vspace{4pt}

\begin{codebox}
zpwr gitconfig         \# edit git config file\\
zpwr gitglobalconfig   \# same thing, explicit name\\
zpwr gitignore         \# edit .git/info/exclude (local)\\
zpwr gitignores        \# edit global git ignores file\\
zpwr gitglobalignores  \# same as gitignores\\
\end{codebox}


\vspace{12pt}
{\bfseries\color{neonyellow} STASH AND DIFF (VIA FORGIT)}\vspace{4pt}

\textbf{ZPWR} delegates interactive stash and diff to forgit:

\begin{codebox}
zpwr forgitdiff        \# fzf-powered git diff viewer\\
zpwr forgitstash       \# fzf-powered stash browser\\
\end{codebox}

See page 13 for the full forgit arsenal.


\vspace{12pt}
{\bfseries\color{neonyellow} REMOTE OPERATIONS}\vspace{4pt}

\begin{codebox}
zpwr gitremotes        \# list all git remotes\\
zpwr gitdir            \# check if pwd is a git dir\\
\end{codebox}


\vspace{12pt}
{\bfseries\color{neonyellow} TAGS}\vspace{4pt}

\begin{codebox}
zpwr gitedittag        \# print tag to BUFFER\\
zpwr gitupdatetag      \# print tag and latest msg to BUFFER\\
\end{codebox}


\vspace{12pt}
{\bfseries\color{neonyellow} REAL-WORLD EXAMPLE}\vspace{4pt}

\# You just finished a feature. Ship it:

\begin{codebox}
zpwr gitcommit "add quantum flux capacitor"\\
\end{codebox}

\# Stages everything, commits with message, pushes. \# One line. Done. Next feature.

\# Need to check your git config:

\begin{codebox}
zpwr gitconfig\\
\end{codebox}

\# Opens it in \$\textbf{EDITOR}. Fix, save, done.

Git verbs follow the pattern: git + action noun. If you can name it, zpwr probably has a verb for it.


\vspace{16pt}

\section{GitHub Integration}


\vspace{10pt}
{\large\bfseries\color{neoncyan} GITHUB INTEGRATION}\vspace{6pt}

GitHub is not a website. It is an API with a website. \textbf{ZPWR} treats it accordingly.


\vspace{12pt}
{\bfseries\color{neonyellow} CREATE AND DESTROY}\vspace{4pt}

\begin{codebox}
zpwr githubcreate      \# create a remote GitHub repo\\
zpwr hc               \# same thing, two letters\\
\end{codebox}

Creates a repo on GitHub from your current directory. Sets up the remote. Ready to push.

\begin{codebox}
zpwr githubdelete      \# delete a remote GitHub repo\\
zpwr hd               \# same thing, two letters\\
\end{codebox}

This is permanent. It asks for confirmation. But once confirmed, the repo is gone. Poof.


\vspace{12pt}
{\bfseries\color{neonyellow} BROWSE YOUR GITHUB}\vspace{4pt}

\begin{codebox}
zpwr github            \# open your GitHub profile in browser\\
zpwr gh               \# same thing\\
zpwr mygithub          \# same thing\\
\end{codebox}

Opens your profile page. Quick ego check.

\begin{codebox}
zpwr githubzpwr        \# open zpwr repo on GitHub\\
zpwr ghz              \# same thing\\
\end{codebox}

Jump straight to the zpwr source. File issues. Star it. You know you want to.


\vspace{12pt}
{\bfseries\color{neonyellow} CONTRIBUTION STATS}\vspace{4pt}

\begin{codebox}
zpwr githubcontribcount \# count GitHub contribs in last year\\
\end{codebox}

Tells you how green your contribution graph is. The number that keeps you pushing commits at 2 AM.


\vspace{12pt}
{\bfseries\color{neonyellow} REVEAL: THE BROWSER SHORTCUT}\vspace{4pt}

\begin{codebox}
zpwr reveal            \# show current repo in browser\\
\end{codebox}

You are in a git repo. You want to see it on GitHub. One word. Browser opens. You are there.


\vspace{12pt}
{\bfseries\color{neonyellow} WORKFLOW EXAMPLE}\vspace{4pt}

\# Start a new project from scratch:

\begin{codebox}
mkdir new-project \&\& cd new-project\\
git init\\
zpwr hc               \# repo created on GitHub\\
zpwr gitcommit "initial commit"\\
zpwr reveal            \# admire it in the browser\\
\end{codebox}

From zero to pushed repo in four commands. That is the GitHub integration promise.


\vspace{16pt}

\section{Git Archaeology}


\vspace{10pt}
{\large\bfseries\color{neoncyan} GIT ARCHAEOLOGY}\vspace{6pt}

The past is not dead. It is not even past. It is stored in .git/objects and zpwr can dig it up.


\vspace{12pt}
{\bfseries\color{neonyellow} SEARCHING THE LOG}\vspace{4pt}

\begin{codebox}
zpwr gitsearch             \# search for regex in git log\\
\end{codebox}

Uses fzf to grep through your entire commit history. That one commit where you "fixed the thing"? Found.

\begin{codebox}
zpwr gitcommits            \# search git commits with fzf\\
\end{codebox}

Interactive commit browser. Preview diffs inline. Like git log but you can actually find things.


\vspace{12pt}
{\bfseries\color{neonyellow} COUNTING COMMITS}\vspace{4pt}

\begin{codebox}
zpwr gitcommitcount        \# total number of commits\\
\end{codebox}

Just a number. But what a number. The sum total of your decisions in this repo.


\vspace{12pt}
{\bfseries\color{neonyellow} CONTRIBUTION ANALYSIS}\vspace{4pt}

\begin{codebox}
zpwr gitcontribcount       \# contributions by author\\
\end{codebox}

Who committed the most? The leaderboard of blame.

\begin{codebox}
zpwr gitcontribcountdirs   \# contributions across dirs\\
\end{codebox}

Same analysis but across multiple repositories. Useful for org-wide contribution audits.

\begin{codebox}
zpwr gitcontribcountlines  \# lines contributed by author\\
\end{codebox}

Lines of code per author. The metric that lies to you. But it is a fun lie.


\vspace{12pt}
{\bfseries\color{neonyellow} SIZE ANALYSIS}\vspace{4pt}

\begin{codebox}
zpwr gittotallines         \# total line count of git files\\
\end{codebox}

How big is this repo, really? Now you know.

\begin{codebox}
zpwr gitlargest            \# largest files in git history\\
\end{codebox}

Shows the biggest blobs stored by git, descending. That 50MB binary someone committed in 2019? There it is. Now you know what to purge (see page 12).


\vspace{12pt}
{\bfseries\color{neonyellow} DIG SESSION EXAMPLE}\vspace{4pt}

\# Who wrote the most code in this repo?

\begin{codebox}
zpwr gitcontribcount\\
\end{codebox}

\# What is eating all the disk space?

\begin{codebox}
zpwr gitlargest\\
\end{codebox}

\# Find that commit about the database migration:

\begin{codebox}
zpwr gitsearch\\
\end{codebox}

\# Type "migration" in fzf. There it is.


\vspace{16pt}

\section{Git Surgery}


\vspace{10pt}
{\large\bfseries\color{neoncyan} GIT SURGERY}\vspace{6pt}

Sometimes you need to rewrite history. Not the geopolitical kind. The git kind. These are destructive operations. Use with intent.


\vspace{12pt}
{\bfseries\color{neonyellow} CLEARING THE CACHE}\vspace{4pt}

\begin{codebox}
zpwr gitclearcache     \# clear old git refs and objects\\
\end{codebox}

Runs garbage collection on git's internal state. Prunes unreachable objects, packs loose refs. Your .git directory gets a deep clean.


\vspace{12pt}
{\bfseries\color{neonyellow} REMOVING COMMITS}\vspace{4pt}

\begin{codebox}
zpwr gitclearcommit    \# remove matching commits from history\\
\end{codebox}

This is filter-branch territory. Rewrites history. Force push required after. Use when someone committed passwords, API keys, or that embarrassing commented-out code block.


\vspace{12pt}
{\bfseries\color{neonyellow} REMOVING FILES FROM ALL HISTORY}\vspace{4pt}

\begin{codebox}
zpwr gitclearfile      \# rm file from all git refs and objects\\
\end{codebox}

Nuclear option for files. That 200MB video someone committed in 2020? This removes it from every commit that ever existed. Your repo size drops. Your clone time drops. Your blood pressure drops.


\vspace{12pt}
{\bfseries\color{neonyellow} EMAIL SURGERY}\vspace{4pt}

\begin{codebox}
zpwr gitemail          \# change email with git filter-branch\\
\end{codebox}

Changes both author and committer email across history.

\begin{codebox}
zpwr gitaemail         \# change author email only\\
\end{codebox}

When you committed with your personal email at work.

\begin{codebox}
zpwr gitcemail         \# change committer email only\\
\end{codebox}

When the committer field has the wrong identity.


\vspace{12pt}
{\bfseries\color{neonyellow} WHEN TO USE THESE}\vspace{4pt}

Scenario: you committed with wrong@email.com.

\begin{codebox}
zpwr gitemail\\
\end{codebox}

\# Follow the prompts to fix every commit.

Scenario: .env file in git history with secrets.

\begin{codebox}
zpwr gitclearfile\\
\end{codebox}

\# Purges it from every ref and object.

REMEMBER: These rewrite history. Coordinate with your team before force pushing. Or do not coordinate and watch the chaos unfold.


\vspace{16pt}

\section{Forgit: Git on Steroids}


\vspace{10pt}
{\large\bfseries\color{neoncyan} FORGIT: GIT ON STEROIDS}\vspace{6pt}

Forgit wraps common git operations with fzf previews. Every action gets a live preview pane. Every choice is informed. No more blind git commands.


\vspace{12pt}
{\bfseries\color{neonyellow} STAGING AND UNSTAGING}\vspace{4pt}

\begin{codebox}
zpwr forgitadd            \# fzf-powered git add\\
\end{codebox}

Shows unstaged files with inline diffs in the preview. Select files with tab. Enter to stage. Visual staging.

\begin{codebox}
zpwr forgitreset          \# fzf-powered git reset HEAD\\
\end{codebox}

The reverse of forgitadd. Unstage files visually.


\vspace{12pt}
{\bfseries\color{neonyellow} VIEWING CHANGES}\vspace{4pt}

\begin{codebox}
zpwr forgitdiff           \# fzf-powered git diff viewer\\
\end{codebox}

Browse changed files. Preview shows colorized diffs. Side-by-side comparison in your terminal.

\begin{codebox}
zpwr forgitlog            \# fzf-powered git log\\
\end{codebox}

Commit history with preview. See the diff for each commit without leaving the log view. Beautiful.


\vspace{12pt}
{\bfseries\color{neonyellow} STASH OPERATIONS}\vspace{4pt}

\begin{codebox}
zpwr forgitstash          \# fzf-powered stash browser\\
\end{codebox}

Lists stashes. Preview shows what each stash contains. Pick one, apply it. No more "was it stash@\{3\}?"


\vspace{12pt}
{\bfseries\color{neonyellow} CLEANUP}\vspace{4pt}

\begin{codebox}
zpwr forgitclean          \# fzf-powered git clean\\
\end{codebox}

Shows untracked files. Select which to delete. Preview before you destroy. Civilized cleanup.

\begin{codebox}
zpwr forgitrestore        \# fzf-powered git restore\\
\end{codebox}

Discard changes to selected files. With preview.


\vspace{12pt}
{\bfseries\color{neonyellow} GITIGNORE MANAGEMENT}\vspace{4pt}

\begin{codebox}
zpwr forgitignore         \# browse gitignore templates\\
zpwr forgitignoreget      \# download a gitignore template\\
zpwr forgitignorelist     \# list available templates\\
zpwr forgitignoreupdate   \# update gitignore templates\\
zpwr forgitignoreclean    \# clean gitignore cache\\
\end{codebox}


\vspace{12pt}
{\bfseries\color{neonyellow} EXTRAS}\vspace{4pt}

\begin{codebox}
zpwr forgitinfo           \# show git repo information\\
zpwr forgitwarn           \# show forgit warnings\\
\end{codebox}

Forgit turns git into a visual experience. Every forgit* verb is a regular git command with eyes.


\vspace{16pt}

\section{Git Battalion Operations}


\vspace{10pt}
{\large\bfseries\color{neoncyan} GIT BATTALION OPERATIONS}\vspace{6pt}

One repo is easy. Twenty repos is a war. These verbs let you run commands across all of them.


\vspace{12pt}
{\bfseries\color{neonyellow} THE GENERAL: FORALLDIR}\vspace{4pt}

\begin{codebox}
zpwr gitforalldir <cmd>   \# run cmd in ALL git dirs\\
\end{codebox}

Walks every git directory it can find and runs your command inside each one. The nuclear option.


\vspace{12pt}
{\bfseries\color{neonyellow} CACHED REPOS}\vspace{4pt}

\begin{codebox}
zpwr gitreposexec <cmd>    \# run cmd in cached git dirs\\
\end{codebox}

Uses the cached list of known repos. Faster than foralldir because it does not rescan the filesystem.


\vspace{12pt}
{\bfseries\color{neonyellow} CLEAN VS DIRTY TARGETING}\vspace{4pt}

\begin{codebox}
zpwr gitreposcleanexec <cmd>       \# refresh + run in clean repos\\
zpwr gitreposcleancacheexec <cmd>  \# run in cached clean repos\\
zpwr gitreposdirtyexec <cmd>       \# refresh + run in dirty repos\\
zpwr gitreposdirtycacheexec <cmd>  \# run in cached dirty repos\\
\end{codebox}

Clean repos have no uncommitted changes. Dirty repos have pending work. Target your commands accordingly.


\vspace{12pt}
{\bfseries\color{neonyellow} MASS UPDATE}\vspace{4pt}

\begin{codebox}
zpwr gitupdatefordir      \# run git updates in all git dirs\\
\end{codebox}

Pulls latest, updates submodules, the full sync. Run this on a Monday morning. Coffee. Wait.


\vspace{12pt}
{\bfseries\color{neonyellow} BRANCH WIPE OPERATIONS}\vspace{4pt}

\begin{codebox}
zpwr gitfordirmain        \# wipe to main branch in git dirs\\
zpwr gitfordirdevelop     \# wipe to develop branch in git dirs\\
\end{codebox}

\begin{codebox}
zpwr gitzfordir <branch> [z]       \# wipe to any branch + z jump\\
zpwr gitzfordirmain [z]            \# z jump + wipe to main\\
zpwr gitzfordirdevelop [z]         \# z jump + wipe to develop\\
\end{codebox}


\vspace{12pt}
{\bfseries\color{neonyellow} THE REPOS FILE}\vspace{4pt}

\begin{codebox}
zpwr gitreposfile         \# edit the git repos file\\
\end{codebox}

This file lists all known repositories. Edit it to add or remove repos from batch operations.


\vspace{12pt}
{\bfseries\color{neonyellow} EXAMPLE: STATUS CHECK ALL REPOS}\vspace{4pt}

\begin{codebox}
zpwr gitreposexec git status -s\\
\end{codebox}

Runs git status -s in every cached repo. Instantly see which repos have uncommitted work.

Battalion ops turn you from a developer into a fleet commander.


\vspace{16pt}


\chapter{THE SEARCH ENGINE}


\begin{center}
\includegraphics[width=0.85\textwidth]{/Users/wizard/.zpwr/docs/screenshots/figlet_scan.png}
\end{center}

\section{Finding Files}


\vspace{10pt}
{\large\bfseries\color{neoncyan} FINDING FILES}\vspace{6pt}

The filesystem is vast. You are small. But you have fzf and a verb for every search strategy.


\vspace{12pt}
{\bfseries\color{neonyellow} FILESEARCH: LOCAL TREE}\vspace{4pt}

\begin{codebox}
zpwr filesearch          \# find and open a file in a sub dir\\
\end{codebox}

Recursively lists files under \$\textbf{PWD} into fzf. Type fragments of the filename. Fuzzy match. Enter. Opens the file in your pager.

\begin{codebox}
zpwr filesearchedit      \# same but puts path on command line\\
\end{codebox}

Instead of opening immediately, the selected path lands on your \$BUFFER for further editing. Append commands, pipe it, whatever you want.


\vspace{12pt}
{\bfseries\color{neonyellow} DIRSEARCH: DIRECTORIES ONLY}\vspace{4pt}

\begin{codebox}
zpwr dirsearch           \# cd to a sub directory via fzf\\
\end{codebox}

Same idea but only directories. Pick one, land in it. Faster than cd-tab-tab-tab-tab.


\vspace{12pt}
{\bfseries\color{neonyellow} FINDSEARCH: FULL DISK}\vspace{4pt}

\begin{codebox}
zpwr findsearch          \# find across the entire drive\\
\end{codebox}

Uses find(1) to search the whole filesystem. Results feed into fzf. Pick one, open it. Slow on first run. Worth the wait.

\begin{codebox}
zpwr findsearchedit      \# same but result goes to BUFFER\\
\end{codebox}


\vspace{12pt}
{\bfseries\color{neonyellow} LOCATESEARCH: INDEX-BASED}\vspace{4pt}

\begin{codebox}
zpwr locatesearch        \# search the locate database\\
\end{codebox}

Uses locate(1) which searches a pre-built index. Blazing fast. But only as fresh as your last updatedb.

\begin{codebox}
zpwr locatesearchedit    \# same but result goes to BUFFER\\
\end{codebox}


\vspace{12pt}
{\bfseries\color{neonyellow} THE SEARCH HIERARCHY}\vspace{4pt}

Speed    Scope          Verb -----    -----          ---- Fast ... current tree .. filesearch Fast ... dirs only ..... dirsearch Slow ... entire disk ... findsearch Instant  indexed disk .. locatesearch

Start with filesearch for local work. Escalate to locatesearch when you need the whole disk. Fall back to findsearch when the index is stale.


\vspace{16pt}

\section{Searching Content}


\vspace{10pt}
{\large\bfseries\color{neoncyan} SEARCHING CONTENT}\vspace{6pt}

Finding files is one thing. Finding the line inside the file? That is the real game.


\vspace{12pt}
{\bfseries\color{neonyellow} WORDSEARCH: CONTENT IN LOCAL TREE}\vspace{4pt}

\begin{codebox}
zpwr wordsearch          \# find file by content in sub dir\\
\end{codebox}

Greps through files in the current tree. Results show matching lines with filenames. Fzf preview shows the file with match highlighted. Select one, jump straight to it.

\begin{codebox}
zpwr wordsearchedit      \# same but result goes to BUFFER\\
\end{codebox}

Puts the selected file path on your command line. Useful when you want to pipe or chain commands.


\vspace{12pt}
{\bfseries\color{neonyellow} GREP: AG INTO FZF}\vspace{4pt}

\begin{codebox}
zpwr grep               \# grep through pwd with ag into fzf\\
\end{codebox}

This is the heavy artillery. Uses the silver searcher (ag) piped into fzf for interactive filtering. Respects .gitignore. Skips binaries. Fast.


\vspace{12pt}
{\bfseries\color{neonyellow} FILESEARCH VS WORDSEARCH}\vspace{4pt}

Know the filename?     Use filesearch Know the content?      Use wordsearch Know nothing?          Use grep and improvise.


\vspace{12pt}
{\bfseries\color{neonyellow} THE EDIT SUFFIX PATTERN}\vspace{4pt}

Every search verb has an edit variant: filesearch     -> filesearchedit wordsearch     -> wordsearchedit findsearch     -> findsearchedit locatesearch   -> locatesearchedit.

Without edit: opens the result immediately. With edit: puts the path on your command line. This lets you compose: vim \$(zpwr filesearchedit)


\vspace{12pt}
{\bfseries\color{neonyellow} REAL-WORLD EXAMPLES}\vspace{4pt}

\# Find where "\textbf{ZPWR\_VERBS}" is defined:

\begin{codebox}
zpwr wordsearch\\
\end{codebox}

\# Type \textbf{ZPWR\_VERBS} in fzf, browse the matches.

\# Quick grep for TODO comments:

\begin{codebox}
zpwr grep\\
\end{codebox}

\# Type TODO, see every file that has one.

The search verbs are your eyes inside the codebase. Use them before you reach for a GUI.


\vspace{16pt}

\section{The FZF Dimension}


\vspace{10pt}
{\large\bfseries\color{neoncyan} THE FZF DIMENSION}\vspace{6pt}

FZF is the nervous system of zpwr. It is wired into search, selection, navigation, git, vim, environment browsing, and verb dispatch. If zpwr has a UI, fzf IS that UI.


\vspace{12pt}
{\bfseries\color{neonyellow} VERB DISCOVERY}\vspace{4pt}

\begin{codebox}
zpwr verbs     \# fuzzy search all zpwr verbs\\
\end{codebox}

Every verb, every description, in one fzf window. This is the god-mode cheat sheet. Learn it first.


\vspace{12pt}
{\bfseries\color{neonyellow} ENVIRONMENT BROWSING}\vspace{4pt}

\begin{codebox}
zpwr envaccept          \# accept from fzf env\\
\end{codebox}

Browse all shell environment entries (aliases, functions, parameters, builtins) in fzf. Select one, accept it onto the command line.

\begin{codebox}
zpwr envedit            \# edit from fzf env\\
\end{codebox}

Same browsing but puts the result on \$BUFFER for editing.

\begin{codebox}
zpwr envcachesearch     \# search cached aliases, functions, etc.\\
\end{codebox}

Searches the pre-built cache of your entire environment. Faster than envedit because it uses the cache file.


\vspace{12pt}
{\bfseries\color{neonyellow} FZF PATTERNS ACROSS \textbf{ZPWR}}\vspace{4pt}

FZF appears in dozens of verbs. The pattern is always: 1. A data source generates lines (files, commits, dirs) 2. Lines pipe into fzf with a preview command 3. You fuzzy-filter, select, and take an action.

Verbs with fzf built in: filesearch     file list -> fzf -> open wordsearch     grep results -> fzf -> open dirsearch      dir list -> fzf -> cd gitcommits     git log -> fzf -> show diff forgitadd      unstaged files -> fzf -> stage forgitlog      commit log -> fzf -> preview zcd            z database -> fzf -> cd subcommandsfzf all verbs -> fzf -> execute.


\vspace{12pt}
{\bfseries\color{neonyellow} THE FZF KEYBINDINGS}\vspace{4pt}

\begin{codebox}
zpwr vimsearch          \# search vim keybindings in fzf\\
zpwr allsearch          \# search ALL keybindings in fzf\\
\end{codebox}

When you want to find a keybinding but do not remember where it is defined. Search them all.

FZF is not optional in zpwr. It is the interface. Master fzf and you master zpwr.


\vspace{16pt}


\chapter{EDITOR MASTERY}


\begin{center}
\includegraphics[width=0.85\textwidth]{/Users/wizard/.zpwr/docs/screenshots/art_neural_link.png}
\end{center}

\section{Vim Integration Overview}


\vspace{10pt}
{\large\bfseries\color{neoncyan} VIM INTEGRATION OVERVIEW}\vspace{6pt}

Vim is not just an editor in zpwr. It is a first-class citizen with its own verb family. Over 20 verbs, each opening vim in a specific context.


\vspace{12pt}
{\bfseries\color{neonyellow} THE PHILOSOPHY}\vspace{4pt}

Every vim verb answers one question: "What do I want to edit right now?" The answer is always one verb away.


\vspace{12pt}
{\bfseries\color{neonyellow} THE vim* VERB FAMILY}\vspace{4pt}

\begin{codebox}
vimall .............. open ALL zpwr files (for :argdo)\\
vimalledit .......... pick and edit 1+ zpwr configs\\
vimautoload ......... open all autoloads (for :argdo)\\
vimcd ............... vim edit, then cd to first dir\\
vimconfig ........... edit all zpwr config files\\
vimctags ............ vim zpwr ctags\\
vimemacsconfig ...... vim edit emacs zpwr configs\\
vimfilesearch ....... vim a file found via fzf\\
vimfilesearchedit ... vim edit a file found via fzf\\
vimfindsearch ....... vim a file from find(1) via fzf\\
vimfindsearchedit ... vim edit find result\\
vimgtags ............ vim zpwr GNU global tags\\
vimgtagsedit ........ vim edit zpwr GNU global tags\\
vimlocatesearch ..... vim a file from locate via fzf\\
vimlocatesearchedit . vim edit locate result\\
vimplugincount ...... count vim plugins\\
vimpluginlist ....... list all vim plugins\\
vimrecent ........... edit most recent files\\
vimrecentcd ......... cd + edit most recent files\\
vimrecentsudo ....... sudo edit most recent files\\
vimrecentsudocd ..... cd + sudo edit recent files\\
vimscripts .......... open all zpwr scripts (for :argdo)\\
vimscriptedit ....... pick and edit 1+ scripts\\
vimsearch ........... search vim keybindings via fzf\\
vimtests ............ edit all zpwr tests\\
vimtokens ........... vim the .tokens.sh file\\
vimwordsearch ....... vim a file found by content\\
vimwordsearchedit ... vim edit file found by content\\
\end{codebox}


\vspace{12pt}
{\bfseries\color{neonyellow} THE :argdo PATTERN}\vspace{4pt}

Verbs like vimall and vimscripts load many files into vim's arglist. Then you can:

\begin{codebox}
:argdo \%s/old/new/g | update\\
\end{codebox}

Mass find-and-replace across the entire zpwr codebase.

If it can be edited, there is a vim verb for it.


\vspace{16pt}

\section{Vim File Operations}


\vspace{10pt}
{\large\bfseries\color{neoncyan} VIM FILE OPERATIONS}\vspace{6pt}

Opening vim is easy. Opening vim with the right files already loaded? That is where zpwr earns its keep.


\vspace{12pt}
{\bfseries\color{neonyellow} THE BIG OPENS}\vspace{4pt}

\begin{codebox}
zpwr vimall             \# load ALL zpwr files into arglist\\
\end{codebox}

Every .zsh, .sh, .py, .pl file in the zpwr tree. Hundreds of files. Ready for :argdo mass operations. This is the tactical nuke of text editing.

\begin{codebox}
zpwr vimalledit         \# fzf pick 1+ configs to edit\\
\end{codebox}

Like vimall but you choose which files to open. Surgical precision instead of carpet bombing.


\vspace{12pt}
{\bfseries\color{neonyellow} RECENT FILES}\vspace{4pt}

\begin{codebox}
zpwr vimrecent          \# edit most recently opened files\\
\end{codebox}

Pulls from vim's file history. Opens the files you were just working on. Resume where you left off.

\begin{codebox}
zpwr vimrecentcd        \# cd to dir + open recent files\\
\end{codebox}

Same as vimrecent but also changes to the directory. Your shell and editor are now in the same place.

\begin{codebox}
zpwr vimrecentsudo      \# sudo edit recent files\\
\end{codebox}

When you need root to edit those files.

\begin{codebox}
zpwr vimrecentsudocd    \# cd + sudo edit recent files\\
\end{codebox}

The full combo: navigate, elevate, edit.


\vspace{12pt}
{\bfseries\color{neonyellow} \textbf{ZPWR} SUBSYSTEMS}\vspace{4pt}

\begin{codebox}
zpwr vimscripts         \# load all zpwr scripts into arglist\\
zpwr vimscriptedit      \# pick and edit 1+ scripts\\
zpwr vimautoload        \# load all autoloads into arglist\\
zpwr vimtests           \# edit all zpwr test files\\
\end{codebox}


\vspace{12pt}
{\bfseries\color{neonyellow} VIM CD}\vspace{4pt}

\begin{codebox}
zpwr vimcd              \# edit file, cd to its directory\\
\end{codebox}

Opens a file in vim, and when you exit, your shell is cd'd to the directory containing that file. Edit something, then be where it lives. Elegant.


\vspace{12pt}
{\bfseries\color{neonyellow} WORKFLOW EXAMPLE}\vspace{4pt}

\# Resume yesterday's work:

\begin{codebox}
zpwr vimrecentcd\\
\end{codebox}

\# Opens recent files, lands you in the right dir.

\# Mass rename a function across all scripts:

\begin{codebox}
zpwr vimscripts\\
:argdo \%s/oldFunc/newFunc/g | update\\
\end{codebox}

\# Done. Every script updated.


\vspace{16pt}

\section{Vim Search and Tags}


\vspace{10pt}
{\large\bfseries\color{neoncyan} VIM SEARCH AND TAGS}\vspace{6pt}

Finding files is step one. Finding definitions inside those files is step two. Tags are step three: jump to any symbol by name.


\vspace{12pt}
{\bfseries\color{neonyellow} VIM FILE SEARCH}\vspace{4pt}

\begin{codebox}
zpwr vimfilesearch          \# fzf find file -> open in vim\\
\end{codebox}

Lists files in the current tree via fzf. Select one, it opens in vim. No path typing.

\begin{codebox}
zpwr vimfilesearchedit      \# fzf find file -> path to BUFFER\\
\end{codebox}

Same search but puts the path on your command line.


\vspace{12pt}
{\bfseries\color{neonyellow} VIM WORD SEARCH}\vspace{4pt}

\begin{codebox}
zpwr vimwordsearch          \# fzf grep content -> open in vim\\
\end{codebox}

Searches file contents, shows matches in fzf. Select a match, vim opens at that location.

\begin{codebox}
zpwr vimwordsearchedit      \# same but path to BUFFER\\
\end{codebox}


\vspace{12pt}
{\bfseries\color{neonyellow} FULL DISK SEARCH INTO VIM}\vspace{4pt}

\begin{codebox}
zpwr vimfindsearch          \# find(1) -> fzf -> vim\\
zpwr vimfindsearchedit      \# find(1) -> fzf -> BUFFER\\
zpwr vimlocatesearch        \# locate -> fzf -> vim\\
zpwr vimlocatesearchedit    \# locate -> fzf -> BUFFER\\
\end{codebox}

The search escalation ladder from page 15 but with vim as the final destination instead of a pager.


\vspace{12pt}
{\bfseries\color{neonyellow} CTAGS: JUMP TO DEFINITION}\vspace{4pt}

\begin{codebox}
zpwr vimctags               \# open zpwr ctags in vim\\
\end{codebox}

Ctags indexes function and variable definitions. With tags loaded, Ctrl-] on any symbol jumps to its definition. Ctrl-T jumps back. Code navigation at the speed of thought.


\vspace{12pt}
{\bfseries\color{neonyellow} GNU GLOBAL TAGS}\vspace{4pt}

\begin{codebox}
zpwr vimgtags               \# open zpwr GNU global tags in vim\\
zpwr vimgtagsedit           \# edit zpwr GNU global tags\\
\end{codebox}

GNU Global is ctags on steroids. It indexes not just definitions but references too. Find every caller of a function, not just its definition.


\vspace{12pt}
{\bfseries\color{neonyellow} THE SEARCH MATRIX}\vspace{4pt}

What you know       Verb to use -------------       ----------- Filename fragment   vimfilesearch String in a file    vimwordsearch File somewhere      vimlocatesearch Symbol name         vimctags then Ctrl-] All references      vimgtags.


\vspace{16pt}

\section{Vim Config and Sessions}


\vspace{10pt}
{\large\bfseries\color{neoncyan} VIM CONFIG AND SESSIONS}\vspace{6pt}

Configuration is not a one-time event. It is a lifestyle. \textbf{ZPWR} gives you session commands that open the right config files as a group.


\vspace{12pt}
{\bfseries\color{neonyellow} THE SESSION VERBS}\vspace{4pt}

Each of these opens a curated set of config files in vim buffers, ready for synchronized editing:

\begin{codebox}
zpwr vrc          \# vim/nvim config session\\
\end{codebox}

Opens your vimrc or init.vim along with related files. All your editor settings in one session.

\begin{codebox}
zpwr brc          \# shell aliases session\\
\end{codebox}

Opens the tokens file and shell alias configs. Where you tune your personal shell experience.

\begin{codebox}
zpwr zrc          \# zshrc session\\
\end{codebox}

Opens the zshrc and zsh init chain. The root of your shell's boot sequence.

\begin{codebox}
zpwr trc          \# tmux.conf session\\
\end{codebox}

Opens tmux configuration files. Pane layouts, keybindings, status bar.


\vspace{12pt}
{\bfseries\color{neonyellow} CONFIG EDITING}\vspace{4pt}

\begin{codebox}
zpwr vimconfig    \# edit all zpwr configs in vim\\
\end{codebox}

Opens every zpwr configuration file. The master session. Use :bn and :bp to navigate between buffers.

\begin{codebox}
zpwr vimtokens   \# vim the .tokens.sh file\\
zpwr tokens      \# same thing\\
\end{codebox}

Quick jump to your personal aliases and secrets file. This is the file that is gitignored. Your sanctuary.


\vspace{12pt}
{\bfseries\color{neonyellow} PLUGIN INVENTORY}\vspace{4pt}

\begin{codebox}
zpwr vimplugincount  \# how many vim plugins installed\\
zpwr vimpluginlist   \# list all vim plugins\\
\end{codebox}

Know thy plugins. Count them. Name them. If the count scares you, it should.


\vspace{12pt}
{\bfseries\color{neonyellow} TYPICAL CONFIG WORKFLOW}\vspace{4pt}

\# Add a new shell alias:

\begin{codebox}
zpwr brc\\
\end{codebox}

\# Edit .tokens.sh, add your alias, :wq \# Source it or open a new terminal.

\# Tweak a vim keybinding:

\begin{codebox}
zpwr vrc\\
\end{codebox}

\# Navigate to the mapping, change it, :wq.

Three letters to open any config. No file paths. No guessing. Just intent.


\vspace{16pt}

\section{Emacs Mirror Universe}


\vspace{10pt}
{\large\bfseries\color{neoncyan} EMACS MIRROR UNIVERSE}\vspace{6pt}

For every vim verb, there is an emacs verb. Identical semantics. Different editor. Perfect symmetry. \textbf{ZPWR} does not take sides in the editor war. It arms both factions equally.


\vspace{12pt}
{\bfseries\color{neonyellow} THE SYMMETRY TABLE}\vspace{4pt}

vim verb              emacs verb --------              ----------

\begin{codebox}
vimall .............. emacsall\\
vimalledit .......... emacsalledit\\
vimautoload ......... emacsautoload\\
vimcd ............... emacscd\\
vimconfig ........... emacsconfig\\
vimctags ............ emacsctags\\
vimfilesearch ....... emacsfilesearch\\
vimfilesearchedit ... emacsfilesearchedit\\
vimfindsearch ....... emacsfindsearch\\
vimfindsearchedit ... emacsfindsearchedit\\
vimgtags ............ emacsgtags\\
vimgtagsedit ........ emacsgtagsedit\\
vimlocatesearch ..... emacslocatesearch\\
vimlocatesearchedit . emacslocatesearchedit\\
vimrecent ........... emacsrecent\\
vimrecentcd ......... emacsrecentcd\\
vimrecentsudo ....... emacsrecentsudo\\
vimrecentsudocd ..... emacsrecentsudocd\\
vimscripts .......... emacsscripts\\
vimscriptedit ....... emacsscriptedit\\
vimtests ............ emacstests\\
vimtokens ........... emacstokens\\
vimwordsearch ....... emacswordsearch\\
vimwordsearchedit ... emacswordsearchedit\\
vimplugincount ...... emacsplugincount\\
vimpluginlist ....... emacspluginlist\\
\end{codebox}


\vspace{12pt}
{\bfseries\color{neonyellow} EMACS-SPECIFIC}\vspace{4pt}

\begin{codebox}
zpwr emacsallserver       \# emacs all zpwr files (server mode)\\
zpwr emacsemacsconfig     \# emacs edit emacs-specific configs\\
zpwr emacszpwr            \# emacs the zpwr directory\\
\end{codebox}

The lesson: zpwr is editor-agnostic by design. Replace "vim" with "emacs" in any verb. It works. Both sides get the same firepower.


\vspace{16pt}


\chapter{PERFORMANCE CULT}


\begin{center}
\includegraphics[width=0.85\textwidth]{/Users/wizard/.zpwr/docs/screenshots/zpwr_bench_help.png}
\end{center}

\section{Benchmarking Startup}


\begin{center}
\includegraphics[width=0.8\textwidth]{/Users/wizard/.zpwr/docs/screenshots/zpwr_bench_help.png}
\end{center}


\vspace{10pt}
{\large\bfseries\color{neoncyan} BENCHMARKING STARTUP}\vspace{6pt}

A shell that takes 2 seconds to start is a shell that has already lost. \textbf{ZPWR} takes startup time seriously enough to build a benchmarking verb.


\vspace{12pt}
{\bfseries\color{neonyellow} THE BENCH VERB}\vspace{4pt}

\begin{codebox}
zpwr bench\\
\end{codebox}

Runs multiple shell startup iterations and reports: - Mean startup time - Percentiles (p50, p95, p99) - Baseline diff (zpwr overhead vs bare zsh)


\vspace{12pt}
{\bfseries\color{neonyellow} WHAT THE NUMBERS MEAN}\vspace{4pt}

< 100ms .... imperceptible. You are golden. 100-200ms .. noticeable if you open many terminals. 200-500ms .. getting sluggish. Time to investigate. > 500ms .... something is wrong. Check your plugins.


\vspace{12pt}
{\bfseries\color{neonyellow} HOW IT WORKS}\vspace{4pt}

The bench verb spawns N fresh zsh instances, each running through the full init chain. It measures wall-clock time from fork to prompt. Then it does the same for bare zsh (no zpwr). The diff tells you exactly what zpwr costs.


\vspace{12pt}
{\bfseries\color{neonyellow} WHAT SLOWS YOU DOWN}\vspace{4pt}

Common culprits in order of likelihood: 1. Too many vim/emacs plugins loading at shell init 2. Nvm/pyenv/rbenv hooking into every prompt 3. Uncompiled completion files (run zpwr compile) 4. Network calls in your init chain (stop that) 5. Massive \textbf{HISTFILE} being loaded synchronously.


\vspace{12pt}
{\bfseries\color{neonyellow} OPTIMIZATION CHECKLIST}\vspace{4pt}

\begin{codebox}
zpwr compile        \# recompile all cache and completions\\
zpwr compilefpath    \# recompile all fpath files\\
zpwr compilefiles    \# recompile all files\\
\end{codebox}

Compiled .zwc files load faster than source. This is the single biggest performance lever.


\vspace{12pt}
{\bfseries\color{neonyellow} BENCHMARK REGULARLY}\vspace{4pt}

Run zpwr bench after any config change. If the numbers go up, you broke something. If they go down, you are a performance wizard.

Your shell starts thousands of times a year. Every millisecond is multiplied by every terminal open. Measure. Optimize. Measure again.


\vspace{16pt}

\section{Flame Charts}


\begin{center}

\includegraphics[width=0.8\textwidth]{/Users/wizard/.zpwr/docs/screenshots/zpwr_flame_help.png}
\end{center}


\vspace{10pt}
{\large\bfseries\color{neoncyan} FLAME CHARTS}\vspace{6pt}

Bench tells you HOW LONG startup takes. Flame tells you WHERE the time goes.


\vspace{12pt}
{\bfseries\color{neonyellow} THE FLAME VERB}\vspace{4pt}

\begin{codebox}
zpwr flame\\
\end{codebox}

Generates a startup flame chart from zprof data. Shows you exactly which functions and files are consuming your startup time, in descending order.


\vspace{12pt}
{\bfseries\color{neonyellow} READING THE OUTPUT}\vspace{4pt}

The flame chart shows a ranked list of functions: Each entry shows: - Function/file name - Time consumed (milliseconds) - Percentage of total startup time.

The top of the list is where optimization effort pays off the most. If one function takes 40\% of startup, that is your target.


\vspace{12pt}
{\bfseries\color{neonyellow} HOW ZPROF WORKS}\vspace{4pt}

Zsh has a built-in profiler: zprof. When enabled, it instruments every function call during shell init and records the wall-clock time. The flame verb wraps this into a readable report.


\vspace{12pt}
{\bfseries\color{neonyellow} COMMON HOTSPOTS}\vspace{4pt}

What you will typically see at the top: compinit ........... completion system init compaudit .......... completion security check compdump ........... writing completion cache \_zpwr\_* ............ zpwr autoload functions nvm/pyenv/rbenv .... version manager hooks


\vspace{12pt}
{\bfseries\color{neonyellow} OPTIMIZATION TIPS}\vspace{4pt}

1. If compinit is slow, run zpwr compile to cache it 2. If a version manager dominates, lazy-load it 3. If an autoload is slow, check for disk I/O in it 4. If compaudit is slow, fix the directory permissions 5. After changes, run zpwr bench to verify improvement


\vspace{12pt}
{\bfseries\color{neonyellow} THE OPTIMIZATION LOOP}\vspace{4pt}

\begin{codebox}
zpwr flame      \# identify the bottleneck\\
\end{codebox}

... fix it ...

\begin{codebox}
zpwr bench      \# measure the improvement\\
zpwr flame      \# find the next bottleneck\\
\end{codebox}

Repeat until you hit diminishing returns.

Flame charts turn gut feelings into data. Stop guessing. Start profiling.


\vspace{16pt}

\section{Live Dashboard}


\begin{center}
\includegraphics[width=0.8\textwidth]{/Users/wizard/.zpwr/docs/screenshots/zpwr_bench_help.png}
\end{center}


\vspace{10pt}
{\large\bfseries\color{neoncyan} LIVE DASHBOARD}\vspace{6pt}

Bench is a snapshot. Flame is a profile. Top is a live feed of your shell's vital signs.


\vspace{12pt}
{\bfseries\color{neonyellow} THE TOP VERB}\vspace{4pt}

\begin{codebox}
zpwr top\\
\end{codebox}

Launches a live shell resource dashboard. Think htop, but for your zsh environment.


\vspace{12pt}
{\bfseries\color{neonyellow} WHAT EACH SECTION SHOWS}\vspace{4pt}

Shell Environment - Number of aliases currently defined - Number of functions loaded - Number of parameters (variables) in scope - Builtin and keyword counts.

Completion System - Completion functions registered - Cache file sizes and freshness.

Resource Usage - Memory footprint of the shell process - History file size and entry count - Autoload function count.


\vspace{12pt}
{\bfseries\color{neonyellow} INTERPRETING THE DATA}\vspace{4pt}

High alias count?    Normal for zpwr. 1000+ is fine. High function count? Also normal. Autoloaded functions are lazy-loaded, they cost nothing until called. Large \textbf{HISTFILE}?      Consider zpwr backuphistory and trim. Stale cache?         Run zpwr compile to refresh.


\vspace{12pt}
{\bfseries\color{neonyellow} COMPANION VERBS}\vspace{4pt}

\begin{codebox}
zpwr environmentcounts   \# one-shot environment stats\\
zpwr autoloadcount       \# total number of autoloads\\
zpwr autoloadlist        \# list all autoloads\\
zpwr vimplugincount      \# count vim plugins\\
zpwr emacsplugincount    \# count emacs plugins\\
\end{codebox}

Use these for quick spot checks without the full dashboard.


\vspace{12pt}
{\bfseries\color{neonyellow} THE PERFORMANCE QUARTET}\vspace{4pt}

zpwr bench    -> how fast is total startup? zpwr flame    -> ASCII flame chart of function hotspots zpwr startup  -> per-function zprof breakdown zpwr top      -> what is loaded right now?

Four verbs. Four perspectives. Complete observability. Your shell is no longer a black box. It is an open book.


\vspace{16pt}


\chapter{HEALTH \& DIAGNOSTICS}


\begin{center}
\includegraphics[width=0.85\textwidth]{/Users/wizard/.zpwr/docs/screenshots/zpwr_doctor_help.png}
\end{center}

\section{The Doctor Is In (1/2)}


\begin{center}
\includegraphics[width=0.8\textwidth]{/Users/wizard/.zpwr/docs/screenshots/zpwr_doctor_help.png}
\end{center}


\vspace{10pt}
{\large\bfseries\color{neoncyan} THE DOCTOR IS IN}\vspace{6pt}

{\small\itshape\color{dimtext} // every system needs a physician. yours lives at zpwr doctor. //}

Your shell is a living organism. It accumulates cruft, grows tumors in \textbf{PATH}, and occasionally develops broken symlinks. The doctor runs a full-body scan and tells you what hurts.


\vspace{12pt}
{\bfseries\color{neonyellow} INVOKE THE DOCTOR}\vspace{4pt}

\begin{codebox}
zpwr doctor              \# full diagnostic scan\\
zpwr doctor -h           \# help with detailed check list\\
\end{codebox}


\vspace{12pt}
{\bfseries\color{neonyellow} WHAT IT CHECKS}\vspace{4pt}

1. STALE .zwc FILES Compiled zsh files older than their source. Fix: zpwr recompile

2. ZCOMPDUMP FRESHNESS Completion dump older than 1 week. Fix: zpwr regenzsh

3. DUPLICATE \textbf{PATH} ENTRIES Every duplicate is a wasted hash lookup. Fix: zpwr deduppaths

4. DUPLICATE/MISSING \textbf{FPATH} ENTRIES Same disease, different organ.

5. BROKEN SYMLINKS Dangling symlinks in \$\textbf{ZPWR}.

Continued on next page...


\vspace{16pt}

\section{The Doctor Is In (2/2)}


\begin{center}
\includegraphics[width=0.8\textwidth]{/Users/wizard/.zpwr/docs/screenshots/zpwr_doctor_help.png}
\end{center}


\vspace{10pt}
{\large\bfseries\color{neoncyan} THE DOCTOR (continued)}\vspace{6pt}

6. NAMED DIRECTORIES Invalid named dirs and p10k pollution from auto\_name\_dirs leaking gitstatus vars.

7. DEPENDENCIES Checks: git, eza, fzf, tmux, perl, python3, zinit.

8. HISTORY \& BACKUPS Is \textbf{HISTFILE} present? How many backups exist?

9. CONFIG SYMLINKS Checks .zshrc, .vimrc, .tmux.conf are symlinks.


\vspace{12pt}
{\bfseries\color{neonyellow} READING THE RESULTS}\vspace{4pt}

Green checkmark  = healthy yellow warning   = degraded but functional red X            = broken, fix now.


\vspace{12pt}
{\bfseries\color{neonyellow} THE SIGNAL BAR}\vspace{4pt}

\#\#\#\#\#\#\#\#\#\#  = HEALTHY \#\#\#\#\#\#----  = WARNINGS \#\#--------  = CRITICAL.


\vspace{12pt}
{\bfseries\color{neonyellow} COMPANION VERBS}\vspace{4pt}

\begin{codebox}
zpwr stale                \# stale .zwc, orphan caches\\
zpwr pathaudit             \# deep PATH/FPATH audit\\
zpwr lint                  \# shellcheck + zsh -n\\
zpwr startup               \# zprof breakdown\\
zpwr ports                 \# listening ports\\
\end{codebox}


\vspace{12pt}
{\bfseries\color{neonyellow} PRO MOVE}\vspace{4pt}

Add zpwr doctor to weekly cron. Catch rot early.


\vspace{16pt}

\section{Dependency Cartography}


\begin{center}
\includegraphics[width=0.8\textwidth]{/Users/wizard/.zpwr/docs/screenshots/zpwr_deps_help.png}
\end{center}


\vspace{10pt}
{\large\bfseries\color{neoncyan} DEPENDENCY CARTOGRAPHY}\vspace{6pt}

{\small\itshape\color{dimtext} // mapping the call graph. who calls whom. who dies alone. //}

zpwr has hundreds of functions. They call each other in a tangled web of dependencies. The deps verb maps that web. Think of it as Google Maps for your function call graph.


\vspace{12pt}
{\bfseries\color{neonyellow} BASIC USAGE}\vspace{4pt}

\begin{codebox}
zpwr deps                          \# summary overview\\
zpwr deps zpwrClearList             \# what does zpwrClearList call?\\
zpwr deps -h                        \# full help\\
\end{codebox}


\vspace{12pt}
{\bfseries\color{neonyellow} THE FLAGS}\vspace{4pt}

-r, --reverse    Who calls this function? (reverse lookup) -o, --orphans    Functions called by nobody internally -s, --summary    Overview stats (default when no args) -d, --dot        Output dot format for graphviz rendering


\vspace{12pt}
{\bfseries\color{neonyellow} FORWARD LOOKUP: WHAT DOES X CALL?}\vspace{4pt}

\begin{codebox}
zpwr deps zpwrDoctor\\
\end{codebox}

Shows every function zpwrDoctor calls, plus depth-2 callees. Output:  zpwrDoctor -> zpwrPrettyPrint -> [zpwrLogDebug, ...]


\vspace{12pt}
{\bfseries\color{neonyellow} REVERSE LOOKUP: WHO CALLS X?}\vspace{4pt}

\begin{codebox}
zpwr deps -r zpwrPrettyPrint\\
\end{codebox}

Answers: "Who depends on zpwrPrettyPrint?" If the answer is 30 functions, that function is load-bearing. Touch it carefully.


\vspace{12pt}
{\bfseries\color{neonyellow} ORPHAN DETECTION}\vspace{4pt}

\begin{codebox}
zpwr deps --orphans\\
\end{codebox}

Lists functions that no other zpwr function calls. These are either entry points (user-invoked) or dead code. Orphans are not necessarily bad -- many are top-level verbs or zinit hooks. But if you wrote zpwrFoo and nobody calls it, maybe zpwrFoo was a fever dream.


\vspace{12pt}
{\bfseries\color{neonyellow} GRAPHVIZ EXPORT}\vspace{4pt}

\begin{codebox}
zpwr deps --dot > graph.dot\\
dot -Tsvg graph.dot -o graph.svg\\
\end{codebox}

Produces a visual dependency graph. Nodes are functions, edges are call relationships. The dot output uses cyberpunk colors: dark background, cyan nodes, magenta edges. Open graph.svg in a browser and behold the nervous system.


\vspace{12pt}
{\bfseries\color{neonyellow} THE SUMMARY VIEW}\vspace{4pt}

With no args, deps shows two ranked lists: TOP 10 MOST DEPENDENCIES (functions that call the most others) TOP 10 MOST DEPENDED ON  (functions called by the most others) Plus stats: total functions, orphan count, total call edges. The bar charts use Unicode block characters. Very aesthetic.


\vspace{12pt}
{\bfseries\color{neonyellow} HOW IT WORKS}\vspace{4pt}

Scans autoload dirs: common, darwin, linux, comp\_utils. Reads each function body and string-matches against all other function names. Also scans \textbf{\$ZPWR\_SCRIPTS/lib.sh.} It is O(n\textasciicircum{}2) and proud of it. Your function count is not big enough to matter.


\vspace{16pt}

\section{Command Autopsy}


\begin{center}
\includegraphics[width=0.8\textwidth]{/Users/wizard/.zpwr/docs/screenshots/zpwr_trace_help.png}
\end{center}


\vspace{10pt}
{\large\bfseries\color{neoncyan} COMMAND AUTOPSY}\vspace{6pt}

{\small\itshape\color{dimtext} // when you need to know why a command is slow, fat, or dead //}

The trace verb is a command profiler. It wraps any command in instrumentation and reports wall time, CPU time, memory, context switches, page faults, and more.


\vspace{12pt}
{\bfseries\color{neonyellow} BASIC USAGE}\vspace{4pt}

\begin{codebox}
zpwr trace ls -la                   \# full trace of ls -la\\
zpwr trace -q make                  \# quick: timing + resources only\\
zpwr trace -h                       \# full help\\
\end{codebox}


\vspace{12pt}
{\bfseries\color{neonyellow} THE FLAGS}\vspace{4pt}

-q, --quick      Timing + resource stats only (runs once) -s, --strace     Syscall counts only (Linux; no ltrace) -l, --ltrace     C library calls only (Linux; no strace) -a, --all        Both strace + ltrace (default behavior)


\vspace{12pt}
{\bfseries\color{neonyellow} WHAT YOU GET: THE RESOURCE CARD}\vspace{4pt}

Exit code          green = 0, red = nonzero wall time          real clock time with rating: INSTANT (<100ms), FAST (<1s), NOMINAL (<5s), SLOW (<30s), HEAVY (30s+) user/sys time      CPU time in user/kernel space peak RSS           max resident memory (auto-scaled B/KB/MB/GB) ctx switches       voluntary + involuntary context switches page reclaims      soft faults (memory mapped in from cache) page faults        hard faults (actual disk I/O for memory) block I/O          filesystem reads and writes signals recv       number of signals delivered IPC messages       sent/received inter-process messages shell delta        memory change of the parent shell child procs        before/after child process count.


\vspace{12pt}
{\bfseries\color{neonyellow} PLATFORM DIFFERENCES}\vspace{4pt}

MacOS:  \textbf{\textbf{/usr/bin/time}} -l for resource stats. SIP blocks dtruss/dtrace on most binaries. You get timing + memory + kernel stats. Linux:  strace -c for syscall counts, ltrace -c for C library calls. Command re-runs for each tool.


\vspace{12pt}
{\bfseries\color{neonyellow} PRACTICAL EXAMPLES}\vspace{4pt}

\begin{codebox}
zpwr trace git status              \# is git status slow in this repo?\\
zpwr trace -q rm -rf node\_modules  \# quick mode for destructive ops\\
zpwr trace -s npm install          \# syscalls only (Linux)\\
zpwr trace zsh -c 'source \textasciitilde{}/.zshrc' \# profile shell startup\\
\end{codebox}


\vspace{12pt}
{\bfseries\color{neonyellow} THE SIGNAL BAR}\vspace{4pt}

Bottom of output: a 10-block performance signal. \#\#\#\#\#\#\#\#\#\#  score 10 = under 100ms. Lightning. \#\#\#\#\#\#----  score 6  = under 2s. Acceptable. \#\#--------  score 2  = under 30s. Go get coffee.

WARNING: Default mode re-runs the command for strace/ltrace. Use -q for destructive commands (rm, truncate, etc).


\vspace{16pt}

\section{Usage Intelligence}


\vspace{10pt}
{\large\bfseries\color{neoncyan} USAGE INTELLIGENCE}\vspace{6pt}

{\small\itshape\color{dimtext} // your shell keeps a diary. these verbs read it. //}

You have \$\{\#\textbf{ZPWR\_VERBS}\} verbs and thousands of aliases. Which ones do you actually use? Funcrank and aliasrank mine your shell history to answer that question.


\vspace{12pt}
{\bfseries\color{neonyellow} FUNCRANK: FUNCTION \& VERB PROFILER}\vspace{4pt}

\begin{codebox}
zpwr funcrank                       \# top 20 functions and verbs\\
zpwr funcrank 40                    \# top 40\\
zpwr funcrank --reset               \# nuke cache, rebuild from scratch\\
zpwr funcrank --reset 50            \# rebuild then show top 50\\
\end{codebox}

Funcrank scans your \textbf{HISTFILE} for occurrences of known function names and zpwr verbs. It produces two ranked lists: 1. Top functions by history occurrence count 2. Top zpwr verbs by how often you type "zpwr <verb>" Each entry gets a Unicode bar chart. The biggest user wins.


\vspace{12pt}
{\bfseries\color{neonyellow} ALIASRANK: REVERSE EXPANSION PROFILER}\vspace{4pt}

\begin{codebox}
zpwr aliasrank                      \# top 20 by expanded usage\\
zpwr aliasrank 50                   \# top 50\\
zpwr aliasrank --reset              \# force cache rebuild\\
\end{codebox}

Here is the trick: zpwr uses zsh-expand to expand aliases before they hit history. So history contains the expanded form, not the alias name. Aliasrank does a reverse lookup: it matches history lines against alias values to figure out which alias you would have typed.


\vspace{12pt}
{\bfseries\color{neonyellow} THE CACHE SYSTEM}\vspace{4pt}

\textbf{\$ZPWR\_LOCAL/zpwr-aliasrank-cache.zsh}   command counts \textbf{\$ZPWR\_LOCAL/zpwr-verbrank-cache.zsh}    verb counts \textbf{\$ZPWR\_LOCAL/zpwr-rank-stamp}            history line watermark.

The cache auto-rebuilds when history grows by 1000+ lines. The watermark file stores the last-seen history line count. No need to rebuild every time -- incremental is the way. Use --reset to force a full rebuild if cache seems stale.


\vspace{12pt}
{\bfseries\color{neonyellow} INTERPRETING THE RESULTS}\vspace{4pt}

If your top verb is help, you are still learning. Good. If your top verb is cleanall, you break things a lot. If your top function is cd, you need more named dirs. If an alias never appears, maybe it is too long to remember.


\vspace{12pt}
{\bfseries\color{neonyellow} PRACTICAL INTELLIGENCE}\vspace{4pt}

Run funcrank monthly. Notice patterns: - Verbs you never use might need better names or keybindings - Functions with high counts are candidates for optimization - Low-count zpwr verbs might benefit from shorter aliases.

Your history is a behavioral log. These verbs turn it into actionable intelligence. Know thyself, operative.


\vspace{16pt}


\chapter{ENVIRONMENT MASTERY}


\begin{center}

\includegraphics[width=0.85\textwidth]{/Users/wizard/.zpwr/docs/screenshots/art_system_core.png}
\end{center}

\section{Know Your Shell}


\vspace{10pt}
{\large\bfseries\color{neoncyan} KNOW YOUR SHELL}\vspace{6pt}

{\small\itshape\color{dimtext} // you cannot master what you cannot measure //}

Three verbs to audit your environment. Each one reveals a different cross-section of the machine you live inside.


\vspace{12pt}
{\bfseries\color{neonyellow} ENVVARS: THE \textbf{ZPWR} VARIABLES}\vspace{4pt}

\begin{codebox}
zpwr envvars                        \# list all ZPWR\_* env vars\\
zpwr ev                              \# same thing, shorter\\
\end{codebox}

Shows every environment variable starting with \textbf{ZPWR}\_. These are the knobs and dials of the framework: paths, flags, plugin manager config, colorscheme, etc. If something is misconfigured, this is where you look.


\vspace{12pt}
{\bfseries\color{neonyellow} ENVVARSALL: THE FULL PICTURE}\vspace{4pt}

\begin{codebox}
zpwr envvarsall                     \# every env var, all of them\\
zpwr eva                             \# same thing, shorter\\
\end{codebox}

Dumps every exported variable. This includes \textbf{ZPWR}\_*, \textbf{PATH}, \textbf{HOME}, \textbf{EDITOR}, \textbf{LANG}, and everything else. Pipe it to fzf when you are hunting for a specific var:

\begin{codebox}
zpwr eva | fzf\\
\end{codebox}


\vspace{12pt}
{\bfseries\color{neonyellow} ENVIRONMENTCOUNTS: THE NUMBERS GAME}\vspace{4pt}

\begin{codebox}
zpwr environmentcounts              \# full environment statistics\\
zpwr envcounts                      \# alias\\
zpwr e                               \# shortest form\\
\end{codebox}

This is the census. It counts everything loaded in your current shell session:

Aliases         regular + global + suffix aliases functions       autoloaded + defined functions builtins        zsh builtin commands parameters      shell variables (scalar, array, assoc) completions     registered completion functions modules         loaded zsh modules named dirs      hash -d entries \textbf{PATH} entries    directories in \textbf{PATH} \textbf{FPATH} entries   directories in \textbf{FPATH}

The numbers tell you how armed you are. A fresh zpwr install might show thousands of aliases, hundreds of functions, and many registered completions. If the numbers are low, something failed to load.


\vspace{12pt}
{\bfseries\color{neonyellow} PRO TIP: PAIR THEM}\vspace{4pt}

\begin{codebox}
zpwr doctor                   \# is anything broken?\\
zpwr e                         \# how much loaded?\\
zpwr ev                        \# are my zpwr vars right?\\
\end{codebox}

Run all three after installing or updating zpwr. They are your post-deployment smoke test.


\vspace{16pt}

\section{The Option Matrix}


\vspace{10pt}
{\large\bfseries\color{neoncyan} THE OPTION MATRIX}\vspace{6pt}

{\small\itshape\color{dimtext} // zsh has 200+ options. most people know 5. time to fix that. //}

The zsh option system is a maze of boolean switches that control everything from globbing to history to error handling. These verbs help you navigate the labyrinth.


\vspace{12pt}
{\bfseries\color{neonyellow} OPTIONS: THE BOOLEAN SWITCHES}\vspace{4pt}

\begin{codebox}
zpwr options                        \# list all set options\\
\end{codebox}

Shows every zsh option that is currently ON. The important ones for zpwr: extendedglob      (\#, \textasciitilde{}, \textasciicircum{} in globs)  -- zpwr needs this autopushd         (cd pushes to dirstack) histignorealldups (no duplicate history) sharehistory      (history shared across sessions) completeinword    (tab complete mid-word) promptsubst       (parameter expansion in prompts)

If extendedglob is off, half of zpwr's glob tricks die. If sharehistory is off, your sessions are islands.


\vspace{12pt}
{\bfseries\color{neonyellow} MODULES: THE LOADED EXTENSIONS}\vspace{4pt}

\begin{codebox}
zpwr modules                        \# list loaded zsh modules\\
\end{codebox}

Zsh modules are C libraries that add features: zsh/datetime      EPOCHREALTIME, strftime zsh/stat          zstat builtin for file metadata zsh/zutil         zparseopts for flag parsing zsh/complist       menuselect completion listing zsh/parameter      \$functions, \$aliases, \$parameters hashes zsh/terminfo       terminal capability strings.

If zsh/parameter is not loaded, zpwr env verbs break. If zsh/datetime is missing, timing-based verbs fail. Load a missing module: zmodload zsh/modulename.


\vspace{12pt}
{\bfseries\color{neonyellow} ZSTYLE: THE FINE-TUNING KNOBS}\vspace{4pt}

\begin{codebox}
zpwr zstyle                         \# list all zstyle settings\\
\end{codebox}

zstyle is zsh's configuration registry. It controls: :completion:*     how tab completion behaves :vcs\_info:*       git/svn info in prompts :zpwr:*           zpwr-specific settings.

Example styles zpwr sets: zstyle ':completion:*' menu select zstyle ':completion:*' matcher-list 'm:\{a-z\}=\{A-Z\}' The first enables arrow-key menu selection in tab complete. The second makes completion case-insensitive.


\vspace{12pt}
{\bfseries\color{neonyellow} DETECTIVE WORK}\vspace{4pt}

Completions acting weird? Check:

\begin{codebox}
zpwr options    \# is completeinword on?\\
zpwr modules    \# is zsh/complist loaded?\\
zpwr zstyle     \# what completion styles are set?\\
\end{codebox}

The answer is usually hiding in one of these three.


\vspace{16pt}

\section{Cache Search}


\vspace{10pt}
{\large\bfseries\color{neoncyan} CACHE SEARCH}\vspace{6pt}

{\small\itshape\color{dimtext} // the environment is a haystack. these are your magnets. //}

zpwr maintains an environment cache for fast lookup. These verbs let you search, inspect, and edit that cache without grepping through raw variable dumps.


\vspace{12pt}
{\bfseries\color{neonyellow} ENVIRONMENTCACHESEARCH: FZF ENV NAVIGATION}\vspace{4pt}

\begin{codebox}
zpwr environmentcachesearch         \# fzf through cached env\\
\end{codebox}

Opens an fzf interface over the environment cache. Type a few characters to filter. The cache includes: - All exported environment variables - \textbf{ZPWR}\_* configuration variables - \textbf{PATH} and \textbf{FPATH} entries - Named directory mappings.

This is faster than env | grep because it uses the pre-built cache file rather than querying the live env.


\vspace{12pt}
{\bfseries\color{neonyellow} ENVACCEPT: ACCEPT A CACHE ENTRY}\vspace{4pt}

\begin{codebox}
zpwr envaccept                      \# accept/apply cached env entry\\
\end{codebox}

After finding a variable in the cache, this verb applies it to your current session. Useful when the cache has a value you want but the live env does not.


\vspace{12pt}
{\bfseries\color{neonyellow} ENVEDIT: EDIT THE ENVIRONMENT}\vspace{4pt}

\begin{codebox}
zpwr envedit                        \# open env cache in EDITOR\\
\end{codebox}

Opens the environment cache file in your editor. Make changes, save, and the cache is updated. This is the surgical approach: when you know exactly which variable needs tweaking.


\vspace{12pt}
{\bfseries\color{neonyellow} THE CACHE LIFECYCLE}\vspace{4pt}

1. zpwr regenenvcache   builds the cache from live env 2. environmentcachesearch fuzzy-searches the cache 3. envaccept             applies an entry 4. envedit               manual edits

The cache lives at \$\textbf{ZPWR\_LOCAL}. It is regenerated during shell startup and when you run zpwr regenenvcache.


\vspace{12pt}
{\bfseries\color{neonyellow} WORKFLOW EXAMPLE}\vspace{4pt}

"What is \textbf{ZPWR\_EDITOR} set to?"

\begin{codebox}
zpwr environmentcachesearch    \# type 'EDITOR' in fzf\\
\end{codebox}

Found it: \textbf{ZPWR\_EDITOR}=nvim.

"I want to change a bunch of \textbf{ZPWR}\_ vars."

\begin{codebox}
zpwr envedit                       \# opens in EDITOR\\
\end{codebox}

Edit the vars, save, done.

"Something is off, rebuild the cache."

\begin{codebox}
zpwr regenenvcache                 \# fresh cache from live env\\
\end{codebox}


\vspace{12pt}
{\bfseries\color{neonyellow} PRO TIP}\vspace{4pt}

If you are debugging a missing env var, check two places: zpwr environmentcachesearch (cached) and zpwr eva (live). If cache has it but live does not, run zpwr regenenvcache.


\vspace{16pt}

\section{Snapshots \& Time Travel}


\begin{center}
\includegraphics[width=0.8\textwidth]{/Users/wizard/.zpwr/docs/screenshots/zpwr_snapshot_help.png}
\end{center}


\vspace{10pt}
{\large\bfseries\color{neoncyan} SNAPSHOTS \& TIME TRAVEL}\vspace{6pt}

{\small\itshape\color{dimtext} // save your universe. restore it when everything goes wrong. //}

The snapshot verb captures your entire terminal environment to a portable directory. The restore verb brings it back. This is git for your shell state.


\vspace{12pt}
{\bfseries\color{neonyellow} TAKING A SNAPSHOT}\vspace{4pt}

\begin{codebox}
zpwr snapshot                       \# timestamp-named snapshot\\
zpwr snapshot myproject              \# named snapshot\\
zpwr snapshot before-refactor        \# meaningful name\\
\end{codebox}


\vspace{12pt}
{\bfseries\color{neonyellow} WHAT GETS CAPTURED}\vspace{4pt}

Hash dirs         named directory mappings (hash -d) aliases           regular + global + suffix aliases env vars          all exported variables (secrets excluded) zpwr vars         \textbf{ZPWR\_VARS} associative array functions         user function count tracking history           full shell history dump tmux             sessions, windows, panes, resurrect state vim sessions      saved vim session files dir stack         current directory and pushd stack git status        branch, commit, dirty state of cwd.

Storage: \textbf{\$ZPWR\_LOCAL/snapshots/}<name>/ Each snapshot is a directory of plain zsh files. Human-readable. Sourceable. Portable.


\vspace{12pt}
{\bfseries\color{neonyellow} RESTORING A SNAPSHOT}\vspace{4pt}

\begin{codebox}
zpwr restore                        \# list available snapshots\\
zpwr restore myproject               \# restore everything\\
zpwr restore myproject aliases env    \# restore only aliases+env\\
zpwr restore myproject dirs           \# just cd to saved pwd\\
\end{codebox}


\vspace{12pt}
{\bfseries\color{neonyellow} SELECTIVE RESTORE COMPONENTS}\vspace{4pt}

All       everything (default when no component given) hash      named directories only aliases   regular, global, suffix aliases env       environment variables tmux      stage tmux-resurrect file (press prefix+Ctrl-r) vim       vim session files to \textbf{\textasciitilde{}/.vim/sessions} history   merge snapshot history into current session dirs      cd to the saved working directory.


\vspace{12pt}
{\bfseries\color{neonyellow} WHEN TO USE SNAPSHOTS}\vspace{4pt}

Before a big update:     zpwr snapshot before-update Before experiments:      zpwr snapshot stable-state Switching projects:      zpwr snapshot projectA New machine setup:       copy snapshot dir, restore on new host.

NOTE: Sensitive vars (PASSWORD, SECRET, TOKEN, AWS\_SECRET*, GITHUB\_TOKEN) are automatically excluded from env snapshots. Your secrets stay secret.

History merge is additive -- it does not overwrite your current history. It uses fc -R to load entries.


\vspace{16pt}


\chapter{THE JANITOR'S CLOSET}


\begin{center}

\includegraphics[width=0.85\textwidth]{/Users/wizard/.zpwr/docs/screenshots/art_override.png}
\end{center}

\section{Cleanup Operations}


\vspace{10pt}
{\large\bfseries\color{neoncyan} CLEANUP OPERATIONS}\vspace{6pt}

{\small\itshape\color{dimtext} // entropy is the enemy. the janitor is your friend. //}

Shell environments accumulate garbage like a hoarder's attic. Cache files grow stale. Temp dirs fill up. Logs balloon. These verbs are the Marie Kondo of your \$\textbf{ZPWR}.


\vspace{12pt}
{\bfseries\color{neonyellow} CLEANALL: THE THERMONUCLEAR BROOM}\vspace{4pt}

\begin{codebox}
zpwr cleanall                       \# clean everything at once\\
\end{codebox}

Calls all the individual clean verbs in sequence. This is the "I do not know what is wrong but I want a fresh start" button. Safe but thorough.


\vspace{12pt}
{\bfseries\color{neonyellow} INDIVIDUAL CLEAN VERBS}\vspace{4pt}

CLEANCACHE.

\begin{codebox}
zpwr cleancache                     \# clear general caches\\
\end{codebox}

Wipes cached data that zpwr generates during operation. Safe to run any time. Everything rebuilds on demand.

CLEANCOMPCACHE.

\begin{codebox}
zpwr cleancompcache                 \# clear completion cache\\
\end{codebox}

Removes the zsh completion cache directory. Run this when tab completions are showing stale results or when a new command does not appear in completions. Follow up with zpwr regenzsh to rebuild zcompdump.

CLEANDIRS.

\begin{codebox}
zpwr cleandirs                      \# clear cached dir listings\\
\end{codebox}

Removes directory cache files used by navigation commands. If directory completions point to dirs that no longer exist, this verb fixes that.

CLEANTEMP.

\begin{codebox}
zpwr cleantemp                      \# clear temp files\\
\end{codebox}

Removes temporary files zpwr created during operation. These are transient by nature -- safe to nuke.

CLEANDIRSANDTEMP.

\begin{codebox}
zpwr cleandirsandtemp               \# dirs + temp combined\\
\end{codebox}

Convenience combo of cleandirs and cleantemp.

CLEANLOG.

\begin{codebox}
zpwr cleanlog                       \# truncate log files\\
\end{codebox}

The zpwr logfile grows forever if you never clean it. This verb truncates or removes old log entries. Check log size: ls -lh \$\textbf{ZPWR\_LOGFILE}


\vspace{12pt}
{\bfseries\color{neonyellow} THE CLEAN WORKFLOW}\vspace{4pt}

Light cleanup (weekly):

\begin{codebox}
zpwr cleantemp \&\& zpwr cleanlog\\
\end{codebox}

Medium cleanup (monthly):

\begin{codebox}
zpwr cleancache \&\& zpwr cleancompcache \&\& zpwr regenzsh\\
\end{codebox}

Full cleanup (when things feel sluggish):

\begin{codebox}
zpwr cleanall\\
\end{codebox}

None of these delete your configuration or history. They only remove regenerable cache and temp data.


\vspace{16pt}

\section{Git Cache Hygiene}


\vspace{10pt}
{\large\bfseries\color{neoncyan} GIT CACHE HYGIENE}\vspace{6pt}

{\small\itshape\color{dimtext} // zpwr tracks your git repos. the cache needs a bath sometimes. //}

zpwr maintains caches of git repository locations so that repo-aware verbs (navigation, status, batch operations) do not have to scan the filesystem every time. Three verbs manage these caches.


\vspace{12pt}
{\bfseries\color{neonyellow} CLEANGITCACHE: THE GENERAL REPO CACHE}\vspace{4pt}

\begin{codebox}
zpwr cleangitcache                  \# clear all git repo caches\\
\end{codebox}

Wipes the master list of known git repositories. This cache is built by scanning for .git directories and is used by verbs like zpwr fordir and repo navigation commands.

When to use it: - You moved or deleted git repos - Navigation is suggesting repos that no longer exist - You cloned new repos and they are not showing up.

After clearing, run zpwr regengitrepocache to rebuild.


\vspace{12pt}
{\bfseries\color{neonyellow} CLEANGITCLEANCACHE: CLEAN REPOS ONLY}\vspace{4pt}

\begin{codebox}
zpwr cleangitcleancache             \# clear clean repo cache\\
\end{codebox}

zpwr separately tracks which repos have a clean working tree (no uncommitted changes). This cache speeds up operations that only care about clean repos.

Clear this when: - The "clean repos" list seems wrong - You committed in a repo but it still shows as dirty.


\vspace{12pt}
{\bfseries\color{neonyellow} CLEANGITDIRTYCACHE: DIRTY REPOS ONLY}\vspace{4pt}

\begin{codebox}
zpwr cleangitdirtycache             \# clear dirty repo cache\\
\end{codebox}

The mirror of the clean cache. Tracks repos with uncommitted changes. Clear this when: - A repo shows as dirty but you already committed - The dirty count in the banner seems wrong.

Rebuild with: zpwr regengitdirtyrepocache.


\vspace{12pt}
{\bfseries\color{neonyellow} THE GIT CACHE ARCHITECTURE}\vspace{4pt}

Regengitrepocache        scans for all git repos regengitdirtyrepocache   classifies clean vs dirty cleangitcache            wipes the master list cleangitcleancache       wipes the clean-only list cleangitdirtycache       wipes the dirty-only list.

The caches live in \$\textbf{ZPWR\_LOCAL} as plain text files. One repo path per line. Grep-friendly. Human-readable.


\vspace{12pt}
{\bfseries\color{neonyellow} THE FULL REFRESH}\vspace{4pt}

\begin{codebox}
zpwr cleangitcache\\
zpwr regengitrepocache\\
zpwr regengitdirtyrepocache\\
\end{codebox}

Three commands. Fresh git cache. Takes a few seconds depending on how many repos you have scattered around.


\vspace{16pt}

\section{Path Deduplication}


\vspace{10pt}
{\large\bfseries\color{neoncyan} PATH DEDUPLICATION}\vspace{6pt}

{\small\itshape\color{dimtext} // every duplicate entry is a tiny crime against performance //}

\textbf{PATH} and \textbf{FPATH} grow duplicates like weeds. Every subshell, every sourced script, every tool that helpfully "adds itself to your \textbf{PATH}" can insert duplicates. The deduppaths verb removes them.


\vspace{12pt}
{\bfseries\color{neonyellow} THE VERB}\vspace{4pt}

\begin{codebox}
zpwr deduppaths                     \# deduplicate PATH and FPATH\\
\end{codebox}

Removes duplicate entries from both \textbf{PATH} and \textbf{FPATH} while preserving the order of first occurrence. If \textbf{\textbf{/usr/bin}} appears 3 times, only the first one survives.


\vspace{12pt}
{\bfseries\color{neonyellow} WHY DUPLICATES HAPPEN}\vspace{4pt}

1. Shell inheritance Parent shell sets \textbf{PATH}. Child shell re-sources .zshrc. \textbf{PATH} grows. Nested tmux sessions make it worse.

2. Multiple init scripts Homebrew, nvm, pyenv, rbenv, cargo -- each one adds to \textbf{PATH} in their init hook. Source them twice, \textbf{PATH} grows.

3. \textbf{\textbf{/etc/zprofile}} on macOS macOS ships with \textbf{\textbf{/etc/zprofile}} that runs path\_helper. path\_helper reads \textbf{\textbf{/etc/paths}} and \textbf{\textbf{/etc/paths.d/}} and prepends entries. This runs on every login shell, BEFORE your .zshrc. Duplicates guaranteed.

4. Manual \textbf{PATH} manipulation export \textbf{PATH}="/my/dir:\$\textbf{PATH}" in .zshrc. Every time .zshrc is sourced, another copy appears.


\vspace{12pt}
{\bfseries\color{neonyellow} HOW TO PREVENT DUPLICATES}\vspace{4pt}

Zsh has a built-in mechanism: the typeset -U flag.

\begin{codebox}
typeset -U path fpath\\
\end{codebox}

This marks the arrays as "unique" -- zsh automatically removes duplicates when entries are added. zpwr sets this early in startup. But if something adds to \textbf{PATH} before zpwr loads, duplicates can sneak in.


\vspace{12pt}
{\bfseries\color{neonyellow} DIAGNOSING THE PROBLEM}\vspace{4pt}

\begin{codebox}
zpwr pathaudit                       \# full audit: dupes, missing, writable\\
zpwr doctor                         \# checks for duplicate PATH/FPATH\\
print -l path | sort | uniq -d     \# show only duplicated entries\\
echo \#path                         \# count PATH entries\\
\end{codebox}

A healthy \textbf{PATH} has 10-30 entries. If you see 50+, something is adding entries in a loop.


\vspace{12pt}
{\bfseries\color{neonyellow} THE FIX}\vspace{4pt}

\begin{codebox}
zpwr deduppaths                     \# instant dedup\\
\end{codebox}

Run it. Then check: echo \$\#path If the count dropped, you had duplicates. If it drops a lot, investigate the source.

For a permanent fix, find the script that adds duplicates and guard it with a check:

\begin{codebox}
{[}[ -z \{path[(r)/my/dir]\} ]] \&\& path=(/my/dir path)\\
\end{codebox}


\vspace{16pt}

\section{The Nuclear Options}


\vspace{10pt}
{\large\bfseries\color{neoncyan} THE NUCLEAR OPTIONS}\vspace{6pt}

{\small\itshape\color{dimtext} // when gentle cleanup is not enough. when you need to nuke from orbit. //}

Most clean verbs are surgical. These two are orbital strikes. They clean, regenerate, rebuild, and update in one shot. Use them when you want a fresh start without reinstalling.


\vspace{12pt}
{\bfseries\color{neonyellow} CLEANREFRESHUPDATE: THE BIG RED BUTTON}\vspace{4pt}

\begin{codebox}
zpwr cleanrefreshupdate             \# clean + regen + update\\
\end{codebox}

This is the full-cycle maintenance verb. It: 1. Cleans all caches (equivalent to cleanall) 2. Regenerates compiled files, completions, env cache 3. Runs the update cycle for zpwr and plugins.

Think of it as "factory reset without losing data." Your configuration, history, and local files survive. Everything else gets rebuilt from source.

When to use it: - Shell startup feels slow and you have tried individual fixes - After a major zsh or macOS upgrade - When individual clean verbs did not solve the problem - When you want to feel powerful.

TIME COST: Several minutes. This is not instant. It downloads plugin updates and recompiles everything.


\vspace{12pt}
{\bfseries\color{neonyellow} CLEANUPDATEDEPS: THE DEPENDENCY PURGE}\vspace{4pt}

\begin{codebox}
zpwr cleanupdatedeps                \# clean deps + update them\\
\end{codebox}

Focused on external dependencies rather than caches. This verb: 1. Cleans stale dependency state 2. Updates all managed dependencies.

When to use it: - Plugin versions are stuck at old commits - A dependency update broke something - You want the latest of everything.


\vspace{12pt}
{\bfseries\color{neonyellow} DECISION MATRIX}\vspace{4pt}

Symptom                          Verb.

Stale completions                zpwr cleancompcache Slow startup                     zpwr recompile Missing new repos in nav         zpwr cleangitcache General weirdness                zpwr cleanall Everything is broken              zpwr cleanrefreshupdate Plugins are ancient               zpwr cleanupdatedeps.


\vspace{12pt}
{\bfseries\color{neonyellow} THE ESCALATION \textbf{PATH}}\vspace{4pt}

Always start small and escalate:

\begin{codebox}
zpwr cleantemp                      \# level 1: temp files\\
zpwr cleancache                     \# level 2: all caches\\
zpwr cleanall                       \# level 3: everything\\
zpwr cleanrefreshupdate             \# level 4: clean + rebuild + update\\
\end{codebox}

Do not start with level 4. That is like calling an airstrike because you saw a spider.


\vspace{16pt}


\chapter{BUILD SYSTEM}


\begin{center}
\includegraphics[width=0.85\textwidth]{/Users/wizard/.zpwr/docs/screenshots/figlet_exec.png}
\end{center}

\section{Compilation 101}


\vspace{10pt}
{\large\bfseries\color{neoncyan} COMPILATION 101}\vspace{6pt}

{\small\itshape\color{dimtext} // .zwc files: compiled zsh. faster loading. dark magic. //}

Zsh can compile scripts to an internal bytecode format stored in .zwc (Zsh Word Code) files. Loading compiled files is faster than parsing source text every time.


\vspace{12pt}
{\bfseries\color{neonyellow} WHAT IS A .ZWC FILE?}\vspace{4pt}

A .zwc file is the compiled form of a .zsh script or an autoload function directory. When zsh finds foo.zwc next to foo.zsh, it loads the compiled version -- skipping the parsing step entirely.

For individual files: zcompile foo.zsh -> foo.zsh.zwc For directories:      zcompile dir.zwc dir/* -> dir.zwc.

The speed difference is measurable on large files and during startup when many files are sourced.


\vspace{12pt}
{\bfseries\color{neonyellow} RECOMPILE: BUILD ALL .ZWC FILES}\vspace{4pt}

\begin{codebox}
zpwr recompile                      \# compile all zpwr files\\
zpwr compile                        \# same thing\\
\end{codebox}

Scans all zpwr source files and compiles them to .zwc. This includes: - Autoload function files in \$\textbf{ZPWR\_AUTOLOAD} - Shell scripts in \$\textbf{ZPWR\_ENV} - Any .zsh file that has changed since last compile.


\vspace{12pt}
{\bfseries\color{neonyellow} RELATED COMPILE VERBS}\vspace{4pt}

\begin{codebox}
zpwr recompilefiles                 \# compile specific file set\\
zpwr compilefiles                   \# alias for above\\
zpwr recompilefpath                 \# compile FPATH directories\\
zpwr compilefpath                   \# alias for above\\
\end{codebox}


\vspace{12pt}
{\bfseries\color{neonyellow} DECOMPILE: REMOVE .ZWC FILES}\vspace{4pt}

\begin{codebox}
zpwr decompile                      \# remove all .zwc files\\
zpwr uncompile                      \# alias for above\\
\end{codebox}

Deletes all .zwc files in zpwr. After this, zsh falls back to parsing source files directly. Useful when: - Debugging a function that behaves differently than source - You suspect a stale .zwc file is loading old code - zpwr doctor reports stale .zwc files.


\vspace{12pt}
{\bfseries\color{neonyellow} REFRESHZWC: SMART REBUILD}\vspace{4pt}

\begin{codebox}
zpwr refreshzwc                     \# refresh only stale .zwc files\\
\end{codebox}

Only recompiles .zwc files that are older than their source. Faster than a full recompile when only a few files changed.


\vspace{12pt}
{\bfseries\color{neonyellow} THE COMPILATION WORKFLOW}\vspace{4pt}

After editing zpwr source files:

\begin{codebox}
zpwr recompile                    \# rebuild .zwc files\\
\end{codebox}

If something is weird after an update:

\begin{codebox}
zpwr decompile \&\& zpwr recompile  \# clean rebuild\\
\end{codebox}

Quick check for stale files:

\begin{codebox}
zpwr doctor                       \# reports stale .zwc files\\
\end{codebox}

The .zwc files are safe to delete. They are purely a performance cache. Source files are always the authority.


\vspace{16pt}

\section{Compsys Cache}


\vspace{10pt}
{\large\bfseries\color{neoncyan} COMPSYS CACHE}\vspace{6pt}

{\small\itshape\color{dimtext} // tab completion is a religion. the zcompdump is its scripture. //}

The zsh completion system (compsys) is the most complex part of zsh. It maintains a dump file (zcompdump) that caches all registered completions. When it breaks, tab completion goes haywire. Here is how to fix it.


\vspace{12pt}
{\bfseries\color{neonyellow} WHAT IS ZCOMPDUMP?}\vspace{4pt}

A file (usually \textbf{\textasciitilde{}/.zcompdump} or in \$\textbf{ZPWR\_LOCAL}) that contains all compdef registrations. When zsh starts, it sources this file instead of re-scanning all completion functions. This saves 100-500ms of startup time.

Location: \$\textbf{ZSH\_COMPDUMP} The file starts with \#files:TIMESTAMP and contains compdef calls for every registered completion.


\vspace{12pt}
{\bfseries\color{neonyellow} REGENZSH: THE COMPLETION RESET}\vspace{4pt}

\begin{codebox}
zpwr regenzsh                       \# regenerate zcompdump\\
\end{codebox}

Deletes the old zcompdump and forces compinit to rebuild it. This is the \#1 fix for completion problems:

Symptom: New command installed but tab complete ignores it Fix:     zpwr regenzsh.

Symptom: Tab shows wrong completions for a command Fix:     zpwr cleancompcache \&\& zpwr regenzsh.

Symptom: "command not found: compdef" errors on startup Fix:     zpwr regenzsh.

Symptom: Startup is slow (500ms+ for completion init) Fix:     zpwr regenzsh (forces fresh, optimized dump)


\vspace{12pt}
{\bfseries\color{neonyellow} THE COMPLETION CACHE DIRECTORY}\vspace{4pt}

Separate from zcompdump, zsh also caches individual completion results (like package names, hostnames, etc.) in a cache directory.

\begin{codebox}
zpwr cleancompcache                 \# wipe the cache directory\\
\end{codebox}

This forces completions to re-query their sources. Useful when package lists or hostnames changed.


\vspace{12pt}
{\bfseries\color{neonyellow} WHY COMPLETIONS BREAK}\vspace{4pt}

1. Stale zcompdump -- new compdef not registered 2. Corrupted zcompdump -- syntax error in dump file 3. Missing \textbf{FPATH} entry -- completion function not found 4. Wrong function loaded -- .zwc has old version 5. Plugin conflict -- two plugins register for same cmd


\vspace{12pt}
{\bfseries\color{neonyellow} THE NUCLEAR FIX}\vspace{4pt}

\begin{codebox}
zpwr cleancompcache\\
zpwr decompile\\
zpwr regenzsh\\
zpwr recompile\\
\end{codebox}

Four commands: nuke cache, remove .zwc, regen dump, recompile. This fixes 99\% of completion problems. The remaining 1\% is probably a bug in the completion function itself.


\vspace{16pt}

\section{Tag Generation}


\vspace{10pt}
{\large\bfseries\color{neoncyan} TAG GENERATION}\vspace{6pt}

{\small\itshape\color{dimtext} // ctrl-] in vim. M-. in emacs. the tags file makes it work. //}

Tags files are indexes of symbol definitions: functions, variables, classes, methods. Your editor uses them to jump to definitions with a single keystroke. zpwr generates them.


\vspace{12pt}
{\bfseries\color{neonyellow} CTAGS: THE CLASSIC}\vspace{4pt}

\begin{codebox}
zpwr regenctags                     \# regenerate ctags file\\
\end{codebox}

Runs exuberant-ctags or universal-ctags over the zpwr codebase. Produces a tags file that vim and emacs read.

In vim: Ctrl-] jumps to the definition under cursor. In vim: Ctrl-T jumps back to where you were. In vim: :tag zpwrDoctor jumps to zpwrDoctor.

Ctags understands shell functions, C, Python, Ruby, and dozens of other languages out of the box.


\vspace{12pt}
{\bfseries\color{neonyellow} GTAGS: THE GNU GLOBAL SYSTEM}\vspace{4pt}

\begin{codebox}
zpwr regengtags                     \# regenerate GNU Global tags\\
\end{codebox}

GNU Global is a more powerful tagging system. It supports: - Forward references (where is this defined?) - Backward references (where is this called from?) - Symbol completion.

Produces GTAGS, GRTAGS, GPATH files in the project root. Query from command line: global -x zpwrDoctor Query references: global -rx zpwrDoctor.


\vspace{12pt}
{\bfseries\color{neonyellow} GTAGS WITH PYGMENTS PARSER}\vspace{4pt}

\begin{codebox}
zpwr regengtagspygments             \# gtags with Pygments backend\\
\end{codebox}

Uses the Python Pygments parser instead of the built-in one. Pygments understands more languages and file types. Requires: pip install pygments.


\vspace{12pt}
{\bfseries\color{neonyellow} GTAGS BY FILE TYPE}\vspace{4pt}

\begin{codebox}
zpwr regengtagstype                 \# gtags filtered by type\\
\end{codebox}

Generate tags for a specific file type only. Faster than a full regen when you only care about one language.


\vspace{12pt}
{\bfseries\color{neonyellow} HOW TAGS WORK WITH YOUR \textbf{EDITOR}}\vspace{4pt}

vim    reads 'tags' file. :set tags=./tags;,tags Ctrl-] to jump, Ctrl-T to return :tselect for multiple matches.

emacs  reads TAGS file (etags format) or uses gtags M-. To find definition, M-* to return ggtags-mode or helm-gtags for GNU Global.

nvim   same as vim, plus LSP which often replaces tags but tags still work as a fast fallback.


\vspace{12pt}
{\bfseries\color{neonyellow} WHEN TO REGENERATE}\vspace{4pt}

- After adding or renaming functions - After a zpwr update that adds new verbs - When "tag not found" errors appear for known functions - Use zpwr watch '*.zsh' 'zpwr regenctags' for auto-regen


\vspace{16pt}

\section{Cache Regeneration}


\vspace{10pt}
{\large\bfseries\color{neoncyan} CACHE REGENERATION}\vspace{6pt}

{\small\itshape\color{dimtext} // the opposite of cleaning. building the caches back up. //}

Cleaning removes caches. Regeneration rebuilds them. These verbs construct the indexes and precomputed data that make zpwr fast.


\vspace{12pt}
{\bfseries\color{neonyellow} REGEN: THE STANDARD REBUILD}\vspace{4pt}

\begin{codebox}
zpwr regen                          \# regenerate common caches\\
\end{codebox}

Rebuilds the most frequently used caches: environment cache, config links, keybindings. This is the verb you reach for after editing config files.


\vspace{12pt}
{\bfseries\color{neonyellow} REGENALL: THE FULL REBUILD}\vspace{4pt}

\begin{codebox}
zpwr regenall                       \# regenerate everything\\
\end{codebox}

Every cache. Every index. Every precomputed file. This is the thorough version of regen. Takes longer but leaves nothing stale.


\vspace{12pt}
{\bfseries\color{neonyellow} INDIVIDUAL REGEN VERBS}\vspace{4pt}

Regenenvcache.

\begin{codebox}
zpwr regenenvcache                  \# rebuild environment cache\\
\end{codebox}

Snapshots all env vars to a cache file for fast lookup. Used by environmentcachesearch and other env verbs.

Regenkeybindings.

\begin{codebox}
zpwr regenkeybindings               \# rebuild keybinding index\\
\end{codebox}

Re-registers all zpwr keybindings. Run this after editing keybinding configuration or after a plugin clobbered your bindings.

Regenconfiglinks.

\begin{codebox}
zpwr regenconfiglinks               \# rebuild config symlinks\\
\end{codebox}

Re-creates the symlinks from \textbf{\textasciitilde{}/.zshrc}, \textbf{\textasciitilde{}/.vimrc}, etc. Into the zpwr directory tree. Run this after a fresh clone or when zpwr doctor reports broken config links.

Regenhistory.

\begin{codebox}
zpwr regenhistory                   \# clean up history file\\
\end{codebox}

Processes the \textbf{HISTFILE} to remove corrupted entries. Zsh history files can accumulate broken lines from concurrent writes or improper shutdowns.

Regenpowerline.

\begin{codebox}
zpwr regenpowerline                 \# rebuild powerline config\\
\end{codebox}

Regenerates powerline/p10k configuration. Run when the prompt looks wrong after a terminal font change.

Regengitrepocache.

\begin{codebox}
zpwr regengitrepocache              \# scan and cache git repos\\
\end{codebox}

Walks the filesystem finding git repos and builds a master list. Used by repo navigation and batch verbs.


\vspace{12pt}
{\bfseries\color{neonyellow} THE REGEN CHEAT SHEET}\vspace{4pt}

After editing config:     zpwr regen After zpwr update:        zpwr regenall After env var changes:    zpwr regenenvcache After repo changes:       zpwr regengitrepocache After completion issues:  zpwr regenzsh.


\vspace{16pt}

\section{Updates \& Dependencies}


\vspace{10pt}
{\large\bfseries\color{neoncyan} UPDATES \& DEPENDENCIES}\vspace{6pt}

{\small\itshape\color{dimtext} // staying current without breaking everything. a balancing act. //}

zpwr is a living system that depends on external tools, plugins, and its own git repository. These verbs keep everything synchronized.


\vspace{12pt}
{\bfseries\color{neonyellow} UPDATE: THE CORE UPDATE}\vspace{4pt}

\begin{codebox}
zpwr update                         \# update zpwr itself\\
\end{codebox}

Pulls the latest zpwr from the git remote. This is a git pull on \$\textbf{ZPWR}. Your local config in \$\textbf{ZPWR\_LOCAL} is not touched.


\vspace{12pt}
{\bfseries\color{neonyellow} UPDATEALL: THE EVERYTHING UPDATE}\vspace{4pt}

\begin{codebox}
zpwr updateall                      \# update zpwr + all plugins\\
\end{codebox}

Updates zpwr, zinit plugins, and other managed components. This is the weekly maintenance verb. Run it when you want the latest of everything.


\vspace{12pt}
{\bfseries\color{neonyellow} UPDATEDEPS: DEPENDENCY UPDATES}\vspace{4pt}

\begin{codebox}
zpwr updatedeps                     \# update external dependencies\\
\end{codebox}

Updates the external tools zpwr depends on: zinit plugins, zsh-autosuggestions, zsh-syntax-highlighting, etc. Does not touch zpwr's own code.


\vspace{12pt}
{\bfseries\color{neonyellow} UPDATEDEPSCLEAN: CLEAN + UPDATE DEPS}\vspace{4pt}

\begin{codebox}
zpwr updatedepsclean                \# fresh dependency install\\
\end{codebox}

Cleans stale dependency state first, then updates. Use when updatedeps alone does not fix plugin issues.


\vspace{12pt}
{\bfseries\color{neonyellow} UPDATEZINIT: ZINIT ONLY}\vspace{4pt}

\begin{codebox}
zpwr updatezinit                    \# update zinit plugin manager\\
\end{codebox}

Updates just the zinit plugin manager binary, not the plugins it manages. Use when zinit itself has a bug or you need a newer zinit feature.


\vspace{12pt}
{\bfseries\color{neonyellow} UPDATEPULL: GIT PULL ALL REPOS}\vspace{4pt}

\begin{codebox}
zpwr updatepull                     \# pull all tracked git repos\\
\end{codebox}

Iterates over all cached git repos and runs git pull. Useful when you have many repos that track remotes and you want them all current.


\vspace{12pt}
{\bfseries\color{neonyellow} THE UPDATE HIERARCHY}\vspace{4pt}

Update              just zpwr (git pull) updatedeps          just plugins and tools updatezinit         just the plugin manager updateall           zpwr + plugins + tools updatepull          all tracked git repos updatedepsclean     clean slate + update deps.


\vspace{12pt}
{\bfseries\color{neonyellow} RECOMMENDED UPDATE SCHEDULE}\vspace{4pt}

Weekly:      zpwr updateall After update: zpwr recompile If broken:   zpwr updatedepsclean.

Always recompile after updating. New source files need new .zwc files. The old compiled versions will load stale code until you recompile.


\vspace{16pt}


\chapter{MONITORING \& AUTOMATION}


\begin{center}

\includegraphics[width=0.85\textwidth]{/Users/wizard/.zpwr/docs/screenshots/zpwr_watch_help.png}
\end{center}

\section{File Watching}


\begin{center}

\includegraphics[width=0.8\textwidth]{/Users/wizard/.zpwr/docs/screenshots/zpwr_watch_help.png}
\end{center}


\vspace{10pt}
{\large\bfseries\color{neoncyan} FILE WATCHING}\vspace{6pt}

{\small\itshape\color{dimtext} // your files change. your command runs. automatically. //}

The watch verb monitors files for changes and auto-executes a command when something is modified. It is the poor man's CI pipeline, running locally in your terminal.


\vspace{12pt}
{\bfseries\color{neonyellow} USAGE}\vspace{4pt}

\begin{codebox}
zpwr watch '<GLOB>' '<COMMAND>'\\
\end{codebox}

Two arguments: what to watch, and what to run. Quote both arguments to prevent premature expansion.


\vspace{12pt}
{\bfseries\color{neonyellow} PRACTICAL EXAMPLES}\vspace{4pt}

\begin{codebox}
zpwr watch '*.zsh' 'zpwr recompile'\\
\end{codebox}

Recompile zpwr whenever a .zsh file changes. Edit a function, save, and .zwc files auto-rebuild.

\begin{codebox}
zpwr watch 'src/**/*.ts' 'npm run build'\\
\end{codebox}

Build your TypeScript project on every source change.

\begin{codebox}
zpwr watch 'Makefile' 'make'\\
\end{codebox}

Run make whenever the Makefile changes.

\begin{codebox}
zpwr watch '.' 'echo changed'\\
\end{codebox}

Watch the entire current directory.

\begin{codebox}
zpwr watch '*.py' 'python -m pytest'\\
\end{codebox}

Run tests on every Python file change. TDD mode.


\vspace{12pt}
{\bfseries\color{neonyellow} PLATFORM BACKEND}\vspace{4pt}

MacOS:   Uses fswatch (brew install fswatch) Linux:   Uses inotifywait (apt install inotify-tools)

The verb auto-detects which tool is available. If neither is found, it tells you what to install.


\vspace{12pt}
{\bfseries\color{neonyellow} THE OUTPUT FORMAT}\vspace{4pt}

Each trigger shows a box with: - Trigger count (how many times it has fired) - Timestamp      (when the change was detected) - Filename       (which file changed) - Exit status    (green ok or red exit code) - Elapsed time   (how long the command took)


\vspace{12pt}
{\bfseries\color{neonyellow} HOW TO STOP}\vspace{4pt}

Press Ctrl-C. The watcher process terminates cleanly.


\vspace{12pt}
{\bfseries\color{neonyellow} GLOB HANDLING}\vspace{4pt}

The glob is split into directory and file filter: 'src/**/*.ts' -> watch src/ recursively, filter *.ts '.'           -> watch current dir, no filter '*.zsh'       -> watch ./, filter *.zsh.

The glob-to-regex conversion handles * and ? Patterns. For fswatch, it uses -e '.*' -i to exclude-then-include. For inotifywait, it uses zsh's glob matching on filenames.


\vspace{16pt}

\section{Command Replay}


\begin{center}
\includegraphics[width=0.8\textwidth]{/Users/wizard/.zpwr/docs/screenshots/zpwr_replay_help.png}
\end{center}


\vspace{10pt}
{\large\bfseries\color{neoncyan} COMMAND REPLAY}\vspace{6pt}

{\small\itshape\color{dimtext} // your history is a script. replay it on demand. //}

The replay verb plays back the last N commands from your shell history. Three modes: look, execute, or tmux pane. It is a VCR for your terminal.


\vspace{12pt}
{\bfseries\color{neonyellow} BASIC USAGE}\vspace{4pt}

\begin{codebox}
zpwr replay                         \# preview last 10 commands\\
zpwr replay 20                      \# preview last 20\\
zpwr replay 5 exec                  \# execute last 5 with timing\\
zpwr replay 10 tmux                 \# replay in a tmux split pane\\
\end{codebox}


\vspace{12pt}
{\bfseries\color{neonyellow} THE THREE MODES}\vspace{4pt}

PREVIEW MODE (default, safe)

\begin{codebox}
zpwr replay\\
zpwr replay 20 preview\\
\end{codebox}

Lists the commands with line numbers. Nothing executes. Use this to review what you did before replaying it. The output shows helpful hints for switching to exec/tmux.

EXEC MODE (runs in current shell)

\begin{codebox}
zpwr replay 5 exec\\
\end{codebox}

Executes each command sequentially in your current shell. Each command shows: - Command number [1/5] - The command being run - Exit status (green ok or red exit code) - Elapsed time in milliseconds.

WARNING: This runs real commands. If your last 5 commands included rm -rf something, it will do it again. Always preview first.

TMUX MODE (runs in a new pane)

\begin{codebox}
zpwr replay 10 tmux\\
\end{codebox}

Opens a new tmux split pane and replays commands there. Your current pane is untouched. You can watch the replay happen side-by-side with your current work. The pane waits for Enter when done, then cleans up. Requires an active tmux session.


\vspace{12pt}
{\bfseries\color{neonyellow} USE CASES}\vspace{4pt}

Reproduce a bug:       zpwr replay 3 exec Demo your workflow:     zpwr replay 20 tmux Audit what you did:     zpwr replay 50 Re-run a build chain:   zpwr replay 4 exec.


\vspace{12pt}
{\bfseries\color{neonyellow} HOW IT WORKS}\vspace{4pt}

Uses fc -rl -N to extract the last N history entries. Reverses them to chronological order. In exec mode, each command is eval'd with timing via EPOCHREALTIME. In tmux mode, a temp script is built and executed in a split pane with a 0.5s delay between commands for readability. The script self-destructs after completion.


\vspace{16pt}

\section{Log Surveillance}


\vspace{10pt}
{\large\bfseries\color{neoncyan} LOG SURVEILLANCE}\vspace{6pt}

{\small\itshape\color{dimtext} // watching logs in glorious technicolor. the ccze way. //}

zpwr logs everything: startup, errors, plugin loads, debug messages. The taillog verb tails these logs with optional ccze colorization for maximum readability.


\vspace{12pt}
{\bfseries\color{neonyellow} BASIC USAGE}\vspace{4pt}

\begin{codebox}
zpwr taillog                        \# follow log with colors\\
zpwr taillog -n 50                  \# last 50 lines, then follow\\
zpwr taillog -f                     \# colorize without following\\
zpwr taillog -r                     \# raw output with line numbers\\
\end{codebox}


\vspace{12pt}
{\bfseries\color{neonyellow} THE FLAGS}\vspace{4pt}

-n, --lines N     Show last N lines before following -f, --no-follow   Just colorize, do not tail -f -r, --raw         No ccze, plain output with line numbers


\vspace{12pt}
{\bfseries\color{neonyellow} CUSTOM LOGFILE}\vspace{4pt}

\begin{codebox}
zpwr taillog /var/log/system.log\\
zpwr taillog /var/log/syslog\\
\end{codebox}

The last argument can be any logfile path. Default: \$\textbf{ZPWR\_LOGFILE} (the zpwr logfile).


\vspace{12pt}
{\bfseries\color{neonyellow} CCZE: THE SECRET WEAPON}\vspace{4pt}

ccze is a log colorizer that understands dozens of log formats: syslog, Apache, Postfix, PHP, kernel messages, and more. It auto-detects the format and colorizes:

Timestamps     in one color warnings       in yellow errors         in red info           in cyan.

\begin{codebox}
hostnames      in magenta\\
\end{codebox}

Install: brew install ccze  or  apt install ccze If ccze is not found, taillog falls back to plain output with a helpful install hint.


\vspace{12pt}
{\bfseries\color{neonyellow} THE \textbf{ZPWR} LOGGING VERBS}\vspace{4pt}

Taillog consumes logs. These verbs produce them:

\begin{codebox}
zpwr logdebug <msg>                 \# debug level\\
zpwr loginfo <msg>                  \# info level\\
zpwr logerror <msg>                 \# error level\\
zpwr logtrace <msg>                 \# trace level (most verbose)\\
\end{codebox}

The *console variants also print to stderr:

\begin{codebox}
zpwr logdebugconsole <msg>\\
zpwr loginfoconsole <msg>\\
zpwr logerrorconsole <msg>\\
zpwr logtraceconsole <msg>\\
\end{codebox}


\vspace{12pt}
{\bfseries\color{neonyellow} LOG MAINTENANCE}\vspace{4pt}

Logs grow forever. Clean them periodically:

\begin{codebox}
zpwr cleanlog                       \# truncate the logfile\\
\end{codebox}

Or check the size first: ls -lh \$\textbf{ZPWR\_LOGFILE}

A 100MB logfile means your debug logging is too chatty. A 1KB logfile means nothing is being logged. Both are worth investigating.


\vspace{16pt}

\section{Backup \& Restore}


\vspace{10pt}
{\large\bfseries\color{neoncyan} BACKUP \& RESTORE}\vspace{6pt}

{\small\itshape\color{dimtext} // your history is priceless. these verbs are the insurance policy. //}

Your shell history is years of accumulated knowledge: every command you ran, every fix you found at 2am, every pipeline you crafted. Losing it is losing institutional memory. These verbs prevent that.


\vspace{12pt}
{\bfseries\color{neonyellow} BACKUP: THE GENERAL BACKUP}\vspace{4pt}

\begin{codebox}
zpwr backup                         \# backup zpwr configuration\\
\end{codebox}

Creates a timestamped backup of zpwr configuration files. This captures your local settings, tokens, and any customizations in \$\textbf{ZPWR\_LOCAL}.

Backup destination: \textbf{\$ZPWR\_LOCAL/rcBackups/} Each backup is timestamped so you can have many.


\vspace{12pt}
{\bfseries\color{neonyellow} BACKUPHISTORY: THE HISTORY VAULT}\vspace{4pt}

\begin{codebox}
zpwr backuphistory                  \# backup HISTFILE\\
\end{codebox}

Copies your \textbf{HISTFILE} to the backup directory with a timestamp. This is the most important backup in zpwr.

Your \textbf{HISTFILE} can be corrupted by: - Power loss during write - Concurrent sessions fighting over the file - Disk full causing truncated writes - Bad .zsh\_history format from import/merge.

A weekly zpwr backuphistory gives you a safety net.


\vspace{12pt}
{\bfseries\color{neonyellow} RESTOREHISTORY: THE TIME MACHINE}\vspace{4pt}

\begin{codebox}
zpwr restorehistory                 \# restore from backup\\
\end{codebox}

Restores \textbf{HISTFILE} from the most recent backup. Use when your history is corrupted or accidentally truncated.

The restore process: 1. Finds the latest backup in \textbf{\$ZPWR\_LOCAL/rcBackups} 2. Backs up the current (possibly broken) \textbf{HISTFILE} 3. Copies the backup over \textbf{HISTFILE} 4. Reloads history with fc -R.


\vspace{12pt}
{\bfseries\color{neonyellow} SNAPSHOTS VS BACKUPS}\vspace{4pt}

Snapshot          captures full env state (aliases, vars, tmux) backup            copies config files to safety backuphistory     specifically targets \textbf{HISTFILE}

Use snapshots for environment state. Use backups for file safety. They complement each other.


\vspace{12pt}
{\bfseries\color{neonyellow} THE SAFETY PROTOCOL}\vspace{4pt}

Before a big change:

\begin{codebox}
zpwr snapshot before-change\\
zpwr backuphistory\\
zpwr backup\\
\end{codebox}

Now you have three layers of protection: env snapshot, history backup, and config backup. If the change goes sideways, you can recover everything.


\vspace{12pt}
{\bfseries\color{neonyellow} CRON IT}\vspace{4pt}

Add to your crontab (crontab -e):

\begin{codebox}
0 3 * * 0 zsh -c 'source \textasciitilde{}/.zshrc; zpwr backuphistory'\\
\end{codebox}

Every Sunday at 3am. Automatic insurance.


\vspace{16pt}


\chapter{UTILITY ARSENAL}


\begin{center}
\includegraphics[width=0.85\textwidth]{/Users/wizard/.zpwr/docs/screenshots/art_mainframe.png}
\end{center}

\section{File \& Process Operations}


\vspace{10pt}
{\large\bfseries\color{neoncyan} FILE \& PROCESS OPERATIONS}\vspace{6pt}

{\small\itshape\color{dimtext} // the everyday tools. sharpened and color-coded. //}

zpwr wraps common file and process operations with enhanced output, color, and fzf integration. These are the verbs you reach for twenty times a day.


\vspace{12pt}
{\bfseries\color{neonyellow} FILE LISTING}\vspace{4pt}

\begin{codebox}
zpwr clearlist                      \# clear screen + list files\\
zpwr clearls                        \# alias for clearlist\\
zpwr info                           \# alias for clearlist\\
zpwr list                           \# list files (no clear)\\
zpwr ls                             \# alias for list\\
\end{codebox}

Clearlist is the most-used verb after help. It clears the terminal and shows a fresh directory listing. Uses eza if available, falls back to ls.


\vspace{12pt}
{\bfseries\color{neonyellow} FILE VIEWING}\vspace{4pt}

\begin{codebox}
zpwr cat <file>                     \# colorized cat\\
zpwr c <file>                       \# shortest form\\
zpwr catcd <dir>                    \# cat files in a directory\\
\end{codebox}

Cat/c uses bat or pygmentize for syntax highlighting when available. Plain cat as fallback. Catcd lists the contents of each file in a directory.


\vspace{12pt}
{\bfseries\color{neonyellow} PROCESS MANAGEMENT}\vspace{4pt}

\begin{codebox}
zpwr ps                             \# enhanced process listing\\
zpwr kill                           \# fzf-powered process kill\\
zpwr killedit                       \# kill then edit related file\\
\end{codebox}

zpwr kill is the crown jewel here. It pipes process listings through fzf, letting you fuzzy-search for the process you want to terminate. Select it, confirm, dead. No more guessing PIDs.


\vspace{12pt}
{\bfseries\color{neonyellow} FILE DESCRIPTORS}\vspace{4pt}

\begin{codebox}
zpwr lsof                           \# list open files (enhanced)\\
zpwr lsofedit                       \# lsof then edit result\\
\end{codebox}

zpwr lsof wraps lsof with fzf filtering. Useful when tracking down which process holds a file or which files a process has open.


\vspace{12pt}
{\bfseries\color{neonyellow} FILE OPENING}\vspace{4pt}

\begin{codebox}
zpwr open                           \# open file/dir with system\\
\end{codebox}

Opens the target with the system handler: Finder on macOS, xdg-open on Linux. Works with files, directories, URLs.


\vspace{12pt}
{\bfseries\color{neonyellow} THE COMBO MOVES}\vspace{4pt}

Clear terminal, list files, check what is running:

\begin{codebox}
zpwr clearlist \&\& zpwr ps\\
\end{codebox}

Find and kill a zombie process:

\begin{codebox}
zpwr kill     \# fzf search -> select -> terminate\\
\end{codebox}

Find what is holding a port:

\begin{codebox}
zpwr lsof     \# fzf through open files and ports\\
\end{codebox}


\vspace{16pt}

\section{Batch Execution}


\vspace{10pt}
{\large\bfseries\color{neoncyan} BATCH EXECUTION}\vspace{6pt}

{\small\itshape\color{dimtext} // one command, many targets. the industrial strength loop. //}

When you need to run the same operation across multiple files, directories, or git repos, these verbs replace hand-rolled for loops with tested, colorized execution.


\vspace{12pt}
{\bfseries\color{neonyellow} FOR: THE BASIC ITERATOR}\vspace{4pt}

\begin{codebox}
zpwr for <command>                  \# run command in current dir context\\
\end{codebox}

Executes a command with zpwr's enhanced output formatting. Shows timing, exit status, and colorized output.


\vspace{12pt}
{\bfseries\color{neonyellow} FOR10: THE STRESS TEST}\vspace{4pt}

\begin{codebox}
zpwr for10 <command>                \# run command 10 times\\
\end{codebox}

Executes a command 10 times. Useful for: - Benchmarking (is the timing consistent?) - Flaky test detection (does it fail intermittently?) - Load testing (can it handle rapid repeated execution?)


\vspace{12pt}
{\bfseries\color{neonyellow} FORDIR: ITERATE OVER GIT REPOS}\vspace{4pt}

\begin{codebox}
zpwr fordir <command>               \# run command in each git repo\\
\end{codebox}

Iterates over all cached git repos, cd's into each one, and runs the command. Extremely powerful for bulk operations:

\begin{codebox}
zpwr fordir 'git status -s'         \# dirty files in all repos\\
zpwr fordir 'git pull'              \# pull all repos\\
zpwr fordir 'git stash list'        \# find forgotten stashes\\
\end{codebox}

Each repo's output is labeled with its path so you know which repo produced which output.


\vspace{12pt}
{\bfseries\color{neonyellow} FORDIRUPDATE: PULL ALL REPOS}\vspace{4pt}

\begin{codebox}
zpwr fordirupdate                   \# git pull in every repo\\
\end{codebox}

Specialized version of fordir that runs git pull in every cached repo. Same as zpwr fordir 'git pull' but optimized for the common case.


\vspace{12pt}
{\bfseries\color{neonyellow} EXECGLOB: GLOB-BASED EXECUTION}\vspace{4pt}

\begin{codebox}
zpwr execglob <glob> <command>      \# run command on glob matches\\
\end{codebox}

Finds files matching a glob and runs a command on each:

\begin{codebox}
zpwr execglob '*.log' 'wc -l'      \# line count each log file\\
zpwr execglob '*.bak' 'rm'         \# delete all .bak files\\
\end{codebox}


\vspace{12pt}
{\bfseries\color{neonyellow} EXECGLOBPARALLEL: PARALLEL GLOB EXECUTION}\vspace{4pt}

\begin{codebox}
zpwr execglobparallel <glob> <cmd>  \# parallel execution\\
\end{codebox}

Same as execglob but runs commands in parallel. Useful when processing many files that are I/O-bound:

\begin{codebox}
zpwr execglobparallel '*.png' 'optipng'  \# compress images\\
\end{codebox}


\vspace{12pt}
{\bfseries\color{neonyellow} THE POWER COMBOS}\vspace{4pt}

Find all repos with uncommitted work:

\begin{codebox}
zpwr fordir 'git diff --stat HEAD'\\
\end{codebox}

Count lines in all Python files:

\begin{codebox}
zpwr execglob '**/*.py' 'wc -l'\\
\end{codebox}

Stress test a flaky command:

\begin{codebox}
zpwr for10 'curl -s localhost:3000/health'\\
\end{codebox}


\vspace{16pt}

\section{Text Operations}


\vspace{10pt}
{\large\bfseries\color{neoncyan} TEXT OPERATIONS}\vspace{6pt}

{\small\itshape\color{dimtext} // find and replace. rename. count. transform. //}

Text manipulation is the bread and butter of shell work. These verbs handle the common cases: renaming files, replacing strings, counting lines, and transforming text.


\vspace{12pt}
{\bfseries\color{neonyellow} RENAME: FILE RENAMING}\vspace{4pt}

\begin{codebox}
zpwr rename                         \# rename files with patterns\\
\end{codebox}

Renames files using pattern matching. Safer than raw mv because it shows what will happen before doing it.


\vspace{12pt}
{\bfseries\color{neonyellow} REPLACER: FIND AND REPLACE IN FILES}\vspace{4pt}

\begin{codebox}
zpwr replacer                       \# search and replace in files\\
\end{codebox}

Performs find-and-replace across multiple files. The codebase refactoring tool. When you renamed a function and need to update every call site.


\vspace{12pt}
{\bfseries\color{neonyellow} CHANGENAME: RENAME WITH HISTORY}\vspace{4pt}

\begin{codebox}
zpwr changename                     \# rename with tracking\\
\end{codebox}

Renames files or directories with additional tracking. More sophisticated than basic rename for complex renaming operations.


\vspace{12pt}
{\bfseries\color{neonyellow} CHANGENAMEZPWR: \textbf{ZPWR}-SPECIFIC RENAME}\vspace{4pt}

\begin{codebox}
zpwr changenamezpwr                 \# rename within zpwr codebase\\
\end{codebox}

Specialized rename for zpwr function and file names. Handles the autoload naming convention (file name must match function name) and updates references.


\vspace{12pt}
{\bfseries\color{neonyellow} LINECOUNT: COUNT \textbf{LINES}}\vspace{4pt}

\begin{codebox}
zpwr linecount                      \# count lines in files\\
\end{codebox}

Enhanced line counting with formatting and totals. Shows per-file counts and a grand total. More informative than raw wc -l.


\vspace{12pt}
{\bfseries\color{neonyellow} PRE \& POST: TEXT INJECTION}\vspace{4pt}

\begin{codebox}
zpwr pre                            \# prepend text to input\\
zpwr post                           \# append text to input\\
\end{codebox}

Pipe filters that add text before or after each line. Useful in pipelines for formatting:

\begin{codebox}
ls | zpwr pre '  - '               \# bullet-point file list\\
cat urls.txt | zpwr post '/api/v1'  \# append to each URL\\
\end{codebox}


\vspace{12pt}
{\bfseries\color{neonyellow} URLSAFE: URL ENCODING}\vspace{4pt}

\begin{codebox}
zpwr urlsafe                        \# URL-encode a string\\
\end{codebox}

Converts special characters to percent-encoded form. Spaces become \%20, ampersands become \%26, etc. Essential when constructing URLs from user input.


\vspace{12pt}
{\bfseries\color{neonyellow} THE TEXT PIPELINE}\vspace{4pt}

zpwr text verbs compose naturally in pipelines:

\begin{codebox}
cat data.txt | zpwr pre 'ROW: ' | zpwr post ' ;'\\
\end{codebox}

Result: ROW: line1 ; ROW: line2 ;

Rename all .txt to .md:

\begin{codebox}
zpwr rename                         \# interactive rename\\
\end{codebox}

Replace 'oldFunc' with 'newFunc' across the codebase:

\begin{codebox}
zpwr replacer                       \# interactive replace\\
\end{codebox}


\vspace{16pt}

\section{Display \& Banners}


\vspace{10pt}
{\large\bfseries\color{neoncyan} DISPLAY \& BANNERS}\vspace{6pt}

{\small\itshape\color{dimtext} // because aesthetics matter. even in a terminal. //}

zpwr ships with a fleet of display verbs that produce ASCII art, colorized output, and terminal animations. Some are informational. Some are just cool to look at.


\vspace{12pt}
{\bfseries\color{neonyellow} THE BANNER FAMILY}\vspace{4pt}

\begin{codebox}
zpwr about                          \# the welcome banner\\
zpwr banner                         \# same as about\\
zpwr bannercounts                   \# banner + environment stats\\
\end{codebox}

The core banner shows the zpwr ASCII logo with version, author, and system info. Bannercounts adds the environment census (alias count, function count, etc.).

\begin{codebox}
zpwr bannerlolcat                   \# banner piped through lolcat\\
\end{codebox}

Rainbow gradient banner. Requires: gem install lolcat For when monochrome is not enough and you need to assert terminal dominance with a rainbow.

\begin{codebox}
zpwr bannerpony                     \# banner with ponysay\\
\end{codebox}

Displays the banner inside an ASCII pony speech bubble. Requires: pip install ponysay Peak terminal culture.

\begin{codebox}
zpwr bannernopony                   \# banner without the pony\\
\end{codebox}

For the pony-averse. All the info, none of the horse.


\vspace{12pt}
{\bfseries\color{neonyellow} PRETTY PRINTING}\vspace{4pt}

\begin{codebox}
zpwr prettyprint                    \# formatted text output\\
zpwr altprettyprint                 \# alternate format\\
zpwr printmap                       \# print associative array\\
\end{codebox}

Prettyprint is used internally by many zpwr verbs to format output with consistent color and spacing. Printmap formats associative arrays as aligned key-value tables. Useful for debugging hash tables.


\vspace{12pt}
{\bfseries\color{neonyellow} COLOR TESTING}\vspace{4pt}

\begin{codebox}
zpwr colortest                      \# test basic terminal colors\\
zpwr colortest256                   \# full 256-color palette\\
\end{codebox}

Run these after changing terminals, fonts, or color schemes. Colortest shows the 16 ANSI colors. Colortest256 shows the full 256-color grid. If colors look wrong, your \textbf{TERM} or terminal emulator config needs adjustment.


\vspace{12pt}
{\bfseries\color{neonyellow} ANIMATION}\vspace{4pt}

\begin{codebox}
zpwr animate                        \# terminal animation\\
\end{codebox}

Plays an ASCII animation in the terminal. Because sometimes you need to entertain yourself while waiting for a build.


\vspace{12pt}
{\bfseries\color{neonyellow} THE \textbf{DISPLAY} ARSENAL AT A GLANCE}\vspace{4pt}

Impress a colleague:      zpwr bannerlolcat Debug your terminal:       zpwr colortest256 Count your weapons:        zpwr bannercounts Maximum absurdity:         zpwr bannerpony Format a hash table:       zpwr printmap Kill time aesthetically:   zpwr animate.


\vspace{16pt}


\chapter{EMACS DIMENSION}


\begin{center}
\includegraphics[width=0.85\textwidth]{/Users/wizard/.zpwr/docs/screenshots/art_neural_link.png}
\end{center}

\section{Emacs Integration Overview}


\vspace{10pt}
{\large\bfseries\color{neoncyan} ENTERING THE EMACS DIMENSION}\vspace{6pt}

Where M-x is a way of life and pinky fingers go to die.

You thought there were only vim verbs? Think again, operative. \textbf{ZPWR} mirrors every vim operation with an emacs equivalent. Same power, different religion.


\vspace{12pt}
{\bfseries\color{neonyellow} THE 28 EMACS VERBS:}\vspace{4pt}

FILE OPS      emacsall, emacsalledit, emacsautoload RECENT        emacsrecent, emacsrecentcd, emacsrecentsudo SEARCH        emacsfilesearch, emacsfilesearchedit emacsfindsearch, emacsfindsearchedit emacslocatesearch, emacslocatesearchedit emacswordsearch, emacswordsearchedit TAGS          emacsctags, emacsgtags, emacsgtagsedit CONFIG        emacsconfig, emacsemacsconfig, emacstokens SERVER        emacsallserver SCRIPTS       emacsscripts, emacsscriptedit META          emacsplugincount, emacspluginlist, emacszpwr.

THE MIRROR PRINCIPLE:

Every vim verb has an emacs twin. The architecture is symmetrical by design. You can swap editors without losing any \textbf{ZPWR} functionality.

vim verb              emacs twin zpwr vimall     -->   zpwr emacsall zpwr vimrecent  -->   zpwr emacsrecent zpwr vimscripts -->   zpwr emacsscripts.


\vspace{12pt}
{\bfseries\color{neonyellow} PHILOSOPHY:}\vspace{4pt}

\textbf{ZPWR} does not care about your editor war allegiance. It arms both sides equally. The real enemy is notepad.

{\small\itshape\color{dimtext} // the emacs dimension runs parallel to vim-space}

{\small\itshape\color{dimtext} // choose your weapon, the framework adapts}

{\small\itshape\color{dimtext} // \$\textbf{ZPWR\_EDITOR} controls which universe you inhabit}


\vspace{16pt}

\section{Emacs File Operations}


\vspace{10pt}
{\large\bfseries\color{neoncyan} EMACS FILE OPS: OPENING THINGS SINCE 1976}\vspace{6pt}

emacs: an operating system lacking only a decent editor.


\vspace{12pt}
{\bfseries\color{neonyellow} zpwr emacsall}\vspace{4pt}

Opens all \textbf{ZPWR} framework files in emacs via fzf. Every autoload, script, config -- the whole kingdom.

\begin{codebox}
zpwr emacsall              \# browse all files\\
zpwr emacsalledit          \# same but editable\\
\end{codebox}


\vspace{12pt}
{\bfseries\color{neonyellow} zpwr emacsrecent}\vspace{4pt}

Opens recently edited files. Emacs remembers where you have been, like a browser history for code.

\begin{codebox}
zpwr emacsrecent            \# fzf recent files\\
zpwr emacsrecentcd          \# cd to recent dir\\
zpwr emacsrecentsudo        \# recent with sudo\\
zpwr emacsrecentsudocd      \# sudo + cd combo\\
\end{codebox}


\vspace{12pt}
{\bfseries\color{neonyellow} zpwr emacsscripts}\vspace{4pt}

Browse and open \textbf{ZPWR} scripts directory. Where the shell scripts live and breathe.

\begin{codebox}
zpwr emacsscripts           \# browse scripts\\
zpwr emacsscriptedit        \# edit a script\\
\end{codebox}


\vspace{12pt}
{\bfseries\color{neonyellow} zpwr emacsautoload}\vspace{4pt}

Open autoloaded zsh functions in emacs. The beating heart of \textbf{ZPWR}'s function library.

\begin{codebox}
zpwr emacsautoload          \# browse autoloads\\
\end{codebox}

{\small\itshape\color{dimtext} // protip: emacsrecent is your daily driver}

{\small\itshape\color{dimtext} // it learns your habits and surfaces what matters}

{\small\itshape\color{dimtext} // sudo variants for when /etc needs surgery}


\vspace{16pt}

\section{Emacs Search \& Tags}


\vspace{10pt}
{\large\bfseries\color{neoncyan} SEARCH \& TAGS: FIND ANYTHING, ANYWHERE}\vspace{6pt}

grep is for amateurs, we have fzf-powered search beams.


\vspace{12pt}
{\bfseries\color{neonyellow} FILE SEARCH (by filename):}\vspace{4pt}

\begin{codebox}
zpwr emacsfilesearch        \# fzf filename search\\
zpwr emacsfilesearchedit    \# search + edit\\
\end{codebox}

Searches file NAMES across the \textbf{ZPWR} tree. Fast, fuzzy, and feeds results into emacs.


\vspace{12pt}
{\bfseries\color{neonyellow} WORD SEARCH (by content):}\vspace{4pt}

\begin{codebox}
zpwr emacswordsearch        \# search file content\\
zpwr emacswordsearchedit    \# content search + edit\\
\end{codebox}

Searches INSIDE files for patterns. Like ag/rg but with fzf preview and direct-to-emacs piping.


\vspace{12pt}
{\bfseries\color{neonyellow} FIND \& LOCATE SEARCH:}\vspace{4pt}

\begin{codebox}
zpwr emacsfindsearch        \# find-based search\\
zpwr emacsfindsearchedit    \# find + edit\\
zpwr emacslocatesearch      \# locate db search\\
zpwr emacslocatesearchedit  \# locate + edit\\
\end{codebox}


\vspace{12pt}
{\bfseries\color{neonyellow} TAGS -- CODE NAVIGATION:}\vspace{4pt}

\begin{codebox}
zpwr emacsctags             \# ctags generation\\
zpwr emacsgtags             \# GNU Global tags\\
zpwr emacsgtagsedit         \# gtags + jump to def\\
\end{codebox}

Tags let emacs jump to function definitions instantly. Ctags for basic tagging, gtags for cross-reference.

{\small\itshape\color{dimtext} // the search hierarchy:}

{\small\itshape\color{dimtext} //   wordsearch  = grep inside files (content)}

{\small\itshape\color{dimtext} //   filesearch  = find by filename (path)}

{\small\itshape\color{dimtext} //   locatesearch = system locate db (everywhere)}

{\small\itshape\color{dimtext} //   findsearch  = find command (directory tree)}

{\small\itshape\color{dimtext} //   tags        = jump to definitions (symbols)}


\vspace{16pt}

\section{Emacs Config \& Server}


\vspace{10pt}
{\large\bfseries\color{neoncyan} CONFIG \& SERVER: THE EMACS MOTHERSHIP}\vspace{6pt}

Because starting emacs fresh every time is a war crime.


\vspace{12pt}
{\bfseries\color{neonyellow} zpwr emacsconfig}\vspace{4pt}

Edit \textbf{ZPWR}'s emacs configuration files. The bridge between the shell framework and emacs-land.

\begin{codebox}
zpwr emacsconfig            \# edit zpwr emacs cfg\\
\end{codebox}


\vspace{12pt}
{\bfseries\color{neonyellow} zpwr emacsemacsconfig}\vspace{4pt}

Edit your personal emacs init files (init.el, etc). Deeper than emacsconfig -- this is YOUR emacs setup.

\begin{codebox}
zpwr emacsemacsconfig       \# edit \textasciitilde{}/.emacs.d\\
\end{codebox}


\vspace{12pt}
{\bfseries\color{neonyellow} zpwr emacsallserver}\vspace{4pt}

Opens files using emacsclient, connecting to a running emacs daemon. Instant file opens, no startup penalty.

\begin{codebox}
zpwr emacsallserver         \# open via daemon\\
\end{codebox}

HOW THE SERVER WORKFLOW WORKS: 1. Start emacs daemon:  emacs --daemon 2. Use emacsallserver to open files via emacsclient 3. Files open instantly in the running emacs instance 4. No more 3-second startup tax per file.


\vspace{12pt}
{\bfseries\color{neonyellow} zpwr emacstokens}\vspace{4pt}

Edit your tokens/secrets file in emacs. API keys, env vars, the classified stuff.

\begin{codebox}
zpwr emacstokens            \# edit ZPWR\_LOCAL/.tokens.sh\\
\end{codebox}


\vspace{12pt}
{\bfseries\color{neonyellow} zpwr emacszpwr}\vspace{4pt}

Open the entire \textbf{ZPWR} directory tree in emacs.

\begin{codebox}
zpwr emacszpwr              \# the whole framework\\
\end{codebox}

{\small\itshape\color{dimtext} // emacsclient + daemon = the cheat code}

{\small\itshape\color{dimtext} // startup once, open files forever}

{\small\itshape\color{dimtext} // this is how emacs users actually survive}


\vspace{16pt}


\chapter{TMUX WARFARE}


\begin{center}
\includegraphics[width=0.85\textwidth]{/Users/wizard/.zpwr/docs/screenshots/art_hack_the_planet.png}
\end{center}

\section{Tmux Session Management}


\vspace{10pt}
{\large\bfseries\color{neoncyan} TMUX WARFARE: TERMINAL MULTIPLEXER COMBAT}\vspace{6pt}

One terminal to rule them all, split them, and in panes bind them.


\vspace{12pt}
{\bfseries\color{neonyellow} zpwr attach}\vspace{4pt}

Attach to an existing tmux session. If no session exists, one will be summoned.

\begin{codebox}
zpwr attach                 \# rejoin the matrix\\
\end{codebox}


\vspace{12pt}
{\bfseries\color{neonyellow} zpwr detach}\vspace{4pt}

Detach from current tmux session gracefully. Your processes keep running in the background.

\begin{codebox}
zpwr detach                 \# leave but don't kill\\
\end{codebox}


\vspace{12pt}
{\bfseries\color{neonyellow} zpwr killmux}\vspace{4pt}

Kill the tmux server. Nuclear option. Everything dies. WARNING: all sessions, windows, and panes -- gone.

\begin{codebox}
zpwr killmux                \# scorched earth\\
\end{codebox}


\vspace{12pt}
{\bfseries\color{neonyellow} zpwr start}\vspace{4pt}

Start a new tmux session with \textbf{ZPWR}'s custom layout. Preconfigured panes, windows, and keybindings.

\begin{codebox}
zpwr start                  \# bootstrap tmux\\
\end{codebox}


\vspace{12pt}
{\bfseries\color{neonyellow} zpwr starttabs}\vspace{4pt}

Start tmux with multiple pre-configured tabs. Each tab gets its own purpose -- dev, logs, git, etc.

\begin{codebox}
zpwr starttabs              \# multi-tab setup\\
\end{codebox}

{\small\itshape\color{dimtext} // the tmux lifecycle:}

{\small\itshape\color{dimtext} //   start --> attach --> detach --> attach --> killmux}

{\small\itshape\color{dimtext} //   sessions persist through SSH disconnects}

{\small\itshape\color{dimtext} //   your terminal state is immortal}


\vspace{16pt}

\section{Key Duplication -- Multi-Pane Workflow}


\vspace{10pt}
{\large\bfseries\color{neoncyan} KEY DUPLICATION: TYPE ONCE, EXECUTE EVERYWHERE}\vspace{6pt}

The closest thing to being in two places at once.


\vspace{12pt}
{\bfseries\color{neonyellow} zpwr startsend}\vspace{4pt}

Send keystrokes to ALL panes in the current window. Whatever you type goes to every pane simultaneously.

\begin{codebox}
zpwr startsend              \# enable key broadcast\\
\end{codebox}

USE CASE: You have 4 panes SSHed into 4 servers. You need to run the same command on all of them. Startsend + type once = executed on all 4.


\vspace{12pt}
{\bfseries\color{neonyellow} zpwr startsendfull}\vspace{4pt}

Like startsend, but broadcasts across ALL windows and ALL panes in the entire tmux session.

\begin{codebox}
zpwr startsendfull          \# global broadcast\\
\end{codebox}

DANGER: every pane in every window receives input. Make sure you know what you are doing.


\vspace{12pt}
{\bfseries\color{neonyellow} zpwr stopsend}\vspace{4pt}

Turn off key broadcasting. Back to single-pane input.

\begin{codebox}
zpwr stopsend               \# disable broadcast\\
\end{codebox}


\vspace{12pt}
{\bfseries\color{neonyellow} THE MULTI-PANE WORKFLOW:}\vspace{4pt}

1. zpwr start                   \# launch tmux 2. Split into 4 panes (Ctrl-b \%) 3. SSH into different servers in each 4. zpwr startsend               \# enable broadcast 5. Type: sudo apt update \&\& sudo apt upgrade -y 6. All 4 servers update simultaneously 7. zpwr stopsend                \# back to normal

{\small\itshape\color{dimtext} // sysadmins call this poor mans ansible}

{\small\itshape\color{dimtext} // but honestly it is faster for quick ops}


\vspace{16pt}

\section{Clipboard \& Buffer Operations}


\vspace{10pt}
{\large\bfseries\color{neoncyan} CLIPBOARD OPS: COPY-PASTE ON STEROIDS}\vspace{6pt}

Because Ctrl-C Ctrl-V is so last millennium.


\vspace{12pt}
{\bfseries\color{neonyellow} zpwr pastebuffer}\vspace{4pt}

Paste from the tmux paste buffer into the terminal. Tmux has its own clipboard, separate from the system.

\begin{codebox}
zpwr pastebuffer            \# paste tmux buffer\\
\end{codebox}


\vspace{12pt}
{\bfseries\color{neonyellow} zpwr pastestring}\vspace{4pt}

Paste a specific string into the active pane. Programmatic pasting -- useful in scripts.

\begin{codebox}
zpwr pastestring 'echo hello' \# paste literal text\\
\end{codebox}


\vspace{12pt}
{\bfseries\color{neonyellow} zpwr pastecommand}\vspace{4pt}

Paste and execute a command in the active pane. Like pastestring but hits enter for you.

\begin{codebox}
zpwr pastecommand 'ls -la'  \# paste + execute\\
\end{codebox}


\vspace{12pt}
{\bfseries\color{neonyellow} zpwr google}\vspace{4pt}

Open a Google search in your default browser. Because sometimes you need to leave the terminal.

\begin{codebox}
zpwr google 'zsh completion' \# search the web\\
\end{codebox}


\vspace{12pt}
{\bfseries\color{neonyellow} zpwr openurl}\vspace{4pt}

Open any URL in the default browser.

\begin{codebox}
zpwr openurl https://github.com \# open URL\\
\end{codebox}

{\small\itshape\color{dimtext} // the clipboard pipeline:}

{\small\itshape\color{dimtext} //   terminal output --> tmux buffer --> pastebuffer}

{\small\itshape\color{dimtext} //   system clipboard --> pastestring --> pane}

{\small\itshape\color{dimtext} //   brain --> google --> browser --> back to terminal}


\vspace{16pt}

\section{Autocomplete Integration}


\vspace{10pt}
{\large\bfseries\color{neoncyan} AUTOCOMPLETE: THE TAB KEY'S FINAL FORM}\vspace{6pt}

Zsh-autocomplete + tmux = sentient tab completion.


\vspace{12pt}
{\bfseries\color{neonyellow} zpwr startauto}\vspace{4pt}

Enable the \textbf{ZPWR} autocomplete integration. Completions appear as you type, no Tab needed.

\begin{codebox}
zpwr startauto              \# enable live complete\\
\end{codebox}


\vspace{12pt}
{\bfseries\color{neonyellow} zpwr stopauto}\vspace{4pt}

Disable autocomplete. Sometimes you want silence.

\begin{codebox}
zpwr stopauto               \# disable live complete\\
\end{codebox}


\vspace{12pt}
{\bfseries\color{neonyellow} HOW IT WORKS WITH TMUX:}\vspace{4pt}

The autocomplete system hooks into zsh's completion engine and renders suggestions in real-time.

FEATURES: * Live suggestions as you type * History-based predictions * Path completion with preview * Command argument completion * Works across all tmux panes independently.


\vspace{12pt}
{\bfseries\color{neonyellow} TMUX + AUTOCOMPLETE TIPS:}\vspace{4pt}

* Each tmux pane runs its own zsh instance * Autocomplete state is per-pane * startauto/stopauto affects current pane only * If completions feel slow, check \$\textbf{ZPWR\_COMPS}\_CACHE

\begin{codebox}
zpwr stopauto               \# when it gets noisy\\
zpwr startauto              \# when you miss it\\
\end{codebox}

{\small\itshape\color{dimtext} // autocomplete is like a copilot for your shell}

{\small\itshape\color{dimtext} // it watches, it learns, it suggests}

{\small\itshape\color{dimtext} // toggle it off when you need to think in silence}


\vspace{16pt}

\section{Process Monitoring}


\vspace{10pt}
{\large\bfseries\color{neoncyan} PROCESS MONITORING: WATCHING THE WATCHERS}\vspace{6pt}

Every PID tells a story, most of them tragic.


\vspace{12pt}
{\bfseries\color{neonyellow} zpwr pstreemonitor}\vspace{4pt}

Displays a live process tree. Watch your processes branch, fork, and occasionally die unexpectedly.

\begin{codebox}
zpwr pstreemonitor          \# live process tree\\
\end{codebox}

Uses pstree under the hood with continuous refresh. Like htop's tree view but from the \textbf{ZPWR} command line.


\vspace{12pt}
{\bfseries\color{neonyellow} zpwr tty}\vspace{4pt}

Show info about the current terminal/tty.

\begin{codebox}
zpwr tty                    \# current tty info\\
\end{codebox}


\vspace{12pt}
{\bfseries\color{neonyellow} TMUX + MONITORING WORKFLOW:}\vspace{4pt}

The power move: dedicate a tmux pane to monitoring.

LAYOUT EXAMPLE: +-------------------+--------------------+ |                   |                    | |   your work       |  pstreemonitor     | |                   |                    | +-------------------+--------------------+ |         zpwr taillog                   | +-----------------------------------------+


\vspace{12pt}
{\bfseries\color{neonyellow} PROCESS HUNTING:}\vspace{4pt}

\textbf{ZPWR} integrates fzf with ps and lsof for killing.

Ctrl-Y Ctrl-K                  \# fzf kill by ps Ctrl-Y Ctrl-L                  \# fzf kill by lsof.

{\small\itshape\color{dimtext} // pstreemonitor in a side pane = instant situational awareness}

{\small\itshape\color{dimtext} // you will see zombie processes before they see you}


\vspace{16pt}


\chapter{NETWORK \& CONTAINERS}


\begin{center}

\includegraphics[width=0.85\textwidth]{/Users/wizard/.zpwr/docs/screenshots/art_firewall.png}
\end{center}

\section{Network Recon}


\vspace{10pt}
{\large\bfseries\color{neoncyan} NETWORK RECON: SCANNING THE WIRE}\vspace{6pt}

Packets don't lie, but DNS sometimes does.


\vspace{12pt}
{\bfseries\color{neonyellow} zpwr digs}\vspace{4pt}

Enhanced DNS lookup. Wraps dig with cleaner output. When you need to know where a domain really points.

\begin{codebox}
zpwr digs github.com         \# DNS lookup\\
zpwr digs example.com        \# check any domain\\
\end{codebox}

Shows A records, CNAME chains, TTLs -- the works. Faster than remembering dig +short +noall +answer.


\vspace{12pt}
{\bfseries\color{neonyellow} zpwr pi / zpwr ping}\vspace{4pt}

Ping a host. Simple, reliable, essential. 'pi' is the short alias because typing matters.

\begin{codebox}
zpwr pi google.com           \# quick ping\\
zpwr ping 8.8.8.8           \# verbose ping\\
\end{codebox}


\vspace{12pt}
{\bfseries\color{neonyellow} zpwr pir}\vspace{4pt}

Ping with extended options. More detailed network diagnostics for when simple ping is not enough.

\begin{codebox}
zpwr pir 192.168.1.1        \# extended ping\\
\end{codebox}


\vspace{12pt}
{\bfseries\color{neonyellow} RECON WORKFLOW:}\vspace{4pt}

1. zpwr digs suspect-domain.com  \# where does it resolve? 2. zpwr pi <ip-from-dig>         \# is it alive? 3. zpwr pir <ip>                 \# deeper probe

{\small\itshape\color{dimtext} // three verbs, complete network recon}

{\small\itshape\color{dimtext} // digs for DNS, pi for heartbeat, pir for details}

{\small\itshape\color{dimtext} // when in doubt, ping first -- dead hosts tell no tales}


\vspace{16pt}

\section{Tor Integration}


\vspace{10pt}
{\large\bfseries\color{neoncyan} TOR INTEGRATION: ANONYMOUS MODE ENGAGED}\vspace{6pt}

The onion router: because privacy is not a crime.


\vspace{12pt}
{\bfseries\color{neonyellow} zpwr torip}\vspace{4pt}

Check your current Tor exit node IP address. Confirms you're routing through the onion network.

\begin{codebox}
zpwr torip                  \# show tor exit IP\\
\end{codebox}

If Tor is running, this shows the exit node's IP. If not, it'll tell you Tor isn't connected.


\vspace{12pt}
{\bfseries\color{neonyellow} zpwr toriprenew}\vspace{4pt}

Request a new Tor circuit. Get a fresh exit node and a new apparent IP address.

\begin{codebox}
zpwr toriprenew             \# new identity\\
\end{codebox}

BEFORE YOU START:

Tor must be installed and running as a service.

MacOS:   brew install tor \&\& brew services start tor Linux:   sudo apt install tor \&\& sudo systemctl start tor.


\vspace{12pt}
{\bfseries\color{neonyellow} THE ANONYMOUS WORKFLOW:}\vspace{4pt}

1. Start tor service 2. zpwr torip            \# confirm you're connected 3. Do your thing 4. zpwr toriprenew       \# rotate identity 5. zpwr torip            \# verify new IP

DISCLAIMER: Tor provides anonymity, not invincibility. DNS leaks, browser fingerprinting, and OPSEC failures can still expose you. Use responsibly.

{\small\itshape\color{dimtext} // "I am become anonymous, destroyer of logs"}

{\small\itshape\color{dimtext} //   -- J. Robert Oppenheimer, probably not}


\vspace{16pt}

\section{Docker Management}


\vspace{10pt}
{\large\bfseries\color{neoncyan} DOCKER MANAGEMENT: CONTAINER WARFARE}\vspace{6pt}

It works on my machine --> we ship the machine.


\vspace{12pt}
{\bfseries\color{neonyellow} zpwr dockerwipe}\vspace{4pt}

The nuclear option for Docker. Removes EVERYTHING: containers, images, volumes, networks -- all of it.

\begin{codebox}
zpwr dockerwipe             \# purge all docker\\
\end{codebox}

WARNING: This is not gentle. It will: - Stop all running containers - Remove all containers - Remove all images - Remove all volumes - Prune everything.

WHEN TO USE: * Disk space emergency (docker images love eating disk) * Fresh start after a botched experiment * When 'docker system df' shows 50GB of shame.


\vspace{12pt}
{\bfseries\color{neonyellow} zpwr dfimage}\vspace{4pt}

Reverse-engineer a Docker image to show its Dockerfile. See how an image was built without the source.

\begin{codebox}
zpwr dfimage nginx:latest   \# reverse Dockerfile\\
zpwr dfimage myapp:v2       \# inspect any image\\
\end{codebox}

Useful for auditing third-party images. Trust but verify what's in that container.


\vspace{12pt}
{\bfseries\color{neonyellow} DOCKER TRIAGE:}\vspace{4pt}

1. \$ docker system df              \# how bad is it? 2. zpwr dfimage suspect:latest   \# what's inside? 3. zpwr dockerwipe               \# burn it all down 4. \$ docker system df              \# ahh, fresh disk

{\small\itshape\color{dimtext} // dockerwipe: for when docker system prune isn't angry enough}

{\small\itshape\color{dimtext} // dfimage: for when you don't trust the README}


\vspace{16pt}

\section{Clone \& GitHub Workflow}


\vspace{10pt}
{\large\bfseries\color{neoncyan} CLONE \& GITHUB: THE FULL GIT WORKFLOW}\vspace{6pt}

git clone is just the beginning of the relationship.


\vspace{12pt}
{\bfseries\color{neonyellow} zpwr clone}\vspace{4pt}

Clone a git repository with \textbf{ZPWR}'s enhanced setup. Handles auth, directory placement, and post-clone hooks.

\begin{codebox}
zpwr clone user/repo         \# clone from GitHub\\
zpwr clone https://...       \# clone any URL\\
\end{codebox}


\vspace{12pt}
{\bfseries\color{neonyellow} zpwr reveal}\vspace{4pt}

Open the current git repo's GitHub page in browser. Instant context switch from terminal to web.

\begin{codebox}
zpwr reveal                 \# open repo in browser\\
\end{codebox}


\vspace{12pt}
{\bfseries\color{neonyellow} zpwr revealrecurse}\vspace{4pt}

Open GitHub pages for the repo AND its submodules. When you need the full picture of a multi-repo project.

\begin{codebox}
zpwr revealrecurse          \# open all repos\\
\end{codebox}


\vspace{12pt}
{\bfseries\color{neonyellow} THE GITHUB WORKFLOW:}\vspace{4pt}

1. zpwr clone user/repo       \# get the code 2. \$ cd repo \&\& hack hack hack 3. zpwr reveal                \# check CI/issues 4. \$ git push                   \# ship it


\vspace{12pt}
{\bfseries\color{neonyellow} RELATED TOOLS:}\vspace{4pt}

\begin{codebox}
zpwr killremote             \# kill remote sessions\\
\end{codebox}

Killremote terminates remote SSH/tmux sessions that have gone stale. Cleanup for the disciplined operator.

{\small\itshape\color{dimtext} // reveal is the bridge between terminal and browser}

{\small\itshape\color{dimtext} // faster than typing github.com/user/repo}

{\small\itshape\color{dimtext} // your fingers will thank you}


\vspace{16pt}


\chapter{THE LOGGING SYSTEM}


\begin{center}

\includegraphics[width=0.85\textwidth]{/Users/wizard/.zpwr/docs/screenshots/art_data_stream.png}
\end{center}

\section{Log Levels}


\vspace{10pt}
{\large\bfseries\color{neoncyan} LOG LEVELS: SIGNAL VS NOISE}\vspace{6pt}

if a tree falls in the forest and nobody tails the log...

\textbf{ZPWR} has a full structured logging system with four levels. Each level has a file variant and a console variant.


\vspace{12pt}
{\bfseries\color{neonyellow} LEVEL 1: TRACE}\vspace{4pt}

\begin{codebox}
zpwr logtrace 'entering func'  \# most verbose\\
\end{codebox}

The firehose. Every function entry, variable assignment, and internal state change. Use for deep debugging only.


\vspace{12pt}
{\bfseries\color{neonyellow} LEVEL 2: DEBUG}\vspace{4pt}

\begin{codebox}
zpwr logdebug 'cache miss'    \# developer detail\\
\end{codebox}

Detailed info for development. Cache behavior, branch decisions, intermediate values.


\vspace{12pt}
{\bfseries\color{neonyellow} LEVEL 3: INFO}\vspace{4pt}

\begin{codebox}
zpwr loginfo 'update done'    \# normal operations\\
\end{codebox}

Standard operational messages. "Started", "Completed", "Loaded 47 plugins". The daily heartbeat.


\vspace{12pt}
{\bfseries\color{neonyellow} LEVEL 4: ERROR}\vspace{4pt}

\begin{codebox}
zpwr logerror 'file not found' \# something broke\\
\end{codebox}

Things that should not have happened but did. Missing files, failed commands, broken assumptions.


\vspace{12pt}
{\bfseries\color{neonyellow} SEVERITY PYRAMID:}\vspace{4pt}


\vspace{16pt}

\section{Console vs File Logging}


\vspace{10pt}
{\large\bfseries\color{neoncyan} CONSOLE VS FILE: TWO DESTINATIONS}\vspace{6pt}

To stderr or to disk, that is the question.

Every log verb comes in two flavors:


\vspace{12pt}
{\bfseries\color{neonyellow} FILE VARIANTS (write to \$\textbf{ZPWR\_LOGFILE}):}\vspace{4pt}

\begin{codebox}
zpwr logtrace 'msg'          \# --> logfile\\
zpwr logdebug 'msg'          \# --> logfile\\
zpwr loginfo 'msg'           \# --> logfile\\
zpwr logerror 'msg'          \# --> logfile\\
\end{codebox}

These write silently to the log file. No terminal output. Perfect for background operations and scripts.


\vspace{12pt}
{\bfseries\color{neonyellow} CONSOLE VARIANTS (write to stderr AND logfile):}\vspace{4pt}

\begin{codebox}
zpwr logtraceconsole 'msg'   \# --> stderr + file\\
zpwr logdebugconsole 'msg'   \# --> stderr + file\\
zpwr loginfoconsole 'msg'    \# --> stderr + file\\
zpwr logerrorconsole 'msg'   \# --> stderr + file\\
\end{codebox}

These print to your terminal AND write to the log. For interactive debugging when you need to see it live.


\vspace{12pt}
{\bfseries\color{neonyellow} THE LOGFILE:}\vspace{4pt}

\$\textbf{ZPWR\_LOGFILE} = \textbf{\$ZPWR\_LOCAL/zpwr.log}

\begin{codebox}
cat ZPWR\_LOGFILE | tail -20   \# recent entries\\
wc -l ZPWR\_LOGFILE           \# how big is it?\\
\end{codebox}


\vspace{12pt}
{\bfseries\color{neonyellow} CHOOSING YOUR VARIANT:}\vspace{4pt}

SCENARIO                          USE Background script running          loginfo Debugging interactively             logdebugconsole Tracking a rare bug                 logtraceconsole Production error handling           logerror.

{\small\itshape\color{dimtext} // file = silent, persistent, searchable}

{\small\itshape\color{dimtext} // console = loud, immediate, ephemeral}

{\small\itshape\color{dimtext} // use both strategically}


\vspace{16pt}

\section{Reading Logs}


\vspace{10pt}
{\large\bfseries\color{neoncyan} READING LOGS: THE ART OF LOG DIVING}\vspace{6pt}

Somewhere in that wall of text is the answer you seek.


\vspace{12pt}
{\bfseries\color{neonyellow} zpwr taillog}\vspace{4pt}

Live tail of the \textbf{ZPWR} log file with ccze colorization. Watch events stream by in glorious technicolor.

\begin{codebox}
zpwr taillog                \# live colorized tail\\
\end{codebox}

ccze automatically highlights timestamps, levels, paths, and errors. Makes raw logs actually readable.


\vspace{12pt}
{\bfseries\color{neonyellow} WHAT TO LOOK FOR:}\vspace{4pt}

ERROR  lines = something broke, investigate immediately WARN   lines = potential issue, worth noting INFO   lines = normal operations, sanity check DEBUG  lines = details for troubleshooting.

\begin{codebox}
TRACE  lines = deep internals, very verbose\\
\end{codebox}


\vspace{12pt}
{\bfseries\color{neonyellow} FILTERING TECHNIQUES:}\vspace{4pt}

\begin{codebox}
grep ERROR ZPWR\_LOGFILE     \# errors only\\
grep -c ERROR ZPWR\_LOGFILE  \# count errors\\
tail -100 ZPWR\_LOGFILE      \# last 100 lines\\
\end{codebox}


\vspace{12pt}
{\bfseries\color{neonyellow} TMUX LOG MONITORING SETUP:}\vspace{4pt}

Dedicate a tmux pane to zpwr taillog:

1. Split pane: Ctrl-b \% 2. In new pane: zpwr taillog 3. Work in other pane, watch logs stream 4. Errors glow red thanks to ccze


\vspace{12pt}
{\bfseries\color{neonyellow} LOG MAINTENANCE:}\vspace{4pt}

Logs grow forever if you let them. \textbf{ZPWR} handles rotation, but you can manually truncate:

\begin{codebox}
: > ZPWR\_LOGFILE            \# truncate to zero\\
\end{codebox}

{\small\itshape\color{dimtext} // taillog + ccze = the cyberpunk log experience}

{\small\itshape\color{dimtext} // errors in red, info in green, timestamps in cyan}


\vspace{16pt}

\section{Debug Mode}


\vspace{10pt}
{\large\bfseries\color{neoncyan} DEBUG MODE: WHEN THINGS GO SIDEWAYS}\vspace{6pt}

Enabling verbose output is how you stop guessing.


\vspace{12pt}
{\bfseries\color{neonyellow} \textbf{ZPWR\_DEBUG}}\vspace{4pt}

Set this env var to enable debug-level logging. More detail in the log, more signal for your hunt.

\begin{codebox}
export ZPWR\_DEBUG=true       \# enable debug\\
zpwr reload                  \# reload with debug\\
zpwr taillog                 \# watch debug output\\
\end{codebox}


\vspace{12pt}
{\bfseries\color{neonyellow} \textbf{ZPWR\_TRACE}}\vspace{4pt}

Maximum verbosity. Every function entry and exit, every variable, every decision point.

\begin{codebox}
export ZPWR\_TRACE=true       \# enable trace\\
zpwr reload                  \# reload with trace\\
\end{codebox}

WARNING: Trace mode generates MASSIVE log output. Your logfile will grow fast. Use sparingly.


\vspace{12pt}
{\bfseries\color{neonyellow} DEBUGGING WORKFLOW:}\vspace{4pt}

1. Reproduce the problem once to confirm it 2. \$ export \textbf{ZPWR\_DEBUG}=true 3. zpwr reload 4. Reproduce the problem again 5. zpwr taillog              \# find the smoking gun 6. Still lost? Escalate: 7. \$ export \textbf{ZPWR\_TRACE}=true 8. zpwr reload \&\& reproduce


\vspace{12pt}
{\bfseries\color{neonyellow} DISABLING DEBUG:}\vspace{4pt}

\begin{codebox}
unset ZPWR\_DEBUG ZPWR\_TRACE  \# back to normal\\
zpwr reload                  \# reload clean\\
\end{codebox}

{\small\itshape\color{dimtext} // debug mode: "trust but verify" in action}

{\small\itshape\color{dimtext} // trace mode: "verify everything, trust nothing"}

{\small\itshape\color{dimtext} // always disable after debugging or your log eats disk}


\vspace{16pt}


\chapter{INTROSPECTION}


\begin{center}
\includegraphics[width=0.85\textwidth]{/Users/wizard/.zpwr/docs/screenshots/art_deep_dive.png}
\end{center}

\section{Counting Everything}


\vspace{10pt}
{\large\bfseries\color{neoncyan} COUNTING EVERYTHING: THE CENSUS BUREAU}\vspace{6pt}

If you can't measure it, you can't brag about it.

\textbf{ZPWR} tracks the size of every component. These verbs tell you exactly how big your installation is.


\vspace{12pt}
{\bfseries\color{neonyellow} zpwr autoloadcount}\vspace{4pt}

Count of autoloaded zsh functions in the framework.

\begin{codebox}
zpwr autoloadcount           \# how many functions?\\
\end{codebox}


\vspace{12pt}
{\bfseries\color{neonyellow} zpwr scriptcount}\vspace{4pt}

Count of shell scripts in \textbf{\$ZPWR/scripts.}

\begin{codebox}
zpwr scriptcount             \# how many scripts?\\
\end{codebox}


\vspace{12pt}
{\bfseries\color{neonyellow} zpwr zshplugincount}\vspace{4pt}

Count of installed zsh plugins (zinit, oh-my-zsh, etc).

\begin{codebox}
zpwr zshplugincount          \# zsh plugins loaded\\
\end{codebox}


\vspace{12pt}
{\bfseries\color{neonyellow} zpwr vimplugincount}\vspace{4pt}

Count of installed vim/neovim plugins.

\begin{codebox}
zpwr vimplugincount          \# vim plugins loaded\\
\end{codebox}


\vspace{12pt}
{\bfseries\color{neonyellow} zpwr emacsplugincount}\vspace{4pt}

Count of installed emacs packages.

\begin{codebox}
zpwr emacsplugincount        \# emacs pkgs loaded\\
\end{codebox}


\vspace{12pt}
{\bfseries\color{neonyellow} THE FULL CENSUS:}\vspace{4pt}

\begin{codebox}
zpwr autoloadcount \&\& zpwr scriptcount \&\& \textbackslash{}\\
\end{codebox}

zpwr zshplugincount \&\& zpwr vimplugincount \&\& \textbackslash{} zpwr emacsplugincount.

Run all five to get the full picture. \textbf{ZPWR} typically ships with 300+ autoloads, 50+ scripts, and dozens of plugins across all three editors.

{\small\itshape\color{dimtext} // these counts tell you the scope of \textbf{ZPWR}}

{\small\itshape\color{dimtext} // hundreds of functions working in concert}

{\small\itshape\color{dimtext} // like a small city running on your laptop}


\vspace{16pt}

\section{Listing Everything}


\vspace{10pt}
{\large\bfseries\color{neoncyan} LISTING EVERYTHING: THE INVENTORY MANIFEST}\vspace{6pt}

Counting is nice, but names are power.

Counts tell you how many. Lists tell you WHAT. Every count verb has a list counterpart.


\vspace{12pt}
{\bfseries\color{neonyellow} zpwr autoloadlist}\vspace{4pt}

List every autoloaded function by name.

\begin{codebox}
zpwr autoloadlist            \# all function names\\
zpwr autoloadlist | grep fzf \# find fzf funcs\\
\end{codebox}


\vspace{12pt}
{\bfseries\color{neonyellow} zpwr scriptlist}\vspace{4pt}

List every script in the scripts directory.

\begin{codebox}
zpwr scriptlist              \# all scripts\\
\end{codebox}


\vspace{12pt}
{\bfseries\color{neonyellow} zpwr zshpluginlist}\vspace{4pt}

List installed zsh plugins with their sources.

\begin{codebox}
zpwr zshpluginlist           \# all zsh plugins\\
\end{codebox}


\vspace{12pt}
{\bfseries\color{neonyellow} zpwr vimpluginlist}\vspace{4pt}

List installed vim/neovim plugins.

\begin{codebox}
zpwr vimpluginlist           \# all vim plugins\\
\end{codebox}


\vspace{12pt}
{\bfseries\color{neonyellow} zpwr emacspluginlist}\vspace{4pt}

List installed emacs packages.

\begin{codebox}
zpwr emacspluginlist         \# all emacs packages\\
\end{codebox}


\vspace{12pt}
{\bfseries\color{neonyellow} POWER COMBOS:}\vspace{4pt}

\begin{codebox}
zpwr autoloadlist | wc -l    \# quick count\\
zpwr autoloadlist | fzf      \# fuzzy browse\\
zpwr scriptlist | sort       \# alphabetical\\
\end{codebox}

{\small\itshape\color{dimtext} // lists are grep-able, fzf-able, pipe-able}

{\small\itshape\color{dimtext} // the raw material for your own exploration}

{\small\itshape\color{dimtext} // pipe into fzf for interactive discovery}


\vspace{16pt}

\section{History Archaeology}


\vspace{10pt}
{\large\bfseries\color{neoncyan} HISTORY ARCHAEOLOGY: DIGGING UP THE PAST}\vspace{6pt}

Your terminal remembers everything, even the mistakes.


\vspace{12pt}
{\bfseries\color{neonyellow} zpwr hist}\vspace{4pt}

Browse command history with fzf. Search, preview, and re-execute past commands interactively.

\begin{codebox}
zpwr hist                   \# fzf history browser\\
\end{codebox}


\vspace{12pt}
{\bfseries\color{neonyellow} zpwr histedit}\vspace{4pt}

Edit the history file directly. Surgically remove embarrassing commands or fix typos in saved history.

\begin{codebox}
zpwr histedit               \# edit history file\\
\end{codebox}

Use with care. History surgery is delicate work.


\vspace{12pt}
{\bfseries\color{neonyellow} zpwr histfile}\vspace{4pt}

Show the path to the current history file. Useful when you're not sure which histfile is active.

\begin{codebox}
zpwr histfile               \# show histfile path\\
\end{codebox}


\vspace{12pt}
{\bfseries\color{neonyellow} zpwr catviminfo}\vspace{4pt}

Display vim's recent file history from viminfo. What files has vim touched recently?

\begin{codebox}
zpwr catviminfo             \# vim file history\\
\end{codebox}


\vspace{12pt}
{\bfseries\color{neonyellow} zpwr catrecentfviminfo / catnviminforecentf}\vspace{4pt}

Cross-reference vim/nvim recent files with the recentf list. Where vim and shell history intersect.

\begin{codebox}
zpwr catrecentfviminfo      \# vim+recentf cross\\
zpwr catnviminforecentf     \# nvim+recentf cross\\
\end{codebox}

{\small\itshape\color{dimtext} // your history is a goldmine of workflow patterns}

{\small\itshape\color{dimtext} // hist + fzf = instant command recall}

{\small\itshape\color{dimtext} // histedit = revisionist history (use responsibly)}


\vspace{16pt}

\section{Verb Discovery}


\vspace{10pt}
{\large\bfseries\color{neoncyan} VERB DISCOVERY: FINDING WHAT'S POSSIBLE}\vspace{6pt}

You can't use what you don't know exists.


\vspace{12pt}
{\bfseries\color{neonyellow} zpwr subcommands}\vspace{4pt}

List all zpwr subcommands (verbs). The master index of everything zpwr can do.

\begin{codebox}
zpwr subcommands             \# all subcommands\\
\end{codebox}


\vspace{12pt}
{\bfseries\color{neonyellow} zpwr subcommandscount}\vspace{4pt}

How many subcommands exist? The number keeps growing.

\begin{codebox}
zpwr subcommandscount        \# total count\\
\end{codebox}


\vspace{12pt}
{\bfseries\color{neonyellow} zpwr verbs (alias: subcommandsfzf)}\vspace{4pt}

Browse subcommands interactively with fzf. The best way to explore when you're not sure what to try.

\begin{codebox}
zpwr verbs         \# fzf subcommands\\
\end{codebox}


\vspace{12pt}
{\bfseries\color{neonyellow} zpwr verbslist / verbsfzf / verbscount}\vspace{4pt}

Alternative verb-centric views of available commands.

\begin{codebox}
zpwr verbslist              \# list all verbs\\
zpwr verbs               \# fzf browse verbs\\
zpwr verbscount             \# count verbs\\
\end{codebox}


\vspace{12pt}
{\bfseries\color{neonyellow} zpwr wizard / manual / tutorial / docs}\vspace{4pt}

You're reading it right now. The interactive encyclopedia with vim-style navigation. Very meta. All four names work.

\begin{codebox}
zpwr wizard                 \# start the wizard\\
zpwr wizard --toc           \# table of contents\\
zpwr wizard --search docker  \# search pages\\
zpwr wizard 42              \# jump to page 42\\
\end{codebox}

{\small\itshape\color{dimtext} // verbsfzf is the fastest way to discover (alias: subcommandsfzf)}

{\small\itshape\color{dimtext} // type a few letters, fzf narrows it down}

{\small\itshape\color{dimtext} // when in doubt, fzf it out}


\vspace{16pt}

\section{Shell State Inspection}


\vspace{10pt}
{\large\bfseries\color{neoncyan} SHELL STATE: KNOW THY ZSH}\vspace{6pt}

Introspection: when the shell looks in the mirror.


\vspace{12pt}
{\bfseries\color{neonyellow} zpwr exists}\vspace{4pt}

Check if a \textbf{ZPWR} component exists (functions, files, etc).

\begin{codebox}
zpwr exists zpwrReload       \# does this exist?\\
\end{codebox}


\vspace{12pt}
{\bfseries\color{neonyellow} zpwr existscommand}\vspace{4pt}

Check if a command exists in \textbf{PATH}. Cleaner than 'which' or 'command -v' in scripts.

\begin{codebox}
zpwr existscommand git      \# is git installed?\\
zpwr existscommand rg       \# is ripgrep here?\\
\end{codebox}


\vspace{12pt}
{\bfseries\color{neonyellow} zpwr reload}\vspace{4pt}

Reload the entire \textbf{ZPWR} framework without restarting. Re-sources configs, rebuilds caches, refreshes state.

\begin{codebox}
zpwr reload                 \# hot reload everything\\
\end{codebox}


\vspace{12pt}
{\bfseries\color{neonyellow} ZSH INSPECTION TOOLS:}\vspace{4pt}

These are built into zsh, but \textbf{ZPWR} makes them accessible:

\begin{codebox}
setopt                      \# list active options\\
zmodload                    \# list loaded modules\\
zstyle -L                   \# list all zstyles\\
\end{codebox}


\vspace{12pt}
{\bfseries\color{neonyellow} zpwr copycommand}\vspace{4pt}

Copy a command's path or definition to clipboard.

\begin{codebox}
zpwr copycommand ls         \# copy ls path\\
\end{codebox}


\vspace{12pt}
{\bfseries\color{neonyellow} zpwr opencommand}\vspace{4pt}

Open a command's source file in your editor. Works for functions, scripts, and binaries.

\begin{codebox}
zpwr opencommand zpwrReload  \# view the source\\
\end{codebox}

{\small\itshape\color{dimtext} // existscommand before you call, reload after you change}

{\small\itshape\color{dimtext} // opencommand when you need to understand the internals}


\vspace{16pt}


\chapter{INSTALLATION \& RESET}


\begin{center}
\includegraphics[width=0.85\textwidth]{/Users/wizard/.zpwr/docs/screenshots/figlet_init.png}
\end{center}

\section{Config Sessions}


\vspace{10pt}
{\large\bfseries\color{neoncyan} CONFIG SESSIONS: EDIT EVERYTHING FROM HERE}\vspace{6pt}

One verb to open any config, no path memorization.

\textbf{ZPWR} provides shortcut verbs for every config file you'll ever need to edit. No more 'vim \textbf{\textasciitilde{}/.zshrc}'.


\vspace{12pt}
{\bfseries\color{neonyellow} zpwr brc  -- Bash RC}\vspace{4pt}

Edit your .bashrc. For when bash compatibility matters.

\begin{codebox}
zpwr brc                    \# edit .bashrc\\
\end{codebox}


\vspace{12pt}
{\bfseries\color{neonyellow} zpwr zrc  -- Zsh RC}\vspace{4pt}

Edit your .zshrc. The most-touched config in \textbf{ZPWR}.

\begin{codebox}
zpwr zrc                    \# edit .zshrc\\
\end{codebox}


\vspace{12pt}
{\bfseries\color{neonyellow} zpwr vrc  -- Vim RC}\vspace{4pt}

Edit your .vimrc or init.vim/init.lua.

\begin{codebox}
zpwr vrc                    \# edit vim config\\
\end{codebox}


\vspace{12pt}
{\bfseries\color{neonyellow} zpwr trc  -- Tmux RC}\vspace{4pt}

Edit your .tmux.conf. Keybindings, statusbar, the works.

\begin{codebox}
zpwr trc                    \# edit tmux config\\
\end{codebox}


\vspace{12pt}
{\bfseries\color{neonyellow} zpwr tokens  -- Secrets File}\vspace{4pt}

Edit \textbf{\$ZPWR\_LOCAL/.tokens.sh.} API keys, secrets, env vars.

\begin{codebox}
zpwr tokens                 \# edit tokens/secrets\\
zpwr vimtokens              \# same, force vim\\
zpwr emacstokens            \# same, force emacs\\
\end{codebox}


\vspace{12pt}
{\bfseries\color{neonyellow} zpwr zcompdump  -- Completion Cache}\vspace{4pt}

Inspect or regenerate the zsh completion dump file. When completions act weird, this is where you look.

\begin{codebox}
zpwr zcompdump              \# inspect comp cache\\
\end{codebox}

{\small\itshape\color{dimtext} // brc, zrc, vrc, trc -- four letters, four configs}

{\small\itshape\color{dimtext} // tokens for secrets, zcompdump for completion fixes}


\vspace{16pt}

\section{Install \& Uninstall}


\vspace{10pt}
{\large\bfseries\color{neoncyan} INSTALL \& UNINSTALL: SETUP AND TEARDOWN}\vspace{6pt}

Easy to install, clean to remove -- no malware energy.


\vspace{12pt}
{\bfseries\color{neonyellow} zpwr install}\vspace{4pt}

Run the full \textbf{ZPWR} installation process. Installs dependencies, links configs, sets up plugins.

\begin{codebox}
zpwr install                \# full installation\\
\end{codebox}

WHAT IT DOES: 1. Installs system dependencies (brew/apt) 2. Links config files to their proper locations 3. Installs zsh, vim, and emacs plugins 4. Generates caches and completion dumps 5. Sets up tmux configuration.


\vspace{12pt}
{\bfseries\color{neonyellow} zpwr uninstall}\vspace{4pt}

Remove \textbf{ZPWR} configuration. Restores your shell to its pre-\textbf{ZPWR} state. Clean breakup, no hard feelings.

\begin{codebox}
zpwr uninstall              \# remove zpwr config\\
\end{codebox}


\vspace{12pt}
{\bfseries\color{neonyellow} CONFIG LINK MANAGEMENT:}\vspace{4pt}

\textbf{ZPWR} uses symlinks to place configs where tools expect them.

How it works: \textbf{\$ZPWR/configs/.zshrc}  -->  \textbf{\textasciitilde{}/.zshrc}     (symlink) \textbf{\$ZPWR/configs/.vimrc}  -->  \textbf{\textasciitilde{}/.vimrc}     (symlink) \textbf{\$ZPWR/configs/.tmux}   -->  \textbf{\textasciitilde{}/.tmux.conf} (symlink)

The source of truth lives in \$\textbf{ZPWR}, the symlinks just point there. Version controlled, portable, clean.


\vspace{12pt}
{\bfseries\color{neonyellow} PARTIAL OPERATIONS:}\vspace{4pt}

Sometimes you don't want the full install/uninstall. You just need to fix the symlinks:

Regenconfiglinks = recreate all config symlinks rmconfiglinks    = remove symlinks only.

{\small\itshape\color{dimtext} // install is idempotent -- run it again when things drift}

{\small\itshape\color{dimtext} // uninstall is clean -- your system goes back to stock}


\vspace{16pt}

\section{Reset \& Recovery}


\vspace{10pt}
{\large\bfseries\color{neoncyan} RESET \& RECOVERY: WHEN THINGS GO WRONG}\vspace{6pt}

Every system needs a panic button -- here are yours.


\vspace{12pt}
{\bfseries\color{neonyellow} zpwr reset}\vspace{4pt}

Full framework reset. Clears caches, regenerates configs, rebuilds completion dumps, relinks everything.

\begin{codebox}
zpwr reset                  \# the big red button\\
\end{codebox}

WHEN TO USE: * Completions are broken or stale * Config changes aren't taking effect * Something feels "off" and you can't pin it down * After a major \textbf{ZPWR} update.


\vspace{12pt}
{\bfseries\color{neonyellow} zpwr reload}\vspace{4pt}

Hot reload without resetting. Re-sources all configs and refreshes the environment. Faster than reset.

\begin{codebox}
zpwr reload                 \# hot reload\\
\end{codebox}

Use reload when you've edited a config and want to apply changes without nuking caches.


\vspace{12pt}
{\bfseries\color{neonyellow} ESCALATION LADDER:}\vspace{4pt}

Level 1: zpwr reload           \# re-source configs Level 2: zpwr reset            \# rebuild everything Level 3: zpwr install          \# full reinstall Level 4: rm -rf \$\textbf{ZPWR} \&\& clone \# scorched earth.

Start at Level 1. Only escalate if the problem persists. Most issues are fixed by reload. Reset handles the rest.


\vspace{12pt}
{\bfseries\color{neonyellow} COMMON RECOVERY SCENARIOS:}\vspace{4pt}

"Completions are wrong"    --> zpwr reset "alias isn't working"      --> zpwr reload "new plugin not loading"   --> zpwr reload "everything is broken"     --> zpwr reset "still broken after reset" --> zpwr install.

{\small\itshape\color{dimtext} // reload is ctrl-r for your shell framework}

{\small\itshape\color{dimtext} // reset is the factory reset you hope to never need}

{\small\itshape\color{dimtext} // but when you do need it, you'll be glad it exists}


\vspace{16pt}


\chapter{BATCH OPERATIONS}


\begin{center}

\includegraphics[width=0.85\textwidth]{/Users/wizard/.zpwr/docs/screenshots/art_override.png}
\end{center}

\section{Loop Warriors}


\vspace{10pt}
{\large\bfseries\color{neoncyan} LOOP WARRIORS — THE REPETITION TOOLKIT}\vspace{6pt}

Because doing things once is for civilians.

The batch operations verbs are the assault rifles of the \textbf{ZPWR} arsenal. Point them at a problem, pull the trigger, and watch every target fall.


\vspace{12pt}
{\bfseries\color{neonyellow} THE FOR LOOP — YOUR BASIC INFANTRY}\vspace{4pt}

\begin{codebox}
zpwr for <cmd> <args...>\\
\end{codebox}

Executes a command for each argument passed to it. Stdin becomes your ammo belt — pipe lines in, each one fires.

Example: run a command against multiple targets.

\begin{codebox}
echo 'file1.txt\textbackslash{}nfile2.txt' | zpwr for wc -l\\
\end{codebox}


\vspace{12pt}
{\bfseries\color{neonyellow} FOR10 — THE BURST-FIRE VARIANT}\vspace{4pt}

\begin{codebox}
zpwr for10 <cmd>\\
\end{codebox}

Runs your command 10 times. No arguments, no mercy. Perfect for stress testing, benchmarking, or just asserting dominance over your CPU.

Example: benchmark a command 10 times.

\begin{codebox}
zpwr for10 zpwr bench\\
\end{codebox}


\vspace{12pt}
{\bfseries\color{neonyellow} WHILETRUE — THE INFINITE LOOP}\vspace{4pt}

\begin{codebox}
zpwr whiletrue <cmd>\\
\end{codebox}

Runs your command forever. Ctrl-C is your only exit. For monitoring, watching, waiting. The sentinel loop.

Example: watch disk usage forever.

\begin{codebox}
zpwr whiletrue df -h /\\
\end{codebox}

WARNING: there is no break condition. You ARE the break.


\vspace{12pt}
{\bfseries\color{neonyellow} WHILESLEEP — THE PATIENT SNIPER}\vspace{4pt}

\begin{codebox}
zpwr whilesleep <seconds> <cmd>\\
\end{codebox}

Like whiletrue, but sleeps between rounds. Polling with grace. Surveillance with manners.

Example: check a URL every 30 seconds.

\begin{codebox}
zpwr whilesleep 30 curl -s -o /dev/null -w '\%\{http\_code\}' https://example.com\\
\end{codebox}

TACTICAL SUMMARY: for          iterate over stdin lines for10        run 10 times, no questions asked whiletrue    infinite loop, Ctrl-C to escape whilesleep   infinite loop with naptime between rounds.

{\small\itshape\color{dimtext} // in the loop, there is no time, only iterations}

{\small\itshape\color{dimtext} // the loop does not end. the operator ends the loop.}


\vspace{16pt}

\section{Directory Sweeps}


\vspace{10pt}
{\large\bfseries\color{neoncyan} DIRECTORY SWEEPS — CARPET BOMBING YOUR FILESYSTEM}\vspace{6pt}

One directory at a time, or all of them at once.

Single-file loops are amateur hour. Real operators sweep directories. They glob patterns. They parallelize. These verbs are your air support.


\vspace{12pt}
{\bfseries\color{neonyellow} FORDIR — THE DIRECTORY CRAWLER}\vspace{4pt}

\begin{codebox}
zpwr fordir <cmd>\\
\end{codebox}

Iterates into each subdirectory of the current dir and executes the command inside it. Cd in, run, cd out. Like a SWAT team clearing rooms one by one.

Example: check git status in every project dir.

\begin{codebox}
cd \textasciitilde{}/repos \&\& zpwr fordir git status -s\\
\end{codebox}


\vspace{12pt}
{\bfseries\color{neonyellow} FORDIRUPDATE — THE MASS UPDATER}\vspace{4pt}

\begin{codebox}
zpwr fordirupdate\\
\end{codebox}

Enters every subdirectory and runs git pull. Your entire project portfolio, synced in one command. Monday morning? Run this first. Coffee second.

Example: update all repos under \textbf{\textasciitilde{}/repos}

\begin{codebox}
cd \textasciitilde{}/repos \&\& zpwr fordirupdate\\
\end{codebox}


\vspace{12pt}
{\bfseries\color{neonyellow} EXECGLOB — GLOB AND EXECUTE}\vspace{4pt}

\begin{codebox}
zpwr execglob <glob> <cmd>\\
\end{codebox}

Finds files matching a glob pattern and runs your command against each one. Pattern matching meets execution. The regex sniper rifle of file operations.

Example: lint every Python file in the tree.

\begin{codebox}
zpwr execglob '*.py' pylint\\
\end{codebox}


\vspace{12pt}
{\bfseries\color{neonyellow} EXECGLOBPARALLEL — CONCURRENT DEVASTATION}\vspace{4pt}

\begin{codebox}
zpwr execglobparallel <glob> <cmd>\\
\end{codebox}

Same as execglob, but fires all jobs in parallel. Every CPU core gets drafted into service. Your fan will spin up. This is normal. This is power.

Example: format every .rs file in parallel.

\begin{codebox}
zpwr execglobparallel '*.rs' rustfmt\\
\end{codebox}

CHOOSING YOUR WEAPON: fordir             run cmd in each subdirectory fordirupdate       git pull in each subdirectory execglob           glob match + execute sequentially execglobparallel   glob match + execute in parallel.

{\small\itshape\color{dimtext} // sequential is predictable. parallel is fast.}

{\small\itshape\color{dimtext} // choose based on whether output order matters.}

{\small\itshape\color{dimtext} // when in doubt, go parallel. life is short.}


\vspace{16pt}

\section{Polyglot Execution}


\vspace{10pt}
{\large\bfseries\color{neoncyan} POLYGLOT EXECUTION — BEYOND THE SHELL}\vspace{6pt}

zsh is the conductor, but the orchestra speaks many tongues.

Sometimes zsh loops aren't enough. Sometimes you need Python. Sometimes you need background jobs. Sometimes you need to poll a service until it wakes up or dies.


\vspace{12pt}
{\bfseries\color{neonyellow} EXECPY — THE PYTHON BRIDGE}\vspace{4pt}

\begin{codebox}
zpwr execpy <python\_expression>\\
\end{codebox}

Fires a Python expression from the comfort of your shell. No need to open a REPL. No need for a .py file. Quick math, string ops, one-liners — python as a verb.

Example: quick calculation.

\begin{codebox}
zpwr execpy '2**128'\\
\end{codebox}

=> 340282366920938463463374607431768211456.


\vspace{12pt}
{\bfseries\color{neonyellow} BACKGROUND — THE SHADOW OPERATIVE}\vspace{4pt}

\begin{codebox}
zpwr background <cmd>\\
\end{codebox}

Sends a command into the background. Detached. Silent. It runs in the shadows while you keep working. Output goes to the void unless you redirect it.

Example: start a long build in the background.

\begin{codebox}
zpwr background make -j8\\
\end{codebox}


\vspace{12pt}
{\bfseries\color{neonyellow} TIMER — THE STOPWATCH}\vspace{4pt}

\begin{codebox}
zpwr timer <cmd>\\
\end{codebox}

Wraps any command in a high-resolution timer. How long did that take? Now you know. No guessing. Benchmarking made trivial.

Example: time your shell startup.

\begin{codebox}
zpwr timer zsh -ic exit\\
\end{codebox}


\vspace{12pt}
{\bfseries\color{neonyellow} POLL — THE WATCHDOG}\vspace{4pt}

\begin{codebox}
zpwr poll <cmd>\\
\end{codebox}

Runs a command repeatedly until it succeeds (exit 0). Waiting for a server to come up? A build to finish? Poll is your patient, relentless sentinel.

Example: wait for a port to open.

\begin{codebox}
zpwr poll nc -z localhost 8080\\
\end{codebox}

THE BATCH OPS PHILOSOPHY: Loops        for, for10, whiletrue, whilesleep Sweeps       fordir, fordirupdate, execglob* Polyglot     execpy, background, timer, poll.

{\small\itshape\color{dimtext} // the terminal is not a single-threaded mind}

{\small\itshape\color{dimtext} // it is a factory floor. these verbs are the machines.}

{\small\itshape\color{dimtext} // feed them work. they do not tire. they do not complain.}


\vspace{16pt}


\chapter{FORGIT DEEP DIVE}


\begin{center}
\includegraphics[width=0.85\textwidth]{/Users/wizard/.zpwr/docs/screenshots/art_the_grid.png}
\end{center}

\section{Forgit Add \& Diff}


\vspace{10pt}
{\large\bfseries\color{neoncyan} FORGIT ADD \& DIFF — INTERACTIVE STAGING}\vspace{6pt}

git add, but make it beautiful.

Forgit is the fusion of fzf and git. It wraps the most common git operations in fuzzy, interactive, preview-enabled interfaces. You will never go back to raw git add.


\vspace{12pt}
{\bfseries\color{neonyellow} FORGITADD — INTERACTIVE STAGING}\vspace{4pt}

\begin{codebox}
zpwr forgitadd\\
\end{codebox}

Opens an fzf interface showing all unstaged changes. Each file gets a diff preview in the right pane. Tab to select multiple files. Enter to stage them.

The workflow: 1. Run zpwr forgitadd 2. Browse files with arrow keys 3. See the diff preview for each file 4. Tab to multi-select 5. Enter to stage selected files.

No more git add -p scrolling through 500-line diffs. No more accidentally staging that debug print.


\vspace{12pt}
{\bfseries\color{neonyellow} FORGITDIFF — \textbf{VISUAL} DIFF BROWSER}\vspace{4pt}

\begin{codebox}
zpwr forgitdiff\\
\end{codebox}

Opens an fzf browser of all changed files with a live diff preview. Read diffs without committing to anything. Pure reconnaissance.

Example workflow:

\begin{codebox}
zpwr forgitdiff           \# browse all diffs visually\\
zpwr forgitadd            \# stage what you want\\
zpwr gitcommit 'msg'      \# commit with message\\
\end{codebox}

WHY FORGIT > RAW GIT: Raw git add:     type filenames from memory like a caveman forgitadd:       browse, preview, multi-select, confirm Raw git diff:    wall of text scrolling past forgitdiff:      file-by-file preview with fuzzy search.

{\small\itshape\color{dimtext} // staging should not be a memory test}

{\small\itshape\color{dimtext} // it should be a visual, interactive experience}

{\small\itshape\color{dimtext} // forgit turns git into a GUI without leaving the terminal}


\vspace{16pt}

\section{Forgit Log \& Stash}


\vspace{10pt}
{\large\bfseries\color{neoncyan} FORGIT LOG \& STASH — VISUAL HISTORY}\vspace{6pt}

Time travel with a preview window.

Git log is a river of hashes. Git stash is a pile of forgotten context. Forgit turns both into something you can actually navigate without losing your mind.


\vspace{12pt}
{\bfseries\color{neonyellow} FORGITLOG — THE COMMIT BROWSER}\vspace{4pt}

\begin{codebox}
zpwr forgitlog\\
\end{codebox}

Opens your git history in fzf with commit previews. Each commit shows its full diff in the preview pane. Search commits by message, author, or hash.

What you see: Left pane:  commit graph with short messages Right pane: full diff of highlighted commit Search:     fuzzy match on any commit field.

"Who wrote this garbage?" — you, using forgitlog, discovering it was also you, three months ago.


\vspace{12pt}
{\bfseries\color{neonyellow} FORGITSTASH — THE STASH MANAGER}\vspace{4pt}

\begin{codebox}
zpwr forgitstash\\
\end{codebox}

Browse your stash entries with preview diffs. See what is in each stash before applying it. No more git stash pop roulette.

The old way: git stash list              \# cryptic one-liners git stash show -p stash@\{2\}  \# which one was it again? git stash pop               \# pray.

The forgit way:

\begin{codebox}
zpwr forgitstash            \# browse, preview, apply\\
\end{codebox}

LOG + STASH TOGETHER: These two verbs give you eyes into the past. Forgitlog shows committed history. Forgitstash shows uncommitted history. Between them, nothing is lost.

{\small\itshape\color{dimtext} // the commit graph is not a wall of text}

{\small\itshape\color{dimtext} // it is a timeline. forgit lets you walk it.}

{\small\itshape\color{dimtext} // stashes are not forgotten. they are filed.}


\vspace{16pt}

\section{Forgit Ignore \& Clean}


\vspace{10pt}
{\large\bfseries\color{neoncyan} FORGIT IGNORE \& CLEAN — .GITIGNORE MANAGEMENT}\vspace{6pt}

Because node\_modules should never see the inside of a commit.

The .gitignore is your repo's bouncer list. These verbs manage it with the precision of a nightclub velvet rope.


\vspace{12pt}
{\bfseries\color{neonyellow} FORGITIGNORE — THE IGNORE BROWSER}\vspace{4pt}

\begin{codebox}
zpwr forgitignore\\
\end{codebox}

Interactive gitignore template browser powered by gitignore.io. Search for languages, frameworks, IDEs. Never hand-write a .gitignore again.


\vspace{12pt}
{\bfseries\color{neonyellow} THE IGNORE FAMILY:}\vspace{4pt}

Forgitignorelist     list available gitignore templates forgitignoreget      fetch a specific template by name forgitignoreupdate   update your local template cache forgitignoreclean    clean the template cache.

Example: generate a Python + macOS gitignore.

\begin{codebox}
zpwr forgitignoreget python,macos > .gitignore\\
\end{codebox}

Example: browse all available templates.

\begin{codebox}
zpwr forgitignorelist\\
\end{codebox}


\vspace{12pt}
{\bfseries\color{neonyellow} FORGITCLEAN — THE UNTRACKED FILE PURGE}\vspace{4pt}

\begin{codebox}
zpwr forgitclean\\
\end{codebox}

Shows untracked files in an fzf interface with previews. Select which ones to delete. Surgical, not scorched earth. Unlike git clean -fd, you see before you destroy.

The old way: git clean -fd    \# delete everything. Hope for the best.

The forgit way:

\begin{codebox}
zpwr forgitclean  \# browse, select, confirm, then delete.\\
\end{codebox}

IGNORE WORKFLOW: 1. zpwr forgitignorelist    browse templates 2. zpwr forgitignoreget     fetch what you need 3. zpwr forgitclean         purge untracked cruft 4. zpwr forgitadd           stage the clean repo.

{\small\itshape\color{dimtext} // a clean repo is a happy repo}

{\small\itshape\color{dimtext} // an ignored file is a file that cannot hurt you}

{\small\itshape\color{dimtext} // let the templates do the thinking}


\vspace{16pt}

\section{Forgit Reset \& Restore}


\vspace{10pt}
{\large\bfseries\color{neoncyan} FORGIT RESET \& RESTORE — UNDO WITH EYES OPEN}\vspace{6pt}

The nuclear options, but with a safety preview.

Git reset and restore are the most terrifying commands in version control. Forgit wraps them in bubble wrap and gives you a preview before you pull the trigger.


\vspace{12pt}
{\bfseries\color{neonyellow} FORGITRESET — UNSTAGE WITH PREVIEW}\vspace{4pt}

\begin{codebox}
zpwr forgitreset\\
\end{codebox}

Shows staged files in fzf with diff previews. Select files to unstage. Like git reset HEAD but you can see exactly what you are unstaging first.

Use case: you ran git add -A like a maniac and now you need to surgically remove files.

\begin{codebox}
zpwr forgitreset    \# browse and unstage selectively\\
\end{codebox}


\vspace{12pt}
{\bfseries\color{neonyellow} FORGITRESTORE — DISCARD CHANGES SAFELY}\vspace{4pt}

\begin{codebox}
zpwr forgitrestore\\
\end{codebox}

Shows modified files with diff previews. Select files to restore to their last committed state. This discards changes permanently. No undo.

The preview saves you. You see the diff before the changes evaporate into the void.


\vspace{12pt}
{\bfseries\color{neonyellow} FORGITWARN — THE SAFETY NET}\vspace{4pt}

\begin{codebox}
zpwr forgitwarn\\
\end{codebox}

Displays warnings about uncommitted changes, diverged branches, and other hazardous states. Your pre-flight checklist before doing anything destructive.


\vspace{12pt}
{\bfseries\color{neonyellow} FORGITINFO — THE STATUS DASHBOARD}\vspace{4pt}

\begin{codebox}
zpwr forgitinfo\\
\end{codebox}

Comprehensive repo information at a glance. Branch, remotes, recent commits, dirty state. Everything you need before making decisions.

THE FORGIT COMPLETE ARSENAL: forgitadd        interactive stage forgitdiff       browse diffs forgitlog        commit history browser forgitstash      stash manager forgitreset      interactive unstage forgitrestore    discard changes with preview forgitclean      purge untracked files forgitignore*    gitignore template management forgitwarn       hazard warnings forgitinfo       repo dashboard.

{\small\itshape\color{dimtext} // forgit is not a git replacement}

{\small\itshape\color{dimtext} // it is git with eyes. git with a UI. git for humans.}


\vspace{16pt}


\chapter{ZPWR ENVIRONMENT VARIABLES}


\begin{center}

\includegraphics[width=0.85\textwidth]{/Users/wizard/.zpwr/docs/screenshots/art_system_core.png}
\end{center}

\section{Core Paths}


\vspace{10pt}
{\large\bfseries\color{neoncyan} CORE PATHS — THE SKELETON OF THE MACHINE}\vspace{6pt}

Every system has bones. These are \textbf{ZPWR}'s.

\textbf{ZPWR} does not hardcode paths. Every directory is an environment variable. Move the root, and the entire system follows. This is not a dotfiles repo — it is a relocatable operating system.


\vspace{12pt}
{\bfseries\color{neonyellow} \$\textbf{ZPWR} — THE ROOT OF ALL THINGS}\vspace{4pt}

\begin{codebox}
echo ZPWR   \# typically \textasciitilde{}/.zpwr\\
\end{codebox}

The root directory. Every other path derives from this. Change this one variable and the whole tree moves.


\vspace{12pt}
{\bfseries\color{neonyellow} \$\textbf{ZPWR\_LOCAL} — YOUR PRIVATE SANCTUM}\vspace{4pt}

\begin{codebox}
echo ZPWR\_LOCAL   \# ZPWR/local\\
\end{codebox}

Gitignored. Your tokens, secrets, local overrides. Nothing in here ever touches a remote repository. Your .tokens.sh lives here. Guard it with your life.


\vspace{12pt}
{\bfseries\color{neonyellow} \$\textbf{ZPWR\_SCRIPTS} — THE ARMORY}\vspace{4pt}

\begin{codebox}
echo ZPWR\_SCRIPTS   \# ZPWR/scripts\\
\end{codebox}

Executable scripts. Bash, Python, Perl — whatever gets the job done. These are on your \textbf{PATH}.


\vspace{12pt}
{\bfseries\color{neonyellow} \$\textbf{ZPWR\_TMUX} — THE MULTIPLEXER CONFIG}\vspace{4pt}

\begin{codebox}
echo ZPWR\_TMUX   \# ZPWR/tmux\\
\end{codebox}

Tmux configuration, scripts, and status line logic. The control center for your terminal multiplexer.


\vspace{12pt}
{\bfseries\color{neonyellow} \$\textbf{ZPWR\_INSTALL} — THE INSTALLER}\vspace{4pt}

\begin{codebox}
echo ZPWR\_INSTALL   \# ZPWR/install\\
\end{codebox}

OS-specific installation scripts. MacOS, Ubuntu, Fedora, Arch — each gets its own installer. Run zpwr install and the right one fires.

THE \textbf{PATH} MAP: \$\textbf{ZPWR}                \textbf{\textasciitilde{}/.zpwr} (the root) \$\textbf{ZPWR\_LOCAL}          \textbf{\$ZPWR/local} (gitignored secrets) \$\textbf{ZPWR\_SCRIPTS}        \textbf{\$ZPWR/scripts} (executables) \$\textbf{ZPWR\_TMUX}           \textbf{\$ZPWR/tmux} (multiplexer) \$\textbf{ZPWR\_INSTALL}        \textbf{\$ZPWR/install} (OS installers)

{\small\itshape\color{dimtext} // paths are not strings. they are load-bearing structures.}

{\small\itshape\color{dimtext} // hardcode a path and you weld yourself to the floor.}

{\small\itshape\color{dimtext} // use the variables. always use the variables.}


\vspace{16pt}

\section{Autoload System}


\vspace{10pt}
{\large\bfseries\color{neoncyan} THE AUTOLOAD SYSTEM — LAZY LOADING FOR HACKERS}\vspace{6pt}

Why load everything when you can load nothing until needed?

Zsh autoloading is the secret to \textbf{ZPWR}'s fast startup. Functions are not loaded into memory until you call them. The shell knows they exist, but their code stays on disk until the moment of invocation. Maximum laziness. Maximum speed.


\vspace{12pt}
{\bfseries\color{neonyellow} \$\textbf{ZPWR\_AUTOLOAD} — THE FUNCTION WAREHOUSE}\vspace{4pt}

\begin{codebox}
echo ZPWR\_AUTOLOAD   \# ZPWR/autoload\\
\end{codebox}

The root of the autoload tree. Every file in this tree becomes a function name. The filename IS the function.


\vspace{12pt}
{\bfseries\color{neonyellow} \$\textbf{ZPWR\_AUTOLOAD}\_COMMON — CROSS-PLATFORM FUNCTIONS}\vspace{4pt}

\begin{codebox}
echo ZPWR\_AUTOLOAD\_COMMON   \# ZPWR/autoload/common\\
\end{codebox}

Functions that work on every OS. The common denominator. MacOS, Linux, BSD — these functions do not care.


\vspace{12pt}
{\bfseries\color{neonyellow} HOW AUTOLOADING WORKS:}\vspace{4pt}

1. Zsh scans \$fpath directories at startup 2. Each filename is registered as a function stub 3. When you call the function, zsh reads the file 4. The file contents replace the stub 5. The function executes with the loaded code

This means hundreds of autoload stubs cost near-zero at startup. Only the ones you actually use consume memory.


\vspace{12pt}
{\bfseries\color{neonyellow} THE COMPILE ACCELERATION:}\vspace{4pt}

\begin{codebox}
zpwr compile       \# compile autoloads to .zwc\\
zpwr recompile     \# recompile all\\
zpwr uncompile     \# remove compiled files\\
\end{codebox}

Compiled .zwc files load faster than raw source. \textbf{ZPWR} compiles its autoload dirs into wordcode files. It is like bytecode for your shell functions.

INSPECTION VERBS:

\begin{codebox}
zpwr autoload          show the autoload system\\
zpwr autoloadcount     count autoloaded functions\\
zpwr autoloadlist      list all autoloaded functions\\
\end{codebox}

{\small\itshape\color{dimtext} // the fastest code is code that never runs}

{\small\itshape\color{dimtext} // autoloading is the art of deferred execution}

{\small\itshape\color{dimtext} // 400 functions, zero startup cost. that is the deal.}


\vspace{16pt}

\section{Behavior Flags}


\vspace{10pt}
{\large\bfseries\color{neoncyan} BEHAVIOR FLAGS — THE KNOBS AND DIALS}\vspace{6pt}

Tune the machine without cracking open the engine.

\textbf{ZPWR} is not a take-it-or-leave-it monolith. These boolean and string flags let you control startup behavior, logging verbosity, and debug output without editing a single line of source code.


\vspace{12pt}
{\bfseries\color{neonyellow} \$\textbf{ZPWR\_PROFILING} — STARTUP PROFILER}\vspace{4pt}

\begin{codebox}
export ZPWR\_PROFILING=true\\
\end{codebox}

Enables zsh startup profiling. Shows exactly where your milliseconds go during shell init. Pair with.

\begin{codebox}
zpwr bench for the full picture.\\
\end{codebox}

Set it, open a new shell, and read the profiling data. Find the bottleneck. Kill the bottleneck. Repeat.


\vspace{12pt}
{\bfseries\color{neonyellow} \$\textbf{ZPWR\_DEBUG} — DEBUG MODE}\vspace{4pt}

\begin{codebox}
export ZPWR\_DEBUG=true\\
\end{codebox}

Turns on debug-level log messages. When something is broken and you need the machine to talk, flip this. Check output with zpwr logdebug.


\vspace{12pt}
{\bfseries\color{neonyellow} \$\textbf{ZPWR\_TRACE} — FULL TRACE MODE}\vspace{4pt}

\begin{codebox}
export ZPWR\_TRACE=true\\
\end{codebox}

The firehose. Every function call, every expansion, traced to the log. Use zpwr trace to run commands with tracing, or zpwr logtrace to read trace logs. WARNING: produces enormous output. Use surgically.


\vspace{12pt}
{\bfseries\color{neonyellow} \$\textbf{ZPWR\_BANNER}\_CLEARLIST — STARTUP BANNER CONTROL}\vspace{4pt}

\begin{codebox}
export ZPWR\_BANNER\_CLEARLIST=true\\
\end{codebox}

Controls whether the \textbf{ZPWR} banner and clearlist show at shell startup. The banner is glorious, but if you open 50 terminals a day, you might want silence.

Related verbs:

\begin{codebox}
zpwr banner          show the banner manually\\
zpwr clearlist       show clearlist manually\\
zpwr bannerlolcat    banner with rainbow colors\\
\end{codebox}

QUICK REFERENCE: \textbf{ZPWR\_PROFILING}        startup timing data \textbf{ZPWR\_DEBUG}            debug-level logging \textbf{ZPWR\_TRACE}            full execution trace \textbf{ZPWR\_BANNER}\_CLEARLIST startup banner toggle.

{\small\itshape\color{dimtext} // the flags do not change the code. they reveal it.}

{\small\itshape\color{dimtext} // debugging is not failure. it is x-ray vision.}


\vspace{16pt}

\section{Plugin System}


\vspace{10pt}
{\large\bfseries\color{neoncyan} PLUGIN SYSTEM — THE TURBO LOADER}\vspace{6pt}

zinit: loading plugins at the speed of light.

\textbf{ZPWR} uses zinit (formerly zplugin) as its plugin manager. But not the basic way. The turbo way. Plugins load asynchronously after prompt appears. You see your prompt instantly. Plugins catch up.


\vspace{12pt}
{\bfseries\color{neonyellow} \$\textbf{ZPWR\_PLUGIN}\_MANAGER — CHOOSE YOUR LOADER}\vspace{4pt}

\begin{codebox}
echo ZPWR\_PLUGIN\_MANAGER   \# zinit\\
\end{codebox}

Controls which plugin manager \textbf{ZPWR} uses. The default and recommended value is zinit. Other managers exist but zinit's turbo mode is what makes \textbf{ZPWR} fast.


\vspace{12pt}
{\bfseries\color{neonyellow} TURBO LOADING EXPLAINED:}\vspace{4pt}

Normal plugin loading: source plugin1.zsh  \# blocks source plugin2.zsh  \# blocks source plugin3.zsh  \# blocks \# prompt finally appears after ALL plugins load.

Zinit turbo loading: prompt appears immediately plugin1 loads in background ... done plugin2 loads in background ... done plugin3 loads in background ... done \# user never notices the delay


\vspace{12pt}
{\bfseries\color{neonyellow} PLUGIN INSPECTION:}\vspace{4pt}

\begin{codebox}
zpwr zshplugincount     count loaded zsh plugins\\
zpwr zshpluginlist      list all loaded zsh plugins\\
zpwr vimplugincount     count vim plugins\\
zpwr vimpluginlist      list all vim plugins\\
zpwr emacsplugincount   count emacs plugins\\
zpwr emacspluginlist    list all emacs plugins\\
\end{codebox}


\vspace{12pt}
{\bfseries\color{neonyellow} UPDATE AND MAINTAIN:}\vspace{4pt}

\begin{codebox}
zpwr updatezinit        update zinit itself\\
zpwr updatedeps          update all plugin deps\\
zpwr plugins             plugin information\\
\end{codebox}

{\small\itshape\color{dimtext} // plugins are not bloat. they are force multipliers.}

{\small\itshape\color{dimtext} // but only if they load in zero perceived time.}

{\small\itshape\color{dimtext} // turbo mode makes 50 plugins feel like zero.}


\vspace{16pt}

\section{Display \& Input}


\vspace{10pt}
{\large\bfseries\color{neoncyan} DISPLAY \& INPUT — THE HUMAN INTERFACE LAYER}\vspace{6pt}

How the machine talks to you, and how you talk back.

These variables control what you see and how fast the shell responds to your keystrokes. Milliseconds matter when your fingers are moving at light speed.


\vspace{12pt}
{\bfseries\color{neonyellow} \$\textbf{ZPWR\_EXA}\_COMMAND — THE LS REPLACEMENT}\vspace{4pt}

\begin{codebox}
echo ZPWR\_EXA\_COMMAND   \# exa or eza\\
\end{codebox}

Controls which modern ls replacement \textbf{ZPWR} uses. Exa is the original, eza is the maintained fork. Both give you colored, git-aware directory listings.

\textbf{ZPWR} verbs like zpwr ls and zpwr clearls use this. Set it in your tokens file to override the default.


\vspace{12pt}
{\bfseries\color{neonyellow} \$\textbf{ZPWR\_KEYTIMEOUT} — KEYSTROKE LATENCY}\vspace{4pt}

\begin{codebox}
echo ZPWR\_KEYTIMEOUT   \# in hundredths of a second\\
\end{codebox}

How long zsh waits for the next key in a multi-key sequence. Lower = snappier. Too low = missed sequences. Default is tuned for fast typists.

If vi-mode feels laggy on Esc, lower this. If key combos misfire, raise this. Sweet spot is personal. Experiment.


\vspace{12pt}
{\bfseries\color{neonyellow} \$\textbf{ZPWR\_EXPAND}\_PRE\_EXEC — ALIAS EXPANSION}\vspace{4pt}

\begin{codebox}
echo ZPWR\_EXPAND\_PRE\_EXEC\\
\end{codebox}

Controls whether \textbf{ZPWR} expands aliases before execution. When enabled, you type the alias and see the full command in your history. Transparency mode.

With expansion:    you type ga  -> history shows git add Without expansion: you type ga  -> history shows ga.


\vspace{12pt}
{\bfseries\color{neonyellow} \$\textbf{ZPWR\_SEND}\_KEYS\_FULL — TMUX KEY PASSTHROUGH}\vspace{4pt}

\begin{codebox}
echo ZPWR\_SEND\_KEYS\_FULL\\
\end{codebox}

Controls full send-keys behavior in tmux integration. When \textbf{ZPWR} manages tmux panes, this determines how keystrokes are forwarded to target panes.

TUNING GUIDE: Feeling laggy?       Lower \textbf{ZPWR\_KEYTIMEOUT} Missing combos?      Raise \textbf{ZPWR\_KEYTIMEOUT} Want transparency?   Enable \textbf{ZPWR\_EXPAND}\_PRE\_EXEC Using eza?           Set \textbf{ZPWR\_EXA}\_COMMAND=eza.

{\small\itshape\color{dimtext} // the interface is not decoration. it is velocity.}

{\small\itshape\color{dimtext} // every millisecond of latency is a papercut.}

{\small\itshape\color{dimtext} // tune these until your shell feels like thought.}


\vspace{16pt}

\section{Runtime State}


\vspace{10pt}
{\large\bfseries\color{neoncyan} RUNTIME STATE — THE MACHINE'S SELF-AWARENESS}\vspace{6pt}

The shell knows what it is. Do you?

These variables hold runtime state — what OS you are on, what TTY you are using, where logs go. They are set at startup and consulted throughout the session.


\vspace{12pt}
{\bfseries\color{neonyellow} \$\textbf{ZPWR\_TTY} — YOUR TERMINAL IDENTITY}\vspace{4pt}

\begin{codebox}
echo ZPWR\_TTY   \# e.g., /dev/ttys003\\
\end{codebox}

The TTY device for the current shell session. Used by \textbf{ZPWR} to track which terminal is which, especially when managing multiple tmux panes.

The zpwr tty verb shows your current TTY info.


\vspace{12pt}
{\bfseries\color{neonyellow} \$\textbf{ZPWR\_OS}\_TYPE — PLATFORM DETECTION}\vspace{4pt}

\begin{codebox}
echo ZPWR\_OS\_TYPE   \# Darwin, Linux, etc.\\
\end{codebox}

Detected at startup. Used throughout \textbf{ZPWR} to branch on platform-specific behavior. MacOS vs Linux vs BSD, the right code path fires automatically.

This is why \textbf{ZPWR} works on macOS and Linux without you changing a single setting. The OS type decides.


\vspace{12pt}
{\bfseries\color{neonyellow} \$\textbf{ZPWR\_LOGFILE} — THE BLACK BOX}\vspace{4pt}

\begin{codebox}
echo ZPWR\_LOGFILE\\
\end{codebox}

Where \textbf{ZPWR} writes its log output. Debug, info, error, trace — all levels go here when enabled.

Read the logs:

\begin{codebox}
zpwr loginfo         info-level messages\\
zpwr logdebug        debug-level messages\\
zpwr logerror        error-level messages\\
zpwr logtrace        trace-level messages\\
zpwr taillog         live tail with colorized output\\
\end{codebox}


\vspace{12pt}
{\bfseries\color{neonyellow} \$\textbf{ZPWR\_COMPSYS}\_CACHE — COMPLETION CACHE}\vspace{4pt}

\begin{codebox}
echo ZPWR\_COMPSYS\_CACHE\\
\end{codebox}

Controls the completion system cache. Zsh completions are expensive to regenerate. This flag keeps them cached.


\vspace{12pt}
{\bfseries\color{neonyellow} \$\textbf{ZPWR\_HISTFILE} — HISTORY FILE}\vspace{4pt}

\begin{codebox}
echo ZPWR\_HISTFILE\\
\end{codebox}

Your shell history file location. \textbf{ZPWR} manages history with deduplication and backup. Your commands are sacred.

Related: zpwr backuphistory and zpwr restorehistory.

THE FULL ENV VAR MAP: Core paths:      \textbf{ZPWR}, \textbf{ZPWR\_LOCAL}, \textbf{ZPWR\_SCRIPTS}, \textbf{ZPWR\_TMUX} Autoload:        \textbf{ZPWR\_AUTOLOAD}, \textbf{ZPWR\_AUTOLOAD}\_COMMON Flags:           \textbf{ZPWR\_PROFILING}, \textbf{ZPWR\_DEBUG}, \textbf{ZPWR\_TRACE} Display:         \textbf{ZPWR\_EXA}\_COMMAND, \textbf{ZPWR\_KEYTIMEOUT} Runtime:         \textbf{ZPWR\_TTY}, \textbf{ZPWR\_OS}\_TYPE, \textbf{ZPWR\_LOGFILE}

{\small\itshape\color{dimtext} // the machine knows itself. every variable is a mirror.}

{\small\itshape\color{dimtext} // read them to understand the current state of reality.}


\vspace{16pt}


\chapter{CREATIVE TOOLS}


\begin{center}
\includegraphics[width=0.85\textwidth]{/Users/wizard/.zpwr/docs/screenshots/figlet_zpwr.png}
\end{center}

\section{Media Creation}


\vspace{10pt}
{\large\bfseries\color{neoncyan} MEDIA CREATION — ART FROM THE COMMAND LINE}\vspace{6pt}

Who needs photoshop when you have a terminal?

The terminal is not just for text. \textbf{ZPWR} includes verbs for generating GIFs, rendering ASCII art, and creating visual content without ever opening a GUI application.


\vspace{12pt}
{\bfseries\color{neonyellow} CREATEGIF — SCREEN TO GIF}\vspace{4pt}

\begin{codebox}
zpwr creategif\\
\end{codebox}

Creates animated GIFs from your terminal sessions. Record your workflow, your demos, your flex. Perfect for README demos and pull request previews.

Your terminal is the camera. The GIF is the film. No screen recording software needed.


\vspace{12pt}
{\bfseries\color{neonyellow} FIGLETFONTS — THE ASCII ART GALLERY}\vspace{4pt}

\begin{codebox}
zpwr figletfonts\\
\end{codebox}

Previews all available figlet fonts on your system. figlet turns text into massive ASCII art banners. This verb shows you every font so you can pick one.

Ever wondered what your name looks like in 47 different ASCII art styles? Now you know.

Example output (one of many fonts): \_\_\_\_\_\_\_\_  \_\_\_\_      \_\_\_\_  \_\_\_\_  \_\_\_\_ |\_\_\_  \_\_\_||  \_ \textbackslash{}    |  \_ \textbackslash{}|  \_ \textbackslash{}|  \_ \textbackslash{} / /    | |\_) |   | |\_) | |\_) | |\_) | / /     |  \_\_/    |  \_\_/|  \_ <|  \_ < / /      | |       | |   | | \textbackslash{} | | \textbackslash{} \textbackslash{} /\_/       |\_|       |\_|   |\_|  \textbackslash{}|\_|  \textbackslash{}\_\textbackslash{}


\vspace{12pt}
{\bfseries\color{neonyellow} ANIMATE — TERMINAL ANIMATION}\vspace{4pt}

\begin{codebox}
zpwr animate\\
\end{codebox}

Runs terminal animations. Because sometimes the command line needs to look as cool as it feels. Eye candy for the operator who has everything.

THE CREATIVE TOOLKIT: creategif       record terminal to GIF figletfonts     preview all ASCII art fonts animate         terminal animations banner          the \textbf{ZPWR} startup banner bannerlolcat    rainbow-colored banner.

{\small\itshape\color{dimtext} // the terminal is a canvas. most people see a spreadsheet.}

{\small\itshape\color{dimtext} // hackers see a gallery. these verbs are your brushes.}


\vspace{16pt}

\section{Color Systems}


\vspace{10pt}
{\large\bfseries\color{neoncyan} COLOR SYSTEMS — PAINTING THE VOID}\vspace{6pt}

256 colors. true color. the terminal is not monochrome.

Color is not cosmetic. It is information density. A colored diff tells you more in one glance than a monochrome one does in ten seconds. \textbf{ZPWR} gives you tools to explore, test, and master terminal color.


\vspace{12pt}
{\bfseries\color{neonyellow} COLORTEST — THE 16-COLOR PALETTE}\vspace{4pt}

\begin{codebox}
zpwr colortest\\
\end{codebox}

Displays the basic 16-color ANSI palette. The foundation. Black, red, green, yellow, blue, magenta, cyan, white — plus their bright variants. 16 soldiers in formation.

Use this to verify your terminal theme is sane. If red looks orange, your theme needs fixing.


\vspace{12pt}
{\bfseries\color{neonyellow} COLORTEST256 — THE FULL SPECTRUM}\vspace{4pt}

\begin{codebox}
zpwr colortest256\\
\end{codebox}

Displays all 256 extended colors. A grid of every color code your terminal supports. From 0 to 255. Useful for picking colors for prompts and scripts.

"What number is that teal?" Run this and find out.


\vspace{12pt}
{\bfseries\color{neonyellow} PYGMENTCOLORS — SYNTAX HIGHLIGHTING PALETTE}\vspace{4pt}

\begin{codebox}
zpwr pygmentcolors\\
\end{codebox}

Shows available Pygments color schemes. Pygments is the syntax highlighting engine used by many \textbf{ZPWR} tools. Pick a scheme that matches your aesthetic.


\vspace{12pt}
{\bfseries\color{neonyellow} PRETTYPRINT \& ALTPRETTYPRINT — FORMATTED OUTPUT}\vspace{4pt}

\begin{codebox}
zpwr prettyprint <file>\\
zpwr altprettyprint <file>\\
\end{codebox}

Syntax-highlighted file display. Prettyprint uses one engine, altprettyprint uses another. Both turn raw source code into technicolor glory in your terminal.

Like cat, but for people who respect their eyes.

COLOR EXPLORATION WORKFLOW: 1. zpwr colortest        check basic 16 colors 2. zpwr colortest256     explore full 256 palette 3. zpwr pygmentcolors    pick a syntax theme 4. zpwr prettyprint      view files in color.

{\small\itshape\color{dimtext} // monochrome terminals are a war crime}

{\small\itshape\color{dimtext} // color is data. learn to read it.}

{\small\itshape\color{dimtext} // 256 colors and you are still using white-on-black?}


\vspace{16pt}

\section{Script Management}


\vspace{10pt}
{\large\bfseries\color{neoncyan} SCRIPT MANAGEMENT — THE FORGE}\vspace{6pt}

Create scripts. Count scripts. Print scripts to PDF. Ship scripts.

\textbf{ZPWR} manages its own script collection and gives you tools to create new ones, inventory them, and even export them to PDF for documentation or review.


\vspace{12pt}
{\bfseries\color{neonyellow} SCRIPTNEW — THE SCRIPT GENERATOR}\vspace{4pt}

\begin{codebox}
zpwr scriptnew <name>\\
\end{codebox}

Creates a new script with proper boilerplate. Shebang, error handling, argument parsing — all pre-wired. Stop writing boilerplate. Start writing logic.

The script lands in \$\textbf{ZPWR\_SCRIPTS}, already executable, already on your \textbf{PATH}. Zero friction from idea to tool.


\vspace{12pt}
{\bfseries\color{neonyellow} SCRIPTTOPDF — THE DOCUMENTER}\vspace{4pt}

\begin{codebox}
zpwr scripttopdf\\
\end{codebox}

Converts scripts to syntax-highlighted PDFs. For code reviews, documentation, or printing out and framing on your wall like the art it is.

Send your code to a printer. Assert dominance.


\vspace{12pt}
{\bfseries\color{neonyellow} SCRIPTCOUNT — THE CENSUS}\vspace{4pt}

\begin{codebox}
zpwr scriptcount\\
\end{codebox}

How many scripts live in \$\textbf{ZPWR\_SCRIPTS}? This tells you. Track the growth of your personal armory over time.


\vspace{12pt}
{\bfseries\color{neonyellow} SCRIPTLIST — THE INVENTORY}\vspace{4pt}

\begin{codebox}
zpwr scriptlist\\
\end{codebox}

Lists every script in the collection. Names, paths, the full manifest of your executable arsenal.

THE SCRIPT LIFECYCLE: 1. zpwr scriptnew myutil     create the script 2. zpwr vimscriptedit        edit scripts with vim/fzf 3. zpwr scriptcount          check your inventory 4. zpwr scriptlist           browse the collection 5. zpwr scripttopdf          export for documentation.


\vspace{12pt}
{\bfseries\color{neonyellow} ALSO RELATED:}\vspace{4pt}

\begin{codebox}
zpwr makedir <name>     create a directory\\
zpwr makefile <name>    create a file\\
\end{codebox}

Simple creation verbs. No boilerplate, just mkdir/touch with \textbf{ZPWR}'s logging and validation baked in.

{\small\itshape\color{dimtext} // a script unwritten is a problem unsolved}

{\small\itshape\color{dimtext} // a script uncounted is a tool forgotten}

{\small\itshape\color{dimtext} // the forge never closes. the armory always grows.}


\vspace{16pt}


\chapter{EXERCISM \& EXTERNAL}

\section{Exercism Integration}


\vspace{10pt}
{\large\bfseries\color{neoncyan} EXERCISM INTEGRATION — LEVEL UP YOUR CODE}\vspace{6pt}

Practice makes perfect. \textbf{ZPWR} makes practice frictionless.

Exercism.org is a free coding practice platform with mentor-reviewed exercises in 50+ languages. \textbf{ZPWR} wraps the exercism CLI so you never leave your terminal.


\vspace{12pt}
{\bfseries\color{neonyellow} EXERCISM — THE LAUNCHER}\vspace{4pt}

\begin{codebox}
zpwr exercism\\
\end{codebox}

The base exercism command integration. Manages your exercism workspace, configuration, and submission workflow from within \textbf{ZPWR}.

Your exercism workspace lives under \$\textbf{ZPWR} and is accessible via zpwr homeexercism for quick navigation.


\vspace{12pt}
{\bfseries\color{neonyellow} EXERCISMDOWNLOAD — FETCH EXERCISES}\vspace{4pt}

\begin{codebox}
zpwr exercismdownload\\
\end{codebox}

Downloads an exercise to your local workspace. The exercise appears, ready to solve. Tests included. No browser clicking. No manual downloads.

The workflow: 1. Browse exercises on exercism.org 2. zpwr exercismdownload to fetch it 3. Solve the exercise in your editor 4. Run the tests locally 5. Submit when you pass.


\vspace{12pt}
{\bfseries\color{neonyellow} EXERCISMEDIT — SOLVE IN YOUR \textbf{EDITOR}}\vspace{4pt}

\begin{codebox}
zpwr exercismedit\\
\end{codebox}

Opens exercism exercises directly in your configured editor. Download, edit, test, submit — all from the terminal. The browser is optional.

WHY EXERCISM + \textbf{ZPWR}:

The exercism CLI is good. \textbf{ZPWR} makes it seamless. No context switching. No browser tabs. No friction.

Languages supported by exercism: Python, Go, Rust, Elixir, JavaScript, TypeScript, C, C++, Java, Ruby, Haskell, Clojure, Kotlin, Swift, Bash, and 35+ more.

{\small\itshape\color{dimtext} // coding practice should be as easy as}

{\small\itshape\color{dimtext} // opening a terminal and typing three words}

{\small\itshape\color{dimtext} // zpwr exercismdownload. done. now solve.}


\vspace{16pt}

\section{Travis CI \& External Services}


\vspace{10pt}
{\large\bfseries\color{neoncyan} TRAVIS CI \& EXTERNAL SERVICES}\vspace{6pt}

CI/CD from your terminal, not your browser.

\textbf{ZPWR} integrates with external services so you can monitor builds, check statuses, and manage CI/CD pipelines without opening a single browser tab.


\vspace{12pt}
{\bfseries\color{neonyellow} TRAVIS — THE CI DASHBOARD}\vspace{4pt}

\begin{codebox}
zpwr travis\\
\end{codebox}

Travis CI integration. View build status, logs, and pipeline information from your terminal.


\vspace{12pt}
{\bfseries\color{neonyellow} TRAVISBUILD — BUILD STATUS}\vspace{4pt}

\begin{codebox}
zpwr travisbuild\\
\end{codebox}

Check the current build status. Green, red, or pending. No need to refresh a web page in a loop.


\vspace{12pt}
{\bfseries\color{neonyellow} TRAVISBRANCH — BRANCH BUILDS}\vspace{4pt}

\begin{codebox}
zpwr travisbranch\\
\end{codebox}

View build status for a specific branch. When you push to a feature branch and want to know if CI passed.


\vspace{12pt}
{\bfseries\color{neonyellow} TRAVISPR — PULL REQUEST BUILDS}\vspace{4pt}

\begin{codebox}
zpwr travispr\\
\end{codebox}

Check build status for pull requests. The merge button is only as good as the CI behind it.


\vspace{12pt}
{\bfseries\color{neonyellow} OTHER EXTERNAL INTEGRATIONS:}\vspace{4pt}

\begin{codebox}
zpwr goclean          clean Go build cache\\
zpwr google <query>    search Google from terminal\\
zpwr digs <domain>     DNS lookups\\
zpwr torip            check your Tor IP address\\
zpwr toriprenew       renew your Tor circuit\\
zpwr curl             enhanced curl operations\\
\end{codebox}

THE EXTERNAL SERVICES PHILOSOPHY: The browser is a context switch. Every tab you open is a distraction vector. \textbf{ZPWR} brings the services to you, in the terminal, where you already are.

{\small\itshape\color{dimtext} // the terminal is the cockpit}

{\small\itshape\color{dimtext} // external services are instruments on the dashboard}

{\small\itshape\color{dimtext} // you read them. you do not visit them.}


\vspace{16pt}


\chapter{TIPS, TRICKS \& EASTER EGGS}


\begin{center}
\includegraphics[width=0.85\textwidth]{/Users/wizard/.zpwr/docs/screenshots/art_hack_the_planet.png}
\end{center}

\section{The One-Letter Verbs}


\vspace{10pt}
{\large\bfseries\color{neoncyan} THE ONE-LETTER VERBS — MAXIMUM POWER, MINIMUM KEYSTROKES}\vspace{6pt}

Six letters that do the work of six paragraphs.

\textbf{ZPWR} has \$\{\#\textbf{ZPWR\_VERBS}\} verbs. Six of them are single characters. These are not abbreviations — they are the distilled essence of the most common operations. One keystroke each. The speed limit of human-shell interaction.


\vspace{12pt}
{\bfseries\color{neonyellow} a — AUTOLOAD}\vspace{4pt}

\begin{codebox}
zpwr a\\
\end{codebox}

Quick access to the autoload system. When you need to check on autoloaded functions fast.


\vspace{12pt}
{\bfseries\color{neonyellow} c — CAT}\vspace{4pt}

\begin{codebox}
zpwr c\\
\end{codebox}

The cat operation. View file contents. The most basic read operation, compressed to one letter.


\vspace{12pt}
{\bfseries\color{neonyellow} e — EDIT}\vspace{4pt}

\begin{codebox}
zpwr e\\
\end{codebox}

Open your editor. The gateway to modification. Uses your configured \$\textbf{EDITOR}.


\vspace{12pt}
{\bfseries\color{neonyellow} f — FILESEARCH}\vspace{4pt}

\begin{codebox}
zpwr f\\
\end{codebox}

File search. Find things. The single most common operation after listing and editing.


\vspace{12pt}
{\bfseries\color{neonyellow} r — REPLAY}\vspace{4pt}

\begin{codebox}
zpwr r\\
\end{codebox}

Replay operations. Redo what you did before. History is not just for reading.


\vspace{12pt}
{\bfseries\color{neonyellow} z — ZCD}\vspace{4pt}

\begin{codebox}
zpwr z\\
\end{codebox}

Z-jump. Frecency-based directory navigation. Type a partial path, z figures out the rest. The single most addictive feature in \textbf{ZPWR}.

THE ONE-LETTER ARSENAL: a   autoload system c   cat / view files e   edit files f   file search r   replay z   z-jump directory navigation.

{\small\itshape\color{dimtext} // every extra character is latency between thought and action}

{\small\itshape\color{dimtext} // one letter is the event horizon of efficiency}

{\small\itshape\color{dimtext} // below this, there is only telepathy}


\vspace{16pt}

\section{Alias Chains}


\vspace{10pt}
{\large\bfseries\color{neoncyan} ALIAS CHAINS — COMPOUND OPERATIONS}\vspace{6pt}

One verb triggers another. The cascade begins.

\textbf{ZPWR} verbs do not exist in isolation. They chain. They compose. They call each other. This is not an accident — it is the architecture.


\vspace{12pt}
{\bfseries\color{neonyellow} THE SHORTHAND ALIASES:}\vspace{4pt}

\textbf{ZPWR} is aliased to shorter forms for speed:

\begin{codebox}
zpwr   the full command\\
zp     shorter\\
zpz    alternate short form\\
\end{codebox}

All three are equivalent:

\begin{codebox}
zpwr help  ==  zp help  ==  zpz help\\
\end{codebox}


\vspace{12pt}
{\bfseries\color{neonyellow} COMPOUND VERB PATTERNS:}\vspace{4pt}

Many verbs are compounds that do two things:

Filesearch + edit     = zpwr filesearchedit findsearch + edit     = zpwr findsearchedit locatesearch + edit   = zpwr locatesearchedit wordsearch + edit     = zpwr wordsearchedit vimrecent + cd        = zpwr vimrecentcd vimrecent + sudo      = zpwr vimrecentsudo.

The pattern: action + modifier = compound verb Search, then edit. Browse, then cd. Open, then sudo.


\vspace{12pt}
{\bfseries\color{neonyellow} THE CLEAN CHAIN:}\vspace{4pt}

\begin{codebox}
zpwr cleanall             clean everything\\
zpwr cleancache           clean caches only\\
zpwr cleancompcache       clean completion cache\\
zpwr cleangitcache        clean git repo cache\\
zpwr cleangitcleancache   clean git clean cache\\
zpwr cleangitdirtycache   clean git dirty cache\\
zpwr cleantemp            clean temp files\\
zpwr cleanrefreshupdate   clean + refresh + update\\
\end{codebox}


\vspace{12pt}
{\bfseries\color{neonyellow} THE GIT CHAIN:}\vspace{4pt}

\begin{codebox}
zpwr gitforalldir     run git cmd in all dirs\\
zpwr gitfordirmain    run git cmd in dirs on main\\
zpwr gitzfordir       z-style git for all dirs\\
\end{codebox}

{\small\itshape\color{dimtext} // verbs are words. chains are sentences.}

{\small\itshape\color{dimtext} // the grammar of \textbf{ZPWR} is combinatorial.}

{\small\itshape\color{dimtext} // learn the words, and the sentences write themselves.}


\vspace{16pt}

\section{The Token System}


\vspace{10pt}
{\large\bfseries\color{neoncyan} THE TOKEN SYSTEM — YOUR PRIVATE OVERRIDE LAYER}\vspace{6pt}

Secrets, aliases, and local config that never touch git.

\textbf{ZPWR} has sensible defaults. But you are not sensible. You have API keys, custom aliases, machine-specific paths, and opinions. The token system is where all of that lives — safely outside version control.


\vspace{12pt}
{\bfseries\color{neonyellow} THE TOKEN FILES:}\vspace{4pt}

\textbf{\$ZPWR\_LOCAL/.tokens.sh} Your main tokens file. Aliases, exports, secrets. This is sourced during shell init. Gitignored.

\textbf{\$ZPWR\_LOCAL/.tokens-pre.sh} Sourced BEFORE \textbf{ZPWR}'s main configuration. Override defaults before they are set. Set \textbf{ZPWR}\_* variables here to change behavior.

\textbf{\$ZPWR\_LOCAL/.tokens-post.sh} Sourced AFTER \textbf{ZPWR}'s main configuration. Override anything \textbf{ZPWR} set. Final word. Your aliases here trump everything.


\vspace{12pt}
{\bfseries\color{neonyellow} THE OVERRIDE PATTERN:}\vspace{4pt}

Want to change a \textbf{ZPWR} default? Do not edit \textbf{ZPWR} source. Put it in tokens-pre.sh (before) or tokens-post.sh (after).

Example: override in tokens-pre.sh.

\begin{codebox}
export ZPWR\_EXA\_COMMAND=eza\\
export ZPWR\_KEYTIMEOUT=1\\
export ZPWR\_BANNER\_CLEARLIST=false\\
\end{codebox}

Example: add custom aliases in tokens-post.sh.

\begin{codebox}
alias k=kubectl\\
alias tf=terraform\\
export GITHUB\_TOKEN=ghp\_xxxxx\\
\end{codebox}


\vspace{12pt}
{\bfseries\color{neonyellow} MANAGING TOKENS:}\vspace{4pt}

\begin{codebox}
zpwr tokens           view your tokens files\\
zpwr vimtokens        edit tokens in vim\\
zpwr emacstokens      edit tokens in emacs\\
\end{codebox}

THE LOAD ORDER: 1. Tokens-pre.sh      (your early overrides) 2. \textbf{ZPWR} core config    (the framework defaults) 3. Tokens-post.sh     (your final overrides) 4. .tokens.sh         (your secrets and aliases)

NEVER commit token files. They are gitignored for a reason. API keys in git history haunt you forever.

{\small\itshape\color{dimtext} // the token system is your private dimension}

{\small\itshape\color{dimtext} // \textbf{ZPWR} updates without touching your config}

{\small\itshape\color{dimtext} // your overrides survive every git pull}


\vspace{16pt}

\section{FZF Everywhere}


\vspace{10pt}
{\large\bfseries\color{neoncyan} FZF EVERYWHERE — THE FUZZY FINDER CONSPIRACY}\vspace{6pt}

fzf is not a tool. It is a religion. \textbf{ZPWR} is its cathedral.

fzf (fuzzy finder) is woven into the fabric of \textbf{ZPWR}. Dozens of verbs pipe their output through fzf for interactive selection. If a verb ends in 'fzf', it is the fuzzy-interactive variant. But many others use it too.


\vspace{12pt}
{\bfseries\color{neonyellow} DEDICATED FZF VERBS:}\vspace{4pt}

\begin{codebox}
zpwr verbs     browse subcommands in fzf\\
zpwr verbs           browse all verbs in fzf\\
\end{codebox}


\vspace{12pt}
{\bfseries\color{neonyellow} VERBS THAT USE FZF INTERNALLY:}\vspace{4pt}

Forgit* verbs: forgitadd, forgitdiff, forgitlog, forgitstash, forgitreset, forgitrestore, forgitclean All use fzf for interactive selection with previews.

Search + edit verbs: filesearchedit, findsearchedit, locatesearchedit, wordsearchedit — search with fzf, open selection.

Recent file verbs: vimrecent, vimrecentcd, emacsrecent, emacsrecentcd Browse recent files with fzf preview.

Editor verbs: vimalledit, emacsalledit — select files to edit.

Git tag verbs: vimgtagsedit, emacsgtagsedit — browse tags in fzf.


\vspace{12pt}
{\bfseries\color{neonyellow} THE FZF PATTERN:}\vspace{4pt}

The pattern is always the same: 1. Generate a list of candidates 2. Pipe through fzf with a preview command 3. Take the selection and act on it.

This pattern turns every list into a menu. Every menu into a decision. Every decision into action.


\vspace{12pt}
{\bfseries\color{neonyellow} FZF CUSTOMIZATION:}\vspace{4pt}

Set FZF\_DEFAULT\_OPTS in your tokens file:

\begin{codebox}
export FZF\_DEFAULT\_OPTS='--height 80\% --reverse --border'\\
\end{codebox}

{\small\itshape\color{dimtext} // if you can list it, you can fzf it}

{\small\itshape\color{dimtext} // if you can fzf it, you can act on it}

{\small\itshape\color{dimtext} // fzf is the universal adapter between data and decisions}


\vspace{16pt}


\chapter{THE COMPLETE REFERENCE}


\begin{center}
\includegraphics[width=0.85\textwidth]{/Users/wizard/.zpwr/docs/screenshots/figlet_docs.png}
\end{center}

\section{All Verbs A-G}


\vspace{10pt}
{\large\bfseries\color{neoncyan} ALL \$\{\#ZPWR\_VERBS\} VERBS: A through G (1/3)}\vspace{6pt}

The complete arsenal. Memorize at your own peril.

A                          quick autoload access about                      display \textbf{ZPWR} about info alacritty                  alacritty terminal config aliasrank                  rank aliases by usage frequency all                        show all \textbf{ZPWR} info allsearch                  search across all sources altprettyprint             alternate syntax-highlighted print animate                    terminal animations attach                     attach to tmux session autoload                   show autoload system autoloadcount              count autoloaded functions autoloadlist               list all autoloaded functions background                 run command in background backup                     backup \textbf{ZPWR} configuration backuphistory              backup shell history banner                     display \textbf{ZPWR} banner bannercounts               banner with counts bannerlolcat               rainbow banner via lolcat bannernopony               banner without ponysay bannerpony                 banner with ponysay bench                      benchmark shell startup brc                        source bashrc c                          cat / view file cat                        enhanced cat catcd                      cat file and cd to its dir catnviminforecentf         cat nvim info recent files catrecentfviminfo          cat recent files from viminfo catviminfo                 cat viminfo file cd                         enhanced cd cdup                       cd up one directory cfasd                      cd via fasd frecency changename                 rename files interactively changenamezpwr             rename \textbf{ZPWR} components cleanall                   clean everything cleancache                 clean caches cleancompcache             clean completion cache color2                     turn on stderr color filter environmentvariables       list zpwr env vars (alias: envvars) environmentvariablesall    list all env vars (alias: envvarsall) environmentvars            list zpwr env vars (alias: envvars) environmentvarsall         list all env vars (alias: envvarsall)

... continued. Run zpwr verbs for the complete list.


\vspace{16pt}

\section{All Verbs H-S}


\vspace{10pt}
{\large\bfseries\color{neoncyan} ALL \$\{\#ZPWR\_VERBS\} VERBS: forgitdiff through recompilefpath (2/3)}\vspace{6pt}

The middle of the alphabet is where the real work lives.

Forgitdiff                 interactive diff via fzf forgitignore               browse gitignore templates forgitignoreclean          clean gitignore cache forgitignoreget            fetch gitignore template forgitignorelist           list gitignore templates forgitignoreupdate         update gitignore cache forgitinfo                 repo info dashboard forgitlog                  interactive git log via fzf forgitreset                interactive unstage via fzf forgitrestore              interactive restore via fzf forgitstash                interactive stash via fzf forgitwarn                 warn about repo hazards forked                     manage forked repos forward                    port forwarding fp                         fpath info funcrank                   rank functions by usage fwd                        forward shortcut gh                         GitHub CLI integration ghcontribcount             GitHub contribution count ghz                        GitHub + z integration gitaemail                  set git author email gitcemail                  set git committer email gitclearcache              clear git cache gitclearcommit             clear a git commit gitclearfile               clear file from git gitcommit                  enhanced git commit gitcommitcount             count commits gitcommits                 browse commits gitconfig                  show git config gitcontribcount            count contributions gitcontribcountdirs        count contribs per dir gitcontribcountlines       count contributed lines gitdir                     show git directory gitedittag                 edit a git tag gitemail                   show git email gitforalldir               git cmd in all dirs hometests                  cd to test directory (alias: hometest) loadjenv                   lazy load jenv logincount                 count logins by user regenpdf                   regenerate encyclopedia PDF subcommandsedit            edit subcommand selection subcommandslist             list all subcommands.

... continued. Run zpwr verbs for the complete list.


\vspace{16pt}

\section{All Verbs T-Z}


\vspace{10pt}
{\large\bfseries\color{neoncyan} ALL \$\{\#ZPWR\_VERBS\} VERBS: refreshzwc through zstyle (3/3)}\vspace{6pt}

The final stretch. The end of the alphabet. The beginning of mastery.

Refreshzwc                 refresh compiled wordcode regen                      regenerate single component regenall                   regenerate everything regenconfiglinks           regenerate config symlinks regenctags                 regenerate ctags regenenvcache              regenerate env cache regengitdirtyrepocache     regen dirty repo cache regengitrepocache          regen git repo cache regengtags                 regenerate gtags regengtagspygments         regen gtags + pygments regengtagstype             regen gtags by type regenhistory               regenerate history regenkeybindings           regen keybindings regenpowerline             regenerate powerline regenzsh                   regenerate zsh config reload                     reload shell config relpath                    get relative path rename                     rename files replacer                   find and replace replay                     replay commands repo                       repository info reset                      reset shell state restart                    restart shell restore                    restore configuration restorehistory             restore history backup return2                    return to previous dir reveal                     open repo in browser revealrecurse              reveal recursively reverse                    reverse input rmconfiglinks              remove config symlinks rvs                        reverse shortcut scriptcount                count scripts scriptlist                 list all scripts scriptnew                  create new script scripts                    browse scripts scripttopdf                convert script to PDF testsall                   run all env tests (alias: testall) upload                     upload file with curl verbsedit                  edit verb selection zpwrCloneToForked          clone zpwr to forked dir zpwrgithub                 open zpwr GitHub page.

Run zpwr verbs for the complete live list.


\vspace{16pt}


\chapter{POWERLEVEL10K \& INSTANT PROMPT}


\begin{center}
\includegraphics[width=0.85\textwidth]{/Users/wizard/.zpwr/docs/screenshots/art_neural_link.png}
\end{center}

\section{Instant Prompt: Zero Latency}


\vspace{10pt}
{\large\bfseries\color{neoncyan} INSTANT PROMPT: ZERO LATENCY}\vspace{6pt}

{\small\itshape\color{dimtext} // your prompt loads before .zshrc finishes. time is a flat circle. //}

The single best UX trick in zpwr: the prompt appears INSTANTLY while .zshrc is still loading behind the curtain. Zero perceived latency. The illusion of speed IS speed.


\vspace{12pt}
{\bfseries\color{neonyellow} HOW IT WORKS}\vspace{4pt}

On previous shell exit, p10k cached the prompt to a file:

\begin{codebox}
\{XDG\_CACHE\_HOME:-HOME/.cache\}/p10k-instant-prompt-\{(\%):-\%n\}.zsh\\
\end{codebox}

This file contains the last rendered prompt as raw ANSI escape sequences. On next shell start, it is sourced BEFORE .zshrc loads. Result: the prompt appears in <10ms. The user can start typing. .zshrc continues loading in the background. Nobody notices.


\vspace{12pt}
{\bfseries\color{neonyellow} THE LOADING SEQUENCE IN .zshrc}\vspace{4pt}

1. zmodload zsh/datetime (timestamp) 2. Source zpwr-hash-dirs.zsh (cached hash -d entries) 3. Source p10k-instant-prompt-*.zsh (the magic file) 4. promptTimestamp=\$EPOCHREALTIME (measure prompt time) 5. ...rest of .zshrc loads (plugins, aliases, etc)... 6. p10k finalize (at very end of .zshrc) Step 3 must come before ANY console output. If something prints to stdout before the instant prompt, p10k aborts it.


\vspace{12pt}
{\bfseries\color{neonyellow} THE HASH DIR CACHE}\vspace{4pt}

The prompt uses \%\textasciitilde{} to show named dirs like \textasciitilde{}\textbf{ZPWR} instead of /Users/you/.zpwr. But hash -d entries are set in zpwrBindDirs, which runs AFTER instant prompt. Chicken-and-egg problem. Solution: zpwrBindDirs writes a cache file:

\begin{codebox}
ZPWR\_LOCAL/zpwr-hash-dirs.zsh\\
\end{codebox}

This cache is sourced at line 73 of .zshrc, BEFORE instant prompt, so \%\textasciitilde{} resolves \textasciitilde{}\textbf{ZPWR} in the cached prompt correctly.


\vspace{12pt}
{\bfseries\color{neonyellow} \textbf{ZPWR\_BANNER}\_CLEARLIST SUPPRESSION}\vspace{4pt}

ZpwrClearList prints a banner + file listing on login. But if it runs before instant prompt finalizes, it corrupts the display. Fix: the banner is deferred until AFTER p10k finalize:

\begin{codebox}
if [[ -n "\_\_p9k\_instant\_prompt\_active" \&\& "ZPWR\_BANNER\_CLEARLIST" == true ]]; then\\
zpwrClearList || true\\
fi\\
\end{codebox}


\vspace{12pt}
{\bfseries\color{neonyellow} THE promptTimestamp MEASUREMENT}\vspace{4pt}

\textbf{ZPWR\_VARS}[promptTimeMs] records how fast the prompt appeared. Run zpwr environmentcounts to see it. With instant prompt, this is typically 0.010-0.050 seconds. Without it, 0.5-2.0s.


\vspace{12pt}
{\bfseries\color{neonyellow} p10k finalize}\vspace{4pt}

The last line of .zshrc:

\begin{codebox}
(( ! \{+functions[p10k]\} )) || p10k finalize\\
\end{codebox}

This tells p10k that .zshrc is done. The real prompt replaces the cached one. Any deferred output can now safely print.


\vspace{16pt}

\section{Powerlevel10k Configuration}


\vspace{10pt}
{\large\bfseries\color{neoncyan} POWERLEVEL10K CONFIGURATION}\vspace{6pt}

{\small\itshape\color{dimtext} // the prompt is your cockpit HUD. configure it or fly blind. //}

The p10k config lives at \textbf{\$ZPWR\_ENV/.p10k.zsh} -- a 1700+ line file that controls every pixel of your prompt. The wizard generated it, but you own it now.


\vspace{12pt}
{\bfseries\color{neonyellow} LEFT \textbf{PROMPT} SEGMENTS}\vspace{4pt}

These appear on the left side, top to bottom across 4 lines:

\begin{codebox}
Line 1: custom\_env, custom\_env\_bat, custom\_env\_editor\\
Line 2: custom\_env\_term, ip\\
Line 3: context, dir, history, ssh, root\_indicator, status, vcs\\
Line 4: custom\_tty, prompt\_char\\
\end{codebox}

The dir segment shows your current directory using \%\textasciitilde{} expansion. The vcs segment shows git branch, dirty status, ahead/behind. The prompt\_char shows green on success, red on failure.


\vspace{12pt}
{\bfseries\color{neonyellow} RIGHT \textbf{PROMPT} SEGMENTS}\vspace{4pt}

The right side is information-dense. Across 3 lines:

\begin{codebox}
Line 1: dir\_writable, background\_jobs, virtualenv, node\_version,\\
\end{codebox}

go\_version, rust\_version, java\_version, package, os\_icon...

\begin{codebox}
Line 2: load, disk\_usage, ram, swap, wifi\\
Line 3: battery, time, command\_execution\_time, vi\_mode, custom\_pid\\
\end{codebox}


\vspace{12pt}
{\bfseries\color{neonyellow} THE DIR SEGMENT AND \%\textasciitilde{} EXPANSION}\vspace{4pt}

Zsh's \%\textasciitilde{} prompt expansion replaces known paths with named dirs. If you hash -d \textbf{ZPWR}=/Users/you/.zpwr, then /Users/you/.zpwr shows as \textasciitilde{}\textbf{ZPWR} in your prompt. Shorter. Cleaner. Professional. POWERLEVEL9K\_DIR\_MAX\_LENGTH=80 caps the dir segment width. POWERLEVEL9K\_DIR\_MIN\_COMMAND\_COLUMNS=40 reserves typing space.


\vspace{12pt}
{\bfseries\color{neonyellow} HOW zpwrBindDirs POPULATES NAMED DIRECTORIES}\vspace{4pt}

The function runs hash -d for \textasciitilde{}20 directories:

\begin{codebox}
hash -d ZPWR="ZPWR"\\
hash -d ZPWR\_TMUX="ZPWR\_TMUX"\\
hash -d ZPWR\_SCRIPTS="ZPWR\_SCRIPTS"\\
hash -d TMUXH="TMUX\_HOME"\\
\end{codebox}

Then it writes them all to \textbf{\$ZPWR\_LOCAL/zpwr-hash-dirs.zsh} for instant prompt to use on next startup. The cache iterates \$\{(kv)nameddirs\} to capture every registered named directory.


\vspace{12pt}
{\bfseries\color{neonyellow} THE MULTILINE FRAME}\vspace{4pt}

The prompt uses box-drawing characters for the frame:

\begin{codebox}
First line:  \%244F-----  (top connector)\\
Middle:      \%244F-----  (mid connector)\\
Last line:   \%244F-----  (bottom connector)\\
\end{codebox}

Gap fill character is '---' in grey (color 244). Mode: nerdfont-complete. Separators use Powerline glyphs.


\vspace{12pt}
{\bfseries\color{neonyellow} prompt\_char SYMBOLS}\vspace{4pt}

Insert mode: ---  Command mode: ---  Visual: V  Overwrite: ---


\vspace{16pt}

\section{P10k Prompt Customization}


\vspace{10pt}
{\large\bfseries\color{neoncyan} P10K PROMPT CUSTOMIZATION}\vspace{6pt}

{\small\itshape\color{dimtext} // your prompt is a canvas. paint it with ANSI and Unicode. //}

\textbf{ZPWR} ships custom p10k segments beyond the defaults. These are defined with POWERLEVEL9K\_CUSTOM\_* variables in \textbf{\$ZPWR\_ENV/.p10k.zsh.}


\vspace{12pt}
{\bfseries\color{neonyellow} CUSTOM SEGMENTS}\vspace{4pt}

Custom\_tty       Shows \$\textbf{ZPWR\_TTY} (e.g. \textbf{\textbf{/dev/ttys003}}) Defined as:

\begin{codebox}
POWERLEVEL9K\_CUSTOM\_TTY='echo "\%K\{blue\}\%F\{white\} ZPWR\_TTY \%f\%k\%F\{blue\}\%f"'\\
\end{codebox}

Custom\_env       Environment info (custom shell, OS, etc.) custom\_env\_bat   Shows bat (syntax highlighter) version/status custom\_env\_editor Shows \$\textbf{EDITOR} custom\_env\_term  Shows \$\textbf{TERM}\_PROGRAM custom\_pid       Shows shell PID (right prompt, bottom line)


\vspace{12pt}
{\bfseries\color{neonyellow} ADDING YOUR OWN SEGMENT}\vspace{4pt}

1. Define the command that generates output:

\begin{codebox}
typeset -g POWERLEVEL9K\_CUSTOM\_MYWIDGET='echo "hello"'\\
\end{codebox}

2. Set colors:

\begin{codebox}
typeset -g POWERLEVEL9K\_CUSTOM\_MYWIDGET\_FOREGROUND=255\\
typeset -g POWERLEVEL9K\_CUSTOM\_MYWIDGET\_BACKGROUND=24\\
\end{codebox}

3. Add 'custom\_mywidget' to LEFT or RIGHT \textbf{PROMPT}\_ELEMENTS array. 4. Run p10k reload or source \textbf{\$ZPWR\_ENV/.p10k.zsh} to apply.


\vspace{12pt}
{\bfseries\color{neonyellow} COLORS AND ICONS}\vspace{4pt}

P10k uses 256-color indices. Print the color map:

\begin{codebox}
for i in \{0..255\}; do print -Pn "\%K\{i\}  \%k\%F\{i\}\{(l:3::0:)i\}\%f "; done\\
\end{codebox}

Foreground/background for each segment and state:

\begin{codebox}
POWERLEVEL9K\_DIR\_BACKGROUND=4        (blue)\\
POWERLEVEL9K\_DIR\_FOREGROUND=254      (near-white)\\
\end{codebox}

Anchor dirs (like .git roots) get bold + brighter color (255).


\vspace{12pt}
{\bfseries\color{neonyellow} TRUNCATION}\vspace{4pt}

POWERLEVEL9K\_DIR\_MAX\_LENGTH=80 caps directory width. POWERLEVEL9K\_DIR\_MIN\_COMMAND\_COLUMNS\_PCT=50 means at least half the terminal width is reserved for your command input. SHORTEN\_STRATEGY can be set to truncate\_to\_unique for smart shortening (commented out by default in zpwr).


\vspace{12pt}
{\bfseries\color{neonyellow} THE .p10k.zsh FILE STRUCTURE}\vspace{4pt}

The file is a single anonymous function () \{ ... \} that: - Unsets all POWERLEVEL9K\_* vars (clean slate) - Defines LEFT/RIGHT\_PROMPT\_ELEMENTS arrays - Sets mode, icons, colors, truncation per segment - Configures transient prompt, instant prompt, hot reload


\vspace{12pt}
{\bfseries\color{neonyellow} SEEING \textbf{PROMPT} TIMING}\vspace{4pt}

\begin{codebox}
zpwr environmentcounts\\
\end{codebox}

Shows promptTimeMs, startupTimeMs, and per-phase timestamps. This is how you prove your prompt loaded in 0.015 seconds.


\vspace{16pt}

\section{Transient Prompt \& Visual Tweaks}


\vspace{10pt}
{\large\bfseries\color{neoncyan} TRANSIENT PROMPT \& VISUAL TWEAKS}\vspace{6pt}

{\small\itshape\color{dimtext} // old prompts fade into ghosts. only the present matters. //}

Transient prompt collapses previous prompts to a minimal form, keeping your scrollback clean and focused. In zpwr it is currently set to 'off' but easy to enable.


\vspace{12pt}
{\bfseries\color{neonyellow} TRANSIENT \textbf{PROMPT}}\vspace{4pt}

When enabled, after you press Enter, the full multi-line prompt collapses to just the prompt\_char symbol (--- or ---). All segments, git status, directory -- gone from old lines. Only the current prompt shows full detail.

\begin{codebox}
typeset -g POWERLEVEL9K\_TRANSIENT\_PROMPT=off\\
\end{codebox}

Change to always or same-dir to enable. 'same-dir' only collapses when you stay in the same directory.


\vspace{12pt}
{\bfseries\color{neonyellow} HOT RELOAD}\vspace{4pt}

Edit .p10k.zsh and apply changes without restarting zsh:

\begin{codebox}
p10k reload\\
\end{codebox}

This re-sources .p10k.zsh and rebuilds the prompt in-place. No exec zsh needed. No losing your command history buffer.


\vspace{12pt}
{\bfseries\color{neonyellow} p10k configure}\vspace{4pt}

The interactive configuration wizard. Walks you through: - Font detection (powerline glyphs, nerdfont icons) - Style selection (rainbow, lean, classic, pure) - Separator shapes, line count, frame options - Instant prompt mode (verbose, quiet, off) It generates a fresh .p10k.zsh. \textbf{ZPWR} ships a pre-configured one in \textbf{\$ZPWR\_ENV/.p10k.zsh} (rainbow style, nerdfont-complete). Running the wizard overwrites it, so back up yours first.


\vspace{12pt}
{\bfseries\color{neonyellow} INSTANT \textbf{PROMPT} vs zpwrClearList}\vspace{4pt}

The banner (zpwrClearList) prints colorized file listings on login. Instant prompt redirects all stdout/stderr during init to a buffer. If zpwrClearList runs during this window, the output gets swallowed or garbled. \textbf{ZPWR}'s fix: 1. Check if \_\_p9k\_instant\_prompt\_active is set 2. If so, defer zpwrClearList to AFTER p10k finalize 3. The banner appears cleanly after the prompt settles.


\vspace{12pt}
{\bfseries\color{neonyellow} THE eza --classify=always FIX}\vspace{4pt}

Under instant prompt, fd/stdout gets redirected. eza's --classify flag (append */ to dirs) auto-detects terminal capabilities. When stdout is redirected, it thinks there's no terminal and drops the indicators. Fix: use --classify=always to force classification regardless of fd state. \textbf{ZPWR\_EXA}\_COMMAND uses this:

\begin{codebox}
command eza --git -il --classify=always -H --color-scale -g -a\\
\end{codebox}


\vspace{12pt}
{\bfseries\color{neonyellow} \$TTY IMMUNE TO FD REDIRECTS}\vspace{4pt}

Zsh sets \$TTY before rc files load. Instant prompt redirects fds 0/1/2, but \$TTY still points to the real terminal. \textbf{ZPWR} captures it: export \textbf{ZPWR\_TTY}=\$TTY This is used by the custom\_tty prompt segment.


\vspace{16pt}


\chapter{ZINIT \& PLUGIN SYSTEM}


\begin{center}

\includegraphics[width=0.85\textwidth]{/Users/wizard/.zpwr/docs/screenshots/art_system_core.png}
\end{center}

\section{Zinit: The Plugin Engine}


\vspace{10pt}
{\large\bfseries\color{neoncyan} ZINIT: THE PLUGIN ENGINE}\vspace{6pt}

{\small\itshape\color{dimtext} // the package manager for your shell. apt-get for zsh plugins. //}

Zinit (formerly zplugin) is a flexible zsh plugin manager with turbo mode, ice modifiers, and completion management. \textbf{ZPWR} uses it as the default: \textbf{ZPWR\_PLUGIN}\_MANAGER=zinit.


\vspace{12pt}
{\bfseries\color{neonyellow} WHAT ZINIT DOES}\vspace{4pt}

- Downloads GitHub plugins to \$\textbf{ZPWR\_PLUGIN}\_MANAGER\_HOME - Sources them into your shell (or loads as completions) - Manages load order with turbo mode (deferred loading) - Handles OMZ snippets (individual files from oh-my-zsh) - Compiles plugins to .zwc bytecode for speed


\vspace{12pt}
{\bfseries\color{neonyellow} ZINIT LOAD vs ZINIT SNIPPET}\vspace{4pt}

zinit load user/repo      Clone full GitHub repo, source .plugin.zsh zinit snippet OMZP::git    Download single file from OMZ plugin zinit snippet OMZL::git.zsh Download single OMZ lib file load = full repo clone. Snippet = single file download. \textbf{ZPWR} uses both: load for GitHub plugins, snippet for OMZ.


\vspace{12pt}
{\bfseries\color{neonyellow} THE PLUGIN ARRAYS}\vspace{4pt}

\textbf{ZPWR} organizes plugins into arrays in .zshrc: \textbf{ZPWR\_GH}\_PLUGINS   28+ GitHub repos (fzf-tab, zsh-z, forgit...) \textbf{ZPWR\_OMZ}\_PLUGINS  25+ OMZ snippets (git, python, golang...) \textbf{ZPWR\_OMZ}\_LIBS     4 OMZ lib files (git.zsh, directories.zsh...) \textbf{ZPWR\_OMZ}\_COMPS    12 OMZ completion files (ng, coffee, scala...) Conditional plugins are appended if the command exists:

\begin{codebox}
zpwrCommandExists docker \&\& ZPWR\_GH\_PLUGINS+=( zsh-docker-aliases )\\
\end{codebox}


\vspace{12pt}
{\bfseries\color{neonyellow} PLUGIN STORAGE}\vspace{4pt}

Everything lives under \$\textbf{ZPWR\_PLUGIN}\_MANAGER\_HOME:

\begin{codebox}
ZPWR\_PLUGIN\_MANAGER\_HOME/bin/         zinit itself\\
ZPWR\_PLUGIN\_MANAGER\_HOME/plugins/     cloned repos\\
ZPWR\_PLUGIN\_MANAGER\_HOME/snippets/    OMZ snippets\\
ZPWR\_PLUGIN\_MANAGER\_HOME/completions/  symlinked completions\\
\end{codebox}


\vspace{12pt}
{\bfseries\color{neonyellow} THE 48+ PLUGIN ROSTER}\vspace{4pt}

Total plugins managed by zpwr (from .zshplugins): MenkeTechnologies forks: fzf-tab, forgit, zsh-expand, zsh-autocomplete, zsh-z, fasd-simple, revolver, zunit... zdharma-continuum: fast-syntax-highlighting, zbrowse, zui... zsh-users: completions, autosuggestions, history-substring Plus: hlissner/zsh-autopair, bigH/git-fuzzy, and more.


\vspace{12pt}
{\bfseries\color{neonyellow} FIRST RUN BOOTSTRAP}\vspace{4pt}

If \$\textbf{ZPWR\_PLUGIN}\_MANAGER\_HOME/bin doesn't exist, .zshrc auto-clones zinit from GitHub. After that, OPTIMIZE\_OUT\_DISK\_ACCESSES=1 skips unnecessary stat() calls on subsequent starts.


\vspace{16pt}

\section{Turbo Mode: Deferred Loading}


\vspace{10pt}
{\large\bfseries\color{neoncyan} TURBO MODE: DEFERRED LOADING}\vspace{6pt}

{\small\itshape\color{dimtext} // why load 48 plugins at startup when you can load 1? //}

Turbo mode is zinit's killer feature: plugins load AFTER the prompt appears, using zsh's sched mechanism. Your shell is interactive instantly; plugins trickle in behind the scenes.


\vspace{12pt}
{\bfseries\color{neonyellow} THE wait ICE MODIFIER}\vspace{4pt}

The 'wait' ice defers loading until after the first prompt:

\begin{codebox}
zinit ice wait            Load after prompt (default priority)\\
zinit ice wait'0'         Same as wait (priority 0)\\
zinit ice wait'0a'        Priority 0, sub-order a (loads first)\\
zinit ice wait'0b'        Priority 0, sub-order b (loads second)\\
zinit ice wait'0e'        Priority 0, sub-order e (loads later)\\
\end{codebox}

The letter suffix controls load order WITHIN the same priority. A loads before b loads before c, etc. This is how \textbf{ZPWR} orchestrates 48 plugins into a precise loading sequence.


\vspace{12pt}
{\bfseries\color{neonyellow} THE GRADUATED LOADING SEQUENCE}\vspace{4pt}

Synchronous:  powerlevel10k (must load before prompt) wait:          OMZ completions, libs, plugins, GH plugins wait'0':       zsh-autopair (needs ZLE ready) wait'0a':      zsh-expand + late aliases + alias cache wait'0b':      zsh-autocomplete (if \textbf{ZPWR\_AUTO}\_COMPLETE=true) wait'0c':      zsh-history-substring-search wait'0d':      zsh-autosuggestions + FZF late bindings + verbs wait'0e':      zsh-zinit-final (penultimate, final keybindings) wait'0f':      zsh-more-completions (low priority, last resort) wait'\$\{\textbf{ZPWR\_ZINIT}\_COMPINIT\_DELAY\}g':  compinit + cdreplay wait'\$\{\textbf{ZPWR\_ZINIT}\_COMPINIT\_DELAY\}h':  fast-syntax-highlighting.


\vspace{12pt}
{\bfseries\color{neonyellow} WHY STARTUP IS FAST}\vspace{4pt}

Only ONE plugin loads synchronously: powerlevel10k. Everything else is turbo-deferred. Combined with instant prompt, you see the prompt in \textasciitilde{}15ms. The other 47 plugins load in the next \textasciitilde{}200ms while you're still thinking.


\vspace{12pt}
{\bfseries\color{neonyellow} atinit AND atload HOOKS}\vspace{4pt}

These run code at precise moments during plugin loading: atinit   Runs BEFORE the plugin is sourced atload   Runs AFTER the plugin is sourced \textbf{ZPWR} uses these extensively to bind keymaps and aliases:

\begin{codebox}
zinit ice wait'0a' atload'zpwrBindAliasesLate; zpwrCreateAliasCache'\\
\end{codebox}

The atload fires after zsh-expand is sourced, then late aliases override any OMZ defaults. Precise surgical timing.


\vspace{12pt}
{\bfseries\color{neonyellow} \textbf{ZPWR\_ZINIT}\_COMPINIT\_DELAY}\vspace{4pt}

Default: 0. Controls how long to wait before running compinit. Higher values defer completion system init even further. Set in \textbf{\$ZPWR\_ENV/.zpwr\_env.sh.} The 'g' and 'h' suffixes on the wait ice ensure compinit and syntax highlighting load last.


\vspace{16pt}

\section{Ice Modifiers Deep Dive}


\vspace{10pt}
{\large\bfseries\color{neoncyan} ICE MODIFIERS DEEP DIVE}\vspace{6pt}

{\small\itshape\color{dimtext} // ice modifiers are the spice of zinit. season generously. //}

Every zinit load/snippet can be preceded by 'zinit ice' with modifiers that control HOW the plugin loads. Ice melts after one use -- it applies only to the next command.


\vspace{12pt}
{\bfseries\color{neonyellow} CORE ICE MODIFIERS USED BY \textbf{ZPWR}}\vspace{4pt}

Lucid          Suppress the 'Loaded plugin' message \textbf{ZPWR} uses this on every plugin. No spam. Nocompile      Skip .zwc compilation of the plugin Useful when zpwrRecompileFiles handles it globally. Nocd           Don't cd to plugin dir before sourcing Prevents p10k internals from polluting named dirs. Nocompletions  Don't install completions from this plugin Used for zsh-more-completions (fpath-only approach). As'completion' Load the file as a completion, not a plugin Used for OMZ completion snippets like \_ng, \_coffee. As'program'    Add the pick'd file to \$\textbf{PATH} as executable Used for fzf binary and git-fuzzy. As'null'       Don't source anything -- just run hooks Used for the compinit phase (atinit only). Pick'null'     Don't source any file from the plugin dir Used with as'completion' for OMZ comp snippets. Pick'bin/fzf'  Source or add this specific file With as'program', adds bin/fzf to \textbf{PATH}.


\vspace{12pt}
{\bfseries\color{neonyellow} HOOK ICE MODIFIERS}\vspace{4pt}

Atinit         Code to run BEFORE the plugin sources atload         Code to run AFTER the plugin sources Example from \textbf{ZPWR}:

\begin{codebox}
zinit ice lucid nocompile wait atinit='zpwrBindOverrideOMZ;zpwrBindForGit'\\
zinit load MenkeTechnologies/forgit\\
\end{codebox}

Before forgit loads, zpwrBindOverrideOMZ runs to clear OMZ aliases that would conflict with forgit's git shortcuts.


\vspace{12pt}
{\bfseries\color{neonyellow} HOW \textbf{ZPWR} USES ICE BY CATEGORY}\vspace{4pt}

OMZ completions: lucid nocompile as'completion' pick'null' wait OMZ libs:        lucid nocompile wait atload='zpwrOmzOverrides' OMZ plugins:     lucid nocompile wait GH plugins:      lucid nocompile wait Programs:        as'program' lucid nocompile pick'bin/fzf' wait Syntax highlight:lucid nocompile wait"\$\{\textbf{ZPWR\_ZINIT}\_COMPINIT\_DELAY\}h"


\vspace{12pt}
{\bfseries\color{neonyellow} ICE STACKING}\vspace{4pt}

Multiple ices combine. Order doesn't matter (mostly). The = syntax for atload/atinit works with single quotes:

\begin{codebox}
atload='zpwrBindAliasesLate; zpwrCreateAliasCache'\\
\end{codebox}

Or with quoted strings:

\begin{codebox}
atinit'zpwrBindZdharma' atload'zpwrBindZdharmaPost'\\
\end{codebox}


\vspace{16pt}

\section{Zinit Completions \& Compinit}


\vspace{10pt}
{\large\bfseries\color{neoncyan} ZINIT COMPLETIONS \& COMPINIT}\vspace{6pt}

{\small\itshape\color{dimtext} // compinit is the most expensive call in zsh startup. defer it. //}

The completion system (compsys) needs compinit to initialize. But compinit scans all fpath entries and parses completion files. On a system with many registered completions, this is slow. \textbf{ZPWR} defers it behind turbo mode.


\vspace{12pt}
{\bfseries\color{neonyellow} ZINIT COMPINIT OPTIONS}\vspace{4pt}

\begin{codebox}
ZINIT[COMPINIT\_OPTS]='-C -u'\\
\end{codebox}

-C skips security checks (don't re-verify file ownership). -u suppresses warnings about insecure directories. These flags make compinit fast but trust your fpath.

\begin{codebox}
ZINIT[ZCOMPDUMP\_PATH]="ZSH\_COMPDUMP"\\
\end{codebox}

Points compinit to zpwr's zcompdump location.


\vspace{12pt}
{\bfseries\color{neonyellow} THE DEFERRED COMPINIT}\vspace{4pt}

Compinit runs in the LAST turbo phase:

\begin{codebox}
zinit ice lucid nocompile nocd as'null'\\
wait"\{ZPWR\_ZINIT\_COMPINIT\_DELAY\}g"\\
atinit'zicompinit; zicdreplay; zpwrBindOverrideOMZCompdefs'\\
\end{codebox}

Zicompinit: runs compinit with the configured opts. Zicdreplay: replays all compdef calls that were recorded while compinit wasn't loaded yet. Plugins that called compdef during turbo load had their calls queued -- now they fire.


\vspace{12pt}
{\bfseries\color{neonyellow} THE OMZ COMPLETION SYMLINK HACK}\vspace{4pt}

OMZ completion snippets land in:

\begin{codebox}
ZSH/snippets/OMZP::coffee/\_coffee\\
\end{codebox}

But compinit needs them in the completions dir. \textbf{ZPWR} symlinks them using zsh/files builtin:

\begin{codebox}
zmodload -F zsh/files b:zf\_ln\\
zf\_ln -sfn ZSH/snippets/OMZP::\{p\%/*\}/... ZSH/completions/...\\
\end{codebox}

Zf\_ln is a builtin ln(1), no fork needed. Fast.


\vspace{12pt}
{\bfseries\color{neonyellow} zpwrStaleZcompdump}\vspace{4pt}

The zcompdump cache gets stale when plugins update. zpwrStaleZcompdump checks the cache age. If it is older than \textasciitilde{}1 week, it triggers a full compinit rebuild. This runs automatically every shell startup (cheap check).


\vspace{12pt}
{\bfseries\color{neonyellow} FORCE REBUILD}\vspace{4pt}

\begin{codebox}
zpwr regenzsh\\
\end{codebox}

Nukes the zcompdump, re-runs compinit from scratch, and recompiles everything. Use when completions go missing or after adding new completion files to fpath.


\vspace{12pt}
{\bfseries\color{neonyellow} \textbf{ZPWR\_ZINIT}\_COMPINIT\_DELAY}\vspace{4pt}

Default: 0. Set in \textbf{\$ZPWR\_ENV/.zpwr\_env.sh.} Controls the numeric prefix of the wait ice for compinit. 0 = load immediately after prompt. Higher = more delay. On fast machines, 0 is fine. On slow machines, 1 or 2 lets other plugins finish first.


\vspace{16pt}

\section{Plugin Troubleshooting}


\vspace{10pt}
{\large\bfseries\color{neoncyan} PLUGIN TROUBLESHOOTING}\vspace{6pt}

{\small\itshape\color{dimtext} // when plugins misbehave, you need diagnostic tools. //}

Zinit ships with built-in debugging commands. When something breaks, these are your first responders.


\vspace{12pt}
{\bfseries\color{neonyellow} DIAGNOSTIC COMMANDS}\vspace{4pt}

zinit report <plugin>      Show load status, errors, files.

\begin{codebox}
zinit report zsh-users/zsh-autosuggestions\\
\end{codebox}

Shows: source file, ice modifiers used, compdef calls, functions defined, aliases created, \textbf{PATH} changes. zinit times               Show plugin load times Sorted by load time. Find the slowest plugin. If something takes >50ms, investigate. zinit loaded              List all loaded plugins zinit unload <plugin>      Unload a plugin from current session Useful for debugging conflicts. Does not persist. zinit list                List all registered plugins zinit cd <plugin>          cd into plugin directory.


\vspace{12pt}
{\bfseries\color{neonyellow} COMMON ISSUES}\vspace{4pt}

Plugin not loading? - Check if 'wait' ice is set (turbo mode defers loading) - Run zinit report to see if it errored silently - The 'lucid' ice suppresses messages -- remove it to debug Completion missing after plugin install? - Run zpwr regenzsh to rebuild zcompdump - Check fpath: print -l \$fpath | grep <plugin> - Ensure the plugin uses as'completion' or adds to fpath Keybinding conflicts? - Load order matters. Plugins loaded later override earlier ones. - \textbf{ZPWR} uses wait'0a' through wait'0e' to control this. - zpwrBindOverrideZLE fires in the 0a phase to win conflicts.


\vspace{12pt}
{\bfseries\color{neonyellow} ADDING CUSTOM PLUGINS}\vspace{4pt}

Add to \textbf{ZPWR\_GH}\_PLUGINS in \textbf{\$ZPWR\_LOCAL/.tokens.sh}:

\begin{codebox}
ZPWR\_GH\_PLUGINS+=( your-user/your-plugin )\\
\end{codebox}

The .tokens.sh file is sourced before plugin loading starts (via zpwrTokenPre), so your additions join the array in time. For conditional loading:

\begin{codebox}
zpwrCommandExists myapp \&\& ZPWR\_GH\_PLUGINS+=( user/zsh-myapp )\\
\end{codebox}


\vspace{12pt}
{\bfseries\color{neonyellow} UPDATING PLUGINS}\vspace{4pt}

\begin{codebox}
zpwr updatezinit\\
\end{codebox}

Runs zinit update --all to pull latest from all repos. After updating, run zpwr regenzsh to rebuild completions.


\vspace{12pt}
{\bfseries\color{neonyellow} NUCLEAR OPTION}\vspace{4pt}

If zinit is truly broken, nuke and rebuild:

\begin{codebox}
rm -rf ZPWR\_PLUGIN\_MANAGER\_HOME\\
\end{codebox}

Next shell start auto-clones zinit and re-downloads everything. Slow but thorough. The cockroach approach to debugging.


\vspace{16pt}


\chapter{ZSH INTERNALS}


\begin{center}
\includegraphics[width=0.85\textwidth]{/Users/wizard/.zpwr/docs/screenshots/art_deep_dive.png}
\end{center}

\section{Autoloading: Lazy Function Loading}


\vspace{10pt}
{\large\bfseries\color{neoncyan} AUTOLOADING: LAZY FUNCTION LOADING}\vspace{6pt}

{\small\itshape\color{dimtext} // 443 functions. zero parsing until you call one. //}

Autoloading is zsh's lazy-loading mechanism. Register a name, and zsh only reads the function body from disk on first call. This is why \textbf{ZPWR} can have hundreds of autoloaded functions with negligible startup cost.


\vspace{12pt}
{\bfseries\color{neonyellow} HOW autoload -z WORKS}\vspace{4pt}

1. autoload -z funcname registers 'funcname' as autoloadable. 2. Zsh records the name but does NOT read the file yet. 3. On first call to funcname, zsh searches \$fpath for a file named 'funcname' (exact match, no extension). 4. The file is sourced. The function is now defined in memory. 5. The function executes with the original arguments. 6. Subsequent calls use the in-memory version. No more disk.


\vspace{12pt}
{\bfseries\color{neonyellow} THE GLOB IN .zshrc}\vspace{4pt}

\begin{codebox}
builtin autoload -z ZPWR\_AUTOLOAD\_COMMON/*(.:t)\\
\end{codebox}

The glob *(.:t) means: *    All files in the directory (.)  Qualifier: regular files only (not dirs or symlinks) :t   Modifier: tail only (strip directory, keep filename) So it registers every regular file's name as a function. Same pattern runs for \textbf{ZPWR\_AUTOLOAD}\_DARWIN or \textbf{ZPWR\_AUTOLOAD}\_LINUX.


\vspace{12pt}
{\bfseries\color{neonyellow} THE fpath}\vspace{4pt}

Fpath is the function search path (like \textbf{PATH} but for functions):

\begin{codebox}
fpath=( ZPWR\_AUTOLOAD\_COMMON ZPWR\_AUTOLOAD\_COMP\_UTILS ZPWR\_COMPS fpath )\\
\end{codebox}

Order matters: first match wins. \textbf{ZPWR} dirs come first so they can override system completions if needed. Fpath is NOT exported (typeset +x \textbf{FPATH}) -- child shells don't inherit it.


\vspace{12pt}
{\bfseries\color{neonyellow} THE FUNCTION FILE FORMAT}\vspace{4pt}

Each file in the autoload dir contains one function. The file is named exactly like the function. Inside:

\begin{codebox}
function zpwrSomething() \{\\
\# function body\\
\}\\
zpwrSomething "@"\\
\end{codebox}

The last line calls the function with args passed through. This is the \textbf{ZPWR} convention. The function wraps itself.


\vspace{12pt}
{\bfseries\color{neonyellow} WHY AUTOLOADING IS FAST}\vspace{4pt}

443 function registrations = 443 hash table entries. Cost: microseconds. No file I/O, no parsing, no execution. Compare to sourcing 443 files: hundreds of milliseconds. Autoloading turns O(n) startup into O(1) with lazy O(1) per call.


\vspace{12pt}
{\bfseries\color{neonyellow} autoload -Uz}\vspace{4pt}

-U suppresses alias expansion during loading (safety). -z forces zsh-style function loading (vs ksh-style). \textbf{ZPWR} uses -z for its functions, -Uz for system utilities like zrecompile, zmv, zargs, and compinit.


\vspace{16pt}

\section{ZLE: The Zsh Line Editor}


\vspace{10pt}
{\large\bfseries\color{neoncyan} ZLE: THE ZSH LINE EDITOR}\vspace{6pt}

{\small\itshape\color{dimtext} // the engine behind every keystroke. readline on steroids. //}

ZLE (Zsh Line Editor) is the subsystem that handles your command line input. Every key press triggers a ZLE widget. \textbf{ZPWR} rewires dozens of these widgets for custom behavior.


\vspace{12pt}
{\bfseries\color{neonyellow} WIDGETS: NAMED \textbf{EDITOR} ACTIONS}\vspace{4pt}

A widget is a named action that ZLE can execute:

\begin{codebox}
.accept-line          Built-in: execute the current line\\
.backward-delete-char  Built-in: backspace\\
.self-insert           Built-in: insert the typed character\\
\end{codebox}

The dot prefix means 'built-in widget'. User-defined widgets can override them. \textbf{ZPWR} defines custom widgets like:

\begin{codebox}
zpwrMagicEnter         Custom Enter key behavior\\
zpwrExpandTerminateSpace Custom Space key behavior\\
zpwrAcceptLine          Custom line acceptance\\
\end{codebox}


\vspace{12pt}
{\bfseries\color{neonyellow} bindkey: MAPPING KEYS TO WIDGETS}\vspace{4pt}

bindkey connects key sequences to widgets:

\begin{codebox}
bindkey '\textasciicircum{}M' zpwrMagicEnter    Enter key\\
bindkey ' ' zpwrExpandTerminateSpace  Space key\\
bindkey -v                     Switch to vi mode\\
bindkey -e                     Switch to emacs mode\\
\end{codebox}

\textbf{ZPWR\_BINDKEY}\_VI controls which mode: true=vi, false=emacs.


\vspace{12pt}
{\bfseries\color{neonyellow} HOW zpwrMagicEnter WORKS}\vspace{4pt}

On empty line (just pressing Enter): - If in a git dir: runs git status -u . - If not: runs ls -lh . - Resets the prompt via \_p9k\_precmd\_impl + zle reset-prompt On non-empty line: delegates to zle accept-line (zpwrAcceptLine).


\vspace{12pt}
{\bfseries\color{neonyellow} HOW zpwrAcceptLine WORKS}\vspace{4pt}

The accept-line replacement does serious work: 1. Updates zconvey name with current command + TTY + PID 2. Detects commands that modify files (rm, mv, cp, mkdir...) and sets \textbf{ZPWR\_WILL}\_CLEAR=true for post-exec ls refresh 3. Handles sudo + grc colorization for aliased commands 4. Prevents global alias expansion in first position 5. Optionally expands aliases before execution 6. Calls .accept-line (the real built-in)


\vspace{12pt}
{\bfseries\color{neonyellow} THE Ctrl-F CHORD PREFIX SYSTEM}\vspace{4pt}

\textbf{ZPWR} binds Ctrl-F as a prefix key for multi-key chords:

\begin{codebox}
Ctrl-F Ctrl-D   FZF file search\\
Ctrl-F Ctrl-I   FZF word search (ag/grep)\\
Ctrl-F Ctrl-G   FZF git operations\\
\end{codebox}

This gives you an entire second keymap layer without conflicting with normal editing keys. Like tmux prefix but for your command line.


\vspace{16pt}

\section{Hook Functions: precmd, preexec, chpwd}


\vspace{10pt}
{\large\bfseries\color{neoncyan} HOOK FUNCTIONS: precmd, preexec, chpwd}\vspace{6pt}

{\small\itshape\color{dimtext} // zsh calls your code at key moments. event-driven shell. //}

Zsh provides hook points where you can inject code that runs automatically before prompts, before commands, and when directories change. \textbf{ZPWR} uses all three.


\vspace{12pt}
{\bfseries\color{neonyellow} precmd: BEFORE EACH \textbf{PROMPT}}\vspace{4pt}

Runs after a command finishes, right before the prompt draws. P10k's \_p9k\_precmd is the main user -- it updates git status, execution time, and every prompt segment. \textbf{ZPWR} adds zpwrPrecmd to the hook:

\begin{codebox}
function zpwrPrecmd()\{ export ZPWR\_OPTS="-" \}\\
\end{codebox}

This captures the current shell options before each prompt.


\vspace{12pt}
{\bfseries\color{neonyellow} preexec: BEFORE EACH COMMAND}\vspace{4pt}

Runs after Enter but BEFORE the command executes. Receives the command string as \$1. \textbf{ZPWR}'s zpwrPreexec is minimal (a placeholder hook). P10k uses preexec to record the start timestamp for command\_execution\_time segment.


\vspace{12pt}
{\bfseries\color{neonyellow} chpwd: WHEN DIRECTORY CHANGES}\vspace{4pt}

Runs whenever \$\textbf{PWD} changes (cd, pushd, popd, auto\_cd). \textbf{ZPWR}'s zpwrChpwd hook handles directory-change actions.


\vspace{12pt}
{\bfseries\color{neonyellow} THE \_functions ARRAYS}\vspace{4pt}

Zsh supports multiple hooks via arrays:

\begin{codebox}
precmd\_functions+=( myHook )      Add a precmd hook\\
preexec\_functions+=( myHook )     Add a preexec hook\\
chpwd\_functions+=( myHook )       Add a chpwd hook\\
\end{codebox}

Every function in the array runs at the hook point. Order matters: first in array = first to execute.


\vspace{12pt}
{\bfseries\color{neonyellow} HOW \textbf{ZPWR} REGISTERS HOOKS}\vspace{4pt}

ZpwrBindPrecmd (runs during p10k atload):

\begin{codebox}
precmd\_functions=( zpwrPrecmd \{(@)precmd\_functions:\#zpwrPrecmd\} )\\
\end{codebox}

This puts zpwrPrecmd FIRST, removing any duplicate. The \$\{(@)array:\#pattern\} syntax removes matching elements. zpwrBindPreexecChpwd (runs during syntax-highlighting atload):

\begin{codebox}
preexec\_functions=( \{(@)preexec\_functions:\#zpwrPreexec\} zpwrPreexec )\\
chpwd\_functions=( \{(@)chpwd\_functions:\#zpwrChpwd\} zpwrChpwd )\\
\end{codebox}

These put zpwr hooks LAST (after p10k and other plugins).


\vspace{12pt}
{\bfseries\color{neonyellow} THE HOOK COUNT IN zpwr top}\vspace{4pt}

The zpwr top dashboard shows the total hook count:

\begin{codebox}
nhooks=((\{\#precmd\_functions\}+\{\#preexec\_functions\}+\{\#chpwd\_functions\}))\\
\end{codebox}

Typical value: 5-8 hooks total across all three arrays. If this number grows unexpectedly, a plugin is adding hooks you may not want. Check with: print -l \$precmd\_functions.


\vspace{16pt}

\section{Emulate \& Options}


\vspace{10pt}
{\large\bfseries\color{neoncyan} EMULATE \& OPTIONS}\vspace{6pt}

{\small\itshape\color{dimtext} // zsh has 200+ options. emulate -L zsh resets them all locally. //}

Zsh options control fundamental behavior: globbing, word splitting, quoting, history. \textbf{ZPWR} functions use emulate -L zsh to guarantee a known environment regardless of user settings.


\vspace{12pt}
{\bfseries\color{neonyellow} emulate -L zsh}\vspace{4pt}

\begin{codebox}
builtin emulate -L zsh\\
\end{codebox}

Resets ALL options to zsh defaults for the current function scope. -L means local: options revert when the function returns. Without this, a function might inherit noglob from the caller and break globbing, or inherit shwordsplit and break arrays. \textbf{ZPWR} functions start with this to be immune to the caller's state.


\vspace{12pt}
{\bfseries\color{neonyellow} KEY OPTIONS IN \textbf{ZPWR}}\vspace{4pt}

extendedglob      \#\# = one-or-more, (\#i) = case-insensitive glob \textasciicircum{}pattern = negation. Essential for \textbf{ZPWR} globs. glob\_dots         Dot files matched by * (no need for .*) globstarshort     **.c = **/*.c (recursive glob shorthand) braceccl          \{abcd0-9\} expands as character class rc\_quotes         '' inside single quotes = literal single quote rc\_expand\_param   Array expansion includes prefix: xx\$\{arr\} = xxA xxB nocaseglob        Globs are case-insensitive globally numeric\_glob\_sort  file1 file2 file10 sort numerically prompt\_subst       Allow \$(...) and \$\{...\} in prompt strings transient\_rprompt  Right prompt only on current line


\vspace{12pt}
{\bfseries\color{neonyellow} THE noshwordsplit DEFAULT}\vspace{4pt}

Bash splits unquoted variables on whitespace by default. Zsh does NOT (noshwordsplit). This means:

\begin{codebox}
x="hello world"; print x  -> one argument in zsh, two in bash\\
\end{codebox}

If a function inherits shwordsplit (e.g., from emulate sh), arrays and parameter expansions break silently. Another reason for emulate -L zsh.


\vspace{12pt}
{\bfseries\color{neonyellow} THE local-IN-LOOP BUG}\vspace{4pt}

In zsh, 'local' inside a loop prints the variable value as a side effect when the variable already exists:

\begin{codebox}
for i in 1 2 3; do local x=i; done  \# prints x=1 x=2 x=3\\
\end{codebox}

This is because 'local' on an existing local acts like typeset, which prints. Declare locals BEFORE loops, not inside them.


\vspace{12pt}
{\bfseries\color{neonyellow} auto\_name\_dirs: REMOVED}\vspace{4pt}

This option auto-registers any parameter set to a directory path as a named directory. Sounds nice, but p10k/gitstatus set internal variables to paths, polluting the named dir table. \textbf{ZPWR} removed it. Explicit hash -d in zpwrBindDirs handles it. The comment in .zshrc: \# auto\_name\_dirs removed: pollutes named \# dirs with p10k/gitstatus internals.


\vspace{16pt}

\section{Parameter Expansion Magic}


\vspace{10pt}
{\large\bfseries\color{neoncyan} PARAMETER EXPANSION MAGIC}\vspace{6pt}

{\small\itshape\color{dimtext} // zsh parameter expansion is a programming language in itself. //}

Zsh's \$\{...\} syntax can do transformations that would require sed, awk, basename, dirname, tr, and cut in bash. No forks. No subshells. Pure in-process string manipulation.


\vspace{12pt}
{\bfseries\color{neonyellow} PARAMETER EXPANSION FLAGS (the parenthesized letters)}\vspace{4pt}

\$\{(K)assoc\}          Keys of associative array \$\{(v)assoc\}          Values of associative array \$\{(kv)assoc\}         Key-value pairs interleaved \$\{(s:/:)var\}         Split string on / into array \$\{(j:,:)array\}       Join array elements with , \$\{(L)var\}            Lowercase the entire string \$\{(U)var\}            Uppercase the entire string \$\{(f)var\}            Split on newlines (lines to array) \$\{(F)array\}          Join with newlines (array to lines) \$\{(t)var\}            Type of variable (scalar, array, etc.) \$\{(q)var\}            Quote the value for re-use in eval \$\{(Q)var\}            Remove one level of quoting \$\{(z)var\}            Split into words like the shell parser \$\{(o)array\}          Sort ascending \$\{(O)array\}          Sort descending.


\vspace{12pt}
{\bfseries\color{neonyellow} MODIFIERS (the colon letters)}\vspace{4pt}

\$\{Var:t\}             Tail: basename (/a/b/c.txt -> c.txt) \$\{var:h\}             Head: dirname (/a/b/c.txt -> /a/b) \$\{var:r\}             Root: remove extension (file.txt -> file) \$\{var:e\}             Extension: (file.txt -> txt) \$\{var:A\}             Absolute path (resolve symlinks) \$\{var:l\}             Lowercase (zsh 5.0+) \$\{var:u\}             Uppercase (zsh 5.0+)


\vspace{12pt}
{\bfseries\color{neonyellow} SUBSTITUTION AND PATTERN MATCHING}\vspace{4pt}

{\small\itshape\color{dimtext} \$\{var//pat/rep\}       Global substitution}

\$\{Var/pat/rep\}        First match substitution \$\{var\#pat\}            Remove shortest prefix match \$\{var\#\#pat\}           Remove longest prefix match \$\{var\%pat\}            Remove shortest suffix match \$\{var\%\%pat\}           Remove longest suffix match.


\vspace{12pt}
{\bfseries\color{neonyellow} \textbf{ZPWR} EXAMPLES}\vspace{4pt}

From zpwrBindDirs (writing hash dir cache):

\begin{codebox}
print -r "hash -d \{k\}=\{(q)v\}"  -- (q) quotes the value safely\\
\end{codebox}

From zpwrAcceptLine (parsing command buffer):

\begin{codebox}
mywords=( \{(z)cmdstr\} )           -- (z) shell-word splits\\
cmd=\{mywords[1]\#=\}                -- strip leading = from =cmd\\
\end{codebox}

From zpwrBindPrecmd (removing duplicates from array):

\begin{codebox}
\{(@)precmd\_functions:\#zpwrPrecmd\}  -- remove matching element\\
\end{codebox}

These are not academic exercises. They replace forks.


\vspace{16pt}

\section{The .zwc Compilation System}


\vspace{10pt}
{\large\bfseries\color{neoncyan} THE .zwc COMPILATION SYSTEM}\vspace{6pt}

{\small\itshape\color{dimtext} // bytecode for your shell. because even scripts deserve a compiler. //}

Zsh can compile scripts to .zwc (Zsh Word Code) bytecode. These files skip the parsing phase on load, going straight to execution. \textbf{ZPWR} compiles 50+ files for speed.


\vspace{12pt}
{\bfseries\color{neonyellow} WHAT .zwc FILES ARE}\vspace{4pt}

A .zwc file is a binary representation of parsed zsh code. When zsh sources a .zsh file, it checks for a .zwc sibling. If the .zwc exists and is newer than the source, zsh loads the bytecode directly. No lexing, no parsing. Speedup: 2-5x for large files. Marginal for tiny ones.


\vspace{12pt}
{\bfseries\color{neonyellow} zrecompile}\vspace{4pt}

The zrecompile utility compiles only if the source is newer:

\begin{codebox}
autoload -Uz zrecompile\\
zrecompile -p -U ZPWR\_ENV\_FILE\\
\end{codebox}

-p means 'check pairs' (source vs .zwc timestamps). -U suppresses alias expansion during compilation. If the .zwc is current, zrecompile is a no-op.


\vspace{12pt}
{\bfseries\color{neonyellow} zpwrRecompileFiles: WHAT GETS COMPILED}\vspace{4pt}

The function compiles two categories of files: Local files (\textasciitilde{}47): \textbf{\textasciitilde{}/.zshrc}, \textbf{\textasciitilde{}/.zshenv}, \textbf{\textasciitilde{}/.zlogout}, \textbf{\textasciitilde{}/.zlogin} p10k-instant-prompt-*.zsh, p10k-dump-*.zsh \$\textbf{ZSH\_COMPDUMP}, \$\textbf{ZPWR\_ALIAS}\_FILE, \$\textbf{ZPWR\_PROMPT}\_FILE \textbf{\$ZPWR\_LOCAL/.tokens.sh}, \$\textbf{ZPWR\_ENV}\_FILE, \$\textbf{ZPWR\_RE}\_ENV\_FILE All .zsh/.sh in \$\textbf{ZPWR\_ENV}/, all .zsh in \$\textbf{ZPWR\_SCRIPTS}/ All plugin .zsh files under \$\textbf{ZPWR\_PLUGIN}\_MANAGER\_HOME/bin/ OMZ lib, plugin, theme, and custom files Global files (\textasciitilde{}8): \textbf{\textbf{/etc/profile}}, \textbf{\textbf{/etc/zshrc}}, \textbf{\textbf{/etc/zshenv}}, \textbf{\textbf{/etc/zprofile}}, etc. These need sudo (compiled with sudo zsh -c zrecompile).


\vspace{12pt}
{\bfseries\color{neonyellow} \textbf{ZPWR} VERBS FOR COMPILATION}\vspace{4pt}

\begin{codebox}
zpwr recompile       Compile all known files (zpwrRecompileFiles)\\
zpwr decompile       Remove .zwc files (zpwrUncompile)\\
zpwr refreshzwc      Remove + recompile (clean rebuild)\\
\end{codebox}


\vspace{12pt}
{\bfseries\color{neonyellow} THE zsh/files MODULE}\vspace{4pt}

Zsh can load builtin file operations via zmodload:

\begin{codebox}
zmodload -F zsh/files b:zf\_ln b:zf\_rm b:zf\_mkdir\\
\end{codebox}

These are shell builtins that replace ln, rm, mkdir. No fork, no exec -- pure kernel syscalls from the shell process. \textbf{ZPWR} uses zf\_ln for the OMZ completion symlink hack. In a loop creating 12 symlinks, this saves 12 fork+exec pairs.


\vspace{12pt}
{\bfseries\color{neonyellow} OTHER USEFUL MODULES}\vspace{4pt}

\begin{codebox}
zsh/datetime    EPOCHREALTIME, EPOCHSECONDS\\
zsh/complist    Menu-select completion mode\\
zsh/zprof       Profiler (ZPWR\_PROFILING=true enables it)\\
\end{codebox}

Modules are loaded with: zmodload zsh/datetime.


\vspace{16pt}


\chapter{FZF ARCHITECTURE}


\begin{center}
\includegraphics[width=0.85\textwidth]{/Users/wizard/.zpwr/docs/screenshots/figlet_scan.png}
\end{center}

\section{FZF: The Neural Interface}


\vspace{10pt}
{\large\bfseries\color{neoncyan} FZF: THE NEURAL INTERFACE}\vspace{6pt}

{\small\itshape\color{dimtext} // the universal adapter between data and decisions. //}

fzf (fuzzy finder) is a general-purpose command-line filter. Pipe anything into it, fuzzy-search interactively, get results. \textbf{ZPWR} wraps fzf into dozens of workflows: files, words, git, verbs.


\vspace{12pt}
{\bfseries\color{neonyellow} THE PATTERN: PIPE | FZF | ACTION}\vspace{4pt}

Every fzf integration follows the same architecture:

\begin{codebox}
<data source> | fzf [options] | <action on selection>\\
\end{codebox}

The data source generates lines. fzf presents them. The action consumes the selected line(s). Simple. Composable. Powerful.


\vspace{12pt}
{\bfseries\color{neonyellow} FZF\_DEFAULT\_OPTS}\vspace{4pt}

\textbf{ZPWR} sets global fzf defaults in zpwrBindFZFLate:

\begin{codebox}
FZF\_DEFAULT\_OPTS="ZPWR\_COMMON\_FZF\_ELEM --reverse --border --height 100\%"\\
\end{codebox}

Plus a Solarized Dark color scheme layered on top:

\begin{codebox}
--color fg:-1,bg:-1,hl:blue,fg+:base2,bg+:base02,hl+:blue\\
\end{codebox}

\textbf{ZPWR\_COMMON}\_FZF\_ELEM adds the logo prompt and navigation binds:

\begin{codebox}
--prompt='ZPWR\_FZF\_LOGO ' --bind=ctrl-n:page-down,...\\
\end{codebox}


\vspace{12pt}
{\bfseries\color{neonyellow} THE PREVIEW WINDOW}\vspace{4pt}

fzf's --preview flag runs a command for each highlighted item:

\begin{codebox}
fzf --preview 'bat --color=always \{\}'\\
\end{codebox}

\textbf{ZPWR} configures rich previews per context: Files:   bat with syntax highlighting, line numbers, grid Dirs:    stat output Packages: dpkg -I or rpm -qi Git:     diff preview, commit log If bat is installed, it's the default colorizer. Falls back to pygmentize with terminal256 output.


\vspace{12pt}
{\bfseries\color{neonyellow} fzf-completion (Tab)}\vspace{4pt}

fzf replaces the default Tab completion with a fuzzy finder. Type a command, press Tab, and fzf opens with completions:

\begin{codebox}
cd **<Tab>     fzf through subdirectories\\
kill <Tab>      fzf through running processes\\
ssh **<Tab>     fzf through known hosts\\
\end{codebox}

FZF\_COMPLETION\_OPTS controls the preview for these.


\vspace{12pt}
{\bfseries\color{neonyellow} fzf-tab PLUGIN}\vspace{4pt}

MenkeTechnologies/fzf-tab replaces zsh's menu completion with fzf. Any time the completion system would show a menu, fzf-tab intercepts it and shows fzf instead. This means: - Fuzzy matching for every completion, not just ** triggers - Previews for file completions, git branches, everything - FZF\_TAB\_OPTS configures the behavior (ansi, nth, delimiter) - FZF\_TAB\_COMMAND sets which fzf binary to use.


\vspace{12pt}
{\bfseries\color{neonyellow} GLOBAL ALIASES FOR FZF}\vspace{4pt}

\begin{codebox}
\{ZPWR\_GLOBAL\_ALIAS\_PREFIX\}z   Pipe into fzf: ls | ,,z\\
\{ZPWR\_GLOBAL\_ALIAS\_PREFIX\}zb  Inline fzf selection: vim ,,zb\\
\end{codebox}


\vspace{16pt}

\section{FZF in Practice}


\vspace{10pt}
{\large\bfseries\color{neoncyan} FZF IN PRACTICE}\vspace{6pt}

{\small\itshape\color{dimtext} // the theory is simple. the practice is where it gets good. //}

\textbf{ZPWR} ships dozens of fzf-powered workflows. Each follows the generate | fzf | action pattern. Here are the heaviest hitters.


\vspace{12pt}
{\bfseries\color{neonyellow} FILE SEARCH (zpwr filesearch / Ctrl-F Ctrl-D)}\vspace{4pt}

Pattern: find files | fzf --preview 'bat' | open in editor 1. Find/fd generates a file list (respecting .gitignore) 2. fzf shows files with bat-powered syntax-highlighted preview 3. Selected file opens in \$\textbf{EDITOR} at the matched location FZF\_CTRL\_T\_OPTS configures the preview command.


\vspace{12pt}
{\bfseries\color{neonyellow} WORD SEARCH (zpwr wordsearch / Ctrl-F Ctrl-I)}\vspace{4pt}

Pattern: ag/rg search | fzf --preview 'context' | open at line 1. Ag (silver searcher) or rg searches file contents 2. fzf shows matches with file:line:content format 3. Selection opens the file at the exact line number FZF\_AG\_OPTS: delimiter ':', match on 3rd+ field (skip file:line).


\vspace{12pt}
{\bfseries\color{neonyellow} GIT ADD (forgitadd)}\vspace{4pt}

Pattern: git diff --stat | fzf --preview 'git diff' | git add 1. Lists modified files with diff stats 2. Preview pane shows the actual diff for each file 3. Toggle files with Tab, confirm to stage them Part of forgit: the git + fzf integration plugin.


\vspace{12pt}
{\bfseries\color{neonyellow} VERB SEARCH (zpwr verbsfzf)}\vspace{4pt}

Pattern: \textbf{ZPWR\_VERBS} keys | fzf --preview 'help' | eval 1. All \$\{\#\textbf{ZPWR\_VERBS}\}+ zpwr verbs piped into fzf 2. Preview shows the verb description and command 3. Select a verb to execute it immediately FZF\_ENV\_OPTS\_VERBS configures the preview for verbs.


\vspace{12pt}
{\bfseries\color{neonyellow} CUSTOM FZF KEY BINDINGS}\vspace{4pt}

The Ctrl-F prefix opens \textbf{ZPWR}'s fzf keymap layer:

\begin{codebox}
Ctrl-F Ctrl-D   File search (find + fzf + editor)\\
Ctrl-F Ctrl-I   Word search (ag + fzf + editor)\\
Ctrl-F Ctrl-G   Git search (git log + fzf)\\
Ctrl-F Ctrl-R   History search with fzf\\
\end{codebox}

These are ZLE widgets bound via zpwrBindOverrideZLE.


\vspace{12pt}
{\bfseries\color{neonyellow} zpwrIntoFzf AND zpwrIntoFzfAg}\vspace{4pt}

Two core ZLE widgets that bridge the command line with fzf: zpwrIntoFzf     File-oriented fzf (Ctrl-T style) zpwrIntoFzfAg   Content-oriented fzf (ag/grep style) Both capture fzf output and insert it into the command line buffer (\$BUFFER). You select files/matches, they appear at your cursor. Seamless integration between search and input.


\vspace{12pt}
{\bfseries\color{neonyellow} THE COLORIZER CHAIN}\vspace{4pt}

\textbf{ZPWR\_COLORIZER} preference: bat > pygmentize. bat uses themes (BAT\_THEME=\$\textbf{ZPWR\_BAT}\_THEME). Pygmentize uses PYGMENTIZE\_COLOR style. Language-specific colorizers: FZF\_COLORIZER\_C, \_SH, \_YAML, \_JAVA.


\vspace{16pt}

\section{Building Your Own FZF Verbs}


\vspace{10pt}
{\large\bfseries\color{neoncyan} BUILDING YOUR OWN FZF VERBS}\vspace{6pt}

{\small\itshape\color{dimtext} // once you grok the pattern, everything becomes fzf-able. //}

The generate | fzf | action pattern is infinitely extensible. Here's how to build your own fzf-powered verbs.


\vspace{12pt}
{\bfseries\color{neonyellow} THE PATTERN}\vspace{4pt}

\begin{codebox}
<generator> | fzf [--preview '<preview-cmd>'] | <consumer>\\
\end{codebox}

Generator: anything that outputs lines (find, git, docker, etc.) Preview: command that shows detail for highlighted line (\textbackslash{}\{\textbackslash{}\}) Consumer: action on selected lines (xargs, while read, eval)


\vspace{12pt}
{\bfseries\color{neonyellow} EXAMPLE: FZF THROUGH DOCKER CONTAINERS}\vspace{4pt}

\begin{codebox}
function fzf-docker() \{\\
docker ps -a --format '\{\{.ID\}\} \{\{.Names\}\} \{\{.Status\}\}' |\\
fzf --preview 'docker inspect \{1\} | head -50' |\\
awk '\{print 1\}' |\\
xargs -r docker logs --tail 100\\
\}\\
\end{codebox}

Generator: docker ps. Preview: docker inspect. Consumer: docker logs.


\vspace{12pt}
{\bfseries\color{neonyellow} EXAMPLE: FZF THROUGH TMUX SESSIONS}\vspace{4pt}

\begin{codebox}
function fzf-tmux-session() \{\\
tmux list-sessions -F '\#\{session\_name\}: \#\{session\_windows\} windows' |\\
fzf --preview 'tmux list-windows -t \{1\}' |\\
cut -d: -f1 |\\
xargs tmux switch-client -t\\
\}\\
\end{codebox}

Generator: tmux list-sessions. Consumer: tmux switch-client.


\vspace{12pt}
{\bfseries\color{neonyellow} THE ACCEPT vs EDIT PATTERN}\vspace{4pt}

\textbf{ZPWR} verbs come in two flavors: Accept verbs  Execute the selection immediately Example: verbsfzf selects a verb and runs it. Edit verbs    Insert the selection into \$BUFFER for editing Example: filesearch inserts the path into your command line. The difference is in the consumer. Accept uses eval/exec. Edit uses BUFFER=\$selection + zle redisplay.


\vspace{12pt}
{\bfseries\color{neonyellow} MAKING IT A ZLE WIDGET}\vspace{4pt}

To bind your fzf verb to a key:

\begin{codebox}
function my-fzf-widget() \{\\
BUFFER=(my-generator | fzf)\\
CURSOR=\{\#BUFFER\}\\
zle redisplay\\
\}\\
zle -N my-fzf-widget\\
bindkey '\textasciicircum{}F\textasciicircum{}X' my-fzf-widget\\
\end{codebox}

Now Ctrl-F Ctrl-X triggers your custom fzf workflow.


\vspace{12pt}
{\bfseries\color{neonyellow} TIPS FOR GOOD FZF EXPERIENCES}\vspace{4pt}

- Use --ansi if your generator outputs color codes - Use --delimiter and --nth to control what fzf matches on - Use -m for multi-select (toggle with Tab) - Use --header for context hints at the top - Preview window size: --preview-window=right:50\%:wrap - Test with: echo 'a\textbackslash{}nb\textbackslash{}nc' | fzf --preview 'echo \{\}'


\vspace{16pt}


\chapter{TMUX DEEP DIVE}


\begin{center}
\includegraphics[width=0.85\textwidth]{/Users/wizard/.zpwr/docs/screenshots/art_mainframe.png}
\end{center}

\section{Tmux Architecture in ZPWR}


\vspace{10pt}
{\large\bfseries\color{neoncyan} TMUX ARCHITECTURE IN ZPWR}\vspace{6pt}

{\small\itshape\color{dimtext} // the terminal multiplexer. one screen, infinite shells. //}

Tmux is the backbone of the \textbf{ZPWR} workflow. Multiple shells, persistent sessions, split panes, and code execution across panes -- all orchestrated from \$\textbf{ZPWR\_TMUX}.


\vspace{12pt}
{\bfseries\color{neonyellow} SESSION / WINDOW / PANE MODEL}\vspace{4pt}

Session    Named group of windows (e.g., 'dev', 'zpwr-study') Window     A tab inside a session (like browser tabs) Pane       A split inside a window (like vim splits) Typical \textbf{ZPWR} layout: 1 session, 2-4 windows, 2-4 panes each. Prefix key: Ctrl-B (default). Followed by a command key.


\vspace{12pt}
{\bfseries\color{neonyellow} SESSION PERSISTENCE}\vspace{4pt}

Tmux-resurrect   Manual save/restore of sessions Save: prefix + Ctrl-s Restore: prefix + Ctrl-r Saves: pane layouts, working directories, running programs. @resurrect-processes ':all:' -- saves ALL running programs. Tmux-continuum    Auto-save every 15 minutes Works silently in the background. No manual intervention. Combined with resurrect: your tmux state survives reboots.


\vspace{12pt}
{\bfseries\color{neonyellow} THE POWERLINE STATUS BAR}\vspace{4pt}

Tmux sources powerline bindings for the status bar:

\begin{codebox}
source-file \textasciitilde{}/.tmux/powerline/bindings/tmux/powerline.conf\\
\end{codebox}

The symlink \textbf{\textasciitilde{}/.tmux/powerline} points to pip's powerline-status. Powerline-daemon runs in the background for performance. Status bar shows: session name, window list, hostname, time, battery level (via tmux-battery plugin).


\vspace{12pt}
{\bfseries\color{neonyellow} \$\textbf{ZPWR\_TMUX}: THE CONFIG DIRECTORY}\vspace{4pt}

All tmux config lives under \$\textbf{ZPWR\_TMUX}:

\begin{codebox}
init.conf               Main config (bindings, options)\\
tmux-mac.conf           macOS-specific overrides\\
tmux-linux.conf         Linux-specific overrides\\
tmux\_ge\_21.conf         Tmux >= 2.1 features\\
tmux\_ge\_29.conf         Tmux >= 2.9 features\\
tmux\_ge\_31.conf         Tmux >= 3.1 features\\
four-panes.conf         Layout: 4 panes\\
sixteen-panes.conf      Layout: 16 panes\\
config-files.conf       Opens brc, vrc, trc, zrc in 4 panes\\
learn.conf              Study layout\\
\end{codebox}


\vspace{12pt}
{\bfseries\color{neonyellow} init.conf HIGHLIGHTS}\vspace{4pt}

- 30,000 line scrollback (\textbf{ZPWR\_TMUX}\_HISTORY\_LIMIT) - Vi keybindings in copy mode - Prefix S: synchronize-panes (type in all panes at once) - Prefix q: display pane numbers for 5 seconds - 3-second repeat-time for directional keys


\vspace{16pt}

\section{Vim-Tmux Code Execution}


\vspace{10pt}
{\large\bfseries\color{neoncyan} VIM-TMUX CODE EXECUTION}\vspace{6pt}

{\small\itshape\color{dimtext} // select code in vim. execute in the pane next door. magic. //}

The killer feature of the \textbf{ZPWR} vim+tmux integration: write code in vim, send it to an adjacent tmux pane for execution. No copy-paste. No switching windows.


\vspace{12pt}
{\bfseries\color{neonyellow} THE TmuxRepeat FUNCTION}\vspace{4pt}

Defined in .vimrc, TmuxRepeat has three modes: 'file'    Save the file, execute it in the right pane 'visual'  Save visual selection to temp file, execute it 'repl'    Paste visual selection directly into right pane.


\vspace{12pt}
{\bfseries\color{neonyellow} SUPPORTED LANGUAGES}\vspace{4pt}

TmuxRepeat detects file type by extension:

\begin{codebox}
sh, zsh, py, rb, pl, clj, tcl, vim, lisp, hs, ml,\\
coffee, swift, lua, java, f90, cr\\
\end{codebox}

For known types, it sends the file to zpwrRunner.sh which dispatches to the correct interpreter. For unknown types, it falls back to TmuxRepeatGeneric (re-run previous command).


\vspace{12pt}
{\bfseries\color{neonyellow} HOW IT WORKS}\vspace{4pt}

1. Ctrl-V in normal/insert mode (or leader-vj): Saves the file, then sends:

\begin{codebox}
tmux send-keys -t right 'bash "ZPWR\_SCRIPTS/zpwrRunner.sh" "file"' C-m\\
\end{codebox}

2. leader-vk in visual mode: Yanks selection to temp file \$\textbf{ZPWR\_HIDDEN}\_DIR\_TEMP/.vimTempFile.<ext>, then sends that temp file to the right pane for execution. 3. leader-vl in visual mode (REPL mode): Yanks selection, loads it into tmux buffer 'buffer0099', then pastes directly into the right pane. For REPLs like python, irb, node -- code runs line-by-line in the REPL.


\vspace{12pt}
{\bfseries\color{neonyellow} THE PANE TARGETING}\vspace{4pt}

Standard tmux: sends to pane named 'right' (the next pane). MacVim/GUI: sends to session 'vimmers:0.' (external tmux). The pane\_count check ensures a target pane exists.


\vspace{12pt}
{\bfseries\color{neonyellow} allPanes.zsh: BROADCAST TO ALL PANES}\vspace{4pt}

\textbf{\$ZPWR\_SCRIPTS/allPanes.zsh} sends commands to all tmux panes:

\begin{codebox}
prefix + Space   Run command in all panes (single mode)\\
prefix + v       Run command in all panes (multi mode)\\
prefix + b       Open URL in all panes\\
\end{codebox}


\vspace{12pt}
{\bfseries\color{neonyellow} SYNCHRONIZE-PANES}\vspace{4pt}

Prefix + S   Toggle synchronize-panes When enabled, every keystroke goes to ALL panes in the window. Useful for running identical commands on multiple servers, or setting up parallel environments simultaneously.


\vspace{12pt}
{\bfseries\color{neonyellow} GOOGLE FROM TMUX}\vspace{4pt}

Prefix + g   Google search from tmux Runs \textbf{\$ZPWR\_TMUX/google.sh} with the current pane content. Select text, press prefix+g, and your browser opens with a Google search for the selection.


\vspace{16pt}

\section{Tmux Plugins \& Extensions}


\vspace{10pt}
{\large\bfseries\color{neoncyan} TMUX PLUGINS \& EXTENSIONS}\vspace{6pt}

{\small\itshape\color{dimtext} // tpm: the plugin manager for your multiplexer's plugins. yo dawg. //}

\textbf{ZPWR} manages 9 tmux plugins via tpm (Tmux Plugin Manager). They live in \textbf{\textasciitilde{}/.tmux/plugins/} and add features that tmux does not provide out of the box.


\vspace{12pt}
{\bfseries\color{neonyellow} TPM: TMUX PLUGIN MANAGER}\vspace{4pt}

Installed to \textbf{\textasciitilde{}/.tmux/plugins/tpm.} Managed via prefix keys:

\begin{codebox}
prefix + I       Install plugins listed in .tmux.conf\\
prefix + U       Update all installed plugins\\
prefix + Alt-u   Clean plugins not in config\\
\end{codebox}

Plugins are declared in \textbf{\$ZPWR\_INSTALL/.tmux.conf}:

\begin{codebox}
set-option -g @plugin 'tmux-plugins/tpm'\\
set-option -g @plugin 'tmux-plugins/tmux-resurrect'\\
\end{codebox}


\vspace{12pt}
{\bfseries\color{neonyellow} THE PLUGIN ROSTER (.tmuxplugins)}\vspace{4pt}

Tmux-plugins/tpm                Plugin manager itself tmux-plugins/tmux-sensible      Sane defaults for tmux tmux-plugins/tmux-resurrect     Session save/restore tmux-plugins/tmux-continuum      Auto-save sessions tmux-plugins/tmux-battery        Battery indicator tmux-plugins/tmux-prefix-highlight Highlight when prefix active tmux-plugins/tmux-yank           System clipboard integration MenkeTechnologies/tmux-thumbs     Quick copy with hints MenkeTechnologies/tmux-fzf-url    FZF through URLs in pane.


\vspace{12pt}
{\bfseries\color{neonyellow} tmux-resurrect: WHAT IT SAVES}\vspace{4pt}

- Pane layouts (splits, sizes, positions) - Working directories for each pane - Running programs (@resurrect-processes ':all:') - Command history is NOT saved (that's zsh's job) Save: prefix+Ctrl-s. Restore: prefix+Ctrl-r.


\vspace{12pt}
{\bfseries\color{neonyellow} tmux-continuum: AUTO-SAVE}\vspace{4pt}

Piggybacks on resurrect. Auto-saves every 15 minutes. @continuum-restore is commented out in \textbf{ZPWR} by default. Enable it to auto-restore on tmux server start.


\vspace{12pt}
{\bfseries\color{neonyellow} tmux-fzf-url: URL PICKER}\vspace{4pt}

Scans the visible pane for URLs, opens them in fzf:

\begin{codebox}
prefix + e   FZF through URLs, open selected in browser\\
prefix + u   FZF through URLs, search mode\\
\end{codebox}

Uses \textbf{ZPWR\_TMUX}\_HISTORY\_LIMIT for scrollback depth.


\vspace{12pt}
{\bfseries\color{neonyellow} CUSTOM SCRIPTS IN \$\textbf{ZPWR\_TMUX}}\vspace{4pt}

Layout scripts create multi-pane environments:

\begin{codebox}
prefix + D   control-window.conf (IDE layout)\\
prefix + F   four-panes.conf\\
prefix + G   eight-panes.conf\\
prefix + O   sixteen-panes.conf\\
prefix + T   config-files.conf (brc, vrc, trc, zrc)\\
prefix + M   learn.conf (study layout)\\
\end{codebox}


\vspace{12pt}
{\bfseries\color{neonyellow} zpwr study: THE SPLIT-PANE LEARNING ENVIRONMENT}\vspace{4pt}

Creates a dedicated tmux session 'zpwr-study' with the encyclopedia on the left and a live shell on the right. Read the docs, try commands. Knowledge through practice.


\vspace{16pt}


\chapter{VIM DEEP DIVE}


\begin{center}
\includegraphics[width=0.85\textwidth]{/Users/wizard/.zpwr/docs/screenshots/art_neural_link.png}
\end{center}

\section{Vim Plugin Architecture}


\vspace{10pt}
{\large\bfseries\color{neoncyan} VIM PLUGIN ARCHITECTURE}\vspace{6pt}

{\small\itshape\color{dimtext} // 15 plugins in .vimrc. each one a force multiplier. //}

\textbf{ZPWR}'s vim config uses Vundle as the plugin manager and ships 15 curated plugins that transform vim into a full IDE with completion, snippets, and tmux integration.


\vspace{12pt}
{\bfseries\color{neonyellow} VUNDLE: THE PLUGIN MANAGER}\vspace{4pt}

Plugins live in \textbf{\textasciitilde{}/.vim/bundle/.} Managed by Vundle:

\begin{codebox}
set runtimepath+=\textasciitilde{}/.vim/bundle/Vundle.vim\\
call vundle\#begin()\\
Plugin 'VundleVim/Vundle.vim'\\
... more plugins ...\\
call vundle\#end()\\
\end{codebox}

Install: :PluginInstall  Update: :PluginUpdate  Clean: :PluginClean.


\vspace{12pt}
{\bfseries\color{neonyellow} YouCompleteMe: SEMANTIC COMPLETION}\vspace{4pt}

\begin{codebox}
Plugin 'ycm-core/YouCompleteMe'\\
\end{codebox}

The heavyweight champion. Provides semantic completion for: C/C++ (via clangd), Python (via jedi), and more. Config: completion after 1 char, enabled in all file types, completes in comments. Uses C-n/C-p for selection.


\vspace{12pt}
{\bfseries\color{neonyellow} UltiSnips + vim-snippets: SNIPPET ENGINE}\vspace{4pt}

\begin{codebox}
Plugin 'SirVer/ultisnips'\\
Plugin 'honza/vim-snippets'\\
\end{codebox}

Tab expands snippets, Tab/Shift-Tab navigate tab stops. Hundreds of snippets for every language: type 'for<Tab>' in Python and get a complete for loop with cursor placement.


\vspace{12pt}
{\bfseries\color{neonyellow} SuperTab: TAB COMPLETION GLUE}\vspace{4pt}

\begin{codebox}
Plugin 'ervandew/supertab'\\
\end{codebox}

Unifies Tab behavior: C-n cycles forward through completions. Coordinates with YCM and UltiSnips so Tab does the right thing in every context (snippet vs completion vs indent).


\vspace{12pt}
{\bfseries\color{neonyellow} vim-auto-save: NEVER :w AGAIN}\vspace{4pt}

\begin{codebox}
Plugin '907th/vim-auto-save'\\
\end{codebox}

Saves on InsertLeave, TextChanged, and TextChangedI. Every edit auto-saves silently. No more 'Unsaved changes' panics. G:auto\_save\_silent=1 suppresses the notification.


\vspace{12pt}
{\bfseries\color{neonyellow} rainbow: MATCHING BRACKET COLORS}\vspace{4pt}

\begin{codebox}
Plugin 'luochen1990/rainbow'\\
\end{codebox}

Nested brackets get different colors: parentheses in blue, inner brackets in orange, deeper in cyan, deepest in magenta. G:rainbow\_active=1 enables it globally on startup.


\vspace{12pt}
{\bfseries\color{neonyellow} vim-multiple-cursors: MULTI-CURSOR EDITING}\vspace{4pt}

\begin{codebox}
Plugin 'TerryMa/vim-multiple-cursors'\\
\end{codebox}

Sublime Text-style multi-cursor. Ctrl-N selects next match, then edit all occurrences simultaneously.


\vspace{12pt}
{\bfseries\color{neonyellow} tmux-complete.vim: CROSS-PANE COMPLETION}\vspace{4pt}

\begin{codebox}
Plugin 'wellle/tmux-complete.vim'\\
\end{codebox}

Completes words from ALL tmux panes. Writing a command in vim that references output in another pane? Tab-complete it. Uses async.vim and asyncomplete.vim for non-blocking operation.


\vspace{16pt}

\section{Vim Sessions \& Workflow}


\vspace{10pt}
{\large\bfseries\color{neoncyan} VIM SESSIONS \& WORKFLOW}\vspace{6pt}

{\small\itshape\color{dimtext} // your vim state, frozen in amber. thaw it when you need it. //}

Vim sessions save your entire editing state: open buffers, window layouts, cursor positions. \textbf{ZPWR} integrates sessions with tmux for a seamless development environment.


\vspace{12pt}
{\bfseries\color{neonyellow} VIM SESSION MANAGEMENT}\vspace{4pt}

\textbf{ZPWR} configures vim sessions as manual-save only:

\begin{codebox}
let g:session\_autosave = 'no'\\
let g:session\_autoload = 'no'\\
\end{codebox}

No surprises. Sessions save when you tell them to. Save: Ctrl-D q (mapped to :SaveSession! In normal/insert mode).


\vspace{12pt}
{\bfseries\color{neonyellow} CONFIG SESSIONS: vrc, brc, zrc, trc}\vspace{4pt}

\textbf{ZPWR} defines verbs for editing its own config files: vrc   Open .vimrc in vim brc   Open shell aliases file (.zpwr\_alias.sh) zrc   Open .zshrc in vim trc   Open tmux.conf in vim The tmux layout script config-files.conf opens all four in a 4-pane arrangement. Prefix+T triggers it. Each pane runs one of these verbs automatically.


\vspace{12pt}
{\bfseries\color{neonyellow} THE AUXILIARY VIM PLUGINS}\vspace{4pt}

Vim-autotag           Auto-regenerate ctags on save.

\begin{codebox}
Plugin 'craigemery/vim-autotag'\\
\end{codebox}

G:autotagStartMethod='fork' -- non-blocking tag generation. Vim-online-thesaurus   Look up synonyms (visual mode \textbackslash{}K)

\begin{codebox}
Plugin 'beloglazov/vim-online-thesaurus'\\
\end{codebox}

Vim-zsh-completion     Zsh syntax completion in vim.

\begin{codebox}
Plugin 'Valodim/vim-zsh-completion'\\
\end{codebox}

Vim-minimap            Sublime-style code minimap.

\begin{codebox}
Plugin 'severin-lemaignan/vim-minimap'\\
\end{codebox}


\vspace{12pt}
{\bfseries\color{neonyellow} THE zpwr snapshot/restore INTEGRATION}\vspace{4pt}

\textbf{ZPWR} has verbs for full environment snapshots:

\begin{codebox}
zpwr snapshot    Capture terminal environment state\\
zpwr restore     Restore from snapshot\\
\end{codebox}

ZpwrRestore rebuilds shell state: directory stack, aliases, parameters, and more. This complements tmux-resurrect: resurrect restores the tmux layout, snapshot/restore restores the shell state within each pane.


\vspace{12pt}
{\bfseries\color{neonyellow} THE FULL \textbf{ZPWR} VIM WORKFLOW}\vspace{4pt}

1. Open a tmux session with your project layout 2. Vim in left pane, shell in right pane 3. Edit code in vim, Ctrl-V to execute in right pane 4. Visual select + leader-vl to paste into a REPL 5. tmux-complete.vim grabs words from the REPL output 6. Auto-save means you never lose work 7. Prefix+Ctrl-s saves the entire tmux state This is not an IDE. This is a weapon system.


\vspace{16pt}


\chapter{TEMPRS — TEMPORARY FILE STACK MANAGER}


\begin{center}
\includegraphics[width=0.85\textwidth]{/Users/wizard/.zpwr/docs/screenshots/temprs_help.png}
\end{center}

\section{temprs: The Stack Machine}


\begin{center}
\includegraphics[width=0.8\textwidth]{/Users/wizard/.zpwr/docs/screenshots/temprs_help.png}
\end{center}


\vspace{10pt}
{\large\bfseries\color{neoncyan} TEMPRS: THE STACK MACHINE}\vspace{6pt}

{\small\itshape\color{dimtext} // your filesystem is a stack. you just didn't know it yet. //}

temprs is a Rust-based temporary file stack manager by MenkeTechnologies. Think of it as a LIFO/FIFO clipboard that persists across shell sessions using actual tempfiles on disk. Version 2.9.3. Every operation is instant. No databases, no daemons — just a master record file and a directory of tempfiles.


\vspace{12pt}
{\bfseries\color{neonyellow} THE STACK PARADIGM}\vspace{4pt}

temprs models your scratch data as a stack with both ends open: push (default), pop, shift, unshift — like a Perl array on disk.


\vspace{12pt}
{\bfseries\color{neonyellow} CORE OPERATIONS}\vspace{4pt}

temprs              \# no args: create new tempfile, push to top temprs -o INDEX     \# read contents of tempfile at INDEX to stdout temprs -i INDEX     \# write stdin into tempfile at INDEX temprs -p           \# pop from top of stack (print + remove) temprs -s           \# shift from bottom of stack (print + remove) temprs -u           \# unshift: push stdin to bottom (silent) temprs -a INDEX     \# insert new tempfile at INDEX temprs -r INDEX     \# remove tempfile at INDEX.


\vspace{12pt}
{\bfseries\color{neonyellow} THE MASTER RECORD}\vspace{4pt}

temprs keeps a master record file that maps stack positions to actual tempfile paths on disk. The temprs directory holds the files themselves. Both are in your system's temp directory. temprs -d           \# list the temprs directory contents temprs -m           \# list the master record file path.


\vspace{12pt}
{\bfseries\color{neonyellow} PIPELINE INTEGRATION}\vspace{4pt}

temprs reads stdin when available — perfect for pipelines:

\begin{codebox}
echo 'stash this' | temprs -u  \# silent push to bottom\\
curl -s api.io/data | temprs    \# push response to top\\
temprs -o 0 | jq .             \# read top, pipe to jq\\
\end{codebox}

The -u unshift flag reads stdin and produces no stdout — ideal for mid-pipeline stashing without breaking the flow.


\vspace{12pt}
{\bfseries\color{neonyellow} POSITIONAL ARGUMENT}\vspace{4pt}

temprs myfile.txt   \# read file into new tempfile on stack If both stdin and a file arg are present, stdin takes priority. Your data flows in, gets a stack address, and waits for recall.


\vspace{16pt}

\section{temprs: Advanced Operations}


\begin{center}
\includegraphics[width=0.8\textwidth]{/Users/wizard/.zpwr/docs/screenshots/temprs_help.png}
\end{center}


\vspace{10pt}
{\large\bfseries\color{neoncyan} TEMPRS: ADVANCED OPERATIONS}\vspace{6pt}

{\small\itshape\color{dimtext} // the stack goes deeper. every tempfile has a story. //}


\vspace{12pt}
{\bfseries\color{neonyellow} LISTING \& INSPECTION}\vspace{4pt}

temprs -l           \# list all tempfile paths on the stack temprs -n           \# list numbered (index + path) temprs -L           \# list all tempfile contents temprs -N           \# list numbered with full contents temprs -k           \# count tempfiles on the stack temprs -I INDEX     \# show metadata for tempfile at INDEX.


\vspace{12pt}
{\bfseries\color{neonyellow} EDITING \& TAGGING}\vspace{4pt}

temprs -e INDEX     \# open tempfile at INDEX in \$\textbf{EDITOR} temprs -w NAME      \# tag a tempfile with an alias name temprs -R OLD NEW   \# rename tag from OLD to NEW Tags let you retrieve tempfiles by name instead of index. Your stack positions shift; your tags don't.


\vspace{12pt}
{\bfseries\color{neonyellow} SEARCH \& CONCATENATE}\vspace{4pt}

temprs -g PATTERN   \# grep through all tempfile contents temprs -C 0 2 5    \# concatenate tempfiles at indices to stdout The grep flag searches across every tempfile on the stack — like a personal knowledge base with instant full-text search.


\vspace{12pt}
{\bfseries\color{neonyellow} STACK MANIPULATION}\vspace{4pt}

temprs -D 0 1      \# diff two tempfiles by index or name temprs -M 0 3      \# move tempfile from position 0 to 3 temprs -x INDEX     \# duplicate tempfile onto top of stack temprs -S 0 2      \# swap two tempfiles by index or name temprs --rev        \# reverse the entire stack order temprs -A INDEX     \# append stdin to existing tempfile.


\vspace{12pt}
{\bfseries\color{neonyellow} HOUSEKEEPING}\vspace{4pt}

temprs -c           \# clear/purge ALL tempfiles from stack temprs --expire 24  \# purge tempfiles older than 24 hours temprs --wc INDEX   \# print line count of tempfile temprs --size INDEX \# print byte size of tempfile temprs --path INDEX \# print the actual filesystem path temprs --sort mtime \# sort stack by name, size, or mtime.


\vspace{12pt}
{\bfseries\color{neonyellow} MODE FLAGS}\vspace{4pt}

temprs -q           \# quiet mode: suppress output on create temprs -v           \# verbose mode: extra diagnostic output.


\vspace{12pt}
{\bfseries\color{neonyellow} THE CYBERPUNK -h}\vspace{4pt}

Run temprs -h for the full ASCII art banner with the TEMPRS logo, status bar, and signal meter. Because even help text should look like a console from a Gibson novel. Also: --head INDEX N and --tail INDEX N for partial reads, and --replace INDEX PAT REPL for in-place sed-style edits.


\vspace{16pt}


\chapter{LSOFRS — ENHANCED LSOF}


\begin{center}
\includegraphics[width=0.85\textwidth]{/Users/wizard/.zpwr/docs/screenshots/lsofrs_help.png}
\end{center}

\section{lsofrs: Network Intelligence}


\begin{center}
\includegraphics[width=0.8\textwidth]{/Users/wizard/.zpwr/docs/screenshots/lsofrs_help.png}
\end{center}


\vspace{10pt}
{\large\bfseries\color{neoncyan} LSOFRS: NETWORK INTELLIGENCE}\vspace{6pt}

{\small\itshape\color{dimtext} // every open file descriptor is a signal. lsofrs reads them all. //}

lsofrs is a Rust-based enhanced lsof wrapper v4.7.0 by MenkeTechnologies. It wraps the venerable lsof with colorized output, structured formatting, and cyberpunk-grade aesthetics.


\vspace{12pt}
{\bfseries\color{neonyellow} WHY NOT RAW LSOF?}\vspace{4pt}

Raw lsof dumps walls of monochrome text. lsofrs adds color coding by type, structured columns, and filtering that doesn't require piping through 5 layers of awk and grep.


\vspace{12pt}
{\bfseries\color{neonyellow} SELECTION FLAGS}\vspace{4pt}

lsofrs -i [ADDR]   \# select internet connections ADDR format: [4|6][proto][@host|addr][:svc|port] lsofrs -c COMMAND   \# select by command name (prefix match) lsofrs -p PID       \# select by PID (comma-separated, \textasciicircum{}excludes) lsofrs -u \textbf{USER}      \# select by login name or UID lsofrs -d FD        \# select by file descriptor set lsofrs -g [PGID]    \# select by process group ID lsofrs -a           \# AND selections (default is OR)


\vspace{12pt}
{\bfseries\color{neonyellow} NETWORK FLAGS}\vspace{4pt}

lsofrs -n           \# inhibit hostname resolution (fast mode) lsofrs -P           \# inhibit port-to-name conversion lsofrs -N           \# select NFS files lsofrs -U           \# select UNIX domain socket files.


\vspace{12pt}
{\bfseries\color{neonyellow} \textbf{ZPWR} INTEGRATION}\vspace{4pt}

The zpwr lsof verb pipes lsof output through fzf for interactive process selection and killing:

\begin{codebox}
zpwr lsof           \# fzf over lsof, select PIDs to kill\\
zpwr lsofedit       \# same but edit the kill command first\\
\end{codebox}

The zpwrLsoffzf widget (bound to \textasciicircum{}F\textasciicircum{}H in viins and vicmd) runs sudo lsof -i | fzf -m, extracts PIDs with awk, and inserts them into your command line. Scan the local network like a deck jockey, pick your targets with fzf.


\vspace{12pt}
{\bfseries\color{neonyellow} FILES \& DIRECTORIES}\vspace{4pt}

lsofrs FILE...      \# list processes using these files lsofrs --dir DIR    \# open files in DIR (one level, like +d) lsofrs --dir-recurse DIR \# recursive (like +D)


\vspace{16pt}

\section{lsofrs: Practical Usage}


\vspace{10pt}
{\large\bfseries\color{neoncyan} LSOFRS: PRACTICAL USAGE}\vspace{6pt}

{\small\itshape\color{dimtext} // who's listening on port 8080? lsofrs knows. lsofrs always knows. //}


\vspace{12pt}
{\bfseries\color{neonyellow} FINDING WHAT'S ON A PORT}\vspace{4pt}

\begin{codebox}
lsofrs -i :8080        \# anything on port 8080\\
lsofrs -i TCP:443      \# TCP connections on 443\\
lsofrs -i 4            \# all IPv4 connections\\
lsofrs -i 6TCP         \# IPv6 TCP only\\
\end{codebox}


\vspace{12pt}
{\bfseries\color{neonyellow} \textbf{DISPLAY} \& OUTPUT}\vspace{4pt}

lsofrs -F [FIELDS]  \# machine-readable output fields lsofrs -J / --json  \# JSON output for piping to jq lsofrs -R           \# list parent PIDs lsofrs -t           \# terse: PID only (for scripting) lsofrs -0           \# NUL field terminator (xargs -0 friendly) lsofrs --color MODE \# auto, always, or never.


\vspace{12pt}
{\bfseries\color{neonyellow} LIVE MONITORING}\vspace{4pt}

lsofrs +r 2         \# repeat every 2 seconds lsofrs -W / --monitor \# live full-screen like top lsofrs --delta       \# highlight new/gone FDs in repeat mode lsofrs --follow PID  \# watch one process, highlight opens/closes lsofrs --top 20     \# live top-N processes by FD count lsofrs --watch FILE  \# watch who opens/closes a file over time.


\vspace{12pt}
{\bfseries\color{neonyellow} DIAGNOSTICS}\vspace{4pt}

lsofrs --leak-detect \# detect FD leaks (poll every 5s, 3 rounds) lsofrs --stale       \# find FDs pointing to deleted files lsofrs --ports       \# listening ports summary (like ss -tlnp) lsofrs --tree        \# process tree with FD counts (pstree + lsof) lsofrs --summary     \# aggregate FD summary by type and user.


\vspace{12pt}
{\bfseries\color{neonyellow} COMPARISON}\vspace{4pt}

lsof    — the OG. Every Unix has it. Output is a wall of text. ss      — socket statistics. Fast but sockets only, no files. netstat — deprecated on Linux. Still works on macOS/BSD. lsofrs  — wraps lsof with color, JSON, live monitoring, leak detection, tree views, and the cyberpunk banner you get with lsofrs -h. Built for operators who stare at FD tables and like what they see.


\vspace{12pt}
{\bfseries\color{neonyellow} QUICK RECIPES}\vspace{4pt}

\begin{codebox}
lsofrs -t -i :3000 | xargs kill  \# kill whatever's on 3000\\
lsofrs -c node -a -i TCP          \# node's TCP connections\\
lsofrs -u (whoami) -a -i         \# your network connections\\
\end{codebox}


\vspace{16pt}


\chapter{EZA — THE MODERN LS}


\begin{center}
\includegraphics[width=0.85\textwidth]{/Users/wizard/.zpwr/docs/screenshots/eza_listing.png}
\end{center}

\section{eza: ls Evolved}


\begin{center}
\includegraphics[width=0.8\textwidth]{/Users/wizard/.zpwr/docs/screenshots/eza_listing.png}
\end{center}


\vspace{10pt}
{\large\bfseries\color{neoncyan} EZA: LS EVOLVED}\vspace{6pt}

{\small\itshape\color{dimtext} // ls is dead. long live ls. //}

eza is the modern Rust replacement for ls, formerly known as exa. Maintained by eza-community after exa's author went AWOL. In zpwr, it's the default file lister for everything.


\vspace{12pt}
{\bfseries\color{neonyellow} \$\textbf{ZPWR\_EXA}\_COMMAND — THE CONFIGURED INVOCATION}\vspace{4pt}

\begin{codebox}
command eza --git -il --classify=always -H --color-scale -g -a --colour=always\\
\end{codebox}

Every flag earns its place: --git              \# git status per file (M modified, N new, I ignored) -i                 \# show inode numbers -l                 \# long listing format --classify=always  \# / after dirs, * after executables, always -H                 \# show hard link count --color-scale      \# gradient colors for file sizes and ages -g                 \# show group ownership -a                 \# show hidden files (dotfiles) --colour=always    \# force color even when piped.


\vspace{12pt}
{\bfseries\color{neonyellow} WHY --classify=always?}\vspace{4pt}

Not just -F. Under p10k instant prompt, stdout gets redirected during init. eza's --classify auto-detects terminal capabilities and when stdout is redirected, it drops the indicators. The =always suffix forces classification regardless of fd state. Without it, your login banner loses all the / and * markers.


\vspace{12pt}
{\bfseries\color{neonyellow} WHERE IT'S USED}\vspace{4pt}

ZpwrClearList and zpwrListNoClear both eval \$\textbf{ZPWR\_EXA}\_COMMAND. The fallback chain: eza -> grc + gls -> raw ls. zpwrListNoClear checks for eza first, falls back to gls on macOS (with grc for colorization) or plain ls on Linux.

\begin{codebox}
alias eza="ZPWR\_EXA\_COMMAND"  \# bound in zpwrBindAliasesLate\\
alias lr="ZPWR\_EXA\_COMMAND -R --tree | less -rMN"\\
\end{codebox}

The lr alias gives you a recursive tree piped through less with line numbers — a full filesystem X-ray.


\vspace{12pt}
{\bfseries\color{neonyellow} EXTENDED ATTRIBUTES}\vspace{4pt}

Set \textbf{ZPWR\_EXA}\_EXTENDED=true in .zpwr\_env.sh to add --extended, which shows xattr metadata on macOS (com.apple.quarantine etc).


\vspace{16pt}

\section{eza: Colors \& Integration}


\begin{center}
\includegraphics[width=0.8\textwidth]{/Users/wizard/.zpwr/docs/screenshots/eza_listing.png}
\end{center}


\vspace{10pt}
{\large\bfseries\color{neoncyan} EZA: COLORS \& INTEGRATION}\vspace{6pt}

{\small\itshape\color{dimtext} // a file listing without color is just noise. //}


\vspace{12pt}
{\bfseries\color{neonyellow} EZA\_COLORS ENVIRONMENT VARIABLE}\vspace{4pt}

eza uses EZA\_COLORS (or EXA\_COLORS for compat) to customize colors beyond the defaults. Format is colon-separated key=value:

\begin{codebox}
export EZA\_COLORS="di=34:ln=35:so=32:pi=33:ex=31"\\
\end{codebox}

Keys match LS\_COLORS extensions. eza also adds its own: da=date, sn=file size number, sb=file size unit, uu=current user, un=other user, gu=current group.


\vspace{12pt}
{\bfseries\color{neonyellow} LS\_COLORS COMPATIBILITY}\vspace{4pt}

eza reads LS\_COLORS as a fallback. If you have a vivid or dircolors setup, eza inherits it. EZA\_COLORS overrides.


\vspace{12pt}
{\bfseries\color{neonyellow} TREE VIEW}\vspace{4pt}

\begin{codebox}
eza --tree          \# recursive tree (like tree command)\\
eza --tree -L 3     \# limit depth to 3 levels\\
eza --tree --git-ignore \# respect .gitignore in tree\\
\end{codebox}

The tree integrates with --git, so you see git status markers inline with the tree structure. New files glow green.


\vspace{12pt}
{\bfseries\color{neonyellow} THE FALLBACK CHAIN IN zpwrListNoClear}\vspace{4pt}

1. If eza exists: eval \$\textbf{ZPWR\_EXA}\_COMMAND 2. macOS + grc: grc -c \textbf{\textasciitilde{}/conf.gls} gls -iFlhA --color=always 3. macOS bare: ls -iFlhAO (O = macOS file flags) 4. Linux + grc: grc -c \textbf{\textasciitilde{}/conf.gls} ls -iFlhA --color=always 5. Linux bare: ls -ifhlA grc (Generic Colouriser) wraps ls output with conf.gls rules. It's the middle ground between raw ls and eza.


\vspace{12pt}
{\bfseries\color{neonyellow} \textbf{ZPWR} VERBS USING EZA}\vspace{4pt}

ZpwrClearList clears the terminal then runs the listing. It's called on every cd via the chpwd hook — so you always see what's in the directory you just landed in. zpwrListNoClear does the same listing without clearing.


\vspace{12pt}
{\bfseries\color{neonyellow} SORTING \& FILTERING}\vspace{4pt}

\begin{codebox}
eza --sort=modified  \# sort by modification time\\
eza --sort=size      \# sort by file size\\
eza -r               \# reverse sort order\\
eza --group-directories-first \# dirs at top\\
eza -I '*.pyc|\_\_pycache\_\_'   \# ignore patterns\\
\end{codebox}

Combine with --git-ignore to hide everything in .gitignore. Your file listings become as curated as your code.


\vspace{16pt}


\chapter{BAT — CAT WITH WINGS}


\begin{center}
\includegraphics[width=0.85\textwidth]{/Users/wizard/.zpwr/docs/screenshots/bat_zshrc.png}
\end{center}

\section{bat: Syntax Highlighting for Everything}


\begin{center}
\includegraphics[width=0.8\textwidth]{/Users/wizard/.zpwr/docs/screenshots/bat_sample.png}
\end{center}


\vspace{10pt}
{\large\bfseries\color{neoncyan} BAT: SYNTAX HIGHLIGHTING FOR EVERYTHING}\vspace{6pt}

{\small\itshape\color{dimtext} // cat is for reading files. bat is for understanding them. //}

bat is a cat clone with syntax highlighting, git integration, and automatic paging. Written in Rust. In zpwr it's the default colorizer: \textbf{ZPWR\_COLORIZER}=bat in .zpwr\_env.sh.


\vspace{12pt}
{\bfseries\color{neonyellow} CORE FEATURES}\vspace{4pt}

Line numbers in the gutter. Git diff markers (+/-/\textasciitilde{}) in the margin. Syntax highlighting for 200+ languages. Automatic paging via less when output exceeds terminal height.


\vspace{12pt}
{\bfseries\color{neonyellow} DISCOVERY}\vspace{4pt}

bat --list-themes    \# see all available color themes bat --list-languages \# every syntax bat can highlight.


\vspace{12pt}
{\bfseries\color{neonyellow} HOW \textbf{ZPWR} USES BAT}\vspace{4pt}

bat is the preview engine for fzf throughout zpwr:

\begin{codebox}
fzf --preview 'bat --color=always \{\}'\\
\end{codebox}

The full zpwr colorizer invocation:

\begin{codebox}
bat --paging never --wrap character --color always\\
--style="numbers,grid,changes,header"\\
\end{codebox}

This is stored in \$\textbf{ZPWR\_COLORIZER} and expanded into: FZF\_COLORIZER          — preview for generic fzf FZF\_COLORIZER\_FILE     — preview for file-specific fzf FZF\_COLORIZER\_FILE\_TEXT — text file preview (uses -l ASP) When bat is unavailable, zpwr falls back to pygmentize.


\vspace{12pt}
{\bfseries\color{neonyellow} \textbf{ZPWR} VERBS}\vspace{4pt}

zpwr filesearch and zpwr wordsearch both use bat in their fzf preview windows. Vimfilesearch, vimwordsearch, and the emacs variants (emacsfilesearch, emacswordsearch) all inherit the same bat-powered preview pipeline.


\vspace{12pt}
{\bfseries\color{neonyellow} STANDALONE USAGE}\vspace{4pt}

\begin{codebox}
bat src/main.rs              \# highlighted source view\\
bat -l json data.txt         \# force language detection\\
bat --diff                    \# show only git-changed lines\\
bat -A                        \# show non-printable chars\\
bat --plain                   \# no line numbers, no border\\
bat -r 10:20 file.zsh        \# show only lines 10-20\\
\end{codebox}

The -r (--line-range) flag is perfect for surgical code review.


\vspace{16pt}

\section{bat: Configuration \& Tricks}


\begin{center}
\includegraphics[width=0.8\textwidth]{/Users/wizard/.zpwr/docs/screenshots/bat_zshrc.png}
\end{center}


\vspace{10pt}
{\large\bfseries\color{neoncyan} BAT: CONFIGURATION \& TRICKS}\vspace{6pt}

{\small\itshape\color{dimtext} // configure your cat. teach it new tricks. make it purr. //}


\vspace{12pt}
{\bfseries\color{neonyellow} BAT\_THEME}\vspace{4pt}

\begin{codebox}
export BAT\_THEME="Dracula"\\
\end{codebox}

Set this in your env to change the default theme globally. Popular themes: Dracula, Monokai Extended, Nord, OneHalfDark, Solarized (Dark), TwoDark, ansi. Preview them all:

\begin{codebox}
bat --list-themes | fzf --preview='bat --theme=\{\} src/main.rs'\\
\end{codebox}


\vspace{12pt}
{\bfseries\color{neonyellow} STYLE COMPONENTS}\vspace{4pt}

\begin{codebox}
bat --style=numbers,changes,header\\
\end{codebox}

Components: full, auto, plain, numbers, changes, header, header-filename, header-filesize, grid, rule, snip. Combine with commas. 'full' enables everything. 'plain' disables everything — just colored text, no chrome.


\vspace{12pt}
{\bfseries\color{neonyellow} BAT AS MANPAGE COLORIZER}\vspace{4pt}

\begin{codebox}
export MANPAGER="sh -c 'col -bx | bat -l man -p'"\\
\end{codebox}

This makes every man page syntax-highlighted. The -l man forces manpage language detection, -p is --plain (no line numbers cluttering man output). Col -bx strips backspace formatting.


\vspace{12pt}
{\bfseries\color{neonyellow} SYNTAX CACHE}\vspace{4pt}

\begin{codebox}
bat cache --build    \# rebuild syntax/theme cache\\
bat cache --clear    \# clear and rebuild from scratch\\
\end{codebox}

After adding custom .sublime-syntax files to \textbf{\textasciitilde{}/.config/bat/syntaxes/}, rebuild the cache to pick them up.


\vspace{12pt}
{\bfseries\color{neonyellow} DEBIAN/UBUNTU GOTCHA}\vspace{4pt}

On Debian-based systems, the binary is batcat not bat due to a name conflict with an unrelated package. Fix:

\begin{codebox}
alias bat=batcat     \# or symlink: ln -s batcat \textasciitilde{}/.local/bin/bat\\
\end{codebox}


\vspace{12pt}
{\bfseries\color{neonyellow} BAT + DELTA FOR GIT DIFFS}\vspace{4pt}

delta is a syntax-highlighting pager for git diffs. It uses bat's syntax themes internally. Configure in .gitconfig:

\begin{codebox}
{[}core] pager = delta\\
{[}delta] syntax-theme = Dracula\\
\end{codebox}

delta adds line numbers, side-by-side view, and word-level diff highlighting. bat does files, delta does diffs.


\vspace{12pt}
{\bfseries\color{neonyellow} BAT\_PAGER}\vspace{4pt}

\begin{codebox}
export BAT\_PAGER="less -RFX"\\
\end{codebox}

Customize the pager bat uses. -R for color, -F to quit if one screen, -X to not clear screen on exit. Or set to '' to disable paging entirely — zpwr does this with --paging never.


\vspace{16pt}


\chapter{FD-FIND — FIND REPLACEMENT}


\begin{center}
\includegraphics[width=0.85\textwidth]{/Users/wizard/.zpwr/docs/screenshots/fd_zsh_files.png}
\end{center}

\section{fd: find for Humans}


\begin{center}
\includegraphics[width=0.8\textwidth]{/Users/wizard/.zpwr/docs/screenshots/fd_help.png}
\end{center}


\vspace{10pt}
{\large\bfseries\color{neoncyan} FD: FIND FOR HUMANS}\vspace{6pt}

{\small\itshape\color{dimtext} // find is powerful. fd is powerful AND usable. //}

fd is a fast, user-friendly alternative to find, written in Rust by David Peter (sharkdp — also the author of bat). It respects .gitignore by default and uses regex, not globs.


\vspace{12pt}
{\bfseries\color{neonyellow} BASIC USAGE}\vspace{4pt}

\begin{codebox}
fd PATTERN           \# search current dir recursively\\
fd PATTERN path/     \# search specific directory\\
fd '\textasciicircum{}test.*\textbackslash{}.zsh'  \# full regex support\\
\end{codebox}

No -name, no -type f, no quoting gymnastics. Just the pattern.


\vspace{12pt}
{\bfseries\color{neonyellow} KEY DIFFERENCES FROM FIND}\vspace{4pt}

1. Regex by default (find uses glob with -name) 2. Ignores .git, node\_modules, .gitignore entries automatically 3. 5-10x faster than find due to parallel directory traversal 4. Smart case: case-insensitive unless pattern has uppercase 5. Colorized output by default


\vspace{12pt}
{\bfseries\color{neonyellow} COMMON FLAGS}\vspace{4pt}

fd -e zsh           \# filter by extension fd -t f             \# files only fd -t d             \# directories only fd -t l             \# symlinks only fd -t x             \# executables only fd -H               \# include hidden files fd -I               \# don't respect .gitignore fd -u               \# unrestricted (-H -I combined) fd -g GLOB          \# use glob pattern instead of regex.


\vspace{12pt}
{\bfseries\color{neonyellow} EXECUTION}\vspace{4pt}

fd -x cmd \{\}        \# execute command per result (parallel) fd -X cmd\{\}         \# execute once with all results as args.

{\small\itshape\color{dimtext} Placeholders: \{\} = path, \{/\} = basename, \{//\} = parent dir,}

\{.\} = Path without extension, \{/.\} = basename without extension.

\begin{codebox}
fd -e jpg -x convert \{\} \{.\}.png  \# convert all jpgs to png\\
\end{codebox}

The -x flag runs in parallel by default — your cores earn their keep. -X collects all results and passes them as one batch. This is the find | xargs pattern without the pipe.


\vspace{16pt}

\section{fd: Integration with zpwr \& fzf}


\begin{center}
\includegraphics[width=0.8\textwidth]{/Users/wizard/.zpwr/docs/screenshots/fd_zsh_files.png}
\end{center}


\vspace{10pt}
{\large\bfseries\color{neoncyan} FD: INTEGRATION WITH ZPWR \& FZF}\vspace{6pt}

{\small\itshape\color{dimtext} // fd feeds fzf. fzf feeds your brain. the pipeline of knowing. //}


\vspace{12pt}
{\bfseries\color{neonyellow} FZF\_DEFAULT\_COMMAND}\vspace{4pt}

zpwr sets FZF\_DEFAULT\_COMMAND to use find + ag by default:

\begin{codebox}
export FZF\_DEFAULT\_COMMAND='find * | ag -v ".git/"'\\
\end{codebox}

You can override this with fd for faster results:

\begin{codebox}
export FZF\_DEFAULT\_COMMAND='fd --type f --hidden --follow'\\
\end{codebox}

fd is faster because it parallelizes traversal and skips .gitignore entries without needing ag to filter them out.


\vspace{12pt}
{\bfseries\color{neonyellow} FD + FZF PATTERNS}\vspace{4pt}

\begin{codebox}
fd -e zsh | fzf                 \# fuzzy find zsh files\\
fd -t f | fzf --preview 'bat \{\}'  \# find + preview\\
vim (fd -t f | fzf)             \# find, select, edit\\
\end{codebox}

This three-tool combo (fd + fzf + bat) is the holy trinity of file navigation. fd finds, fzf filters, bat previews.


\vspace{12pt}
{\bfseries\color{neonyellow} TEMPORAL \& SIZE FILTERS}\vspace{4pt}

fd --changed-within 1d  \# modified in last 24 hours fd --changed-before 1w  \# modified more than a week ago fd --size +1m           \# files larger than 1 megabyte fd --size -10k          \# files smaller than 10 kilobytes Units: b (bytes), k (kilobytes), m (megabytes), g (gigabytes). Time: s (seconds), m (minutes), h (hours), d (days), w (weeks).


\vspace{12pt}
{\bfseries\color{neonyellow} FD + BAT}\vspace{4pt}

\begin{codebox}
fd -e rs --exec bat\{\}   \# colorized find results\\
fd -e md -X bat          \# cat all markdown files at once\\
\end{codebox}


\vspace{12pt}
{\bfseries\color{neonyellow} FD vs FIND vs LOCATE}\vspace{4pt}

Find   — universal. Every Unix has it. Slow on large trees. Supports -exec, -newer, -perm, -prune. The Swiss Army. fd     — fast parallel traversal, sane defaults, .gitignore aware. Best for daily use in code repositories. Locate — instant results from a pre-built database (updatedb). Stale until the next cron run. No real-time accuracy.


\vspace{12pt}
{\bfseries\color{neonyellow} DEPTH \& EXCLUSION}\vspace{4pt}

fd -d 3              \# max depth of 3 levels fd -E node\_modules   \# exclude a directory by name fd -E '*.min.js'     \# exclude by glob pattern Stack -E flags to build precise exclusion lists on the fly.


\vspace{16pt}


\chapter{RIPGREP — GREP AT LIGHT SPEED}


\begin{center}
\includegraphics[width=0.85\textwidth]{/Users/wizard/.zpwr/docs/screenshots/rg_search.png}
\end{center}

\section{ripgrep: grep Reimagined}


\begin{center}
\includegraphics[width=0.8\textwidth]{/Users/wizard/.zpwr/docs/screenshots/rg_help.png}
\end{center}


\vspace{10pt}
{\large\bfseries\color{neoncyan} RIPGREP: GREP REIMAGINED}\vspace{6pt}

{\small\itshape\color{dimtext} // grep -r is a campfire. ripgrep is a directed energy weapon. //}

Ripgrep (rg) is a blazing-fast recursive grep written in Rust by Andrew Gallant (BurntSushi). It respects .gitignore, skips binary files, and uses Rust's regex engine for raw speed. 5-100x faster than grep -r on large codebases. Not a typo.


\vspace{12pt}
{\bfseries\color{neonyellow} BASIC USAGE}\vspace{4pt}

\begin{codebox}
rg PATTERN           \# search current dir recursively\\
rg PATTERN path/     \# search specific directory\\
rg 'fn main'         \# literal or regex — rg auto-detects\\
\end{codebox}

Output is colorized by default: filenames in magenta, line numbers in green, matches highlighted. Grouped by file.


\vspace{12pt}
{\bfseries\color{neonyellow} KEY FLAGS}\vspace{4pt}

rg -i PATTERN       \# case insensitive search rg -l PATTERN       \# files with matches only (paths) rg -c PATTERN       \# count matches per file rg -w PATTERN       \# whole word matches only rg -F PATTERN       \# fixed string (no regex interpretation) rg -v PATTERN       \# invert match (lines NOT matching)


\vspace{12pt}
{\bfseries\color{neonyellow} FILE TYPE FILTERING}\vspace{4pt}

rg -t zsh PATTERN   \# search only zsh files rg -T js PATTERN    \# exclude javascript files rg -g '*.zsh' PAT   \# glob filter (include) rg -g '!*.min.js' PAT \# glob filter (exclude) rg --type-list shows all built-in type definitions.


\vspace{12pt}
{\bfseries\color{neonyellow} CONTEXT \textbf{LINES}}\vspace{4pt}

rg -A 3 PATTERN     \# 3 lines after each match rg -B 3 PATTERN     \# 3 lines before each match rg -C 3 PATTERN     \# 3 lines before AND after.


\vspace{12pt}
{\bfseries\color{neonyellow} MULTILINE MATCHING}\vspace{4pt}

rg -U PATTERN       \# --multiline: match across lines.

\begin{codebox}
rg -U 'struct \textbackslash{}\{[\textbackslash{}s\textbackslash{}S]*?\textbackslash{}\}' \# match entire struct blocks\\
\end{codebox}

Without -U, patterns match within single lines only.


\vspace{12pt}
{\bfseries\color{neonyellow} THE COMPETITION}\vspace{4pt}

grep -r  — universal, slow, no .gitignore, no parallelism ag      — the silver searcher. Fast, .gitignore aware. rg is 2-5x faster on most benchmarks. Ack     — Perl-based, pioneered the "better grep" category. Slower than both ag and rg. Still has loyal users.


\vspace{16pt}

\section{ripgrep: Advanced \& zpwr Integration}


\begin{center}
\includegraphics[width=0.8\textwidth]{/Users/wizard/.zpwr/docs/screenshots/rg_search.png}
\end{center}


\vspace{10pt}
{\large\bfseries\color{neoncyan} RIPGREP: ADVANCED \& ZPWR INTEGRATION}\vspace{6pt}

{\small\itshape\color{dimtext} // rg --json | jq — structured grep for the data-obsessed. //}


\vspace{12pt}
{\bfseries\color{neonyellow} STRUCTURED OUTPUT}\vspace{4pt}

rg --json PATTERN   \# JSON output: one object per line Each line is a JSON object with type (begin, match, end, summary), path, line number, match offsets. Pipe to jq:

\begin{codebox}
rg --json TODO | jq 'select(.type=="match") | .data.path'\\
\end{codebox}


\vspace{12pt}
{\bfseries\color{neonyellow} SEARCH \& REPLACE PREVIEW}\vspace{4pt}

rg --replace 'NEW' PATTERN  \# preview replacements This prints what the output WOULD look like with replacements but does NOT modify files. It's a dry-run. Capture groups work:

\begin{codebox}
rg 'foo(\textbackslash{}w+)' --replace 'bar1'  \# foo\_x -> bar\_x\\
\end{codebox}


\vspace{12pt}
{\bfseries\color{neonyellow} STATISTICS}\vspace{4pt}

rg --stats PATTERN   \# show match stats after results Prints: N matches, N matched lines, N files searched, N files matched, bytes searched, duration. Quick audit.


\vspace{12pt}
{\bfseries\color{neonyellow} CUSTOM FILE TYPES}\vspace{4pt}

\begin{codebox}
rg --type-add 'zpwr:*.zsh' -t zpwr PATTERN\\
\end{codebox}

Define a custom type inline and use it immediately. Persist custom types in \textbf{\textasciitilde{}/.ripgreprc} (set RIPGREP\_CONFIG\_PATH).


\vspace{12pt}
{\bfseries\color{neonyellow} \textbf{ZPWR} INTEGRATION}\vspace{4pt}

zpwr grep uses ag (silver searcher) piped into fzf:

\begin{codebox}
zpwr grep            \# ag into fzf for interactive search\\
\end{codebox}

zpwr wordsearch searches file contents with fzf preview:

\begin{codebox}
zpwr wordsearch      \# search by word in subdirectory files\\
\end{codebox}

Both use bat-powered previews in the fzf window.


\vspace{12pt}
{\bfseries\color{neonyellow} FZF + RG LIVE RELOAD}\vspace{4pt}

The ultimate interactive grep: rg feeds fzf, bat previews:

\begin{codebox}
rg --line-number PATTERN |\\
fzf --delimiter : --preview 'bat --highlight-line \{2\} \{1\}'\\
\end{codebox}

rg outputs file:line:match, fzf splits on :, bat highlights the matching line. Three tools, one seamless experience.


\vspace{12pt}
{\bfseries\color{neonyellow} CONFIGURATION}\vspace{4pt}

\begin{codebox}
export RIPGREP\_CONFIG\_PATH=\textasciitilde{}/.ripgreprc\\
\end{codebox}

Put default flags in this file, one per line: --smart-case, --hidden, --glob=!.git, --max-columns=200 Also: .rgignore in any directory works like .gitignore — project-specific patterns that rg will skip. .ignore files (used by ag too) are also respected by rg.


\vspace{16pt}


\chapter{NEOVIM}


\begin{center}
\includegraphics[width=0.85\textwidth]{/Users/wizard/.zpwr/docs/screenshots/art_neural_link.png}
\end{center}

\section{Neovim: The Next Generation}


\vspace{10pt}
{\large\bfseries\color{neoncyan} NEOVIM: THE NEXT GENERATION}\vspace{6pt}

{\small\itshape\color{dimtext} // vim, but the timeline where Bram had unlimited funding. //}

Neovim is the modernized fork of vim. Refactored internals, async everything, built-in LSP, Lua plugin API. In zpwr, the default editor: \textbf{ZPWR\_VIM}='nvim' in .zpwr\_env.sh.


\vspace{12pt}
{\bfseries\color{neonyellow} THE INIT.VIM BRIDGE}\vspace{4pt}

\textbf{\$ZPWR\_INSTALL/init.vim} is just three lines:

\begin{codebox}
set runtimepath\textasciicircum{}=\textasciitilde{}/.vim runtimepath+=\textasciitilde{}/.vim/after\\
let \&packpath = \&runtimepath\\
source \textasciitilde{}/.vimrc\\
\end{codebox}

This means neovim sources the same .vimrc as vim. One config to rule them both. All 80+ vim plugins work in either editor. The symlink is created by plugins\_install.sh:

\begin{codebox}
ln -sf ZPWR\_INSTALL/init.vim \textasciitilde{}/.config/nvim/init.vim\\
\end{codebox}


\vspace{12pt}
{\bfseries\color{neonyellow} NEOVIM-SPECIFIC FEATURES}\vspace{4pt}

Async job control  — run external commands without blocking the editor. vim 8 added this too, but nvim's API is cleaner (jobstart/jobstop). Built-in LSP       — nvim 0.5+ has a native LSP client. Configure with lua: vim.lsp.start(). No plugin needed for basic LSP. Lua plugin API     — write plugins in Lua, not just VimScript. Lua runs 10-100x faster than VimScript. :terminal          — built-in terminal emulator. Better than vim's :terminal — proper job control, scrollback, and buffer integration. True color         — 24-bit color out of the box. No termguicolors hacks needed (though the option exists).


\vspace{12pt}
{\bfseries\color{neonyellow} PROVIDER SUPPORT}\vspace{4pt}

nvim uses providers for python3 and node integration:

\begin{codebox}
pip3 install pynvim          \# python3 provider\\
npm install -g neovim        \# node.js provider\\
\end{codebox}

Run :checkhealth to verify all providers are working.


\vspace{12pt}
{\bfseries\color{neonyellow} \textbf{ZPWR} VIM INFO}\vspace{4pt}

zpwr catviminfo and zpwr catrecentfviminfo read the nvim shada (shared data) file — the nvim equivalent of .viminfo.


\vspace{16pt}

\section{Neovim: Configuration \& Tips}


\vspace{10pt}
{\large\bfseries\color{neoncyan} NEOVIM: CONFIGURATION \& TIPS}\vspace{6pt}

{\small\itshape\color{dimtext} // one vimrc. two editors. zero compromises. //}


\vspace{12pt}
{\bfseries\color{neonyellow} THE SHARED CONFIG PHILOSOPHY}\vspace{4pt}

zpwr's init.vim just sources \textbf{\textasciitilde{}/.vimrc.} This means: 1. Every plugin in .vimbundle works in both vim and nvim 2. All keybindings, settings, and autocmds are shared 3. One file to edit: zpwr vrc opens .vimrc 4. No divergent configs to maintain across editors The .vimrc is the single source of truth. Init.vim is a shim.


\vspace{12pt}
{\bfseries\color{neonyellow} NEOVIM-SPECIFIC SETTINGS}\vspace{4pt}

Guard nvim-only settings with has('nvim'):

\begin{codebox}
if has('nvim')\\
set inccommand=nosplit   " live substitution preview\\
set pumblend=15          " popup menu transparency\\
set winblend=15          " floating window transparency\\
endif\\
\end{codebox}

Inccommand is a killer feature: :s/old/new shows replacements live as you type. No more guessing what your regex will do.


\vspace{12pt}
{\bfseries\color{neonyellow} :checkhealth}\vspace{4pt}

Run this in nvim to diagnose your setup. It checks: - python3/node providers (pynvim, neovim npm package) - clipboard support (pbcopy on macOS, xclip on Linux) - terminal capabilities, true color support - plugin health (many plugins register health checks) First thing to run when something feels wrong.


\vspace{12pt}
{\bfseries\color{neonyellow} SESSION TRACKING}\vspace{4pt}

Tpope/vim-obsession (in .vimbundle) tracks sessions:

\begin{codebox}
:Obsession           \# start tracking session to Session.vim\\
:Obsession!          \# stop tracking and delete Session.vim\\
nvim -S Session.vim  \# restore session\\
\end{codebox}

Works in both vim and nvim since it's a VimScript plugin.


\vspace{12pt}
{\bfseries\color{neonyellow} WHY \textbf{ZPWR} SUPPORTS BOTH}\vspace{4pt}

nvim is better locally — async, LSP, Lua, true color. But: - Remote servers often only have vim (or vi) - Docker containers ship vim, not nvim - macOS ships vim in \textbf{\textbf{/usr/bin/vim}} The shared .vimrc means your muscle memory works everywhere. nvim when you can, vim when you must, same config either way.


\vspace{12pt}
{\bfseries\color{neonyellow} \textbf{ZPWR} \textbf{EDITOR} VERBS}\vspace{4pt}

\begin{codebox}
zpwr vrc             \# edit .vimrc (shared by both)\\
\end{codebox}

Both editors, one config, zero excuses. The way it should be.


\vspace{16pt}


\chapter{PERL IN ZPWR}


\begin{center}
\includegraphics[width=0.85\textwidth]{/Users/wizard/.zpwr/docs/screenshots/perl_version.png}
\end{center}

\section{Perl Oneliners: The Secret Weapon}


\begin{center}
\includegraphics[width=0.8\textwidth]{/Users/wizard/.zpwr/docs/screenshots/perl_version.png}
\end{center}


\vspace{10pt}
{\large\bfseries\color{neoncyan} PERL ONELINERS: THE SECRET WEAPON}\vspace{6pt}

sed wishes it could do this. awk cries in a corner.

Perl is zpwr's scalpel for text surgery. Where sed and awk stumble on complex patterns, perl oneliners slice clean. zpwr uses perl throughout for text transformation, parsing, delimiter handling, and file manipulation.


\vspace{12pt}
{\bfseries\color{neonyellow} THE KEY FLAGS}\vspace{4pt}

-e 'code'     execute code -n            loop over input lines (like sed) -p            loop and print each line (like sed -p) -l            chomp newlines, add them back on print -a            autosplit into @F array (like awk) -F'sep'       set field separator for -a -i            in-place edit (like sed -i) -E            enable say, state, switch (modern perl)


\vspace{12pt}
{\bfseries\color{neonyellow} WHITESPACE CLEANUP (scripts/lib.sh)}\vspace{4pt}

\begin{codebox}
perl -pi -e 's@\textbackslash{}s+@\textbackslash{}n@g; s@\textbackslash{}x09@    @g'\\
\end{codebox}

Strips trailing whitespace, converts trailing tabs to spaces. The -pi combo: -p loops, -i edits in place. zpwr's file cleaner.


\vspace{12pt}
{\bfseries\color{neonyellow} PRINT SEPARATOR \textbf{LINES}}\vspace{4pt}

\begin{codebox}
perl -le "print '\#'x80"\\
\end{codebox}

Prints 80 \# characters. The x operator repeats a string. Used in zpwrPrettyPrint for visual separators.


\vspace{12pt}
{\bfseries\color{neonyellow} DELIMITER TRANSFORMATION}\vspace{4pt}

\begin{codebox}
perl -F"ZPWR\_DELIMITER\_CHAR" -anE 'say ...'\\
\end{codebox}

Splits lines on \textbf{ZPWR\_DELIMITER}\_CHAR into @F fields. Used heavily in zpwrContribCount, zpwrTotalLines for reformatting git contribution data into tabular output.


\vspace{12pt}
{\bfseries\color{neonyellow} FILTER BLANK \textbf{LINES}}\vspace{4pt}

\begin{codebox}
perl -ne 'print if /\textbackslash{}S/'\\
\end{codebox}

Only prints lines containing non-whitespace. Used in zpwrRegenAllKeybindingsCache to clean keybinding dumps.


\vspace{12pt}
{\bfseries\color{neonyellow} THE PERL REPL ALIAS}\vspace{4pt}

\begin{codebox}
plr (alias for rlwrap perl -wnE'say eval()//@')\\
\end{codebox}

Interactive perl REPL with readline. Type expressions, get results.


\vspace{16pt}

\section{Perl Oneliners: Advanced Patterns}


\vspace{10pt}
{\large\bfseries\color{neoncyan} PERL ONELINERS: ADVANCED PATTERNS}\vspace{6pt}

When you need a chainsaw, not a butter knife.


\vspace{12pt}
{\bfseries\color{neonyellow} \textbf{PATH} QUOTING (zpwrGetFound)}\vspace{4pt}

\begin{codebox}
perl -pe 's@\textasciicircum{}([\textasciitilde{}]*)([\textasciicircum{}\textasciitilde{}].*)@1"2"@;s@\textbackslash{}s+@ @g'\\
\end{codebox}

Wraps non-tilde paths in double quotes and collapses spaces. Handles filenames with spaces for safe shell consumption.


\vspace{12pt}
{\bfseries\color{neonyellow} FZF RESULT FORMATTING (zpwrAgIntoFzf)}\vspace{4pt}

\begin{codebox}
perl -pe 's@\textasciicircum{}(.*?):(\textbackslash{}d+):(.*)@+2 "PWD/1"@'\\
\end{codebox}

Converts grep output (file:line:match) into vim args (+line "file"). The brain behind zpwr grep -> fzf -> vim pipeline.


\vspace{12pt}
{\bfseries\color{neonyellow} VERB LIST FORMATTING (zpwrVerbsFZF)}\vspace{4pt}

\begin{codebox}
perl -e '@a=<>;for(@a)\{print "zpwr 1" if m\{\textasciicircum{}(\textbackslash{}S+)\textbackslash{}s+\}\}'\\
\end{codebox}

Reads verb list, extracts first word, prepends 'zpwr '. Builds the command string that fzf executes on selection.


\vspace{12pt}
{\bfseries\color{neonyellow} IN-PLACE RENAME (zpwrChangeNameZpwr)}\vspace{4pt}

\begin{codebox}
perl -i -pe 's@\textbackslash{}bOLDNAME\textbackslash{}b@NEWNAME@g' **/*(.)'\\
\end{codebox}

Renames all occurrences across all files. The \textbackslash{}b word boundary prevents partial matches. The **(.) glob hits all regular files.


\vspace{12pt}
{\bfseries\color{neonyellow} LOGIN COUNTING (zpwrLoginCount)}\vspace{4pt}

\begin{codebox}
perl -e '"'"'print qx(last -f "\_") for </var/log/wtmp*>'"'"'\\
\end{codebox}

Reads all wtmp files and runs last on each. The diamond operator <> globs and iterates. Then pipes to:

\begin{codebox}
perl -lane '"'"'print [0] if /\textbackslash{}S+/ \&\& !/wtmp/'"'"'\\
\end{codebox}

Extracts username (first field) from non-empty, non-header lines.


\vspace{12pt}
{\bfseries\color{neonyellow} FILE COUNT DEDUP (zpwrEnvCounts)}\vspace{4pt}

\begin{codebox}
perl -lne '@l=<>;@u=do\{my \%seen;grep\{!seen\{\_\}++\}@l\};'\\
\end{codebox}

Reads all lines, deduplicates with a hash, then checks if each unique path is a real file. The densest oneliner in zpwr. Counts unique zpwr files across all tracked locations.


\vspace{12pt}
{\bfseries\color{neonyellow} IP ADDRESS EXTRACTION (ipa alias)}\vspace{4pt}

\begin{codebox}
ifconfig | perl -lane 'do \{print \textbackslash{}[1]=\textasciitilde{}s/addr//r;exit0\} if /inet\textbackslash{}s/ and !/127/'\\
\end{codebox}

{\small\itshape\color{dimtext} Finds the first non-loopback inet address. The s///r flag}

Returns the modified string without changing the original.


\vspace{12pt}
{\bfseries\color{neonyellow} WHY PERL OVER SED/AWK?}\vspace{4pt}

1. Perl regex is more powerful (lookahead, backrefs, /x) 2. -i flag for in-place edit is reliable across OS 3. Complex logic in one line (if/unless/for/do blocks) 4. @F autosplit is awk-like but with perl's full stdlib 5. Perl is installed everywhere (even minimal containers)


\vspace{16pt}


\chapter{HIDDEN SUBSYSTEMS}


\begin{center}
\includegraphics[width=0.85\textwidth]{/Users/wizard/.zpwr/docs/screenshots/art_mainframe.png}
\end{center}

\section{Zconvey: Terminal Communication}


\vspace{10pt}
{\large\bfseries\color{neoncyan} ZCONVEY: TERMINAL COMMUNICATION}\vspace{6pt}

Every terminal screams its identity into the void.


\vspace{12pt}
{\bfseries\color{neonyellow} WHAT IS ZCONVEY?}\vspace{4pt}

MenkeTechnologies/zconvey - a zinit plugin Allows sending commands between terminal sessions. Your terminals can talk to each other. Yes, really. Loaded via zinit in \textbf{\$ZPWR\_INSTALL/.zshplugins}


\vspace{12pt}
{\bfseries\color{neonyellow} HOW \textbf{ZPWR} USES IT}\vspace{4pt}

Every time you press Enter, zpwrAcceptLine fires. It calls zc-rename to label the current terminal with: TTY:\$\{TTY\} PID:\$\{\$\} CMD:\$BUFFER \textbf{PWD}:\$\{PWD\} DATE:\$(date) This broadcasts your terminal's state to the zconvey bus. Every pane, every tab, every session - all labeled.


\vspace{12pt}
{\bfseries\color{neonyellow} THE CONVEY VARIABLES}\vspace{4pt}

\textbf{ZPWR\_CONVEY}\_NAME      the full label string sent to zconvey \textbf{ZPWR\_CONVEY}\_LAST\_CMD  the last command (for zconvey commands) When you run a zc-* command, it preserves the previous label so your terminal doesn't get renamed to "zc-send ls".


\vspace{12pt}
{\bfseries\color{neonyellow} THE RENAME LOGIC}\vspace{4pt}

\begin{codebox}
if ! [[ (zpwrExpandAliases BUFFER) = zc* ]]; then\\
ZPWR\_CONVEY\_NAME="TTY:\{TTY\} PID:\{\} CMD:..."\\
zc-rename ZPWR\_CONVEY\_NAME\\
\end{codebox}

If BUFFER starts with zc*, keep the old label. Otherwise, stamp the terminal with the new command.


\vspace{12pt}
{\bfseries\color{neonyellow} WHY THIS MATTERS}\vspace{4pt}

tmux status bar can show what each pane is running. You can identify panes by their last command and \textbf{PWD}. Remote debugging: see exactly what terminal ran what. It's like process accounting, but for your shell sessions.


\vspace{12pt}
{\bfseries\color{neonyellow} SENDING COMMANDS BETWEEN TERMINALS}\vspace{4pt}

Zc-send 2 'cd /tmp \&\& ls'  send command to terminal 2 zc-all 'echo hello'        broadcast to all terminals Your terminals are a mesh network. Act accordingly. SIGNAL: zpwrAcceptLine is the heartbeat of zpwr. Every Enter keypress: rename terminal, check for clear, detect file-modifying commands, send tmux keys. Zconvey makes that heartbeat visible across sessions.


\vspace{16pt}

\section{The 24-Hour Updater Loop}


\vspace{10pt}
{\large\bfseries\color{neoncyan} THE 24-HOUR UPDATER LOOP}\vspace{6pt}

Your system updates itself while you sleep.


\vspace{12pt}
{\bfseries\color{neonyellow} THE INFINITE LOOP}\vspace{4pt}

\textbf{\$ZPWR\_SCRIPTS/autoUpdater.sh} runs forever. Every 24 hours it wakes up and updates everything.

\begin{codebox}
while true; do\\
bash -l updater.sh -e\\
\# sleep in 10-minute chunks until 24h elapsed\\
done\\
\end{codebox}

Uses Perl Time::Piece to calculate the next update time. Logs the time diff on each wake to \$\textbf{ZPWR\_LOGFILE}.


\vspace{12pt}
{\bfseries\color{neonyellow} WHAT GETS UPDATED}\vspace{4pt}

Updater.sh orchestrates the full update: zpwrUpdater      runs the main zpwr update sequence zpwrUpdateDeps   updates all package manager dependencies Homebrew, npm, pip, gem, go - everything gets refreshed. History files backed up. Logs rotated. Deps pinned.


\vspace{12pt}
{\bfseries\color{neonyellow} SYSTEMD SERVICE (LINUX)}\vspace{4pt}

On Linux, autoUpdater.sh runs as a systemd service. zpwr restart      restart the zpwr systemd service zpwr serviceup    start and enable the systemd service zpwr servicedown  stop and disable the systemd service On macOS, \textbf{\$ZPWR\_SCRIPTS/updaterForLaunchCtl.sh} is used.


\vspace{12pt}
{\bfseries\color{neonyellow} EMAIL NOTIFICATIONS}\vspace{4pt}

\textbf{\$ZPWR\_SCRIPTS/updaterEmail.sh} Pipes the update log to email via mutt. Wake up to a full report of what changed overnight. Your system admin is a bash script. It never sleeps.


\vspace{12pt}
{\bfseries\color{neonyellow} THE SLEEP STRATEGY}\vspace{4pt}

Doesn't sleep for 24h straight (that would drift). Sleeps in 10-minute intervals, checks elapsed time:

\begin{codebox}
((timediff > ((3600 * 24)))) \&\& break\\
\end{codebox}

This handles laptop suspends and clock adjustments. WARNING: The updater keeps zpwr and all deps current. If something breaks after an update, check the log: zpwr taillog  to see what the updater did last. The machine maintains itself. You just write code.


\vspace{16pt}

\section{The 32-Pane REPL Layout}


\vspace{10pt}
{\large\bfseries\color{neoncyan} THE 32-PANE REPL LAYOUT}\vspace{6pt}

Every language, one screen, zero excuses.


\vspace{12pt}
{\bfseries\color{neonyellow} THE TMUX KEYBINDINGS}\vspace{4pt}

Prefix R  source \textbf{\$ZPWR\_TMUX/thirtytwo-panes-repl.conf} prefix O  source \textbf{\$ZPWR\_TMUX/sixteen-panes.conf} prefix G  source \textbf{\$ZPWR\_TMUX/eight-panes.conf} prefix F  source \textbf{\$ZPWR\_TMUX/four-panes.conf} One keystroke splits your terminal into a REPL army.


\vspace{12pt}
{\bfseries\color{neonyellow} THE 32 REPLs (prefix R)}\vspace{4pt}

Pane 0:  lua          pane 1:  swift        pane 2:  ghci pane 3:  clisp        pane 4:  ocaml        pane 5:  clojure pane 6:  pry (ruby)   pane 7:  bpython      pane 8:  php -a pane 9:  groovysh     pane 10: csharp       pane 11: kotlinc pane 12: node         pane 13: coffee       pane 14: tclsh pane 15: ksh          pane 16: bash         pane 17: csh pane 18: perl -dE1    pane 19: gore (go)    pane 20: jshell pane 21: mycli        pane 22: pgcli        pane 23: python2 pane 24: nslookup     pane 25: mongo        pane 26: irb pane 27: scala        pane 28: utop         pane 29: iex pane 30: bc           pane 31: gforth 32 languages. 32 panes. One tmux window. Absolute madness.


\vspace{12pt}
{\bfseries\color{neonyellow} THE VIM-TMUX EXECUTION PIPELINE}\vspace{4pt}

Write code in vim, send it to a REPL pane: Supported languages for vim-to-tmux execution: zsh, bash, python, ruby, perl, node, go java, scala, R, octave Select code in vim, fire it into the matching pane. Edit-run-edit-run without ever leaving vim.


\vspace{12pt}
{\bfseries\color{neonyellow} PROGRESSIVE LAYOUTS}\vspace{4pt}

Don't need 32 languages? Scale down: prefix G  8 panes   - quick polyglot session prefix O  16 panes  - serious multi-language work prefix R  32 panes  - full REPL matrix (needs big screen) Each layout is a tmux conf file in \$\textbf{ZPWR\_TMUX}/ NOTE: prefix R on a 13" laptop is an act of hubris. You'll need at least a 27" monitor to read anything. But hey, it works. And it's glorious.


\vspace{16pt}

\section{Cross-Platform Installer}


\vspace{10pt}
{\large\bfseries\color{neoncyan} CROSS-PLATFORM INSTALLER}\vspace{6pt}

One repo to rule them all, one script to install them.


\vspace{12pt}
{\bfseries\color{neonyellow} INSTALLER SCRIPTS IN \$\textbf{ZPWR\_INSTALL}}\vspace{4pt}

Formulae.sh        Homebrew formulae (macOS/Linux) casks.sh           Homebrew casks (macOS GUI apps) mincasks.sh        minimal cask set for lean installs taps.sh            Homebrew tap repositories npm\_install.sh     Node.js packages via npm npm.sh             npm package list gems.sh            Ruby gems pip\_install.sh     Python packages via pip pip.sh / pip3.sh   pip package lists go\_install.sh      Go packages rustupinstall.sh   Rust toolchain via rustup.


\vspace{12pt}
{\bfseries\color{neonyellow} \textbf{EDITOR} \& PLUGIN INSTALLERS}\vspace{4pt}

Neovim\_install.sh       Neovim setup vim\_install.sh          Vim from source vim\_plugins\_install.sh  Vim plugin manager + plugins zsh\_plugins\_install.sh  Zsh plugins via zinit tmux\_plugins\_install.sh Tmux plugin manager + plugins spacemacs\_install.sh    Spacemacs (Emacs) setup.


\vspace{12pt}
{\bfseries\color{neonyellow} SHARED INFRASTRUCTURE}\vspace{4pt}

Common.sh  shared installer functions used by all scripts zpwrInstall.sh  the master installer orchestrator Supports 20+ Linux distros, macOS, and FreeBSD. Detects your OS and calls the right package manager.


\vspace{12pt}
{\bfseries\color{neonyellow} \textbf{ZPWR} VERBS FOR INSTALLATION}\vspace{4pt}

zpwr install            run configure, make, make install zpwr uninstall          uninstall all zpwr configs zpwr regenconfiglinks   regen symlinks from \$\textbf{ZPWR\_INSTALL} to \$\textbf{HOME} regenconfiglinks recreates dotfile symlinks after a reinstall.


\vspace{12pt}
{\bfseries\color{neonyellow} THE SYMLINK STRATEGY}\vspace{4pt}

zpwr doesn't copy config files to \$\textbf{HOME}. It symlinks them. .vimrc, .tmux.conf, .gitconfig - all point back to \$\textbf{ZPWR\_INSTALL}. Edit once in the repo, git commit, and every machine gets it. zpwr regenconfiglinks rebuilds them if they break. PHILOSOPHY: Package lists are just text files. Add a line to formulae.sh, and your next machine gets it. Your entire dev environment is version-controlled.


\vspace{16pt}


\chapter{QUALITY ASSURANCE}


\begin{center}
\includegraphics[width=0.85\textwidth]{/Users/wizard/.zpwr/docs/screenshots/figlet_trace.png}
\end{center}

\section{The Test Suite}


\vspace{10pt}
{\large\bfseries\color{neoncyan} THE TEST SUITE}\vspace{6pt}

60 test files standing between you and chaos


\vspace{12pt}
{\bfseries\color{neonyellow} TESTING FRAMEWORK}\vspace{4pt}

zpwr uses zunit - a testing framework written in zsh. Config lives in \textbf{\$ZPWR/.zunit.yml} Tests dir: \$\textbf{ZPWR\_TEST}  Output: \textbf{\$ZPWR\_TEST/\_output} Support helpers: \textbf{\$ZPWR\_TEST/\_support}


\vspace{12pt}
{\bfseries\color{neonyellow} TEST CATEGORIES}\vspace{4pt}

T-autoload-*       verify all autoloaded functions load, defined, existence, fpath, whence, type, body Every autoload function must: exist, load, have a body. T-autoload-subdir-* same checks for subdirectory autoloads t-env-*            verify environment variables defined, exported, comprehensive, type checks Every \textbf{ZPWR}\_* var must be defined and exported. T-scripts-*        verify script files executable, readable, non-empty, have shebangs, valid syntax t-verbs-*          verify zpwr verb definitions key format, non-empty values, callable, type checks.


\vspace{12pt}
{\bfseries\color{neonyellow} MORE TEST FILES}\vspace{4pt}

T-functions.zsh     test core shell functions t-callable-fns.zsh  verify functions are callable t-navigation.zsh    test directory navigation t-install-files.zsh verify installer file integrity t-repo-files.zsh    verify repo structure t-plugin-counts.zsh verify plugin load counts.


\vspace{12pt}
{\bfseries\color{neonyellow} RUNNING TESTS}\vspace{4pt}

zpwr test     run all zpwr tests zpwr testall  run ALL env tests (comprehensive sweep) fail\_fast: true means first failure stops the run. Verbose: true means you see every assertion. GUARANTEE: Every autoload loads. Every env var defined. Every script executable. Every verb callable. 60 test files verify zpwr won't break your shell. If tests pass, your environment is sound. Ship it.


\vspace{16pt}


\chapter{COLORIZATION STACK}


\begin{center}

\includegraphics[width=0.85\textwidth]{/Users/wizard/.zpwr/docs/screenshots/art_data_stream.png}
\end{center}

\section{GRC: Generic Colorizer}


\vspace{10pt}
{\large\bfseries\color{neoncyan} GRC: GENERIC COLORIZER}\vspace{6pt}

Because monochrome ls is a crime against humanity.


\vspace{12pt}
{\bfseries\color{neonyellow} WHAT IS GRC?}\vspace{4pt}

grc wraps commands and adds color output via regex rules. It reads a config file that maps patterns to colors. Any command's output can be colorized. Zero modifications.


\vspace{12pt}
{\bfseries\color{neonyellow} GRC IN zpwrListNoClear}\vspace{4pt}

The fallback chain when listing files: 1. eza              preferred (color-scale, git status) 2. grc + ls         fallback (colorized via conf.gls) 3. Plain ls         last resort (no color config)


\vspace{12pt}
{\bfseries\color{neonyellow} PLATFORM-SPECIFIC COMMANDS}\vspace{4pt}

MacOS:

\begin{codebox}
grc -c HOME/conf.gls gls -iFlhA --color=always\\
\end{codebox}

Linux:

\begin{codebox}
grc -c HOME/conf.gls ls -iFlhA --color=always\\
\end{codebox}

MacOS uses gls (GNU ls from coreutils) because BSD ls doesn't support all the flags. Linux uses native ls.


\vspace{12pt}
{\bfseries\color{neonyellow} THE conf.gls CONFIG}\vspace{4pt}

Custom color configuration for ls output, lives at:

\begin{codebox}
HOME/conf.gls\\
\end{codebox}

Maps file types, permissions, sizes to ANSI colors. Directories, executables, symlinks - all distinct colors.


\vspace{12pt}
{\bfseries\color{neonyellow} OTHER GRC TARGETS}\vspace{4pt}

grc can colorize almost any command: ping         latency numbers in color traceroute   hop-by-hop color coding make         errors red, warnings yellow gcc          compiler output colorized netstat      connection states highlighted diff         additions green, deletions red.


\vspace{12pt}
{\bfseries\color{neonyellow} GRC IN LOG COLORIZATION}\vspace{4pt}

\textbf{\$ZPWR\_SCRIPTS/colorLogger.sh} uses grc for system logs. tmux control windows show colorized live log output. TIP: If eza is missing and ls looks dull, install grc. It's the middle tier of zpwr's color fallback chain. One package, instant color for dozens of commands.


\vspace{16pt}

\section{Ponysay, Lolcat \& ASCII Art}


\vspace{10pt}
{\large\bfseries\color{neoncyan} PONYSAY, LOLCAT \& ASCII ART}\vspace{6pt}

Enterprise-grade cartoon horse infrastructure.


\vspace{12pt}
{\bfseries\color{neonyellow} THE ASCII ART TOOLCHAIN}\vspace{4pt}

ponysay    My Little Pony ASCII art (cowsay but fabulous) lolcat     rainbow gradient text output figlet     large ASCII text banners from font files cowsay     classic ASCII art speech bubbles toilet     another ASCII art text renderer zpwr combines all of these for its startup banner.


\vspace{12pt}
{\bfseries\color{neonyellow} \textbf{ZPWR} BANNER VERBS}\vspace{4pt}

zpwr bannerpony      hostname banner with ponysay zpwr bannerlolcat    ponysay + lolcat rainbow gradient zpwr bannernopony    banner without ponysay zpwr figletfonts     preview all installed figlet fonts Each verb calls a different autoloaded function:

\begin{codebox}
zpwrPonyBanner, zpwrBannerLolcat, zpwrNoPonyBanner\\
\end{codebox}


\vspace{12pt}
{\bfseries\color{neonyellow} THE STARTUP BANNER}\vspace{4pt}

On shell startup, zpwrDarwinBanner / zpwrLinuxBanner fires. It picks a random figlet font each time via:

\begin{codebox}
ZPWR\_SCRIPTS/macOnly/figletRandomFontOnce.sh\\
\end{codebox}

\textbf{ZPWR\_BANNER}\_TYPE controls pony vs no-pony: ponies  -> zpwrBannerLolcat (rainbow horse) other   -> zpwrNoPonyBanner (figlet only)


\vspace{12pt}
{\bfseries\color{neonyellow} BANNER CONTROL}\vspace{4pt}

\textbf{ZPWR\_BANNER}\_CLEARLIST  controls startup banner + ls display When true: clear screen, show banner, then zpwrListNoClear. When false: skip the banner entirely (fast startup).


\vspace{12pt}
{\bfseries\color{neonyellow} FONT LOGGING}\vspace{4pt}

The random figlet font name is logged to:

\begin{codebox}
ZPWR\_LOGFILE\\
\end{codebox}

Want to know which font made that gorgeous banner? zpwr taillog  and look for the font name. Found one you love? Set it permanently in your config. REALITY CHECK: A rainbow pony greeting you on every terminal open is either peak engineering or peak insanity. zpwr does not distinguish between the two. Your hostname, rendered in ASCII art, by a cartoon horse. This is what unlimited power looks like.


\vspace{16pt}

\section{GRC, ccze \& the Color Pipeline}


\vspace{10pt}
{\large\bfseries\color{neoncyan} GRC, CCZE \& THE COLOR PIPELINE}\vspace{6pt}

Every byte of output deserves to be beautiful.


\vspace{12pt}
{\bfseries\color{neonyellow} THE FULL COLORIZATION STACK}\vspace{4pt}

zpwr has a color tool for every type of output: FILE LISTING eza (color-scale, git status) -> grc+ls -> plain ls FILE VIEWING bat (syntax highlighting) -> cat \textbf{ZPWR\_COLORIZER}=bat in \textbf{\$ZPWR\_ENV/.zpwr\_env.sh} LOG VIEWING ccze (log colorizer) -> plain tail zpwrTailLog pipes \$\textbf{ZPWR\_LOGFILE} through ccze. GIT DIFFS delta (syntax-aware diff) -> plain diff delta adds line numbers, word-level highlighting, and side-by-side view with Dracula theme.


\vspace{12pt}
{\bfseries\color{neonyellow} CCZE: LOG FILE COLORIZER}\vspace{4pt}

ccze parses log formats and applies semantic colors: timestamps, severity levels, PIDs, IP addresses. zpwr taillog  uses ccze for colorized log tailing. System logs become readable at a glance.


\vspace{12pt}
{\bfseries\color{neonyellow} COLOR TESTING VERBS}\vspace{4pt}

zpwr colortest       test terminal basic 16 colors zpwr colortest256    test full 256-color palette zpwr pygmentcolors   show all pygment color schemes zpwr colorsdiff      colorized side-by-side diff Use these to verify your terminal supports the full palette.


\vspace{12pt}
{\bfseries\color{neonyellow} THE FALLBACK PHILOSOPHY}\vspace{4pt}

Every color tool has a plain fallback: No eza? grc+ls. No grc? Plain ls. No bat? Cat. zpwr degrades gracefully. Color is never a hard dependency. SSH into a minimal server? Everything still works. Just less pretty. Function over form, always. THE RESULT: In a fully-equipped zpwr terminal: File listings glow. Logs are semantic. Diffs are surgical. Your eyes never have to parse raw monochrome again. The color pipeline is zpwr's cosmetic surgery department.


\vspace{16pt}


\chapter{HEALTH \& DIAGNOSTICS}


\begin{center}
\includegraphics[width=0.85\textwidth]{/Users/wizard/.zpwr/docs/screenshots/zpwr_doctor_help.png}
\end{center}

\section{Startup Profile}


\vspace{10pt}
{\large\bfseries\color{neoncyan} STARTUP PROFILE}\vspace{6pt}

{\small\itshape\color{dimtext} // bench tells you HOW FAST. startup tells you WHERE the time goes. //}

zpwr bench measures total wall time over multiple iterations. zpwr startup gives a single-shot, per-function breakdown using zsh's built-in zprof profiler.


\vspace{12pt}
{\bfseries\color{neonyellow} THE VERB}\vspace{4pt}

\begin{codebox}
zpwr startup               \# top 30 functions (default)\\
zpwr startup 50            \# top 50 functions\\
zpwr startup 100           \# full breakdown\\
\end{codebox}


\vspace{12pt}
{\bfseries\color{neonyellow} HOW IT WORKS}\vspace{4pt}

1. Launches a fresh zsh with \textbf{ZPWR\_PROFILING}=true 2. zprof runs at the end of the init sequence 3. Output is captured and colorized by self-time: >= 10\%  RED       — hot path, optimize here >= 3\%   YELLOW    — warm, worth investigating >= 1\%   CYAN      — normal < 1\%    DIM       — negligible


\vspace{12pt}
{\bfseries\color{neonyellow} THE PERFORMANCE QUARTET}\vspace{4pt}

zpwr bench    -> how fast is total startup? zpwr flame    -> ASCII flame chart of function hotspots zpwr startup  -> per-function zprof breakdown zpwr top      -> what is loaded right now?

Four verbs. Four perspectives. Your shell is transparent.


\vspace{16pt}

\section{Shell Lint}


\vspace{10pt}
{\large\bfseries\color{neoncyan} SHELL LINT}\vspace{6pt}

{\small\itshape\color{dimtext} // if it parses, it ships. but does it REALLY parse? //}

zpwr lint runs two levels of analysis on every zpwr file: syntax checking with zsh -n, and static analysis with shellcheck.


\vspace{12pt}
{\bfseries\color{neonyellow} THE VERB}\vspace{4pt}

\begin{codebox}
zpwr lint                  \# run all checks\\
zpwr lint -z               \# zsh -n syntax check only\\
zpwr lint -s               \# shellcheck only\\
\end{codebox}


\vspace{12pt}
{\bfseries\color{neonyellow} WHAT IT CHECKS}\vspace{4pt}

zsh -n (syntax check) Every file in autoload/common, autoload/darwin, autoload/linux, env/*.sh, env/*.zsh, install/.zshrc. This catches unmatched brackets, bad redirections, and broken quoting. Fast — no execution, just parsing.

shellcheck (static analysis) Every .sh and .zsh file in scripts/ and install/. Catches quoting bugs, unused variables, subshell pitfalls. Zsh-incompatible warnings are suppressed (SC1071 etc).


\vspace{12pt}
{\bfseries\color{neonyellow} WHEN TO RUN IT}\vspace{4pt}

After editing any zpwr file. Before committing. Before pushing. Catch it at the gate, not in production.


\vspace{12pt}
{\bfseries\color{neonyellow} COMPANION VERBS}\vspace{4pt}

\begin{codebox}
zpwr doctor               \# broader health check\\
zpwr stale                \# find stale compiled files\\
zpwr recompile            \# rebuild .zwc files after fixes\\
\end{codebox}


\vspace{16pt}

\section{Listening Ports}


\vspace{10pt}
{\large\bfseries\color{neoncyan} LISTENING PORTS}\vspace{6pt}

{\small\itshape\color{dimtext} // who is listening? and on which door? //}

zpwr ports wraps lsof to show every TCP port in LISTEN state with the process name, PID, user, and bind address.


\vspace{12pt}
{\bfseries\color{neonyellow} THE VERB}\vspace{4pt}

\begin{codebox}
zpwr ports                 \# show all listening ports\\
zpwr ports 8080            \# filter to port 8080\\
zpwr ports -k 3000         \# kill whatever is on port 3000\\
\end{codebox}


\vspace{12pt}
{\bfseries\color{neonyellow} COLOR CODING}\vspace{4pt}

Privileged ports (< 1024)    — RED These require root. Sshd on 22, httpd on 80. Common ports (< 10000)       — YELLOW Dev servers, databases. 3000, 5432, 8080. High ports (>= 10000)        — GREEN Ephemeral and user-space. Generally safe.


\vspace{12pt}
{\bfseries\color{neonyellow} THE KILL FLAG}\vspace{4pt}

zpwr ports -k 3000 finds the PID on port 3000 and sends SIGTERM. Shows process name before killing. Use when a dev server refuses to die.


\vspace{12pt}
{\bfseries\color{neonyellow} COMPANION VERBS}\vspace{4pt}

\begin{codebox}
zpwr lsof                 \# interactive fzf over all lsof output\\
zpwr ps                   \# grep through ps output\\
zpwr kill                 \# fzf kill from ps\\
\end{codebox}


\vspace{16pt}

\section{Stale File Finder}


\vspace{10pt}
{\large\bfseries\color{neoncyan} STALE FILE FINDER}\vspace{6pt}

{\small\itshape\color{dimtext} // compiled cruft is invisible cruft. make it visible. //}

zpwr stale scans the entire zpwr tree for artifacts that have gone stale, orphaned, or broken.


\vspace{12pt}
{\bfseries\color{neonyellow} THE VERB}\vspace{4pt}

\begin{codebox}
zpwr stale                 \# scan and report\\
zpwr stale -f              \# delete all stale files\\
\end{codebox}


\vspace{12pt}
{\bfseries\color{neonyellow} WHAT IT FINDS}\vspace{4pt}

STALE .zwc Source file is newer than its compiled .zwc. Your shell is running yesterday's bytecode.

ORPHAN .zwc The source file was deleted but the .zwc remains. Dead weight taking up space and confusing zsh.

BROKEN SYMLINKS Symlinks pointing to files that no longer exist. Source of cryptic 'file not found' errors.

EMPTY CACHE FILES Zero-byte files in cache directories. Usually left behind by interrupted operations.


\vspace{12pt}
{\bfseries\color{neonyellow} THE FIX FLAG}\vspace{4pt}

zpwr stale -f deletes everything it finds in one pass. Follow up with zpwr recompile to rebuild .zwc files.


\vspace{12pt}
{\bfseries\color{neonyellow} WORKFLOW}\vspace{4pt}

\begin{codebox}
zpwr stale                 \# see what is stale\\
zpwr stale -f              \# nuke the cruft\\
zpwr recompile             \# rebuild fresh .zwc files\\
zpwr doctor                \# verify everything is clean\\
\end{codebox}


\vspace{16pt}

\section{Path Audit}


\vspace{10pt}
{\large\bfseries\color{neoncyan} PATH AUDIT}\vspace{6pt}

{\small\itshape\color{dimtext} // your \textbf{PATH} is your attack surface. audit it. //}

zpwr pathaudit inspects \textbf{PATH}, \textbf{FPATH}, and \textbf{MANPATH} for duplicates, missing directories, and world-writable entries.


\vspace{12pt}
{\bfseries\color{neonyellow} THE VERB}\vspace{4pt}

\begin{codebox}
zpwr pathaudit             \# audit all three arrays\\
zpwr pathaudit -v          \# verbose: show every entry\\
\end{codebox}


\vspace{12pt}
{\bfseries\color{neonyellow} WHAT IT CHECKS}\vspace{4pt}

DUPLICATES Same directory appearing more than once. Wastes hash lookups and obscures precedence. Fix: zpwr deduppaths.

MISSING DIRECTORIES \textbf{PATH} entry pointing to a dir that does not exist. Slows command lookup — zsh checks every entry.

WORLD-WRITABLE DIRECTORIES Any user can write to this directory. If it is in your \textbf{PATH}, an attacker can plant a malicious binary that shadows a real command. Security risk.


\vspace{12pt}
{\bfseries\color{neonyellow} ALL THREE ARRAYS}\vspace{4pt}

\textbf{PATH}     — command lookup for executables \textbf{FPATH}    — function autoload lookup for zsh \textbf{MANPATH}  — man page lookup for documentation.


\vspace{12pt}
{\bfseries\color{neonyellow} COMPANION VERBS}\vspace{4pt}

\begin{codebox}
zpwr doctor               \# broader health check (also checks paths)\\
zpwr deduppaths            \# remove duplicate entries\\
zpwr environmentcounts      \# count all env arrays\\
\end{codebox}

Doctor checks paths as part of a broader scan. Pathaudit goes deeper with world-writable detection.


\vspace{16pt}


\chapter{CREATIVE TOOLS}


\begin{center}
\includegraphics[width=0.85\textwidth]{/Users/wizard/.zpwr/docs/screenshots/figlet_zpwr.png}
\end{center}

\section{The Matrix}


\vspace{10pt}
{\large\bfseries\color{neoncyan} THE MATRIX}\vspace{6pt}

{\small\itshape\color{dimtext} // not cmatrix. this is YOUR shell, raining. //}

zpwr matrix is an animated cyberpunk hacker story followed by an infinite ecosystem rain screensaver. Unlike cmatrix, it rains YOUR data — verb names, function names, aliases, env vars, and git repo names. Color-coded by type.


\vspace{12pt}
{\bfseries\color{neonyellow} THE STORY}\vspace{4pt}

ACT I    The Signal — the terminal underground survives ACT II   The Discovery — curl | zsh changes everything ACT III  The Power — bench, flame, doctor, gitwho ACT IV   The Choice — VS Code or zpwr matrix? Between acts: fake SSH sessions, encrypted transmissions, payload downloads with glitchy progress bars.


\vspace{12pt}
{\bfseries\color{neonyellow} THE FX}\vspace{4pt}

CRT scanlines       green sweep top-to-bottom VHS tracking        horizontal distortion bands Pixel corruption    spreading block static Screen tears        offset colored bands Figlet chaos        random fonts, random words, dissolve Progress bars       with random glitch spikes Pulse events        ecosystem facts flash across screen Scrolling ticker    horizontal stat crawl HUD overlay         live verb/alias/func/repo counters.


\vspace{12pt}
{\bfseries\color{neonyellow} THE RAIN COLORS}\vspace{4pt}

Green   verbs         cyan    functions yellow  aliases        magenta env vars red     git repos.


\vspace{12pt}
{\bfseries\color{neonyellow} USAGE}\vspace{4pt}

\begin{codebox}
zpwr matrix                \# full story + rain\\
zpwr matrix -r             \# skip story, rain only\\
zpwr matrix -s 7           \# faster (1-9)\\
zpwr rain                  \# alias\\
zpwr screensaver           \# alias\\
\end{codebox}


\vspace{16pt}


\chapter{INTROSPECTION}


\begin{center}
\includegraphics[width=0.85\textwidth]{/Users/wizard/.zpwr/docs/screenshots/art_deep_dive.png}
\end{center}

\section{Command Resolver}


\vspace{10pt}
{\large\bfseries\color{neoncyan} COMMAND RESOLVER}\vspace{6pt}

{\small\itshape\color{dimtext} // what IS that command, really? //}

zpwr resolve traces a command through its full resolution chain. Unlike 'which' or 'type', it follows the entire path from alias to function to builtin to binary.


\vspace{12pt}
{\bfseries\color{neonyellow} THE VERB}\vspace{4pt}

\begin{codebox}
zpwr resolve ll               \# where does ll go?\\
zpwr resolve ls cd git        \# trace multiple commands\\
zpwr resolve -a grep          \# show all type -a matches\\
zpwr which vim                \# alias for resolve\\
\end{codebox}


\vspace{12pt}
{\bfseries\color{neonyellow} WHAT IT SHOWS}\vspace{4pt}

Aliases          expansion, then resolves first word.

\begin{codebox}
Functions        source file path, line count, verb mapping\\
\end{codebox}

Builtins         identified as builtin or reserved word Binaries         full path, symlink chain, file type Global aliases   detected and shown Suffix aliases   detected and shown.


\vspace{12pt}
{\bfseries\color{neonyellow} WHY NOT JUST 'which'?}\vspace{4pt}

'Which' shows one answer. 'type -a' shows all candidates. Neither follows the chain. If 'll' is an alias for 'eza' and 'eza' is a symlink to '../Cellar/eza/0.18\textbf{\textbf{/bin/eza}}', resolve shows the entire path from ll to the Mach-O binary.


\vspace{12pt}
{\bfseries\color{neonyellow} COMPANION VERBS}\vspace{4pt}

\begin{codebox}
zpwr exists                \# check if identifier is valid\\
zpwr deps                  \# function dependency graph\\
zpwr funcrank              \# rank functions by usage\\
\end{codebox}


\vspace{16pt}


\chapter{CREATIVE TOOLS}


\begin{center}
\includegraphics[width=0.85\textwidth]{/Users/wizard/.zpwr/docs/screenshots/figlet_zpwr.png}
\end{center}

\section{Fortune Cookies}


\vspace{10pt}
{\large\bfseries\color{neoncyan} FORTUNE COOKIES}\vspace{6pt}

{\small\itshape\color{dimtext} // wisdom from the digital underground //}

zpwr fortune displays a random hacker quote rendered in a random figlet font with neon colors. 40 curated quotes from legends and original \textbf{ZPWR} wisdom.


\vspace{12pt}
{\bfseries\color{neonyellow} THE VERB}\vspace{4pt}

\begin{codebox}
zpwr fortune                \# random quote, random font\\
zpwr fortune -p             \# plain text with box frame\\
zpwr fortune -f doom        \# force doom font\\
zpwr wisdom                 \# alias\\
zpwr quote                  \# alias\\
\end{codebox}


\vspace{12pt}
{\bfseries\color{neonyellow} THE QUOTES}\vspace{4pt}

Linus Torvalds    "Talk is cheap. Show me the code." Brian Kernighan   "Controlling complexity is the essence..." Grace Hopper      "The most dangerous phrase is..." \textbf{ZPWR}              "Home is where the .zshrc is." MenkeTechnologies  "Aliases are love letters to your future self." ...and 35 more.


\vspace{12pt}
{\bfseries\color{neonyellow} THE FONTS}\vspace{4pt}

Randomly selected from: cyberlarge, cybermedium, doom, slant, speed, banner, standard. Each quote picks a keyword for large figlet rendering, full quote below.


\vspace{12pt}
{\bfseries\color{neonyellow} USE CASES}\vspace{4pt}

- Add to .zshrc for a fortune on every new shell - Pipe to lolcat: zpwr fortune -p | lolcat - Use as a tmux status-right: too slow, but fun to try


\vspace{16pt}


\chapter{THE SEARCH ENGINE}


\begin{center}
\includegraphics[width=0.85\textwidth]{/Users/wizard/.zpwr/docs/screenshots/figlet_scan.png}
\end{center}

\section{The Learning System}


\vspace{10pt}
{\large\bfseries\color{neoncyan} THE LEARNING SYSTEM}\vspace{6pt}

{\small\itshape\color{dimtext} // your shell remembers what you learn //}

The zsh-learn plugin turns your shell into a persistent knowledge base backed by MySQL. Save commands, search them, quiz yourself, and build a personal learning collection.


\vspace{12pt}
{\bfseries\color{neonyellow} CORE VERBS}\vspace{4pt}

\begin{codebox}
zpwr learn                  \# save a learning entry\\
zpwr searchl                \# search the collection\\
zpwr quiz                   \# quiz yourself\\
zpwr quizall                \# quiz all entries\\
zpwr redo                   \# redo last search as SQL update\\
\end{codebox}


\vspace{12pt}
{\bfseries\color{neonyellow} SEARCH VARIANTS}\vspace{4pt}

\begin{codebox}
zpwr se                     \# basic search\\
zpwr see                    \# category search\\
zpwr seee                   \# timestamp search\\
zpwr ser                    \# random search\\
zpwr sef                    \# search into fzf\\
zpwr searchall              \# entire collection into fzf\\
zpwr searchfull             \# entire collection into fzf\\
zpwr searchle               \# category search (alias for see)\\
zpwr searchlee              \# timestamp search (alias for seee)\\
zpwr learnsearch            \# search the collection\\
zpwr rsql                   \# redo into vim\\
zpwr redosql                \# redo into vim\\
\end{codebox}


\vspace{12pt}
{\bfseries\color{neonyellow} MANAGEMENT}\vspace{4pt}

\begin{codebox}
zpwr learndel               \# delete entry\\
zpwr learndelete            \# delete entry (alias)\\
zpwr delid                  \# delete by ID\\
zpwr editl                  \# edit entry by ID\\
zpwr getlearn               \# get all entries\\
zpwr learnget               \# get all entries\\
zpwr createlearningcollection \# create MySQL table\\
zpwr droplearningcollection   \# drop MySQL table\\
\end{codebox}


\vspace{12pt}
{\bfseries\color{neonyellow} HOW IT WORKS}\vspace{4pt}

Entries are stored in a MySQL database (configurable via \$\textbf{ZPWR\_SCHEMA}\_NAME and \$\textbf{ZPWR\_TABLE}\_NAME). Each entry has a category, timestamp, and content. The quiz system pulls random entries and tests your recall.

Plugin: MenkeTechnologies/zsh-learn.


\vspace{16pt}


\chapter{BATCH OPERATIONS}


\begin{center}

\includegraphics[width=0.85\textwidth]{/Users/wizard/.zpwr/docs/screenshots/art_override.png}
\end{center}

\section{Git Repo Cache System}


\vspace{10pt}
{\large\bfseries\color{neoncyan} GIT REPO CACHE SYSTEM}\vspace{6pt}

{\small\itshape\color{dimtext} // index every repo on your machine. operate on all of them. //}

The zsh-git-repo-cache plugin scans your entire filesystem for git repos, caches the list, and lets you search, filter, and broadcast commands across all of them.


\vspace{12pt}
{\bfseries\color{neonyellow} CACHE GENERATION}\vspace{4pt}

\begin{codebox}
zpwr regengitrepocache      \# scan filesystem with fd/find\\
zpwr regengitdirtyrepocache  \# scan only dirty repos\\
\end{codebox}


\vspace{12pt}
{\bfseries\color{neonyellow} SEARCH \& BROWSE}\vspace{4pt}

\begin{codebox}
zpwr gitrepos               \# all repos into fzf\\
zpwr gitreposclean          \# clean repos into fzf\\
zpwr gitreposdirty          \# dirty repos into fzf\\
zpwr gitreposlist           \# list all repos (no fzf)\\
zpwr gitreposcleanlist      \# list clean repos\\
zpwr gitreposdirtylist      \# list dirty repos\\
\end{codebox}


\vspace{12pt}
{\bfseries\color{neonyellow} CACHED VARIANTS}\vspace{4pt}

\begin{codebox}
zpwr gitreposcleancache     \# clean repos from cache\\
zpwr gitreposdirtycache     \# dirty repos from cache\\
zpwr gitreposcleancachelist  \# list clean from cache\\
zpwr gitreposdirtycachelist  \# list dirty from cache\\
\end{codebox}


\vspace{12pt}
{\bfseries\color{neonyellow} BATCH OPERATIONS}\vspace{4pt}

\begin{codebox}
zpwr gitreposexec <cmd>      \# run cmd in all cached repos\\
zpwr gitreposdirtyexec <cmd> \# run cmd in dirty repos\\
zpwr gitreposcleanexec <cmd> \# run cmd in clean repos\\
zpwr gitwho                 \# contributor stats across all repos\\
\end{codebox}


\vspace{12pt}
{\bfseries\color{neonyellow} CACHE FILES}\vspace{4pt}

\$\textbf{ZPWR\_ALL}\_GIT\_DIRS       — all repos \$\textbf{ZPWR\_ALL}\_GIT\_DIRS\_CLEAN — clean repos only \$\textbf{ZPWR\_ALL}\_GIT\_DIRS\_DIRTY — dirty repos only.

Plugin: MenkeTechnologies/zsh-git-repo-cache.


\vspace{16pt}


\chapter{FZF ARCHITECTURE}


\begin{center}
\includegraphics[width=0.85\textwidth]{/Users/wizard/.zpwr/docs/screenshots/figlet_scan.png}
\end{center}

\section{FZF Brew \& Chrome}


\vspace{10pt}
{\large\bfseries\color{neoncyan} FZF BREW \& CHROME}\vspace{6pt}

{\small\itshape\color{dimtext} // fuzzy finding extends to your package manager and browser //}

The fzf-zsh-plugin adds fzf interfaces to Homebrew package management and Chrome browser data.


\vspace{12pt}
{\bfseries\color{neonyellow} HOMEBREW VERBS}\vspace{4pt}

\begin{codebox}
zpwr brewinstall            \# fzf install formula\\
zpwr brewuninstall          \# fzf uninstall formula\\
zpwr brewcaskinstall        \# fzf install cask\\
zpwr brewcaskuninstall      \# fzf uninstall cask\\
zpwr brewupdate             \# fzf update formula\\
\end{codebox}

Each command pipes the available packages into fzf with multi-select. Pick one or many, then install/remove in one shot. No more googling package names.


\vspace{12pt}
{\bfseries\color{neonyellow} CHROME VERBS}\vspace{4pt}

\begin{codebox}
zpwr chromehistory          \# search Chrome history in fzf\\
zpwr chromebookmarks        \# search Chrome bookmarks in fzf\\
\end{codebox}

Reads Chrome's SQLite databases directly. Select a URL from fzf and open it in the browser.

Plugin: MenkeTechnologies/fzf-zsh-plugin.


\vspace{16pt}


\chapter{THE GIT DIMENSION}


\begin{center}

\includegraphics[width=0.85\textwidth]{/Users/wizard/.zpwr/docs/screenshots/art_data_stream.png}
\end{center}

\section{Git Add-Commit-Push}


\vspace{10pt}
{\large\bfseries\color{neoncyan} GIT ADD-COMMIT-PUSH}\vspace{6pt}

{\small\itshape\color{dimtext} // 318 aliases for the git power user //}

The zsh-git-acp plugin provides 318 aliases that collapse multi-step git workflows into single commands.


\vspace{12pt}
{\bfseries\color{neonyellow} THE GACP PATTERN}\vspace{4pt}

\begin{codebox}
gacp "message"              \# git add -A, commit, push\\
gac "message"               \# git add -A, commit (no push)\\
gcp                          \# git commit, push\\
\end{codebox}


\vspace{12pt}
{\bfseries\color{neonyellow} BRANCH SHORTCUTS}\vspace{4pt}

\begin{codebox}
gcm                          \# git checkout main\\
gcb <branch>                 \# git checkout -b <branch>\\
gbd <branch>                 \# git branch -d <branch>\\
\end{codebox}


\vspace{12pt}
{\bfseries\color{neonyellow} \textbf{ZPWR} INTEGRATION}\vspace{4pt}

\begin{codebox}
zpwr gitcommit "msg"        \# commit and push with message\\
zpwr gitcommitcount          \# count commits in current repo\\
\end{codebox}

The full alias list is enormous — 318 combinations of git operations compressed into 2-4 character shortcuts. Most follow the pattern: g + operation letter(s).

Plugin: MenkeTechnologies/zsh-git-acp.


\vspace{16pt}


\chapter{NETWORK \& CONTAINERS}


\begin{center}

\includegraphics[width=0.85\textwidth]{/Users/wizard/.zpwr/docs/screenshots/art_firewall.png}
\end{center}

\section{Docker \& Container Aliases}


\vspace{10pt}
{\large\bfseries\color{neoncyan} DOCKER \& CONTAINER ALIASES}\vspace{6pt}

{\small\itshape\color{dimtext} // 148 aliases for the containerized world //}

The zsh-docker-aliases plugin provides 148 aliases for Docker and Docker Compose operations.


\vspace{12pt}
{\bfseries\color{neonyellow} DOCKER SHORTCUTS}\vspace{4pt}

\begin{codebox}
dps                          \# docker ps\\
dpsa                         \# docker ps -a\\
di                           \# docker images\\
drm                          \# docker rm\\
drmi                         \# docker rmi\\
dex                          \# docker exec -it\\
dl                           \# docker logs -f\\
\end{codebox}


\vspace{12pt}
{\bfseries\color{neonyellow} DOCKER COMPOSE}\vspace{4pt}

\begin{codebox}
dcu                          \# docker compose up -d\\
dcd                          \# docker compose down\\
dcl                          \# docker compose logs -f\\
dcr                          \# docker compose restart\\
\end{codebox}


\vspace{12pt}
{\bfseries\color{neonyellow} \textbf{ZPWR} INTEGRATION}\vspace{4pt}

\begin{codebox}
zpwr dockerwipe              \# nuke all images/containers\\
zpwr dfimage <ID>            \# extract Dockerfile from image\\
\end{codebox}

Plugin: MenkeTechnologies/zsh-docker-aliases.


\vspace{16pt}


\chapter{ZINIT \& PLUGIN SYSTEM}


\begin{center}

\includegraphics[width=0.85\textwidth]{/Users/wizard/.zpwr/docs/screenshots/art_system_core.png}
\end{center}

\section{The Plugin Ecosystem}


\vspace{10pt}
{\large\bfseries\color{neoncyan} THE PLUGIN ECOSYSTEM}\vspace{6pt}

{\small\itshape\color{dimtext} // 30+ custom plugins from MenkeTechnologies //}

\textbf{ZPWR} loads approximately 95 plugins via zinit. Over 30 are custom MenkeTechnologies creations that extend zpwr with additional verbs, aliases, completions, and integrations.


\vspace{12pt}
{\bfseries\color{neonyellow} KEY PLUGINS}\vspace{4pt}

Zsh-learn              60 verbs, MySQL learning DB zsh-git-repo-cache    20 verbs, filesystem repo index fzf-zsh-plugin        14 verbs, brew/chrome fzf zsh-git-acp           318 aliases, git workflows zsh-docker-aliases    148 aliases, Docker shortcuts zsh-openshift-aliases 106 aliases, OpenShift zsh-expand            spell correct, abbreviation expand zsh-more-completions  many thousands of tab completions zsh-cargo-completion  32 Rust/Cargo aliases zsh-travis            8 functions, CI integration zsh-nginx             nginx management forgit                fzf-powered git (add, diff, log) zsh-z                 frecency directory jumping fzf-tab               fzf-powered tab completion zconvey               cross-session communication.


\vspace{12pt}
{\bfseries\color{neonyellow} CONDITIONAL LOADING}\vspace{4pt}

Plugins for docker, dotnet, kubectl, ngrok, rails, and systemctl only load if those commands exist. No bloat on machines without those tools.


\vspace{12pt}
{\bfseries\color{neonyellow} DEFERRED LOADING}\vspace{4pt}

All plugins load with zinit wait'' ice — deferred after the prompt appears. Only powerlevel10k loads synchronously for instant prompt.


\vspace{16pt}


\chapter{TIPS, TRICKS \& EASTER EGGS}


\begin{center}
\includegraphics[width=0.85\textwidth]{/Users/wizard/.zpwr/docs/screenshots/art_hack_the_planet.png}
\end{center}

\section{Alias Cheat Sheet}


\vspace{10pt}
{\large\bfseries\color{neoncyan} ALIAS CHEAT SHEET}\vspace{6pt}

{\small\itshape\color{dimtext} // the short names you will actually use //}

Many zpwr verbs have short aliases. These are the ones worth memorizing — one or two keystrokes that replace entire workflows.


\vspace{12pt}
{\bfseries\color{neonyellow} DIAGNOSTICS}\vspace{4pt}

\begin{codebox}
zpwr doc          = doctor       zpwr zprof      = startup\\
zpwr check        = lint         zpwr audit      = pathaudit\\
zpwr cruft        = stale        zpwr rank       = funcrank\\
zpwr benchmark    = bench        zpwr dashboard  = top\\
zpwr heatmap      = flame        zpwr instrument = trace\\
zpwr flamechart   = flame        zpwr checkup    = doctor\\
zpwr callers      = deps         zpwr aliasusage = aliasrank\\
\end{codebox}


\vspace{12pt}
{\bfseries\color{neonyellow} MONITORING}\vspace{4pt}

\begin{codebox}
zpwr snap         = snapshot     zpwr load       = restore\\
zpwr rerun        = replay       zpwr onchange   = watch\\
zpwr book         = study\\
\end{codebox}


\vspace{12pt}
{\bfseries\color{neonyellow} CREATIVE}\vspace{4pt}

\begin{codebox}
zpwr rain         = matrix       zpwr wisdom     = fortune\\
zpwr quote        = fortune      zpwr screensaver= matrix\\
\end{codebox}


\vspace{12pt}
{\bfseries\color{neonyellow} GIT}\vspace{4pt}

\begin{codebox}
zpwr who          = gitwho       zpwr gitcontribs= gitwho\\
\end{codebox}


\vspace{12pt}
{\bfseries\color{neonyellow} INTROSPECTION}\vspace{4pt}

\begin{codebox}
zpwr which        = resolve      zpwr whatis     = resolve\\
zpwr today        = timeline     zpwr activity   = timeline\\
\end{codebox}


\vspace{12pt}
{\bfseries\color{neonyellow} LOGGING}\vspace{4pt}

\begin{codebox}
zpwr logs         = taillog\\
\end{codebox}


\vspace{16pt}


\chapter{EMACS DIMENSION}


\begin{center}
\includegraphics[width=0.85\textwidth]{/Users/wizard/.zpwr/docs/screenshots/art_neural_link.png}
\end{center}

\section{Emacs Integration}


\vspace{10pt}
{\large\bfseries\color{neoncyan} EMACS INTEGRATION}\vspace{6pt}

{\small\itshape\color{dimtext} // 29 verbs for the church of emacs //}

Every file operation in \textbf{ZPWR} has an emacs variant. Open, edit, search, and navigate — all wired to emacsclient.


\vspace{12pt}
{\bfseries\color{neonyellow} OPEN \& EDIT}\vspace{4pt}

\begin{codebox}
zpwr emacsall              \# emacs all zpwr files for :argdo\\
zpwr emacsallserver        \# emacs all zpwr files via server\\
zpwr emacsalledit          \# emacs edit 1 or more configs\\
zpwr emacsautoload         \# emacs all autoloads :argdo\\
zpwr emacsscripts          \# emacs all zpwr scripts for :argdo\\
zpwr emacsscriptedit       \# emacs edit 1 or more scripts\\
zpwr emacszpwr             \# emacs zpwr dir\\
\end{codebox}


\vspace{12pt}
{\bfseries\color{neonyellow} CONFIG \& TOKENS}\vspace{4pt}

\begin{codebox}
zpwr emacsconfig           \# emacs all zpwr configs\\
zpwr emacsemacsconfig      \# emacs edit emacs zpwr configs\\
zpwr emacstokens           \# emacs the .tokens.sh file\\
zpwr emacstests            \# emacs edit all zpwr tests\\
\end{codebox}


\vspace{12pt}
{\bfseries\color{neonyellow} NAVIGATION}\vspace{4pt}

\begin{codebox}
zpwr emacscd               \# emacs edit and cd to first dir\\
zpwr emacsrecent           \# emacs edit most recent editor files\\
zpwr emacsrecentcd         \# cd and emacs edit most recent files\\
zpwr emacsrecentsudo       \# sudo emacs edit most recent files\\
zpwr emacsrecentsudocd     \# cd and sudo emacs most recent files\\
\end{codebox}


\vspace{12pt}
{\bfseries\color{neonyellow} SEARCH}\vspace{4pt}

\begin{codebox}
zpwr emacsfilesearch       \# emacs a file in a sub dir\\
zpwr emacsfilesearchedit   \# edit emacs a file in a sub dir\\
zpwr emacsfindsearch       \# emacs find drive for file\\
zpwr emacsfindsearchedit   \# emacs edit find drive for file\\
zpwr emacslocatesearch     \# emacs accept locate drive for file\\
zpwr emacslocatesearchedit \# emacs edit locate drive for file\\
zpwr emacswordsearch       \# emacs a file in a sub dir by word\\
zpwr emacswordsearchedit   \# edit emacs a file by word\\
\end{codebox}


\vspace{12pt}
{\bfseries\color{neonyellow} TAGS \& PLUGINS}\vspace{4pt}

\begin{codebox}
zpwr emacsctags            \# emacs zpwr ctags tags\\
zpwr emacsgtags            \# emacs zpwr GNU global tags\\
zpwr emacsgtagsedit        \# emacs edit zpwr GNU global tags\\
zpwr emacsplugincount      \# total number of emacs plugins\\
zpwr emacspluginlist       \# total list of emacs plugins\\
\end{codebox}


\vspace{16pt}


\chapter{EDITOR MASTERY}


\begin{center}
\includegraphics[width=0.85\textwidth]{/Users/wizard/.zpwr/docs/screenshots/art_neural_link.png}
\end{center}

\section{Vim Integration}


\vspace{10pt}
{\large\bfseries\color{neoncyan} VIM INTEGRATION}\vspace{6pt}

{\small\itshape\color{dimtext} // 35 verbs for the vim faithful //}

The vim verb family mirrors emacs with full parity. Open, edit, search, navigate — all wired to vim/nvim.


\vspace{12pt}
{\bfseries\color{neonyellow} OPEN \& EDIT}\vspace{4pt}

\begin{codebox}
zpwr vimall               \# vim all zpwr files for :argdo\\
zpwr vimalledit            \# vim edit 1 or more configs\\
zpwr vimautoload           \# vim all autoloads :argdo\\
zpwr vimscripts            \# vim all zpwr scripts for :argdo\\
zpwr vimscriptedit         \# vim edit 1 or more scripts\\
\end{codebox}


\vspace{12pt}
{\bfseries\color{neonyellow} CONFIG \& TOKENS}\vspace{4pt}

\begin{codebox}
zpwr vimconfig             \# edit all zpwr configs\\
zpwr vimemacsconfig        \# vim edit emacs zpwr configs\\
zpwr vimtokens             \# vim the .tokens.sh file\\
zpwr vimtests              \# edit all zpwr tests\\
zpwr vrc                   \# vimrc vim session\\
zpwr brc                   \# shell aliases file vim session\\
zpwr trc                   \# tmux.conf vim session\\
zpwr zrc                   \# zshrc vim session\\
zpwr tokens                \# vim the .tokens.sh file\\
\end{codebox}


\vspace{12pt}
{\bfseries\color{neonyellow} NAVIGATION}\vspace{4pt}

\begin{codebox}
zpwr vimcd                 \# vim edit and cd to first dir\\
zpwr vimrecent             \# vim edit most recent editor files\\
zpwr vimrecentcd           \# cd and vim edit most recent files\\
zpwr vimrecentsudo         \# sudo vim edit most recent files\\
zpwr vimrecentsudocd       \# cd and sudo vim most recent files\\
\end{codebox}


\vspace{12pt}
{\bfseries\color{neonyellow} SEARCH}\vspace{4pt}

\begin{codebox}
zpwr vimsearch             \# search vim keybindings\\
zpwr vimfilesearch         \# vim a file in a sub dir\\
zpwr vimfilesearchedit     \# edit vim a file in a sub dir\\
zpwr vimfindsearch         \# vim accept find drive for file\\
zpwr vimfindsearchedit     \# vim edit find drive for file\\
zpwr vimlocatesearch       \# vim accept locate drive for file\\
zpwr vimlocatesearchedit   \# vim edit locate drive for file\\
zpwr vimwordsearch         \# vim a file in a sub dir by word\\
zpwr vimwordsearchedit     \# edit vim a file by word\\
\end{codebox}


\vspace{12pt}
{\bfseries\color{neonyellow} TAGS, PLUGINS \& VIMINFO}\vspace{4pt}

\begin{codebox}
zpwr vimctags              \# vim zpwr ctags tags\\
zpwr vimgtags              \# vim zpwr GNU global tags\\
zpwr vimgtagsedit          \# vim edit zpwr GNU global tags\\
zpwr vimplugincount        \# total number of vim plugins\\
zpwr vimpluginlist         \# total list of vim plugins\\
zpwr catviminfo            \# cat custom vim info file\\
zpwr catrecentfviminfo     \# cat recentf then custom vim info\\
zpwr catnviminforecentf    \# cat custom vim info then recentf\\
\end{codebox}


\vspace{16pt}


\chapter{THE GIT DIMENSION}


\begin{center}

\includegraphics[width=0.85\textwidth]{/Users/wizard/.zpwr/docs/screenshots/art_data_stream.png}
\end{center}

\section{Git Operations Reference}


\vspace{10pt}
{\large\bfseries\color{neoncyan} GIT OPERATIONS REFERENCE}\vspace{6pt}

{\small\itshape\color{dimtext} // 33 verbs for git mastery //}

Core git verbs for commits, config, history rewriting, contributor stats, and multi-repo operations.


\vspace{12pt}
{\bfseries\color{neonyellow} COMMIT \& PUSH}\vspace{4pt}

\begin{codebox}
zpwr gitcommit             \# commit and push with arg message\\
zpwr gitcommitcount        \# count commits in current repo\\
zpwr gitcommits            \# search git commits with fzf\\
zpwr clone                 \# clone and cd to arg\\
\end{codebox}


\vspace{12pt}
{\bfseries\color{neonyellow} CONFIG \& IGNORES}\vspace{4pt}

\begin{codebox}
zpwr gitconfig             \# edit git config file\\
zpwr gitglobalconfig       \# edit git config file\\
zpwr gitignore             \# EDITOR \textasciitilde{}/.git/info/exclude\\
zpwr gitignores            \# edit git ignores file\\
zpwr gitglobalignores      \# edit git ignores file\\
\end{codebox}


\vspace{12pt}
{\bfseries\color{neonyellow} EMAIL \& HISTORY REWRITING}\vspace{4pt}

\begin{codebox}
zpwr gitemail              \# change email with git filter-branch\\
zpwr gitcemail             \# change committer email\\
zpwr gitaemail             \# change author email\\
zpwr gitclearcommit        \# remove matching commits from history\\
zpwr gitclearcache         \# clear old git refs and objects\\
zpwr gitclearfile          \# rm file from all git refs/objects\\
\end{codebox}


\vspace{12pt}
{\bfseries\color{neonyellow} STATS \& INFO}\vspace{4pt}

\begin{codebox}
zpwr gitcontribcount       \# count git contributions by author\\
zpwr gitcontribcountdirs   \# contributions by author for dirs\\
zpwr gitcontribcountlines  \# count lines contributed by author\\
zpwr gittotallines         \# total line count of git files\\
zpwr gitlargest            \# show largest files stored by git\\
zpwr gitremotes            \# list all git remotes\\
zpwr gitdir                \# check if pwd is git dir\\
zpwr gitwho                \# aggregate contributor stats\\
zpwr gitsearch             \# search for regex in git log\\
\end{codebox}


\vspace{12pt}
{\bfseries\color{neonyellow} MULTI-REPO \& TAGS}\vspace{4pt}

\begin{codebox}
zpwr gitforalldir          \# run cmd in all git dirs\\
zpwr gitfordirdevelop      \# git wipe to develop in git dirs\\
zpwr gitfordirmain         \# git wipe to main in git dirs\\
zpwr gitupdatefordir       \# run git updates in all git dirs\\
zpwr gitzfordir            \# git wipe on branch in git dirs\\
zpwr gitzfordirdevelop     \# z and git wipe to develop\\
zpwr gitzfordirmain        \# z and git wipe to main\\
zpwr gitupdatetag          \# print tag and latest msg\\
zpwr gitedittag            \# print tag to BUFFER\\
\end{codebox}


\vspace{16pt}

\section{GitHub \& Travis CI}


\vspace{10pt}
{\large\bfseries\color{neoncyan} GITHUB \& TRAVIS CI}\vspace{6pt}

{\small\itshape\color{dimtext} // 16 verbs for the hosted world //}

Verbs for GitHub repo management, contribution counting, and Travis CI browser integration.


\vspace{12pt}
{\bfseries\color{neonyellow} GITHUB NAVIGATION}\vspace{4pt}

\begin{codebox}
zpwr gh                    \# open github profile\\
zpwr github                \# open github profile\\
zpwr mygithub              \# open github profile\\
zpwr ghz                   \# open zpwr github repo\\
zpwr githubzpwr            \# open zpwr github repo\\
zpwr zpwrgithub            \# open zpwr github\\
\end{codebox}


\vspace{12pt}
{\bfseries\color{neonyellow} GITHUB REPO MANAGEMENT}\vspace{4pt}

\begin{codebox}
zpwr githubcreate          \# create remote github repo\\
zpwr githubdelete          \# delete remote github repo\\
zpwr hc                    \# create remote github repo\\
zpwr hd                    \# delete remote github repo\\
\end{codebox}


\vspace{12pt}
{\bfseries\color{neonyellow} CONTRIBUTION STATS}\vspace{4pt}

\begin{codebox}
zpwr ghcontribcount        \# count github contribs in last year\\
zpwr githubcontribcount    \# count github contribs in last year\\
\end{codebox}


\vspace{12pt}
{\bfseries\color{neonyellow} TRAVIS CI}\vspace{4pt}

\begin{codebox}
zpwr travis                \# open current travis project\\
zpwr travisbranch          \# open current travis branches\\
zpwr travisbuild           \# open current travis builds\\
zpwr travispr              \# open current travis PRs\\
\end{codebox}


\vspace{16pt}


\chapter{NAVIGATION PROTOCOL}


\begin{center}
\includegraphics[width=0.85\textwidth]{/Users/wizard/.zpwr/docs/screenshots/art_the_grid.png}
\end{center}

\section{Navigation Verbs}


\vspace{10pt}
{\large\bfseries\color{neoncyan} NAVIGATION VERBS}\vspace{6pt}

{\small\itshape\color{dimtext} // 39 ways to get where you are going //}

Navigate the filesystem, \textbf{ZPWR} directories, and frecency ranked locations with precision.


\vspace{12pt}
{\bfseries\color{neonyellow} DIRECTORY MOVEMENT}\vspace{4pt}

\begin{codebox}
zpwr cd                    \# cd to directory arg\\
zpwr f                     \# cd to directory arg\\
zpwr forward               \# cd to directory arg\\
zpwr fwd                   \# cd to directory arg\\
zpwr cdup                  \# cd up tree to directory arg\\
zpwr r                     \# cd up tree to directory arg\\
zpwr reverse               \# cd up tree to directory arg\\
zpwr rvs                   \# cd up tree to directory arg\\
zpwr catcd                 \# cat and cd to first dir\\
zpwr clone                 \# clone and cd to arg\\
\end{codebox}


\vspace{12pt}
{\bfseries\color{neonyellow} \textbf{ZPWR} \textbf{HOME} DIRECTORIES}\vspace{4pt}

\begin{codebox}
zpwr home                  \# go to ZPWR\\
zpwr homeautoload          \# go to ZPWR\_AUTOLOAD\\
zpwr homeautoloadcommon    \# go to ZPWR\_AUTOLOAD\_COMMON\\
zpwr homecomps             \# go to ZPWR\_COMPS\\
zpwr homeenv               \# go to ZPWR\_ENV\\
zpwr homeexercism          \# cd to HOME/Exercism\\
zpwr homeinstall           \# go to ZPWR\_INSTALL\\
zpwr homelocal             \# go to ZPWR\_LOCAL\\
zpwr homescripts           \# go to ZPWR\_SCRIPTS\\
zpwr hometest              \# go to ZPWR\_TEST\\
zpwr hometests             \# go to ZPWR\_TEST\\
zpwr hometmux              \# go to ZPWR\_TMUX\\
\end{codebox}


\vspace{12pt}
{\bfseries\color{neonyellow} \textbf{ZPWR} SUBDIRECTORIES}\vspace{4pt}

\begin{codebox}
zpwr autoload              \# cd to autoload directory\\
zpwr comps                 \# cd to completion directory\\
zpwr completions           \# cd to completion directory\\
zpwr forked                \# cd to ZPWR\_FORKED\_DIR\\
zpwr fp                    \# cd to ZPWR\_FORKED\_DIR\\
zpwr hidden                \# go to ZPWR\_HIDDEN\_DIR\\
zpwr plugins               \# cd to ZSH\_CUSTOM/plugins\\
zpwr repo                  \# cd to ZPWR\_REPO\_NAME\\
zpwr scripts               \# cd to scripts directory\\
zpwr web                   \# cd to web dir\\
zpwr zp                    \# cd to ZPWR\_REPO\_NAME\\
zpwr zpwr                  \# cd to ZPWR\_REPO\_NAME\\
zpwr zpz                   \# cd to repo, checkout, rebase, push\\
\end{codebox}


\vspace{12pt}
{\bfseries\color{neonyellow} FRECENCY \& FZF}\vspace{4pt}

\begin{codebox}
zpwr z                     \# cd to z frecency ranked dir\\
zpwr zcd                   \# list then cd to z ranked dir\\
zpwr dirsearch             \# cd to a sub dir via fzf\\
\end{codebox}


\vspace{16pt}


\chapter{BUILD SYSTEM}


\begin{center}
\includegraphics[width=0.85\textwidth]{/Users/wizard/.zpwr/docs/screenshots/figlet_exec.png}
\end{center}

\section{Regen \& Compile}


\vspace{10pt}
{\large\bfseries\color{neoncyan} REGEN \& COMPILE}\vspace{6pt}

{\small\itshape\color{dimtext} // 24 verbs to rebuild the world //}

Regenerate caches, recompile zsh word code, and rebuild tags. The engine room of \textbf{ZPWR} maintenance.


\vspace{12pt}
{\bfseries\color{neonyellow} REGEN CACHES}\vspace{4pt}

\begin{codebox}
zpwr regen                 \# regen caches except git drive search\\
zpwr regenall              \# regen all caches\\
zpwr regenpdf              \# regenerate encyclopedia PDF\\
zpwr regenenvcache         \# regen search env to zpwrEnv files\\
zpwr regengitrepocache     \# regen list of all git repos\\
zpwr regengitdirtyrepocache\# regen list of dirty git repos\\
zpwr regenhistory          \# regen HISTFILE\\
zpwr regenkeybindings      \# regen all keybindings cache\\
zpwr regenconfiglinks      \# regen sym links to HOME\\
zpwr regenpowerline        \# regen powerline sym link\\
zpwr regenzsh              \# regen compsys cache\\
\end{codebox}


\vspace{12pt}
{\bfseries\color{neonyellow} REGEN TAGS}\vspace{4pt}

\begin{codebox}
zpwr regengtags            \# regen gtags files to HOME\\
zpwr regenctags            \# regen ctags files to HOME\\
zpwr regengtagspygments    \# regen GNU gtags with pygments\\
zpwr regengtagstype        \# regen GNU gtags with specified parser\\
\end{codebox}


\vspace{12pt}
{\bfseries\color{neonyellow} COMPILE \& DECOMPILE}\vspace{4pt}

\begin{codebox}
zpwr compile               \# recompile all cache comps\\
zpwr compilefpath          \# recompile all fpath cache comps\\
zpwr compilefiles          \# recompile all files cache comps\\
zpwr recompile             \# recompile all cache comps\\
zpwr recompilefpath        \# recompile all fpath cache comps\\
zpwr recompilefiles        \# recompile all files cache comps\\
zpwr refreshzwc            \# delete then regen compiled zsh word code\\
zpwr decompile             \# delete all cache comps\\
zpwr uncompile             \# delete all cache comps\\
\end{codebox}


\vspace{16pt}


\chapter{THE JANITOR'S CLOSET}


\begin{center}

\includegraphics[width=0.85\textwidth]{/Users/wizard/.zpwr/docs/screenshots/art_override.png}
\end{center}

\section{Cleanup Arsenal}


\vspace{10pt}
{\large\bfseries\color{neoncyan} CLEANUP ARSENAL}\vspace{6pt}

{\small\itshape\color{dimtext} // 14 verbs to sweep the floor //}

Clean directories, caches, temp files, logs, and git state. From surgical precision to nuclear wipe.


\vspace{12pt}
{\bfseries\color{neonyellow} DIRECTORY CLEANING}\vspace{4pt}

\begin{codebox}
zpwr cleandirs             \# clear all ZPWR\_DIRS\_CLEAN\\
zpwr cleandirsandtemp      \# clear dirs and temp\\
zpwr cleanall              \# clear dirs, temp and cache\\
zpwr cleantemp             \# clean all zpwr temp files\\
\end{codebox}


\vspace{12pt}
{\bfseries\color{neonyellow} CACHE CLEANING}\vspace{4pt}

\begin{codebox}
zpwr cleancache            \# clean all zpwr cache files\\
zpwr cleancompcache        \# clean all zpwr compsys cache files\\
zpwr cleangitcache         \# clean all git repo cache files\\
zpwr cleangitdirtycache    \# clean git repo dirty cache files\\
zpwr cleangitcleancache    \# clean git repo clean cache files\\
\end{codebox}


\vspace{12pt}
{\bfseries\color{neonyellow} LOG \& UPDATE CLEANING}\vspace{4pt}

\begin{codebox}
zpwr cleanlog              \# clear ZPWR\_LOGFILE\\
zpwr cleanrefreshupdate    \# clean, refresh zwc, update deps\\
zpwr cleanupdatedeps       \# clean and update deps\\
\end{codebox}


\vspace{12pt}
{\bfseries\color{neonyellow} DO EVERYTHING}\vspace{4pt}

\begin{codebox}
zpwr a                     \# banner, counts, clean, refresh, test, update\\
zpwr all                   \# banner, counts, clean, refresh, test, update\\
\end{codebox}


\vspace{16pt}


\chapter{FORGIT DEEP DIVE}


\begin{center}
\includegraphics[width=0.85\textwidth]{/Users/wizard/.zpwr/docs/screenshots/art_the_grid.png}
\end{center}

\section{Forgit Reference}


\vspace{10pt}
{\large\bfseries\color{neoncyan} FORGIT REFERENCE}\vspace{6pt}

{\small\itshape\color{dimtext} // 14 verbs for fzf-powered git //}

Forgit wraps common git operations in fzf interfaces. Interactive staging, diffing, stashing, and more.


\vspace{12pt}
{\bfseries\color{neonyellow} STAGING \& DIFFING}\vspace{4pt}

\begin{codebox}
zpwr forgitadd             \# fzf interactive git add\\
zpwr forgitdiff            \# fzf interactive git diff\\
zpwr forgitreset           \# fzf interactive git reset\\
zpwr forgitrestore         \# fzf interactive git restore\\
\end{codebox}


\vspace{12pt}
{\bfseries\color{neonyellow} STASH \& CLEAN}\vspace{4pt}

\begin{codebox}
zpwr forgitstash           \# fzf interactive stash show\\
zpwr forgitclean           \# fzf interactive git clean\\
\end{codebox}


\vspace{12pt}
{\bfseries\color{neonyellow} LOG \& INFO}\vspace{4pt}

\begin{codebox}
zpwr forgitlog             \# fzf interactive git log\\
zpwr forgitinfo            \# fzf git info\\
zpwr forgitwarn            \# fzf git warn\\
\end{codebox}


\vspace{12pt}
{\bfseries\color{neonyellow} GITIGNORE MANAGEMENT}\vspace{4pt}

\begin{codebox}
zpwr forgitignore          \# fzf gitignore templates\\
zpwr forgitignoreclean     \# fzf gitignore clean\\
zpwr forgitignoreget       \# fzf gitignore get\\
zpwr forgitignorelist      \# fzf gitignore list\\
zpwr forgitignoreupdate    \# fzf gitignore update\\
\end{codebox}


\vspace{16pt}


\chapter{BUILD SYSTEM}


\begin{center}
\includegraphics[width=0.85\textwidth]{/Users/wizard/.zpwr/docs/screenshots/figlet_exec.png}
\end{center}

\section{Update System}


\vspace{10pt}
{\large\bfseries\color{neoncyan} UPDATE SYSTEM}\vspace{6pt}

{\small\itshape\color{dimtext} // 6 verbs to stay current //}

Keep \textbf{ZPWR}, its dependencies, and zinit plugins up to date. From surgical single-dep pulls to full system updates.


\vspace{12pt}
{\bfseries\color{neonyellow} CONFIG UPDATES}\vspace{4pt}

\begin{codebox}
zpwr update                \# update zpwr custom configs\\
zpwr updatepull            \# pull latest config from repo\\
\end{codebox}


\vspace{12pt}
{\bfseries\color{neonyellow} DEPENDENCY UPDATES}\vspace{4pt}

\begin{codebox}
zpwr updatedeps            \# update all dependencies\\
zpwr updatedepsclean       \# update all dependencies then clean\\
zpwr updatezinit           \# update zinit dependencies\\
\end{codebox}


\vspace{12pt}
{\bfseries\color{neonyellow} FULL UPDATE}\vspace{4pt}

\begin{codebox}
zpwr updateall             \# update zpwr configs and all deps\\
\end{codebox}

Typical workflow: zpwr updateall pulls the latest configs, updates brew/pip/npm/gem/zinit, and refreshes caches. Run zpwr updatedepsclean afterward to reclaim disk space from stale packages.


\vspace{16pt}


\chapter{TMUX WARFARE}


\begin{center}
\includegraphics[width=0.85\textwidth]{/Users/wizard/.zpwr/docs/screenshots/art_hack_the_planet.png}
\end{center}

\section{Clipboard, Send \& Tmux}


\vspace{10pt}
{\large\bfseries\color{neoncyan} CLIPBOARD, SEND \& TMUX}\vspace{6pt}

{\small\itshape\color{dimtext} // 13 verbs for pane orchestration //}

Clipboard access, key duplication across panes, URL handling, and tmux session management.


\vspace{12pt}
{\bfseries\color{neonyellow} CLIPBOARD}\vspace{4pt}

\begin{codebox}
zpwr copycommand           \# get the command to copy with system\\
zpwr pastecommand          \# get the command to paste with system\\
zpwr pastebuffer           \# paste clipboard to BUFFER\\
zpwr pastestring           \# paste clipboard to stdout\\
\end{codebox}


\vspace{12pt}
{\bfseries\color{neonyellow} KEY DUPLICATION}\vspace{4pt}

\begin{codebox}
zpwr startsend             \# dup some keys to pane\\
zpwr startsendfull         \# dup all keys to pane\\
zpwr stopsend              \# stop dup keys to pane\\
\end{codebox}


\vspace{12pt}
{\bfseries\color{neonyellow} URL \& BROWSER}\vspace{4pt}

\begin{codebox}
zpwr google                \# google clipboard contents in tmux\\
zpwr openurl               \# open URL from clipboard in tmux\\
\end{codebox}


\vspace{12pt}
{\bfseries\color{neonyellow} SESSION MANAGEMENT}\vspace{4pt}

\begin{codebox}
zpwr attach                \# attach to tmux session\\
zpwr detach                \# detach from all tmux sessions\\
zpwr killmux               \# kill tmux server\\
zpwr killremote            \# kill tmux server and ssh\\
\end{codebox}


\vspace{16pt}


\chapter{ENVIRONMENT MASTERY}


\begin{center}

\includegraphics[width=0.85\textwidth]{/Users/wizard/.zpwr/docs/screenshots/art_system_core.png}
\end{center}

\section{Environment \& Search}


\vspace{10pt}
{\large\bfseries\color{neoncyan} ENVIRONMENT \& SEARCH}\vspace{6pt}

{\small\itshape\color{dimtext} // 22 verbs to inspect and search the shell //}

List env vars, search keybindings, inspect zsh internals, and navigate the cached environment.


\vspace{12pt}
{\bfseries\color{neonyellow} ENVIRONMENT VARIABLES}\vspace{4pt}

\begin{codebox}
zpwr ev                    \# list all zpwr env vars\\
zpwr envvars               \# list all zpwr env vars\\
zpwr environmentvars       \# list all zpwr env vars\\
zpwr environmentvariables  \# list all zpwr env vars\\
zpwr envvarsall            \# list all env vars\\
zpwr environmentvarsall    \# list all env vars\\
zpwr eva                   \# list all env vars\\
\end{codebox}


\vspace{12pt}
{\bfseries\color{neonyellow} ENVIRONMENT COUNTS \& CACHE}\vspace{4pt}

\begin{codebox}
zpwr e                     \# count all zpwr env\\
zpwr envcounts             \# count all zpwr env\\
zpwr environmentcounts     \# count all zpwr env\\
zpwr envcachesearch        \# fzf search cached aliases, params, builtins\\
zpwr environmentcachesearch\# fzf search cached env\\
zpwr envaccept             \# accept from fzf env\\
zpwr envedit               \# edit from fzf env\\
\end{codebox}


\vspace{12pt}
{\bfseries\color{neonyellow} KEYBINDING SEARCH}\vspace{4pt}

\begin{codebox}
zpwr allsearch             \# search all keybindings\\
zpwr vimsearch             \# search vim keybindings\\
zpwr zshsearch             \# search zsh keybindings\\
\end{codebox}


\vspace{12pt}
{\bfseries\color{neonyellow} ZSH INTERNALS}\vspace{4pt}

\begin{codebox}
zpwr reload                \# reparse env\\
zpwr modules               \# show zsh modules via fzf\\
zpwr options               \# show zsh options via fzf\\
zpwr zstyle                \# fuzzy search zstyle\\
zpwr zcompdump             \# edit zcompdump\\
\end{codebox}


\vspace{16pt}


\chapter{THE LOGGING SYSTEM}


\begin{center}

\includegraphics[width=0.85\textwidth]{/Users/wizard/.zpwr/docs/screenshots/art_data_stream.png}
\end{center}

\section{Logging System}


\vspace{10pt}
{\large\bfseries\color{neoncyan} LOGGING SYSTEM}\vspace{6pt}

{\small\itshape\color{dimtext} // 12 verbs for structured observability //}

\textbf{ZPWR} has a four-level logging system: trace, debug, info, error. Each level writes to file or console.


\vspace{12pt}
{\bfseries\color{neonyellow} LOG TO FILE}\vspace{4pt}

\begin{codebox}
zpwr logtrace              \# log trace to ZPWR\_LOGFILE\\
zpwr logdebug              \# log debug to ZPWR\_LOGFILE\\
zpwr loginfo               \# log info to ZPWR\_LOGFILE\\
zpwr logerror              \# log error to ZPWR\_LOGFILE\\
\end{codebox}


\vspace{12pt}
{\bfseries\color{neonyellow} LOG TO CONSOLE}\vspace{4pt}

\begin{codebox}
zpwr logtraceconsole       \# log trace to console\\
zpwr logdebugconsole       \# log debug to console\\
zpwr loginfoconsole        \# log info to console\\
zpwr logerrorconsole       \# log error to console\\
\end{codebox}


\vspace{12pt}
{\bfseries\color{neonyellow} LOG MANAGEMENT}\vspace{4pt}

\begin{codebox}
zpwr taillog               \# colorized tail of ZPWR\_LOGFILE\\
zpwr logs                  \# colorized tail of ZPWR\_LOGFILE\\
zpwr cleanlog              \# clear ZPWR\_LOGFILE\\
zpwr logincount            \# count of logins by user\\
\end{codebox}

Log levels from most verbose to least: TRACE >   DEBUG >   INFO >   ERROR.


\vspace{16pt}


\chapter{NETWORK \& CONTAINERS}


\begin{center}

\includegraphics[width=0.85\textwidth]{/Users/wizard/.zpwr/docs/screenshots/art_firewall.png}
\end{center}

\section{Process \& Network}


\vspace{10pt}
{\large\bfseries\color{neoncyan} PROCESS \& NETWORK}\vspace{6pt}

{\small\itshape\color{dimtext} // 15 verbs for procs, ports, and packets //}

Kill processes via fzf, monitor process trees, toggle Tor, ping LANs, and run network diagnostics.


\vspace{12pt}
{\bfseries\color{neonyellow} PROCESS MANAGEMENT}\vspace{4pt}

\begin{codebox}
zpwr ps                    \# ps -ef | grep arg\\
zpwr kill                  \# kill from ps output via fzf\\
zpwr killedit              \# edit kill from ps output via fzf\\
zpwr lsof                  \# kill from lsof output via fzf\\
zpwr lsofedit              \# edit kill from lsof output via fzf\\
zpwr pstreemonitor         \# monitor process tree with watch\\
\end{codebox}


\vspace{12pt}
{\bfseries\color{neonyellow} NETWORK DIAGNOSTICS}\vspace{4pt}

\begin{codebox}
zpwr digs                  \# run series of networking cmds on arg\\
zpwr pi                    \# ping all LAN devices\\
zpwr ping                  \# ping all LAN devices\\
zpwr pir                   \# run command on all devices\\
\end{codebox}


\vspace{12pt}
{\bfseries\color{neonyellow} TOR \& IP}\vspace{4pt}

\begin{codebox}
zpwr torip                 \# get IP address via Tor and whois\\
zpwr toriprenew            \# renew Tor identity\\
zpwr to                    \# toggle external ip\\
zpwr toggle                \# toggle external ip\\
\end{codebox}


\vspace{12pt}
{\bfseries\color{neonyellow} POLLING}\vspace{4pt}

\begin{codebox}
zpwr poll                  \# poll git remote and run command\\
\end{codebox}


\vspace{16pt}


\chapter{UTILITY ARSENAL}


\begin{center}
\includegraphics[width=0.85\textwidth]{/Users/wizard/.zpwr/docs/screenshots/art_mainframe.png}
\end{center}

\section{File \& Text Utilities}


\vspace{10pt}
{\large\bfseries\color{neoncyan} FILE \& TEXT UTILITIES}\vspace{6pt}

{\small\itshape\color{dimtext} // 31 verbs for file wrangling //}

Cat, list, grep, search, rename, and transform files. The everyday toolkit for text manipulation.


\vspace{12pt}
{\bfseries\color{neonyellow} VIEW \& LIST}\vspace{4pt}

\begin{codebox}
zpwr cat                   \# zpwr cat args\\
zpwr c                     \# zpwr cat args\\
zpwr ls                    \# list the files with no args\\
zpwr list                  \# list the files with no args\\
zpwr clearlist             \# clear and list the files\\
zpwr clearls               \# clear and list the files\\
zpwr info                  \# get info on command type with args\\
\end{codebox}


\vspace{12pt}
{\bfseries\color{neonyellow} SEARCH}\vspace{4pt}

\begin{codebox}
zpwr grep                  \# grep through pwd with ag into fzf\\
zpwr search                \# recursive grep search with context\\
zpwr filesearch            \# cat a file in a sub dir via fzf\\
zpwr filesearchedit        \# edit a file in a sub dir via fzf\\
zpwr wordsearch            \# cat a file by word via fzf\\
zpwr wordsearchedit        \# edit a file by word via fzf\\
zpwr findsearch            \# cat find drive for file via fzf\\
zpwr findsearchedit        \# edit find drive for file via fzf\\
zpwr locatesearch          \# accept locate drive for file\\
zpwr locatesearchedit      \# edit locate drive for file\\
zpwr linecount             \# get line count of search term\\
zpwr cfasd                 \# cat the fasd frecency ranked file\\
\end{codebox}


\vspace{12pt}
{\bfseries\color{neonyellow} TRANSFORM \& MANAGE}\vspace{4pt}

\begin{codebox}
zpwr replacer              \# find and replace text in files\\
zpwr rename                \# rename files\\
zpwr man                   \# fzf through man pages\\
zpwr open                  \# open with system\\
zpwr opencommand           \# get the command to open with system\\
\end{codebox}


\vspace{12pt}
{\bfseries\color{neonyellow} CONVERT \& UPLOAD}\vspace{4pt}

\begin{codebox}
zpwr scripttopdf           \# convert script to PDF\\
zpwr creategif             \# create GIF from file with ffmpeg\\
zpwr upload                \# upload with curl\\
zpwr urlsafe               \# base64 encode\\
zpwr execglob              \# exec on globbed files\\
zpwr execglobparallel      \# parallel exec on globbed files\\
\end{codebox}


\vspace{16pt}


\chapter{INTROSPECTION}


\begin{center}
\includegraphics[width=0.85\textwidth]{/Users/wizard/.zpwr/docs/screenshots/art_deep_dive.png}
\end{center}

\section{Shell Introspection}


\vspace{10pt}
{\large\bfseries\color{neoncyan} SHELL INTROSPECTION}\vspace{6pt}

{\small\itshape\color{dimtext} // 32 verbs to look inward //}

Inspect \textbf{ZPWR} itself: banners, counts, plugin lists, subcommand discovery, and identity checks.


\vspace{12pt}
{\bfseries\color{neonyellow} EXISTENCE CHECKS}\vspace{4pt}

\begin{codebox}
zpwr exists                \# check if identifier is valid\\
zpwr existscommand         \# check if command is valid\\
\end{codebox}


\vspace{12pt}
{\bfseries\color{neonyellow} BANNERS}\vspace{4pt}

\begin{codebox}
zpwr about                 \# show zpwr banner\\
zpwr banner                \# show zpwr banner\\
zpwr bannercounts          \# show banner and env counts\\
zpwr bannerlolcat          \# banner with ponysay and lolcat\\
zpwr bannernopony          \# show banner without ponysay\\
zpwr bannerpony            \# show banner with ponysay\\
\end{codebox}


\vspace{12pt}
{\bfseries\color{neonyellow} COUNTS \& LISTS}\vspace{4pt}

\begin{codebox}
zpwr autoloadcount         \# total number of autoloads\\
zpwr autoloadlist          \# total list of autoloads\\
zpwr scriptcount           \# total number of scripts\\
zpwr scriptlist            \# total list of scripts\\
zpwr zshplugincount        \# total number of zsh plugins\\
zpwr zshpluginlist         \# total list of zsh plugins\\
zpwr vimplugincount        \# total number of vim plugins\\
zpwr vimpluginlist         \# total list of vim plugins\\
zpwr emacsplugincount      \# total number of emacs plugins\\
zpwr emacspluginlist       \# total list of emacs plugins\\
\end{codebox}


\vspace{12pt}
{\bfseries\color{neonyellow} SUBCOMMANDS \& VERBS}\vspace{4pt}

\begin{codebox}
zpwr subcommands           \# run subcommands for zpwr <tab>\\
zpwr subcommandscount      \# number of subcommands\\
zpwr subcommandsedit       \# edit subcommands for zpwr <tab>\\
zpwr subcommandsfzf        \# fzf subcommands for zpwr <tab>\\
zpwr subcommandslist       \# list subcommands for zpwr <tab>\\
zpwr verbs                 \# run subcommands for zpwr <tab>\\
zpwr verbscount            \# number of verbs\\
zpwr verbsedit             \# edit subcommands for zpwr <tab>\\
zpwr verbsfzf              \# fzf subcommands for zpwr <tab>\\
zpwr verbslist             \# list subcommands for zpwr <tab>\\
zpwr help                  \# list subcommands with color\\
\end{codebox}


\vspace{12pt}
{\bfseries\color{neonyellow} MISC INTROSPECTION}\vspace{4pt}

\begin{codebox}
zpwr tty                   \# print current tty\\
zpwr relpath               \# get relative path to PWD\\
zpwr printmap              \# pretty print maps\\
\end{codebox}


\vspace{16pt}


\chapter{UTILITY ARSENAL}


\begin{center}
\includegraphics[width=0.85\textwidth]{/Users/wizard/.zpwr/docs/screenshots/art_mainframe.png}
\end{center}

\section{Miscellaneous Arsenal (1/2)}


\vspace{10pt}
{\large\bfseries\color{neoncyan} MISCELLANEOUS ARSENAL (1/2)}\vspace{6pt}

{\small\itshape\color{dimtext} // the verbs that defy categorization //}


\vspace{12pt}
{\bfseries\color{neonyellow} TERMINAL \& COLOR}\vspace{4pt}

\begin{codebox}
zpwr animate               \# run terminal animation\\
zpwr alacritty             \# edit alacritty config\\
zpwr color2                \# turn on stderr filter\\
zpwr return2               \# turn off stderr filter\\
zpwr colorsdiff            \# colorized side diff\\
zpwr colortest             \# test colors\\
zpwr colortest256          \# test 256 colors\\
zpwr figletfonts           \# show all figlet fonts\\
zpwr pygmentcolors         \# show all pygment colors\\
\end{codebox}


\vspace{12pt}
{\bfseries\color{neonyellow} BACKUP \& HISTORY}\vspace{4pt}

\begin{codebox}
zpwr backup                \# backup files\\
zpwr backuphistory         \# backup HISTFILE\\
zpwr hist                  \# exec history command via fzf\\
zpwr histedit              \# edit history command via fzf\\
zpwr histfile              \# edit history file\\
\end{codebox}


\vspace{12pt}
{\bfseries\color{neonyellow} FILE \& DIR CREATION}\vspace{4pt}

\begin{codebox}
zpwr makedir               \# make a dir tree\\
zpwr makefile              \# make a dir tree with file at end\\
zpwr scriptnew             \# create new script without editor\\
zpwr changename            \# mv dir into namespace and sed sub\\
zpwr changenamezpwr        \# mv into zpwr namespace and sed sub\\
\end{codebox}


\vspace{12pt}
{\bfseries\color{neonyellow} DOCKER \& GO}\vspace{4pt}

\begin{codebox}
zpwr dfimage               \# cat Dockerfile from IMAGE\_ID\\
zpwr dockerwipe            \# wipe all docker images\\
zpwr goclean               \# clean Go package artifacts\\
\end{codebox}


\vspace{12pt}
{\bfseries\color{neonyellow} EXERCISM}\vspace{4pt}

\begin{codebox}
zpwr exercism              \# go to Exercism home page\\
zpwr exercismedit          \# edit file in HOME/Exercism\\
zpwr exercismdownload      \# download all for track\\
\end{codebox}


\vspace{16pt}

\section{Miscellaneous Arsenal (2/2)}


\vspace{10pt}
{\large\bfseries\color{neoncyan} MISCELLANEOUS ARSENAL (2/2)}\vspace{6pt}

{\small\itshape\color{dimtext} // looping, testing, system, and services //}


\vspace{12pt}
{\bfseries\color{neonyellow} LOOPING \& BACKGROUND}\vspace{4pt}

\begin{codebox}
zpwr background            \# run arg in background\\
zpwr for                   \# run first arg times for command\\
zpwr for10                 \# run 10 times for command\\
zpwr fordir                \# run first arg in following dirs\\
zpwr fordirupdate          \# run git updaters in following dirs\\
zpwr whiletrue             \# run command forever\\
zpwr whilesleep            \# run command and sleep N seconds\\
\end{codebox}


\vspace{12pt}
{\bfseries\color{neonyellow} TESTING}\vspace{4pt}

\begin{codebox}
zpwr test                  \# run all zpwr tests\\
zpwr tests                 \# run all zpwr tests\\
zpwr testall               \# run all env tests\\
zpwr testsall              \# run all env tests\\
\end{codebox}


\vspace{12pt}
{\bfseries\color{neonyellow} SYSTEM \& MISC}\vspace{4pt}

\begin{codebox}
zpwr curl                  \# custom curl with cleaned output\\
zpwr deduppaths            \# deduplicate PATH fpath manpath\\
zpwr edit                  \# exec EDITOR\\
zpwr editor                \# exec EDITOR\\
zpwr install               \# run configure, make, make install\\
zpwr loadjenv              \# lazy load jenv\\
zpwr prettyprint           \# pretty print with color\\
zpwr altprettyprint        \# pretty with alternating color\\
zpwr reset                 \# reset repo to latest main\\
zpwr reveal                \# show remote repo in browser\\
zpwr revealrecurse         \# recursively reveal hidden dirs\\
zpwr rmconfiglinks         \# remove sym links to HOME\\
zpwr startauto             \# start autocomplete\\
zpwr stopauto              \# stop autocomplete\\
zpwr timer                 \# timer one or more commands\\
zpwr uninstall             \# uninstall all zpwr configs\\
zpwr post                  \# postfix all output\\
zpwr pre                   \# prefix all output\\
zpwr zpwrCloneToForked     \# clone repo to ZPWR\_FORKED\_DIR\\
\end{codebox}


\vspace{12pt}
{\bfseries\color{neonyellow} SYSTEMD SERVICES}\vspace{4pt}

\begin{codebox}
zpwr serviceup             \# start and enable systemd service\\
zpwr servicedown           \# stop and disable systemd service\\
zpwr restart               \# restart zpwr systemd service\\
\end{codebox}


\vspace{16pt}


\chapter{APPENDIX A: KEYBINDING REFERENCE}


\begin{center}
\includegraphics[width=0.85\textwidth]{/Users/wizard/.zpwr/docs/screenshots/art_the_grid.png}
\end{center}

\section{Tmux Prefix Keys: Navigation \& Panes}


\vspace{10pt}
{\large\bfseries\color{neoncyan} TMUX PREFIX KEYS: NAVIGATION \& PANES}\vspace{6pt}

pane traversal, splitting, resizing. the grid of power.\\


\vspace{12pt}
{\bfseries\color{neonyellow} PANE SELECTION (hjkl repeatable):}\vspace{4pt}

prefix h              select pane left\\
prefix j              select pane down\\
prefix k              select pane up\\
prefix l              select pane right\\
prefix ;              last-pane (toggle to previous)\\
prefix o              select next pane in order\\
prefix q              display-panes (5s timeout)\\


\vspace{12pt}
{\bfseries\color{neonyellow} PANE RESIZE (HJKL repeatable, 5 cells):}\vspace{4pt}

prefix H              resize pane left 5\\
prefix J              resize pane down 5\\
prefix K              resize pane up 5\\
prefix L              resize pane right 5\\
prefix Ctrl-Up/Down   resize pane by 1\\
prefix Alt-Up/Down    resize pane by 5\\


\vspace{12pt}
{\bfseries\color{neonyellow} PANE SPLITTING:}\vspace{4pt}

prefix -              split vertical (same dir)\\
prefix \textbackslash{}\textbackslash{}              split horizontal (same dir)\\
prefix "              split vertical (default dir)\\
prefix \%              split horizontal (default dir)\\
prefix |              split horizontal\\
prefix \_              split vertical\\


\vspace{12pt}
{\bfseries\color{neonyellow} PANE OPERATIONS:}\vspace{4pt}

prefix z              zoom/unzoom current pane\\
prefix !              break pane into own window\\
prefix x              kill current pane\\
prefix m              mark current pane\\
prefix +              split vertical then close prev pane\\
prefix Y              split vertical then close prev pane\\
prefix \{              rotate window\\
prefix \}              swap pane down\\


\vspace{12pt}
{\bfseries\color{neonyellow} ARROW KEY SELECTION:}\vspace{4pt}

prefix Up/Down        select pane up/down\\
prefix Left/Right     select pane left/right\\
Alt-Up/Down           select pane (root table, no prefix)\\
Alt-Left/Right        select pane (root table, no prefix)\\

All hjkl/HJKL bindings are repeatable (-r flag).\\
Arrow keys in root table work without prefix for speed.\\

\begin{codebox}
// the wizard navigates panes by instinct\\
// hjkl is muscle memory, arrows are for civilians\\
\end{codebox}


\vspace{16pt}

\section{Tmux Prefix Keys: Windows \& Sessions}


\vspace{10pt}
{\large\bfseries\color{neoncyan} TMUX PREFIX KEYS: WINDOWS \& SESSIONS}\vspace{6pt}

window management, session juggling, the multiplexer core.\\


\vspace{12pt}
{\bfseries\color{neonyellow} WINDOW CREATION \& SELECTION:}\vspace{4pt}

prefix c              new window\\
prefix a              last-window (toggle to previous)\\
prefix 0-9            select window by number\\
prefix '              select window by index prompt\\
prefix n              next window (repeatable)\\
prefix p              previous window (repeatable)\\
prefix Ctrl-n         next window\\
prefix Ctrl-p         previous window\\


\vspace{12pt}
{\bfseries\color{neonyellow} WINDOW OPERATIONS:}\vspace{4pt}

prefix ,              rename current window\\
prefix \&              kill window (with confirm)\\
prefix .              move window to target\\
prefix f              find window by name\\
prefix w              choose-tree window view\\
prefix i              display-message (window info)\\
prefix t              clock-mode\\


\vspace{12pt}
{\bfseries\color{neonyellow} SESSION MANAGEMENT:}\vspace{4pt}

prefix \$              rename current session\\
prefix (              switch to previous client (repeatable)\\
prefix )              switch to next client (repeatable)\\
prefix d              detach client\\
prefix s              choose-tree session view\\
prefix Ctrl-z         suspend client\\


\vspace{12pt}
{\bfseries\color{neonyellow} RELOAD \& PERSIST:}\vspace{4pt}

prefix r              reload tmux config + display confirmation\\
prefix Ctrl-s         save session (tmux-resurrect)\\
prefix Ctrl-r         restore session (tmux-resurrect)\\


\vspace{12pt}
{\bfseries\color{neonyellow} ROOT TABLE SHORTCUTS:}\vspace{4pt}

Ctrl-\textbackslash{}\textbackslash{}               switch to next client (no prefix)\\
Ctrl-]               switch to prev client (no prefix)\\


\vspace{12pt}
{\bfseries\color{neonyellow} META:}\vspace{4pt}

prefix ?              list all keys\\
prefix :              command prompt\\
prefix /              describe a key binding\\
prefix \textasciitilde{}              show messages\\

\begin{codebox}
// sessions are parallel universes\\
// the wizard keeps many realities open at once\\
\end{codebox}


\vspace{16pt}

\section{Tmux Prefix Keys: Layouts \& Pane Configs}


\vspace{10pt}
{\large\bfseries\color{neoncyan} TMUX PREFIX KEYS: LAYOUTS \& PANE CONFIGS}\vspace{6pt}

preconfigured pane arrangements. instant workspaces.\\


\vspace{12pt}
{\bfseries\color{neonyellow} \textbf{ZPWR} LAYOUT PRESETS:}\vspace{4pt}

prefix D              control-window layout\\
prefix E              four vertical panes\\
prefix F              four panes layout\\
prefix G              eight panes layout\\
prefix O              sixteen panes layout\\
prefix R              thirty-two panes REPL layout\\
prefix T              config files layout\\
prefix M              learn layout\\

Each layout sources a conf file from \textbf{\$ZPWR/tmux/}:\\
D -> control-window.conf\\
E -> fourVertical.conf\\
F -> four-panes.conf\\
G -> eight-panes.conf\\
O -> sixteen-panes.conf\\
R -> thirtytwo-panes-repl.conf\\
T -> config-files.conf\\
M -> learn.conf\\


\vspace{12pt}
{\bfseries\color{neonyellow} SYNCHRONIZE \& CUSTOMIZE:}\vspace{4pt}

prefix S              toggle synchronize-panes (type in all)\\
prefix C              customize-mode\\


\vspace{12pt}
{\bfseries\color{neonyellow} BUILT-IN LAYOUTS (Alt+number):}\vspace{4pt}

prefix Alt-1          even-horizontal\\
prefix Alt-2          even-vertical\\
prefix Alt-3          main-horizontal\\
prefix Alt-4          main-vertical\\
prefix Alt-5          tiled\\
prefix Alt-6          main-horizontal-mirrored\\
prefix Alt-7          main-vertical-mirrored\\


\vspace{12pt}
{\bfseries\color{neonyellow} WINDOW ROTATION:}\vspace{4pt}

prefix Ctrl-o         rotate window panes\\
prefix Alt-o          rotate window panes (reverse)\\

Synchronize-panes (prefix S) sends keystrokes to ALL panes.\\
Perfect for running the same command across many servers.\\

\begin{codebox}
// sixteen panes is not chaos\\
// sixteen panes is the wizard's dashboard\\
\end{codebox}


\vspace{16pt}

\section{Tmux Prefix Keys: Special Operations}


\vspace{10pt}
{\large\bfseries\color{neoncyan} TMUX PREFIX KEYS: SPECIAL OPERATIONS}\vspace{6pt}

multi-pane execution, URL picking, buffers, plugins.\\


\vspace{12pt}
{\bfseries\color{neonyellow} MULTI-PANE EXECUTION:}\vspace{4pt}

prefix Space          run command in all panes (single)\\
prefix v              run command in all panes (multi)\\
prefix b              open URL/path in all panes\\
prefix g              google search in all panes\\

These source \textbf{\$ZPWR/scripts/allPanes.zsh} with different modes.\\
Space -> single mode (one command per pane)\\
v     -> multi mode (batch operations)\\
b     -> single open (open URL in each pane)\\
g     -> single google (search in each pane)\\


\vspace{12pt}
{\bfseries\color{neonyellow} URL HANDLING:}\vspace{4pt}

prefix e              fzf URL picker -> open\\
prefix u              fzf URL picker -> search\\

Both use tmux-fzf-url plugin with 30s timeout.\\


\vspace{12pt}
{\bfseries\color{neonyellow} PASTE \& BUFFERS:}\vspace{4pt}

prefix P              paste buffer\\
prefix ]              paste buffer (bracketed)\\
prefix =              choose-buffer (fzf style)\\
prefix \#              list buffers\\
prefix Ctrl-v         paste from system clipboard\\


\vspace{12pt}
{\bfseries\color{neonyellow} COPY MODE:}\vspace{4pt}

prefix [              enter copy mode\\
prefix PgUp           enter copy mode and page up\\


\vspace{12pt}
{\bfseries\color{neonyellow} PLUGIN MANAGEMENT (TPM):}\vspace{4pt}

prefix I              install plugins\\
prefix U              update plugins\\
prefix Alt-u          clean unused plugins\\


\vspace{12pt}
{\bfseries\color{neonyellow} SEND PREFIX:}\vspace{4pt}

prefix Ctrl-a         send prefix key to inner tmux\\
prefix Ctrl-]         send literal Ctrl-]\\

Use prefix Ctrl-a when nesting tmux sessions.\\
The outer session captures prefix, Ctrl-a forwards it.\\

\begin{codebox}
// buffers are the wizard's memory\\
// every copy is a spell stored for later\\
\end{codebox}


\vspace{16pt}

\section{Tmux Copy Mode (vi)}


\vspace{10pt}
{\large\bfseries\color{neoncyan} TMUX COPY MODE (VI BINDINGS)}\vspace{6pt}

scrollback navigation, selection, copy, search. pure vi.\\


\vspace{12pt}
{\bfseries\color{neonyellow} ENTER/EXIT:}\vspace{4pt}

prefix [              enter copy mode\\
q                     cancel and exit copy mode\\
Escape                clear selection\\


\vspace{12pt}
{\bfseries\color{neonyellow} SELECTION \& COPY:}\vspace{4pt}

v                     begin selection\\
V                     select entire line\\
Ctrl-v               toggle rectangle selection\\
Space                 begin selection (alternate)\\
y                     copy selection to clipboard (no clear)\\
Enter                 copy selection to clipboard and cancel\\
D                     copy to end of line and cancel\\
A                     append selection and cancel\\


\vspace{12pt}
{\bfseries\color{neonyellow} NAVIGATION:}\vspace{4pt}

h/j/k/l              cursor movement\\
Ctrl-h/j/k/l         move cursor 4 cells at a time\\
w/b/e                word forward/backward/end\\
W/B/E                WORD (space-delimited) movement\\
0 / \$                start/end of line\\
\textasciicircum{}                     first non-blank character\\
g / G                top / bottom of history\\
H / M / L            top / middle / bottom of screen\\


\vspace{12pt}
{\bfseries\color{neonyellow} SCROLLING:}\vspace{4pt}

Ctrl-u / Ctrl-d      half page up / down\\
Ctrl-b / Ctrl-f      full page up / down\\
Ctrl-y / Ctrl-e      scroll up / down one line\\
J / K                scroll down / up one line\\


\vspace{12pt}
{\bfseries\color{neonyellow} SEARCH:}\vspace{4pt}

/                     search forward (incremental)\\
?                     search backward (incremental)\\
n / N                next / previous match\\
*                     search forward for word under cursor\\
\#                     search backward for word under cursor\\


\vspace{12pt}
{\bfseries\color{neonyellow} \textbf{ZPWR} EXTENSIONS:}\vspace{4pt}

s                     copy + google search selection\\
x                     copy + open URL selection\\
z                     copy + google search selection\\
RightClick            copy + google search selection\\

Copy goes to system clipboard via reattach-to-user-namespace pbcopy.\\
\textbf{ZPWR} extends vi copy mode with instant google/open actions.\\

\begin{codebox}
// in copy mode the wizard moves through time\\
// scrollback is archaeology, search is divination\\
\end{codebox}


\vspace{16pt}

\section{Zsh Insert Mode: Basic \& Movement}


\vspace{10pt}
{\large\bfseries\color{neoncyan} ZSH INSERT MODE: BASIC \& MOVEMENT}\vspace{6pt}

the primary editing mode. where commands are born.\\


\vspace{12pt}
{\bfseries\color{neonyellow} LINE MOVEMENT:}\vspace{4pt}

Ctrl-A                beginning-of-line\\
Ctrl-E                end-of-line\\
Ctrl-O                EOL or next tab stop\\


\vspace{12pt}
{\bfseries\color{neonyellow} EDITING:}\vspace{4pt}

Ctrl-U                clear entire line\\
Ctrl-W                delete word backward (autopair)\\
Ctrl-H                delete char backward (autopair)\\
Ctrl-Y                change surrounding quotes\\
Ctrl-Z                undo\\
Ctrl-R                redo\\
Ctrl-Q                duplicate last word\\


\vspace{12pt}
{\bfseries\color{neonyellow} \textbf{ZPWR} SPECIAL:}\vspace{4pt}

Space                 zpwr supernatural expand (aliases+globs)\\
Ctrl-@                zpwr expand + terminate space\\
Ctrl-M (Enter)        zpwr magic enter (show status on empty)\\
.                     zpwr rationalize dot (.. -> ../..)\\
Ctrl-B                zpwr clipboard paste\\
Ctrl-S                git add commit push (no check)\\
Ctrl-K                zsh-learn-Learn (flashcard)\\


\vspace{12pt}
{\bfseries\color{neonyellow} COMPLETION \& FZF:}\vspace{4pt}

Tab (Ctrl-I)          fzf-completion\\
Ctrl-D                list-choices\\
Ctrl-L                clear-screen\\
Ctrl-T                fzf file widget\\
Alt-c                 fzf cd widget\\


\vspace{12pt}
{\bfseries\color{neonyellow} HISTORY:}\vspace{4pt}

Ctrl-N                down-history\\
Ctrl-P                up-history\\
Up/Down arrows        history substring search up/down\\


\vspace{12pt}
{\bfseries\color{neonyellow} MODE SWITCHING:}\vspace{4pt}

Escape                enter vi command mode\\
jf / fj              enter vi command mode (chord escape)\\


\vspace{12pt}
{\bfseries\color{neonyellow} SURROUND \& PAIRS:}\vspace{4pt}

" ' ( [ \{ \textbackslash{}`          auto-pair/surround insert\\
) ] \}                 autopair close\\
Backspace             autopair-aware delete\\

Space is the most powerful key: it expands aliases,\\
global aliases, and abbreviations before inserting.\\

\begin{codebox}
// insert mode is where spells are typed\\
// every keystroke is an incantation\\
\end{codebox}


\vspace{16pt}

\section{Zsh Insert Mode: Ctrl-F \& Ctrl-V Chords}


\vspace{10pt}
{\large\bfseries\color{neoncyan} ZSH INSERT MODE: CTRL-F \& CTRL-V CHORDS}\vspace{6pt}

two-key combos that unlock the full zpwr arsenal.\\


\vspace{12pt}
{\bfseries\color{neonyellow} CTRL-F CHORDS (fzf \& actions):}\vspace{4pt}

Ctrl-F Ctrl-D         pipe buffer into fzf\\
Ctrl-F Ctrl-F         cd + vim fzf file search (accept)\\
Ctrl-F Ctrl-G         google search\\
Ctrl-F Ctrl-H         lsof into fzf\\
Ctrl-F Ctrl-I         ag (silver searcher) into fzf\\
Ctrl-F Ctrl-J         zpwr verbs widget (accept)\\
Ctrl-F Ctrl-K         alternate quotes around buffer\\
Ctrl-F Ctrl-L         list choices\\
Ctrl-F Ctrl-M         zzcomplete\\
Ctrl-F Ctrl-N         zpwr verbs widget (browse)\\
Ctrl-F Ctrl-O         open URL from buffer\\
Ctrl-F Ctrl-P         basic sed substitution\\
Ctrl-F Ctrl-R         wrap buffer as variable\\
Ctrl-F Ctrl-S         git diff check (gacp)\\
Ctrl-F Ctrl-V         edit command line in \$\textbf{EDITOR}\\
Ctrl-F j              zpwr verbs widget (accept)\\
Ctrl-F n              zpwr verbs widget (browse)\\


\vspace{12pt}
{\bfseries\color{neonyellow} CTRL-V CHORDS (vim/fzf/navigation):}\vspace{4pt}

Ctrl-V Ctrl-@         vim fzf (open file in vim)\\
Ctrl-V Space          vim fzf (open file in vim)\\
Ctrl-V Ctrl-F         fasd fzf (frecency file open)\\
Ctrl-V Ctrl-G         fzf cd widget\\
Ctrl-V Ctrl-K         emacs fzf\\
Ctrl-V Ctrl-L         z fzf (frecency dirs)\\
Ctrl-V Ctrl-N         vim fzf sudo\\
Ctrl-V Ctrl-O         fzf tab complete\\
Ctrl-V Ctrl-P         sudo command line toggle\\
Ctrl-V Ctrl-R         history search multi-word\\
Ctrl-V Ctrl-S         z fzf (alternate)\\
Ctrl-V Ctrl-V         vim all files widget (accept)\\
Ctrl-V Ctrl-Z         fzf history widget\\
Ctrl-V ,              fzf environment variables\\
Ctrl-V .              fzf all keybindings\\

{\small\itshape\color{dimtext} Ctrl-V //             locate fzf}

Ctrl-V /.             locate fzf + edit\\
Ctrl-V ;              fzf surround\\
Ctrl-V c              fzf git commits\\
Ctrl-V k              fzf vim keybindings\\
Ctrl-V m              vim all files widget (browse)\\
Ctrl-V v              vim all files widget (accept)\\

Ctrl-F is the action chord: search, edit, google, diff.\\
Ctrl-V is the vim/fzf chord: open files, browse, navigate.\\

\begin{codebox}
// chords are the wizard's combos\\
// two keys unlock what single keys cannot reach\\
\end{codebox}


\vspace{16pt}

\section{Zsh Visual \& Menu Select Modes}


\vspace{10pt}
{\large\bfseries\color{neoncyan} ZSH VISUAL \& MENU SELECT MODES}\vspace{6pt}

text objects, selection, and completion menu navigation.\\


\vspace{12pt}
{\bfseries\color{neonyellow} \textbf{VISUAL} MODE (enter with v/V in vicmd):}\vspace{4pt}

Escape                deactivate region\\
j / k                move selection down / up\\
o                     exchange point and mark\\
p                     put replace selection\\
x                     delete selection\\
u                     downcase selection\\
U                     upcase selection\\
\textasciitilde{}                     swap case of selection\\


\vspace{12pt}
{\bfseries\color{neonyellow} TEXT OBJECTS (a = around, i = inner):}\vspace{4pt}

a" / i"              select around/inner double quotes\\
a' / i'              select around/inner single quotes\\
a\textbackslash{}` / i\textbackslash{}`              select around/inner backticks\\
a( / i(              select around/inner parentheses\\
a[ / i[              select around/inner brackets\\
a\{ / i\{              select around/inner braces\\
a< / i<              select around/inner angle brackets\\
aB / iB              select around/inner block\\
aw / iw              select a word / in word\\
aW / iW              select a blank word / in blank word\\
aa / ia              select a shell word / in shell word\\

Extended quote chars also work as text objects:\\
+ , - . / : ; = @ | \textbackslash{}\textbackslash{} are all valid quote delimiters.\\


\vspace{12pt}
{\bfseries\color{neonyellow} MENU SELECT MODE (during completion):}\vspace{4pt}

Ctrl-J / Ctrl-K       move down / up in menu\\
Ctrl-H / Ctrl-L       move left / right in menu\\
Tab / Ctrl-I          move forward (next item)\\
Ctrl-S / Shift-Tab    reverse menu complete\\
Ctrl-@ (Ctrl-Space)   accept line\\
Ctrl-D                accept and menu complete (stay open)\\
Ctrl-F                accept and infer next\\
Ctrl-M (Enter)        accept line (.accept-line)\\
Ctrl-V                insert mode within menu\\
Ctrl-N / Ctrl-P       forward / backward word in menu\\
Ctrl-X                incremental search forward in menu\\
?                     incremental search backward in menu\\
.                     self-insert (exit menu + type dot)\\
Backspace             undo last menu selection\\

Menu select uses hjkl-style navigation with Ctrl prefix.\\
Ctrl-D is key: accept current match but keep menu open.\\

\begin{codebox}
// visual mode is precision surgery\\
// select only what you need, transform only what you touch\\
\end{codebox}


\vspace{16pt}

\section{Vim Normal Mode Highlights}


\vspace{10pt}
{\large\bfseries\color{neoncyan} VIM NORMAL MODE: ZPWR HIGHLIGHTS}\vspace{6pt}

leader mappings, fzf integration, tmux synergy.\\
leader key is Space. only zpwr-specific bindings listed.\\


\vspace{12pt}
{\bfseries\color{neonyellow} FILE \& BUFFER (Space prefix):}\vspace{4pt}

Space w               :w! (force write)\\
Space e               :q! (force quit)\\
Space q               :qa! (quit all)\\
Space wq              :wq! (write and quit)\\
Space n / Space p     next / prev buffer\\
Space an / Space ap   next / prev in arglist\\
Space Tab             alternate or next buffer\\
Space t               new tab\\
Space s / Space v     horizontal / vertical split\\


\vspace{12pt}
{\bfseries\color{neonyellow} FZF INTEGRATION (Space prefix):}\vspace{4pt}

Space f               :Files! (fzf files)\\
Space b               :Buffers! (fzf buffers)\\
Space j               :Lines! (fzf lines)\\
Space ag              :Ag (silver searcher)\\
Space aa              :Agg (ag variant)\\
Space rg              :Rg (ripgrep)\\
Space z               :Commits! (fzf git commits)\\
Space ;               :Commands! (fzf commands)\\
Space ,               :FZFMaps (fzf all mappings)\\
Space ke / , / m / i  fzf keys / maps / nmap / imap\\
Space ta / ma / env  fzf tags / marks / env vars\\
Space / / .           locate all / colorscheme picker\\


\vspace{12pt}
{\bfseries\color{neonyellow} HISTORY (Space prefix):}\vspace{4pt}

Space hh              :History! (fzf file history)\\
Space h/              :History/! (fzf search history)\\
Space h:              :History:! (fzf command history)\\


\vspace{12pt}
{\bfseries\color{neonyellow} SEARCH \& REPLACE (Space prefix):}\vspace{4pt}

Space r               replace word under cursor (whole file)\\
Space g               global substitute word under cursor\\
Space x               repeat last substitution on line\\
n / N                 next / prev search (centered)\\
z/                    toggle auto-highlight word under cursor\\


\vspace{12pt}
{\bfseries\color{neonyellow} GIT \& TMUX:}\vspace{4pt}

Space hs              GitGutter stage hunk\\
Space hu              GitGutter undo hunk\\
Space hp              GitGutter preview hunk\\
{[}c / ]c              prev / next git hunk\\
Space vj              save + TmuxRepeat file\\
Space os / ox / oc   google / open URL / copy word\\


\vspace{12pt}
{\bfseries\color{neonyellow} MINIMAP \& COMMENTS:}\vspace{4pt}

Space mm              open minimap\\
Space mt              toggle minimap\\
Space cc              NERDCommenter comment\\
Space c Space         NERDCommenter toggle\\
Space cu              NERDCommenter uncomment\\

\begin{codebox}
// the leader key is Space because power needs room\\
// every binding is one thought away from execution\\
\end{codebox}


\vspace{16pt}

\section{Vim Insert \& Visual Mode + Keybinding Search}


\vspace{10pt}
{\large\bfseries\color{neoncyan} VIM INSERT \& VISUAL MODE + KEYBINDING SEARCH}\vspace{6pt}

vim mode highlights and how to find any keybinding.\\


\vspace{12pt}
{\bfseries\color{neonyellow} VIM INSERT MODE HIGHLIGHTS:}\vspace{4pt}

Ctrl-A                jump to first non-blank (like \textasciicircum{})\\
Ctrl-E                end of line or dismiss popup\\
Ctrl-D                delete char forward or default\\
Ctrl-F                move right or scroll\\
Ctrl-S                isurround (surround insert)\\
Ctrl-G s / S          isurround / ISurround (line)\\
Ctrl-Tab              list ultisnips snippets\\
Shift-Tab             SuperTab backward completion\\
Alt-d                 delete word forward\\
" ' ( ) [ ] \textbackslash{}` \{ \}    delimitMate auto-pairing\\


\vspace{12pt}
{\bfseries\color{neonyellow} VIM \textbf{VISUAL} MODE HIGHLIGHTS:}\vspace{4pt}

Space r               replace word in selection range\\
Space g               global substitute in selection range\\
Space x               repeat last sub on selected lines\\
Space " ' \textbackslash{}` ( [ \{    surround selection with quote/bracket\\
Space b               copy selection to tmux buffer\\
Space = / Space -     grow / shrink selection (4 lines)\\
J / K                move selected lines down / up\\
Y                     yank and move to end\\
< / >                indent / outdent (stay in visual)\\
S                     vim-surround visual surround\\
go                    copy selection + open as URL\\
gs                    copy selection + google search\\
Space vl / vk / vj   TmuxRepeat repl / visual / file\\
Space em / ev / ef   extract method / variable / fold\\


\vspace{12pt}
{\bfseries\color{neonyellow} SNEAK \& MOTION:}\vspace{4pt}

s / S                sneak forward / backward (2-char)\\
f / F / t / T        sneak-enhanced find char\\
; / ,                repeat / reverse last sneak/find\\
w / b / e            CamelCaseMotion word movement\\


\vspace{10pt}
{\large\bfseries\color{neoncyan} =}\vspace{6pt}


\vspace{12pt}
{\bfseries\color{neonyellow} HOW TO SEARCH KEYBINDINGS:}\vspace{4pt}


\vspace{10pt}
{\large\bfseries\color{neoncyan} =}\vspace{6pt}

zpwr allsearch        search ALL keybindings in fzf\\
zpwr vimsearch        search vim keybindings in fzf\\
zpwr zshsearch        search zsh keybindings in fzf\\


\vspace{12pt}
{\bfseries\color{neonyellow} KEYBINDING CACHE FILES:}\vspace{4pt}

\textbf{\$ZPWR\_LOCAL/zpwrAllKeybindings.txt}\\
\textbf{\$ZPWR\_LOCAL/zpwr}\{Vim,Zsh,Tmux\}Keybindings.txt\\


\vspace{12pt}
{\bfseries\color{neonyellow} REGENERATE:}\vspace{4pt}

zpwr regenkeybindings  regenerate all keybinding caches\\

Cache files are auto-generated. Run regenkeybindings after\\
adding new bindings to refresh the searchable index.\\

\begin{codebox}
// the wizard does not memorize 4000 keybindings\\
// the wizard searches them at the speed of thought\\
// fzf is the oracle, the cache is the grimoire\\
\end{codebox}


\vspace{16pt}


\chapter{APPENDIX B: COMPLETE KEYBINDING DUMP}


\begin{center}
\includegraphics[width=0.85\textwidth]{/Users/wizard/.zpwr/docs/screenshots/art_the_grid.png}
\end{center}

\section{ALL Tmux Prefix Keys (1/4)}


\vspace{10pt}
{\large\bfseries\color{neoncyan} ALL TMUX PREFIX KEYS (1/4)}\vspace{6pt}

bind-key  -T prefix  Space  run-shell -b '\textbf{ZPWR}/scripts/allPanes.zsh single\\
bind-key  -T prefix  .  break-pane\\
bind-key  -T prefix  \textbackslash{}'  split-window\\
bind-key  -T prefix  \textbackslash{}\#  list-buffers\\
bind-key  -T prefix  \textbackslash{}  command-prompt -I '\#S' \{ rename-session '\%\%' \}\\
bind-key  -T prefix  \textbackslash{}\%  split-window -h\\
bind-key  -T prefix  \&  confirm-before -p 'kill-window \#W? (y/n)' kill-win\\
bind-key  -T prefix  \textbackslash{}'  command-prompt -T window-target -p index \{ selec\\
bind-key -r -T prefix  (  switch-client -p\\
bind-key -r -T prefix  )  switch-client -n\\
bind-key  -T prefix  +  split-window -v -c '\#\{pane\_current\_path\}' \textbackslash{}; sele\\
bind-key  -T prefix  ,  command-prompt -I '\#W' \{ rename-window '\%\%' \}\\
bind-key  -T prefix  -  split-window -v -c '\#\{pane\_current\_path\}'\\
bind-key  -T prefix  .  command-prompt -T target \{ move-window -t '\%\%' \}\\
bind-key  -T prefix  /  command-prompt -k -p key \{ list-keys -1N '\%\%' \}\\
bind-key  -T prefix  0  select-window -t :=0\\
bind-key  -T prefix  1  select-window -t :=1\\
bind-key  -T prefix  2  select-window -t :=2\\
bind-key  -T prefix  3  select-window -t :=3\\
bind-key  -T prefix  4  select-window -t :=4\\
bind-key  -T prefix  5  select-window -t :=5\\
bind-key  -T prefix  6  select-window -t :=6\\
bind-key  -T prefix  7  select-window -t :=7\\
bind-key  -T prefix  8  select-window -t :=8\\
bind-key  -T prefix  9  select-window -t :=9\\
bind-key  -T prefix  :  command-prompt\\
bind-key  -T prefix  \textbackslash{};  last-pane\\
bind-key  -T prefix  <  display-menu -T '\#[align=centre]\#\{window\_index\}:\#\{\\
bind-key  -T prefix  =  choose-buffer -Z\\
bind-key  -T prefix  >  display-menu -T '\#[align=centre]\#\{pane\_index\} (\#\{p\\
bind-key  -T prefix  ?  list-keys -N\\
bind-key  -T prefix  C  customize-mode -Z\\
bind-key  -T prefix  D  source-file \textbf{ZPWR}/tmux/control-window.conf\\
bind-key  -T prefix  E  source-file \textbf{ZPWR}/tmux/fourVertical.conf\\
bind-key  -T prefix  F  source-file \textbf{ZPWR}/tmux/four-panes.conf\\
bind-key  -T prefix  G  source-file \textbf{ZPWR}/tmux/eight-panes.conf\\
bind-key -r -T prefix  H  resize-pane -L 5\\
bind-key  -T prefix  I  run-shell \textbf{HOME}/.tmux/plugins/tpm/bindings/install\_\\
bind-key -r -T prefix  J  resize-pane -D 5\\
bind-key -r -T prefix  K  resize-pane -U 5\\
bind-key -r -T prefix  L  resize-pane -R 5\\
bind-key  -T prefix  M  source-file \textbf{ZPWR}/tmux/learn.conf\\
bind-key  -T prefix  O  source-file \textbf{ZPWR}/tmux/sixteen-panes.conf\\
bind-key  -T prefix  P  paste-buffer\\
bind-key  -T prefix  R  source-file \textbf{ZPWR}/tmux/thirtytwo-panes-repl.conf\\
bind-key  -T prefix  S  set-window-option synchronize-panes\\
bind-key  -T prefix  T  source-file \textbf{ZPWR}/tmux/config-files.conf\\
bind-key  -T prefix  U  run-shell \textbf{HOME}/.tmux/plugins/tpm/bindings/update\_p\\


\vspace{16pt}

\section{ALL Tmux Prefix Keys (2/4)}


\vspace{10pt}
{\large\bfseries\color{neoncyan} ALL TMUX PREFIX KEYS (2/4)}\vspace{6pt}

bind-key  -T prefix  Y  split-window -v -c '\#\{pane\_current\_path\}' \textbackslash{}; sele\\
bind-key  -T prefix  [  copy-mode\\
bind-key  -T prefix  \textbackslash{}\textbackslash{}  split-window -h -c '\#\{pane\_current\_path\}'\\
bind-key  -T prefix  ]  paste-buffer -p\\
bind-key  -T prefix  \_  split-window -v\\
bind-key  -T prefix  a  last-window\\
bind-key  -T prefix  b  run-shell -b '\textbf{ZPWR}/scripts/allPanes.zsh single ope\\
bind-key  -T prefix  c  new-window\\
bind-key  -T prefix  d  detach-client\\
bind-key  -T prefix  e  run-shell -b '\textbf{HOME}/.tmux/plugins/tmux-fzf-url/fzf-\\
bind-key  -T prefix  f  command-prompt \{ find-window -Z '\%\%' \}\\
bind-key  -T prefix  g  run-shell -b '\textbf{ZPWR}/scripts/allPanes.zsh single goo\\
bind-key -r -T prefix  h  select-pane -L\\
bind-key  -T prefix  i  display-message\\
bind-key -r -T prefix  j  select-pane -D\\
bind-key -r -T prefix  k  select-pane -U\\
bind-key -r -T prefix  l  select-pane -R\\
bind-key  -T prefix  m  select-pane -m\\
bind-key -r -T prefix  n  next-window\\
bind-key  -T prefix  o  select-pane -t :.+\\
bind-key -r -T prefix  p  previous-window\\
bind-key  -T prefix  q  display-panes -d 5000\\
bind-key  -T prefix  r  source-file \textbf{ZPWR}/tmux/init.conf \textbackslash{}; display-messag\\
bind-key  -T prefix  s  choose-tree -Zs\\
bind-key  -T prefix  t  clock-mode\\
bind-key  -T prefix  u  run-shell -b '\textbf{HOME}/.tmux/plugins/tmux-fzf-url/fzf-\\
bind-key  -T prefix  v  run-shell -b '\textbf{ZPWR}/scripts/allPanes.zsh multi'\\
bind-key  -T prefix  w  choose-tree -Zw\\
bind-key  -T prefix  x  kill-pane\\
bind-key  -T prefix  z  resize-pane -Z\\
bind-key -r -T prefix  \textbackslash{}\{  rotate-window\\
bind-key  -T prefix  |  split-window -h\\
bind-key  -T prefix  \textbackslash{}\}  swap-pane -D\\
bind-key  -T prefix  \textbackslash{}\textasciitilde{}  show-messages\\
bind-key -r -T prefix  DC  refresh-client -c\\
bind-key  -T prefix  PPage  copy-mode -u\\
bind-key -r -T prefix  Up  select-pane -U\\
bind-key -r -T prefix  Down  select-pane -D\\
bind-key -r -T prefix  Left  select-pane -L\\
bind-key -r -T prefix  Right  select-pane -R\\
bind-key  -T prefix  M-1  select-layout even-horizontal\\
bind-key  -T prefix  M-2  select-layout even-vertical\\
bind-key  -T prefix  M-3  select-layout main-horizontal\\
bind-key  -T prefix  M-4  select-layout main-vertical\\
bind-key  -T prefix  M-5  select-layout tiled\\
bind-key  -T prefix  M-6  select-layout main-horizontal-mirrored\\
bind-key  -T prefix  M-7  select-layout main-vertical-mirrored\\
bind-key  -T prefix  M-n  next-window -a\\


\vspace{16pt}

\section{ALL Tmux Prefix Keys (3/4)}


\vspace{10pt}
{\large\bfseries\color{neoncyan} ALL TMUX PREFIX KEYS (3/4)}\vspace{6pt}

bind-key  -T prefix  M-o  rotate-window -D\\
bind-key  -T prefix  M-p  previous-window -a\\
bind-key  -T prefix  M-u  run-shell \textbf{HOME}/.tmux/plugins/tpm/bindings/clean\_\\
bind-key -r -T prefix  M-Up  resize-pane -U 5\\
bind-key -r -T prefix  M-Down  resize-pane -D 5\\
bind-key -r -T prefix  M-Left  resize-pane -L 5\\
bind-key -r -T prefix  M-Right  resize-pane -R 5\\
bind-key  -T prefix  C-]  send-keys \textasciicircum{}]\\
bind-key  -T prefix  C-a  send-prefix\\
bind-key  -T prefix  C-n  next-window\\
bind-key  -T prefix  C-o  rotate-window\\
bind-key  -T prefix  C-p  previous-window\\
bind-key  -T prefix  C-r  run-shell \textbf{HOME}/.tmux/plugins/tmux-resurrect/scri\\
bind-key  -T prefix  C-s  run-shell \textbf{HOME}/.tmux/plugins/tmux-resurrect/scri\\
bind-key  -T prefix  C-v  run-shell 'tmux set buffer \textbackslash{}'(reattach-to-user-\\
bind-key  -T prefix  C-z  suspend-client\\
bind-key -r -T prefix  C-Up  resize-pane -U\\
bind-key -r -T prefix  C-Down  resize-pane -D\\
bind-key -r -T prefix  C-Left  resize-pane -L\\
bind-key -r -T prefix  C-Right  resize-pane -R\\
bind-key -r -T prefix  S-Up  refresh-client -U 10\\
bind-key -r -T prefix  S-Down  refresh-client -D 10\\
bind-key -r -T prefix  S-Left  refresh-client -L 10\\
bind-key -r -T prefix  S-Right  refresh-client -R 10\\
i  <C-X><C-L>  * fzf\#vim\#complete(fzf\#wrap(\{ 'prefix': '\textasciicircum{}.*', 'source': 'r\\
<Space><Space>  <Plug>(easymotion-prefix)\\
<Plug>(easymotion-prefix)N  <Plug>(easymotion-N)\\
<Plug>(easymotion-prefix)n  <Plug>(easymotion-n)\\
<Plug>(easymotion-prefix)k  <Plug>(easymotion-k)\\
<Plug>(easymotion-prefix)j  <Plug>(easymotion-j)\\
<Plug>(easymotion-prefix)gE  <Plug>(easymotion-gE)\\
<Plug>(easymotion-prefix)ge  <Plug>(easymotion-ge)\\
<Plug>(easymotion-prefix)E  <Plug>(easymotion-E)\\
<Plug>(easymotion-prefix)e  <Plug>(easymotion-e)\\
<Plug>(easymotion-prefix)B  <Plug>(easymotion-B)\\
<Plug>(easymotion-prefix)b  <Plug>(easymotion-b)\\
<Plug>(easymotion-prefix)W  <Plug>(easymotion-W)\\
<Plug>(easymotion-prefix)w  <Plug>(easymotion-w)\\
<Plug>(easymotion-prefix)T  <Plug>(easymotion-T)\\
<Plug>(easymotion-prefix)t  <Plug>(easymotion-t)\\
<Plug>(easymotion-prefix)s  <Plug>(easymotion-s)\\
<Plug>(easymotion-prefix)F  <Plug>(easymotion-F)\\
<Plug>(easymotion-prefix)f  <Plug>(easymotion-f)\\
<Space><Space>  <Plug>(easymotion-prefix)\\
<Plug>(easymotion-prefix)N  <Plug>(easymotion-N)\\
<Plug>(easymotion-prefix)n  <Plug>(easymotion-n)\\
<Plug>(easymotion-prefix)k  <Plug>(easymotion-k)\\
<Plug>(easymotion-prefix)j  <Plug>(easymotion-j)\\


\vspace{16pt}

\section{ALL Tmux Prefix Keys (4/4)}


\vspace{10pt}
{\large\bfseries\color{neoncyan} ALL TMUX PREFIX KEYS (4/4)}\vspace{6pt}

<Plug>(easymotion-prefix)gE  <Plug>(easymotion-gE)\\
<Plug>(easymotion-prefix)ge  <Plug>(easymotion-ge)\\
<Plug>(easymotion-prefix)E  <Plug>(easymotion-E)\\
<Plug>(easymotion-prefix)e  <Plug>(easymotion-e)\\
<Plug>(easymotion-prefix)B  <Plug>(easymotion-B)\\
<Plug>(easymotion-prefix)b  <Plug>(easymotion-b)\\
<Plug>(easymotion-prefix)W  <Plug>(easymotion-W)\\
<Plug>(easymotion-prefix)w  <Plug>(easymotion-w)\\
<Plug>(easymotion-prefix)T  <Plug>(easymotion-T)\\
<Plug>(easymotion-prefix)t  <Plug>(easymotion-t)\\
<Plug>(easymotion-prefix)s  <Plug>(easymotion-s)\\
<Plug>(easymotion-prefix)F  <Plug>(easymotion-F)\\
<Plug>(easymotion-prefix)f  <Plug>(easymotion-f)\\


\vspace{16pt}

\section{ALL Tmux Copy-Mode-Vi (1/2)}


\vspace{10pt}
{\large\bfseries\color{neoncyan} ALL TMUX COPY-MODE-VI (1/2)}\vspace{6pt}

bind-key  -T copy-mode-vi Enter  send-keys -X copy-pipe-and-cancel 'reatta\\
bind-key  -T copy-mode-vi Escape  send-keys -X clear-selection\\
bind-key  -T copy-mode-vi Space  send-keys -X begin-selection\\
bind-key  -T copy-mode-vi \textbackslash{}\#  send-keys -FX search-backward -- '\#\{copy\_cu\\
bind-key  -T copy-mode-vi \textbackslash{}  send-keys -X end-of-line\\
bind-key  -T copy-mode-vi \textbackslash{}\%  send-keys -X next-matching-bracket\\
bind-key  -T copy-mode-vi *  send-keys -FX search-forward -- '\#\{copy\_curso\\
bind-key  -T copy-mode-vi ,  send-keys -X jump-reverse\\
bind-key  -T copy-mode-vi /  command-prompt -i -p '(search down)' 'send -X\\
bind-key  -T copy-mode-vi 0  send-keys -X start-of-line\\
bind-key  -T copy-mode-vi 1  command-prompt -N -I 1 -p (repeat) \{ send-key\\
bind-key  -T copy-mode-vi 2  command-prompt -N -I 2 -p (repeat) \{ send-key\\
bind-key  -T copy-mode-vi 3  command-prompt -N -I 3 -p (repeat) \{ send-key\\
bind-key  -T copy-mode-vi 4  command-prompt -N -I 4 -p (repeat) \{ send-key\\
bind-key  -T copy-mode-vi 5  command-prompt -N -I 5 -p (repeat) \{ send-key\\
bind-key  -T copy-mode-vi 6  command-prompt -N -I 6 -p (repeat) \{ send-key\\
bind-key  -T copy-mode-vi 7  command-prompt -N -I 7 -p (repeat) \{ send-key\\
bind-key  -T copy-mode-vi 8  command-prompt -N -I 8 -p (repeat) \{ send-key\\
bind-key  -T copy-mode-vi 9  command-prompt -N -I 9 -p (repeat) \{ send-key\\
bind-key  -T copy-mode-vi :  command-prompt -p '(goto line)' \{ send-keys -\\
bind-key  -T copy-mode-vi \textbackslash{};  send-keys -X jump-again\\
bind-key  -T copy-mode-vi ?  command-prompt -i -p '(search up)' 'send -X s\\
bind-key  -T copy-mode-vi A  send-keys -X append-selection-and-cancel\\
bind-key  -T copy-mode-vi B  send-keys -X previous-space\\
bind-key  -T copy-mode-vi D  send-keys -X copy-pipe-end-of-line-and-cancel\\
bind-key  -T copy-mode-vi E  send-keys -X next-space-end\\
bind-key  -T copy-mode-vi F  command-prompt -1 -p '(jump backward)' \{ send\\
bind-key  -T copy-mode-vi G  send-keys -X history-bottom\\
bind-key  -T copy-mode-vi H  send-keys -X top-line\\
bind-key  -T copy-mode-vi J  send-keys -X scroll-down\\
bind-key  -T copy-mode-vi K  send-keys -X scroll-up\\
bind-key  -T copy-mode-vi L  send-keys -X bottom-line\\
bind-key  -T copy-mode-vi M  send-keys -X middle-line\\
bind-key  -T copy-mode-vi N  send-keys -X search-reverse\\
bind-key  -T copy-mode-vi P  send-keys -X toggle-position\\
bind-key  -T copy-mode-vi T  command-prompt -1 -p '(jump to backward)' \{ s\\
bind-key  -T copy-mode-vi V  send-keys -X select-line\\
bind-key  -T copy-mode-vi W  send-keys -X next-space\\
bind-key  -T copy-mode-vi X  send-keys -X set-mark\\
bind-key  -T copy-mode-vi \textasciicircum{}  send-keys -X back-to-indentation\\
bind-key  -T copy-mode-vi b  send-keys -X previous-word\\
bind-key  -T copy-mode-vi e  send-keys -X next-word-end\\
bind-key  -T copy-mode-vi f  command-prompt -1 -p '(jump forward)' \{ send-\\
bind-key  -T copy-mode-vi g  send-keys -X history-top\\
bind-key  -T copy-mode-vi h  send-keys -X cursor-left\\
bind-key  -T copy-mode-vi j  send-keys -X cursor-down\\
bind-key  -T copy-mode-vi k  send-keys -X cursor-up\\
bind-key  -T copy-mode-vi l  send-keys -X cursor-right\\


\vspace{16pt}

\section{ALL Tmux Copy-Mode-Vi (2/2)}


\vspace{10pt}
{\large\bfseries\color{neoncyan} ALL TMUX COPY-MODE-VI (2/2)}\vspace{6pt}

bind-key  -T copy-mode-vi n  send-keys -X search-again\\
bind-key  -T copy-mode-vi o  send-keys -X other-end\\
bind-key  -T copy-mode-vi q  send-keys -X cancel\\
bind-key  -T copy-mode-vi r  send-keys -X refresh-from-pane\\
bind-key  -T copy-mode-vi s  send-keys -X copy-pipe-no-clear 'reattach-to-\\
bind-key  -T copy-mode-vi t  command-prompt -1 -p '(jump to forward)' \{ se\\
bind-key  -T copy-mode-vi v  send-keys -X begin-selection\\
bind-key  -T copy-mode-vi w  send-keys -X next-word\\
bind-key  -T copy-mode-vi x  send-keys -X copy-pipe-no-clear 'reattach-to-\\
bind-key  -T copy-mode-vi y  send-keys -X copy-pipe-no-clear 'reattach-to-\\
bind-key  -T copy-mode-vi z  send-keys -X copy-pipe-no-clear 'reattach-to-\\
bind-key  -T copy-mode-vi \textbackslash{}\{  send-keys -X previous-paragraph\\
bind-key  -T copy-mode-vi \textbackslash{}\}  send-keys -X next-paragraph\\
bind-key  -T copy-mode-vi WheelUpPane  select-pane \textbackslash{}; send-keys -X -N 5 s\\
bind-key  -T copy-mode-vi WheelDownPane  select-pane \textbackslash{}; send-keys -X -N 5\\
bind-key  -T copy-mode-vi DoubleClick1Pane  select-pane \textbackslash{}; send-keys -X s\\
bind-key  -T copy-mode-vi TripleClick1Pane  select-pane \textbackslash{}; send-keys -X s\\
bind-key  -T copy-mode-vi BSpace  send-keys -X cursor-left\\
bind-key  -T copy-mode-vi Home  send-keys -X start-of-line\\
bind-key  -T copy-mode-vi End  send-keys -X end-of-line\\
bind-key  -T copy-mode-vi NPage  send-keys -X page-down\\
bind-key  -T copy-mode-vi PPage  send-keys -X page-up\\
bind-key  -T copy-mode-vi Up  send-keys -X cursor-up\\
bind-key  -T copy-mode-vi Down  send-keys -X cursor-down\\
bind-key  -T copy-mode-vi Left  send-keys -X cursor-left\\
bind-key  -T copy-mode-vi Right  send-keys -X cursor-right\\
bind-key  -T copy-mode-vi M-x  send-keys -X jump-to-mark\\
bind-key  -T copy-mode-vi C-b  send-keys -X page-up\\
bind-key  -T copy-mode-vi C-c  send-keys -X cancel\\
bind-key  -T copy-mode-vi C-d  send-keys -X halfpage-down\\
bind-key  -T copy-mode-vi C-e  send-keys -X scroll-down\\
bind-key  -T copy-mode-vi C-f  send-keys -X page-down\\
bind-key  -T copy-mode-vi C-h  send-keys -X -N 4 cursor-left\\
bind-key  -T copy-mode-vi C-j  send-keys -X -N 4 cursor-down\\
bind-key  -T copy-mode-vi C-k  send-keys -X -N 4 cursor-up\\
bind-key  -T copy-mode-vi C-l  send-keys -X -N 4 cursor-right\\
bind-key  -T copy-mode-vi C-n  send-keys -X page-down\\
bind-key  -T copy-mode-vi C-p  send-keys -X page-up\\
bind-key  -T copy-mode-vi C-u  send-keys -X halfpage-up\\
bind-key  -T copy-mode-vi C-v  send-keys -X rectangle-toggle\\
bind-key  -T copy-mode-vi C-y  send-keys -X scroll-up\\
bind-key  -T copy-mode-vi C-Up  send-keys -X scroll-up\\
bind-key  -T copy-mode-vi C-Down  send-keys -X scroll-down\\


\vspace{16pt}

\section{ALL Zsh Insert Mode (1/3)}


\vspace{10pt}
{\large\bfseries\color{neoncyan} ALL ZSH INSERT MODE (1/3)}\vspace{6pt}

bindkey -M viins '\textasciicircum{}@' zpwrExpandTerminateSpace\\
bindkey -M viins '\textasciicircum{}A' beginning-of-line\\
bindkey -M viins '\textasciicircum{}B' zpwrClipboard\\
bindkey -M viins '\textasciicircum{}C' self-insert\\
bindkey -M viins '\textasciicircum{}D' list-choices\\
bindkey -M viins '\textasciicircum{}E' end-of-line\\
bindkey -M viins '\textasciicircum{}F\textasciicircum{}D' zpwrIntoFzf\\
bindkey -M viins '\textasciicircum{}F\textasciicircum{}F' zpwrCdVimFzfFilesearchWidgetAccept\\
bindkey -M viins '\textasciicircum{}F\textasciicircum{}G' zpwrGoogle\\
bindkey -M viins '\textasciicircum{}F\textasciicircum{}H' zpwrLsoffzf\\
bindkey -M viins '\textasciicircum{}F\textasciicircum{}I' zpwrIntoFzfAg\\
bindkey -M viins '\textasciicircum{}F\textasciicircum{}J' zpwrVerbsWidgetAccept\\
bindkey -M viins '\textasciicircum{}F\textasciicircum{}K' zpwrAlternateQuotes\\
bindkey -M viins '\textasciicircum{}F\textasciicircum{}L' list-choices\\
bindkey -M viins '\textasciicircum{}F\textasciicircum{}M' zzcomplete\\
bindkey -M viins '\textasciicircum{}F\textasciicircum{}N' zpwrVerbsWidget\\
bindkey -M viins '\textasciicircum{}F\textasciicircum{}O' zpwrOpen\\
bindkey -M viins '\textasciicircum{}F\textasciicircum{}P' basicSedSub\\
bindkey -M viins '\textasciicircum{}F\textasciicircum{}R' zpwrAsVar\\
bindkey -M viins '\textasciicircum{}F\textasciicircum{}S' zsh-gacp-CheckDiff\\
bindkey -M viins '\textasciicircum{}F\textasciicircum{}V' edit-command-line\\
bindkey -M viins '\textasciicircum{}Fj' zpwrVerbsWidgetAccept\\
bindkey -M viins '\textasciicircum{}Fn' zpwrVerbsWidget\\
bindkey -M viins '\textasciicircum{}G' what-cursor-position\\
bindkey -M viins '\textasciicircum{}H' autopair-delete\\
bindkey -M viins '\textasciicircum{}I' fzf-completion\\
bindkey -M viins '\textasciicircum{}J' accept-line\\
bindkey -M viins '\textasciicircum{}K' zsh-learn-Learn\\
bindkey -M viins '\textasciicircum{}L' clear-screen\\
bindkey -M viins '\textasciicircum{}M' zpwrMagicEnter\\
bindkey -M viins '\textasciicircum{}N' down-history\\
bindkey -M viins '\textasciicircum{}O' zpwrEOLorNextTabStop\\
bindkey -M viins '\textasciicircum{}O\textasciicircum{}I' zc-logo\\
bindkey -M viins '\textasciicircum{}O\textasciicircum{}Z' zui-demo-various\\
bindkey -M viins '\textasciicircum{}P' up-history\\
bindkey -M viins '\textasciicircum{}Q' zpwrLastWordDouble\\
bindkey -M viins '\textasciicircum{}R' redo\\
bindkey -M viins '\textasciicircum{}S' zsh-gacp-NoCheck\\
bindkey -M viins '\textasciicircum{}T' fzf-file-widget\\
bindkey -M viins '\textasciicircum{}U' zpwrClearLine\\
bindkey -M viins '\textasciicircum{}V\textasciicircum{}@' zpwrVimFzf\\
bindkey -M viins '\textasciicircum{}V\textasciicircum{}F' zpwrFasdFZF\\
bindkey -M viins '\textasciicircum{}V\textasciicircum{}G' fzf-cd-widget\\
bindkey -M viins '\textasciicircum{}V\textasciicircum{}K' zpwrEmacsFzf\\
bindkey -M viins '\textasciicircum{}V\textasciicircum{}L' zpwrZFZF\\
bindkey -M viins '\textasciicircum{}V\textasciicircum{}N' zpwrVimFzfSudo\\
bindkey -M viins '\textasciicircum{}V\textasciicircum{}O' fzf-tab-complete\\
bindkey -M viins '\textasciicircum{}V\textasciicircum{}P' sudo-command-line\\


\vspace{16pt}

\section{ALL Zsh Insert Mode (2/3)}


\vspace{10pt}
{\large\bfseries\color{neoncyan} ALL ZSH INSERT MODE (2/3)}\vspace{6pt}

bindkey -M viins '\textasciicircum{}V\textasciicircum{}R' history-search-multi-word\\
bindkey -M viins '\textasciicircum{}V\textasciicircum{}S' zpwrZFZF\\
bindkey -M viins '\textasciicircum{}V\textasciicircum{}V' zpwrVimAllWidgetAccept\\
bindkey -M viins '\textasciicircum{}V\textasciicircum{}Z' fzf-history-widget\\
bindkey -M viins '\textasciicircum{}V ' zpwrVimFzf\\
bindkey -M viins '\textasciicircum{}V,' zpwrFzfEnv\\
bindkey -M viins '\textasciicircum{}V.' zpwrFzfAllKeybind\\
bindkey -M viins '\textasciicircum{}V/.' zpwrLocateFzfEdit\\

{\small\itshape\color{dimtext} bindkey -M viins '\textasciicircum{}V//' zpwrLocateFzf}

bindkey -M viins '\textasciicircum{}V;' zpwrFzfSurround\\
bindkey -M viins '\textasciicircum{}Vc' zpwrFzfCommits\\
bindkey -M viins '\textasciicircum{}Vk' zpwrFzfVimKeybind\\
bindkey -M viins '\textasciicircum{}Vl' zpwrZFZF\\
bindkey -M viins '\textasciicircum{}Vm' zpwrVimAllWidget\\
bindkey -M viins '\textasciicircum{}Vv' zpwrVimAllWidgetAccept\\
bindkey -M viins '\textasciicircum{}W' autopair-delete-word\\
bindkey -M viins '\textasciicircum{}X\textasciicircum{}R' \_read\_comp\\
bindkey -M viins '\textasciicircum{}X\textasciicircum{}X' zpwrBufferXtrace\\
bindkey -M viins '\textasciicircum{}X?' \_complete\_debug\\
bindkey -M viins '\textasciicircum{}XC' \_correct\_filename\\
bindkey -M viins '\textasciicircum{}Xa' \_expand\_alias\\
bindkey -M viins '\textasciicircum{}Xc' \_correct\_word\\
bindkey -M viins '\textasciicircum{}Xd' \_list\_expansions\\
bindkey -M viins '\textasciicircum{}Xe' \_expand\_word\\
bindkey -M viins '\textasciicircum{}Xh' \_complete\_help\\
bindkey -M viins '\textasciicircum{}Xm' \_most\_recent\_file\\
bindkey -M viins '\textasciicircum{}Xn' \_next\_tags\\
bindkey -M viins '\textasciicircum{}Xt' \_complete\_tag\\
bindkey -M viins '\textasciicircum{}X\textasciitilde{}' \_bash\_list-choices\\
bindkey -M viins '\textasciicircum{}Y' zpwrChangeQuotes\\
bindkey -M viins '\textasciicircum{}Z' undo\\
bindkey -M viins '\textasciicircum{}[' vi-cmd-mode\\
bindkey -M viins '\textasciicircum{}[\textasciicircum{}D' capitalize-word\\
bindkey -M viins '\textasciicircum{}[\textasciicircum{}E' zpwrExpandGlobalAliases\\
bindkey -M viins '\textasciicircum{}[\textasciicircum{}F' sub\\
bindkey -M viins '\textasciicircum{}[\textasciicircum{}L' down-case-word\\
bindkey -M viins '\textasciicircum{}[\textasciicircum{}M' self-insert-unmeta\\
bindkey -M viins '\textasciicircum{}[\textasciicircum{}O' zpwrRunner\\
bindkey -M viins '\textasciicircum{}[\textasciicircum{}P' zpwrUp8widget\\
bindkey -M viins '\textasciicircum{}[\textasciicircum{}T' transpose-words\\
bindkey -M viins '\textasciicircum{}[\textasciicircum{}U' up-case-word\\
bindkey -M viins '\textasciicircum{}[ ' zpwrSshRegain\\
bindkey -M viins '\textasciicircum{}[,' \_history-complete-newer\\
bindkey -M viins '\textasciicircum{}[/' \_history-complete-older\\
bindkey -M viins '\textasciicircum{}[OA' history-substring-search-up\\
bindkey -M viins '\textasciicircum{}[OB' history-substring-search-down\\
bindkey -M viins '\textasciicircum{}[OC' vi-forward-char\\
bindkey -M viins '\textasciicircum{}[OD' vi-backward-char\\


\vspace{16pt}

\section{ALL Zsh Insert Mode (3/3)}


\vspace{10pt}
{\large\bfseries\color{neoncyan} ALL ZSH INSERT MODE (3/3)}\vspace{6pt}

bindkey -M viins '\textasciicircum{}[OP' zpwrUp8widget\\
bindkey -M viins '\textasciicircum{}[OQ' sub\\
bindkey -M viins '\textasciicircum{}[OQ\textasciicircum{}[OQ' npm\_toggle\_install\_uninstall\\
bindkey -M viins '\textasciicircum{}[OR' zpwrGetrcWidget\\
bindkey -M viins '\textasciicircum{}[Q' zi-browse-symbol\\
bindkey -M viins '\textasciicircum{}[[1;2D' sub\\
bindkey -M viins '\textasciicircum{}[[1;5A' zsh-gacp-CheckDiff\\
bindkey -M viins '\textasciicircum{}[[1;5B' zpwrUpdater\\
bindkey -M viins '\textasciicircum{}[[1;5C' zpwrTutsUpdate\\
bindkey -M viins '\textasciicircum{}[[1;5D' dbz\\
bindkey -M viins '\textasciicircum{}[[200\textasciitilde{}' bracketed-paste\\
bindkey -M viins '\textasciicircum{}[[5\textasciitilde{}' zpwrClipboard\\
bindkey -M viins '\textasciicircum{}[[A' history-substring-search-up\\
bindkey -M viins '\textasciicircum{}[[B' history-substring-search-down\\
bindkey -M viins '\textasciicircum{}[[C' vi-forward-char\\
bindkey -M viins '\textasciicircum{}[[D' vi-backward-char\\
bindkey -M viins '\textasciicircum{}[[Z' zpwrClipboard\\
bindkey -M viins '\textasciicircum{}[c' fzf-cd-widget\\
bindkey -M viins '\textasciicircum{}[man' man-command-line\\
bindkey -M viins '\textasciicircum{}[\textasciitilde{}' \_bash\_complete-word\\
bindkey -R -M viins '\textasciicircum{}\textbackslash{}\textbackslash{}\textbackslash{}\textbackslash{}'-'\textasciicircum{}\textbackslash{}\textasciicircum{}' self-insert\\
bindkey -M viins '\textasciicircum{}\_' zpwrAsVar\\
bindkey -M viins ' ' zpwrExpandSupernaturalSpace\\
bindkey -M viins '.' self-insert\\
bindkey -M viins '\textbackslash{}'' zpwrInterceptSurround\\
bindkey -R -M viins '\#'-'\&' self-insert\\
bindkey -R -M viins '''-'(' zpwrInterceptSurround\\
bindkey -M viins ')' autopair-close\\
bindkey -R -M viins '*'-'-' self-insert\\
bindkey -M viins '.' zpwrRationalizeDot\\
bindkey -R -M viins '/'-'Z' self-insert\\
bindkey -M viins '[' zpwrInterceptSurround\\
bindkey -M viins '\textbackslash{}\textbackslash{}\textbackslash{}\textbackslash{}' self-insert\\
bindkey -M viins ']' autopair-close\\
bindkey -R -M viins '\textbackslash{}\textasciicircum{}'-'\_' self-insert\\
bindkey -M viins '\textbackslash{}'' zpwrInterceptSurround\\
bindkey -R -M viins 'a'-'f' self-insert\\
bindkey -M viins 'fj' vi-cmd-mode\\
bindkey -R -M viins 'g'-'j' self-insert\\
bindkey -M viins 'jf' vi-cmd-mode\\
bindkey -R -M viins 'k'-'z' self-insert\\
bindkey -M viins '\{' zpwrInterceptSurround\\
bindkey -M viins '|' self-insert\\
bindkey -M viins '\}' autopair-close\\
bindkey -M viins '\textasciitilde{}' self-insert\\
bindkey -M viins '\textasciicircum{}?' zpwrInterceptDelete\\
bindkey -R -M viins '\textbackslash{}M-\textasciicircum{}@'-'\textbackslash{}M-\textasciicircum{}?' self-insert\\


\vspace{16pt}

\section{ALL Zsh Visual Mode (1/2)}


\vspace{10pt}
{\large\bfseries\color{neoncyan} ALL ZSH VISUAL MODE (1/2)}\vspace{6pt}

bindkey -a 'V' visual-line-mode\\
bindkey -a 'v' visual-mode\\
bindkey -M visual '\textasciicircum{}[' deactivate-region\\
bindkey -M visual '\textasciicircum{}[OA' up-line\\
bindkey -M visual '\textasciicircum{}[OB' down-line\\
bindkey -M visual '\textasciicircum{}[[A' up-line\\
bindkey -M visual '\textasciicircum{}[[B' down-line\\
bindkey -M visual '-' vi-up-line-or-history\\
bindkey -M visual 'U' vi-up-case\\
bindkey -M visual 'a\textbackslash{}'' select-quoted\\
bindkey -M visual 'a'' select-quoted\\
bindkey -M visual 'a(' select-bracketed\\
bindkey -M visual 'a)' select-bracketed\\
bindkey -M visual 'a+' select-quoted\\
bindkey -M visual 'a,' select-quoted\\
bindkey -M visual 'a-' select-quoted\\
bindkey -M visual 'a.' select-quoted\\
bindkey -M visual 'a/' select-quoted\\
bindkey -M visual 'a:' select-quoted\\
bindkey -M visual 'a;' select-quoted\\
bindkey -M visual 'a<' select-bracketed\\
bindkey -M visual 'a=' select-quoted\\
bindkey -M visual 'a>' select-bracketed\\
bindkey -M visual 'a@' select-quoted\\
bindkey -M visual 'aB' select-bracketed\\
bindkey -M visual 'aW' select-a-blank-word\\
bindkey -M visual 'a[' select-bracketed\\
bindkey -M visual 'a\textbackslash{}\textbackslash{}\textbackslash{}\textbackslash{}' select-quoted\\
bindkey -M visual 'a]' select-bracketed\\
bindkey -M visual 'a\textbackslash{}'' select-quoted\\
bindkey -M visual 'aa' select-a-shell-word\\
bindkey -M visual 'ab' select-bracketed\\
bindkey -M visual 'aw' select-a-word\\
bindkey -M visual 'a\{' select-bracketed\\
bindkey -M visual 'a|' select-quoted\\
bindkey -M visual 'a\}' select-bracketed\\
bindkey -M visual 'i\textbackslash{}'' select-quoted\\
bindkey -M visual 'i'' select-quoted\\
bindkey -M visual 'i(' select-bracketed\\
bindkey -M visual 'i)' select-bracketed\\
bindkey -M visual 'i+' select-quoted\\
bindkey -M visual 'i,' select-quoted\\
bindkey -M visual 'i-' select-quoted\\
bindkey -M visual 'i.' select-quoted\\
bindkey -M visual 'i/' select-quoted\\
bindkey -M visual 'i:' select-quoted\\
bindkey -M visual 'i;' select-quoted\\
bindkey -M visual 'i<' select-bracketed\\


\vspace{16pt}

\section{ALL Zsh Visual Mode (2/2)}


\vspace{10pt}
{\large\bfseries\color{neoncyan} ALL ZSH VISUAL MODE (2/2)}\vspace{6pt}

bindkey -M visual 'i=' select-quoted\\
bindkey -M visual 'i>' select-bracketed\\
bindkey -M visual 'i@' select-quoted\\
bindkey -M visual 'iB' select-bracketed\\
bindkey -M visual 'iW' select-in-blank-word\\
bindkey -M visual 'i[' select-bracketed\\
bindkey -M visual 'i\textbackslash{}\textbackslash{}\textbackslash{}\textbackslash{}' select-quoted\\
bindkey -M visual 'i]' select-bracketed\\
bindkey -M visual 'i\textbackslash{}'' select-quoted\\
bindkey -M visual 'ia' select-in-shell-word\\
bindkey -M visual 'ib' select-bracketed\\
bindkey -M visual 'iw' select-in-word\\
bindkey -M visual 'i\{' select-bracketed\\
bindkey -M visual 'i|' select-quoted\\
bindkey -M visual 'i\}' select-bracketed\\
bindkey -M visual 'j' down-line\\
bindkey -M visual 'k' up-line\\
bindkey -M visual 'o' exchange-point-and-mark\\
bindkey -M visual 'p' put-replace-selection\\
bindkey -M visual 'u' vi-down-case\\
bindkey -M visual 'x' vi-delete\\
bindkey -M visual '\textasciitilde{}' vi-oper-swap-case\\


\vspace{16pt}

\section{ALL Vim Normal Mode (1/3)}


\vspace{10pt}
{\large\bfseries\color{neoncyan} ALL VIM NORMAL MODE (1/3)}\vspace{6pt}

n  <Esc><C-C>  * wvU\\
n  <Esc><C-T>  * :call TransposeWords()<CR>\\
n  <Space>mt  * :MinimapToggle<CR>\\
n  <Space>mc  * :MinimapClose<CR>\\
n  <Space>mu  * :MinimapUpdate<CR>\\
n  <Space>mm  * :Minimap<CR>\\
n  <Space><Tab>  :call AltOrNextBuffer()<CR>\\
n  <Space>/  * :LocateAll<CR>\\
n  <Space>ma  * :Marks.<CR>\\
n  <Space>]  * <C-W>\}<CR>\\
n  <Space>ta  * :Tags.<CR>\\
n  <Space>h:  * :History:.<CR>\\
n  <Space>h/  * :History/.<CR>\\
n  <Space>hh  * :History.<CR>\\
n  <Space>rg  * :Rg<CR>\\
n  <Space>oa  * :ALEToggle<CR>\\
n  <Space>m  * :Map.<CR>\\
n  <Space>j  * :Lines.<CR>\\
n  <Space>ke  * :FZFKeys<CR>\\
n  <Space>,  * :FZFMaps<CR>\\
n  <Space>aa  * :Agg<CR>\\
n  <Space>i  * :Imap<CR>\\
n  <Space>f  * :Files.<CR>\\
n  <Space>env  * :FZFEnv<CR>\\
n  <Space>;  * :Commands.<CR>\\
n  <Space>.  * :Colors<CR>\\
n  <Space>b  * :Buffers.<CR>\\
n  <Space>ag  * :Ag<CR>\\
n  <Space>vj  * :w.<CR>:call TmuxRepeat('file')<CR>\\
n  <Space>ev  :call ExtractVariable()<CR>\\
n  <Space>os  * :call CopyCWordClip()<CR> :call system('bash \textbf{ZPWR\_TMUX}/goo\\
n  <Space>ox  * :call CopyCWordClip()<CR> :call system('bash \textbf{ZPWR\_TMUX}/goo\\
n  <Space>oc  * :call CopyCWordClip()<CR>\\
n  <Space>x  * :normal mzg\&'zzz<CR>\\
n  <Space>p  * :bprev<CR>\\
n  <Space>n  * :bnext<CR>\\
n  <Space>ap  * :prev<CR>\\
n  <Space>an  * :next<CR>\\
n  <Space>lo  * :lopen<CR>\\
n  <Space>lc  * :lclose<CR>\\
n  <Space>lp  * :lprev<CR>\\
n  <Space>ln  * :lnext<CR>\\
n  <Space>[  * :call Quoter('bracket')<CR>\\
n  <Space>'  * :call Quoter('back')<CR>\\
n  <Space>'  * :call Quoter('single')<CR>\\
n  <Space>'  * :call Quoter('double')<CR>\\
n  <Space>z  * :Commits.<CR>\\
n  <Space>t  * :tabnew<CR>\\


\vspace{16pt}

\section{ALL Vim Normal Mode (2/3)}


\vspace{10pt}
{\large\bfseries\color{neoncyan} ALL VIM NORMAL MODE (2/3)}\vspace{6pt}

n  <Space>s  * :split<CR>\\
n  <Space>v  * :vsplit<CR>\\
n  <Space>w  * :w.<CR>\\
n  <Space>e  * :q.<CR>\\
n  <Space>wq  * :wq.<CR>\\
n  <Space>q  * :qa.<CR>\\
n  <Space>r  * :\%s@\textbackslash{}C\textbackslash{}<<C-R><C-W>\textbackslash{}>@<C-R><C-W>@g<Left><Left>\\
n  <Space>g  * :\%s@\textbackslash{}C\textbackslash{}<<C-R><C-W>\textbackslash{}>@@g<Left><Left>\\
n  <Space>=  * 4+\\
n  <Space>-  * 4-\\
n  \&\&  * :normal mzg\&'zzz<CR>\\
n  \&  * :\&\&<CR>\\
n  <s<Esc>  \& <Nop>\\
n  =s<Esc>  \& <Nop>\\
n  >s<Esc>  \& <Nop>\\
n  N  * :call GoToLastSearch('?')<CR>\\
n  Y  * yy'>\\
n  [o<Esc>  \& <Nop>\\
n  []  k][\%?\}<CR>]\}]]\}]]\\
n  [[  ?\{<CR>w99[\{\\
n  \textbackslash{}K  * :OnlineThesaurusCurrentWord<CR>\\
n  ]o<Esc>  \& <Nop>\\
n  ]]  j0[[\%/\{<CR>\\
n  ][  /\}<CR>b99]\}\\
n  n  * :call GoToLastSearch('/')<CR>\\
n  yo<Esc>  \& <Nop>\\
n  y<C-G>  \& :<C-U>call setreg(v:register, fugitive\#Object(@\%))<CR>\\
n  z/  * :if AutoHighlightToggle()|set hls|endif<CR>\\
n  <SNR>198\_:  * :<C-U><C-R>=v:count ? v:count : ''<CR>\\
n  <C-G>  * :call multiple\_cursors\#new('n', 1)<CR>\\
n  <F11>  * :call conque\_term\#exec\_file()<CR>\\
n  <C-W>\textbackslash{}  * :vsplit<CR>\\
n  <C-W>-  * :split<CR>\\
n  <C-D>,  * :FZFEnv<CR>\\
n  <C-D>.  * :FZFKeys<CR>\\
n  <C-D>/  * :LocateAll<CR>\\
n  <C-D>z  * :TlistAddFiles *<CR>:TlistToggle<CR>\\
n  <C-D>y  * :update<CR>:SyntasticCheck<CR>\\
n  <C-D>x  * :Marks<CR>\\
n  <C-D>w  * :update<CR>\\
n  <C-D>v  * :w.<CR>:call TmuxRepeatGeneric()<CR>\\
n  <C-D>u  * :History:<CR>\\
n  <C-D>]  * <C-W>\}<CR>\\
n  <C-D>t  * :Tags<CR>\\
n  <C-D>s  * :History/<CR>\\
n  <C-D>rr  * :Rg<CR>\\
n  <C-D>rq  * :silent .open -t \%:p:h<CR>:redraw.<CR>\\
n  <C-D>q  * :SaveSession.<CR><Tab>\\


\vspace{16pt}

\section{ALL Vim Normal Mode (3/3)}


\vspace{10pt}
{\large\bfseries\color{neoncyan} ALL VIM NORMAL MODE (3/3)}\vspace{6pt}

n  <C-D>p  * :call PasteClip()<CR>\\
n  <C-D>r  * :call GetRef()<CR>\\
n  <C-D>o  * :ALEToggle<CR>\\
n  <C-D>n  * :Snippets<CR>\\
n  <C-D>m  * :Map.<CR>\\
n  <C-D>l  * :Lines.<CR>\\
n  <C-D>k  * :ALEFix<CR>\\
n  <C-D>j  * :Agg<CR>\\
n  <C-D>i  * :Imap<CR>\\
n  <C-D>h  * :HistoryFiles<CR>\\
n  <C-D>g  * :Commits.<CR>\\
n  <C-D>f  * :Files.<CR>\\
n  <C-D>e  * :ALEInfo<CR>\\
n  <C-D>d  * :Commands<CR>\\
n  <C-D>c  * :Colors<CR>\\
n  <C-D>b  * :Buffers<CR>\\
n  <C-D>a  * :Ag<CR>\\
n  <C-D><C-D>  * :GitGutterUndoHunk<CR>\\
n  <C-V>  * :w.<CR>:call TmuxRepeat('file')<CR>\\
n  <C-Bslash>  * +\\
n  <F7>  * :TTags<CR>\\
n  <F6>  * :SyntasticToggleMode<CR>\\
n  <F5>  * :LOTRToggle<CR>\\
n  <F4>  * :MinimapToggle<CR>\\
n  <F3>  * :TlistAddFiles *<CR>:TlistToggle<CR>\\
n  <F2>  * :UndotreeToggle<CR>\\
n  <F1>  * :NERDTreeToggle<CR>\\
n  <F8>  * :\%s@@@g<Left><Left><Left>\\
n  <C-Up>  * :<C-U>call GoToNextMarker('\{\{\{',1)<CR>\\
n  <C-Down>  * :<C-U>call GoToNextMarker('\{\{\{',0)<CR>\\
n  <End>  * G\\
n  <Home>  * gg\\
n  <C-F>  * :q.<CR>\\
n  <C-C>  * :wq.<CR>:qa.<CR>\\
n  <C-T>  * xp\\
n  <C-L>  * 4l\\
n  <C-H>  * 4h\\
n  <C-K>  * 4k\\
n  <C-J>  * 4j\\
n  <D-v>  '*P\\


\vspace{16pt}

\section{ALL Vim Insert Mode (1/2)}


\vspace{10pt}
{\large\bfseries\color{neoncyan} ALL VIM INSERT MODE (1/2)}\vspace{6pt}

i  <M-d>  * <C-O>dw\\
i  <F31>  * <C-O>dw\\
i  <C-F>  * col('.')>strlen(getline('.'))?'\textbackslash{}<C-F>':'\textbackslash{}<Right>'\\
i  <C-E>  * col('.')>strlen(getline('.'))||pumvisible()?'\textbackslash{}<C-E>':'\textbackslash{}<End>\\
i  <C-D>  * col('.')>strlen(getline('.'))?'\textbackslash{}<C-D>':'\textbackslash{}<Del>'\\
i  <C-X><C-A>  * <C-A>\\
i  <C-A>  * <C-O>\textasciicircum{}\\
i  <C-X>  * <C-R>=<SNR>69\_ManualCompletionEnter()<CR>\\
i  <C-Tab>  * <C-R>=UltiSnips\#ListSnippets()<CR>\\
i  <C-X><C-L>  * fzf\#vim\#complete(fzf\#wrap(\{ 'prefix': '\textasciicircum{}.*', 'source': 'r\\
i  <C-X><C-K>  * fzf\#vim\#complete\#word(\{'left': '15\%'\})\\
i  <F11>  * <C-X><C-T>\\
i  <F10>  * <C-X><C-K>\\
i  <C-D>n  * <C-X><C-O>\\
i  <C-D>\textbackslash{}  * <C-X><C-L>\\
i  <C-D>,  * <C-O>:FZFEnv<CR>\\
i  <C-D>.  * <C-O>:FZFKeys<CR>\\
i  <C-D>/  * <C-O>:LocateAll<CR>\\
i  <C-D>z  * <Esc>:TlistAddFiles * <CR> :TlistToggle<CR>i\\
i  <C-D>y  * <Esc>:update<CR>:SyntasticCheck<CR>a\\
i  <C-D>x  * <C-O>:Marks<CR>\\
i  <C-D>w  * <C-O>:update<CR>\\
i  <C-D>v  * <Esc>:w.<CR>:call TmuxRepeatGeneric()<CR>a\\
i  <C-D>u  * <C-O>:History:<CR>\\
i  <C-D>]  * <C-W>\}<CR>\\
i  <C-D>t  * <C-O>:Tags<CR>\\
i  <C-D>s  * <C-O>:History/<CR>\\
i  <C-D>rr  * <Esc>:Rg<CR>\\
i  <C-D>rq  * <Esc>:silent .open -t \%:p:h<CR>:redraw.<CR>a\\
i  <C-D>q  * <C-O>:SaveSession.<CR><Tab>\\
i  <C-D>p  * <C-O>:call PasteClip()<CR>\\
i  <C-D>r  * <C-O>:call GetRef()<CR>\\
i  <C-D>o  * <C-O>:ALEToggle<CR>\\
i  <C-D>m  * <C-O>:Map<CR>\\
i  <C-D>l  * <C-O>:Lines<CR>\\
i  <C-D>k  * <C-O>:ALEFix<CR>\\
i  <C-D>j  * <C-O>:Agg<CR>\\
i  <C-D>i  * <C-O>:Imap<CR>\\
i  <C-D>h  * <C-O>:HistoryFiles<CR>\\
i  <C-D>g  * <C-O>:Commits.<CR>\\
i  <C-D>f  * <C-O>:Files<CR>\\
i  <C-D>e  * <C-O>:ALEInfo<CR>\\
i  <C-D>d  * <C-O>:Commands<CR>\\
i  <C-D>c  * <C-O>:Colors<CR>\\
i  <C-D>b  * <C-O>:Buffers<CR>\\
i  <C-D>a  * <C-O>:Ag<CR>\\
i  <C-D><C-D>  * <C-O>:GitGutterUndoHunk<CR>\\
i  <C-Bslash>  * <Esc>+\\


\vspace{16pt}

\section{ALL Vim Insert Mode (2/2)}


\vspace{10pt}
{\large\bfseries\color{neoncyan} ALL VIM INSERT MODE (2/2)}\vspace{6pt}

i  <F7>  * <Esc>:TTags<CR>\\
i  <F6>  * <Esc>:SyntasticToggleMode<CR>\\
i  <F5>  * <Esc>:LOTRToggle<CR>\\
i  <F4>  * <Esc>:MinimapToggle<CR>\\
i  <F3>  * <Esc>:TlistAddFiles *<CR>:TlistToggle<CR>\\
i  <F2>  * <Esc>:UndotreeToggle<CR>\\
i  <F1>  * <Esc>:NERDTreeToggle<CR>\\
i  <F8>  * <Esc>:\%s@@@g<Left><Left><Left>\\
i  <C-Up>  * <Esc>:<C-U>call GoToNextMarker('\{\{\{',1)<CR>i\\
i  <C-Down>  * <Esc>:<C-U>call GoToNextMarker('\{\{\{',0)<CR>i\\
i  <C-B><C-N>  * <Esc>\textasciicircum{}2xji\\
i  <C-B>  * getline('.')=\textasciitilde{}'\textasciicircum{}\textbackslash{}s*'\&\&col('.')>strlen(getline('.'))?'0\textbackslash{}<C-D>\\
i  <End>  * <Esc>Gi\\
i  <Home>  * <Esc>ggi\\
i  <C-C>  * <Esc>:wq.<CR>:qa.<CR>\\
i  <C-T>  * i<BS><C-O>:silent. undojoin | normal. xp<CR>\\
i  <C-D><C-T>  * <C-O>:call TransposeWords()<CR>\\
i  <C-L>  * <Esc>mbgg=G'bzza\\
i  <Tab>  * <C-R>=UltiSnips\#ExpandSnippetOrJump()<CR>\\
i  fj  <Esc>\\
i  jf  <Esc>\\


\vspace{16pt}

\section{ALL Vim Visual Mode (1/1)}


\vspace{10pt}
{\large\bfseries\color{neoncyan} ALL VIM VISUAL MODE (1/1)}\vspace{6pt}

v  <Space>vl  * <Esc>:call TmuxRepeat('repl')<CR>gv\\
v  <Space>vk  * <Esc>:call TmuxRepeat('visual')<CR>gv\\
v  <Space>vj  * <Esc>:call TmuxRepeat('file')<CR>gv\\
v  <Space>(  * :call InsertQuoteVisualMode('paren')<CR>\\
v  <Space>\{  * :call InsertQuoteVisualMode('curlybracket')<CR>\\
v  <Space>[  * :call InsertQuoteVisualMode('bracket')<CR>\\
v  <Space>'  * :call InsertQuoteVisualMode('back')<CR>\\
v  <Space>'  * :call InsertQuoteVisualMode('single')<CR<CR>\\
v  <Space>'  * :call InsertQuoteVisualMode('double')<CR>\\
v  <Space>r  * :'<,'>s@\textbackslash{}C\textbackslash{}<<C-R><C-W>\textbackslash{}>@<C-R><C-W>@g<Left><Left>\\
v  <Space>g  * :'<,'>\%s@\textbackslash{}C\textbackslash{}<<C-R><C-W>\textbackslash{}>@@g<Left><Left>\\
v  <Space>b  * :w .tmux set-buffer '(cat)'<CR><CR>\\
v  <Space>=  * 4+\\
v  <Space>-  * 4-\\
v  <  * <gv\\
v  >  * >gv\\
v  J  * :m '> + <CR> gv\\
v  K  * :m '< -- <CR> gv\\
v  Y  * y'>j\\
v  \textbackslash{}K  * y:Thesaurus <C-R>'<CR>\\
v  go  * :call CopyClip()<CR> :call system('bash \textbf{ZPWR\_TMUX}/google.sh open'\\
v  gs  * :call CopyClip()<CR> :call system('bash \textbf{ZPWR\_TMUX}/google.sh googl\\
v  <C-D>d  * :<C-C>:update<CR>\\
v  <C-D>y  * :<C-C>:update<CR>:SyntasticCheck<CR>\\
v  <C-D>,  * :call NERDComment('x','Toggle')<CR>'>\\
v  <RightMouse> * :call CopyClip()<CR>'>\\
v  <C-B>  * :call CopyClip()<CR>'>\\
v  <C-Up>  * :m '< -- <CR> gv\\
v  <C-Down>  * :m '> + <CR> gv\\
v  <C-Left>  * <gv\\
v  <C-Right>  * >gv\\
v  <C-F>  * :<C-C>:q.<CR>\\
v  <C-L>  * 4l\\
v  <C-K>  * 4k\\
v  <C-J>  * 4j\\
v  <D-x>  '*d\\
v  <D-c>  '*y\\
v  <D-v>  '-d'*P\\


\vspace{16pt}


\chapter{TMUX DEEP DIVE}


\begin{center}
\includegraphics[width=0.85\textwidth]{/Users/wizard/.zpwr/docs/screenshots/art_mainframe.png}
\end{center}

\section{Tmux Layout Save \& Restore}


\vspace{10pt}
{\large\bfseries\color{neoncyan} TMUX LAYOUT SAVE \& RESTORE}\vspace{6pt}

Your pane arrangement, preserved forever.


\vspace{12pt}
{\bfseries\color{neonyellow} THE PROBLEM}\vspace{4pt}

You spent 20 minutes splitting panes, resizing them, cd-ing into the right directories. Then tmux crashes. Or you reboot. Or you accidentally kill the session. All that layout work? Gone. Every time.


\vspace{12pt}
{\bfseries\color{neonyellow} THE SOLUTION}\vspace{4pt}

\begin{codebox}
zpwr tmuxsave dev             \# save layout as 'dev'\\
zpwr tmuxload dev             \# restore it later\\
\end{codebox}


\vspace{12pt}
{\bfseries\color{neonyellow} WHAT GETS SAVED}\vspace{4pt}

Windows ............... names and count Pane splits ........... exact geometry via select-layout Directories ........... per-pane working directory Running commands ...... full args from process tree Window options ........ synchronize-panes, auto-rename Session options ....... non-plugin tmux settings Active pane/window .... cursor position restored


\vspace{12pt}
{\bfseries\color{neonyellow} MANAGEMENT}\vspace{4pt}

\begin{codebox}
zpwr tmuxsave                 \# save with timestamp name\\
zpwr tmuxsave --list          \# list saved layouts\\
zpwr tmuxsave --delete dev    \# remove a layout\\
zpwr tmuxload                 \# fzf picker (no args)\\
zpwr tmuxload dev mywork      \# restore as session 'mywork'\\
\end{codebox}


\vspace{12pt}
{\bfseries\color{neonyellow} KEYBINDING}\vspace{4pt}

Prefix W                      \# save layout (tmux keybinding)


\vspace{12pt}
{\bfseries\color{neonyellow} ALIASES}\vspace{4pt}

\begin{codebox}
zpwr layoutsave  = zpwr tmuxsave\\
zpwr layoutload  = zpwr tmuxload\\
\end{codebox}

Layouts are stored as self-contained zsh scripts in \textbf{\$ZPWR\_LOCAL/layouts/.} You can edit them by hand, version control them, or share them with your team.

PROTIP: Save your layout before experimenting. If you break everything, one command brings it back.


\vspace{16pt}


\chapter{MONITORING \& AUTOMATION}


\begin{center}

\includegraphics[width=0.85\textwidth]{/Users/wizard/.zpwr/docs/screenshots/zpwr_watch_help.png}
\end{center}

\section{Live Dashboard \& Sparklines}


\begin{center}
\includegraphics[width=0.8\textwidth]{/Users/wizard/.zpwr/docs/screenshots/zpwr_bench_help.png}
\end{center}


\vspace{10pt}
{\large\bfseries\color{neoncyan} LIVE DASHBOARD \& SPARKLINES}\vspace{6pt}

Your shell's vital signs, in real time.


\vspace{12pt}
{\bfseries\color{neonyellow} zpwr top}\vspace{4pt}

A live-updating dashboard of your shell's resource usage. Runs in alternate screen buffer. Press q to quit.

\begin{codebox}
zpwr top       \# refresh every 2s\\
zpwr top 1     \# fast refresh\\
zpwr top 5     \# slow refresh\\
zpwr dashboard \# same verb\\
zpwr live      \# same verb\\
\end{codebox}


\vspace{12pt}
{\bfseries\color{neonyellow} WHAT IT SHOWS}\vspace{4pt}

Memory ............ RSS bar graph, virtual memory Shell objects ..... functions, aliases, completions, params Delta tracking .... changes between refreshes in yellow Loaded modules .... all zsh modules currently loaded Paths ............ fpath, path, named dirs, dir stack Context .......... git branch, tmux info, startup time Top 5 functions .. largest functions with size bars Top 5 paths ...... longest \textbf{PATH} entries


\vspace{12pt}
{\bfseries\color{neonyellow} SPARKLINES}\vspace{4pt}

Startup time history: last 40 shell starts Color-coded: green=fast yellow=mid red=slow Shows min/max/avg stats below the graph.

Commit velocity: commits per day over 30 days Visualizes your development pace as a mini bar chart.

Startup times are logged automatically on each shell init to \textbf{\$ZPWR\_LOCAL/startup\_history.log.}


\vspace{16pt}


\chapter{CREATIVE TOOLS}


\begin{center}
\includegraphics[width=0.85\textwidth]{/Users/wizard/.zpwr/docs/screenshots/figlet_zpwr.png}
\end{center}

\section{Matrix: The Full Experience}


\vspace{10pt}
{\large\bfseries\color{neoncyan} MATRIX: THE FULL EXPERIENCE}\vspace{6pt}

8 acts, 6 characters, and your ecosystem raining down


\vspace{12pt}
{\bfseries\color{neonyellow} THE STORY}\vspace{4pt}

\begin{codebox}
zpwr matrix               \# 4-act story + rain\\
zpwr matrix --extended     \# full 8-act saga\\
\end{codebox}


\vspace{12pt}
{\bfseries\color{neonyellow} THE CHARACTERS}\vspace{4pt}

\begin{codebox}
zpwr matrix --cast         \# show character bios\\
\end{codebox}

Null Pointer .... protagonist, zpwr operator Segfault ........ bash loyalist in recovery Chmod777 ........ reckless, sudo for everything

\begin{codebox}
DevNull ......... silent assassin\\
\end{codebox}

Chad Electron ... MegaCorp VP, 64GB for Slack The Kernel ...... mythical elder, speaks in dmesg


\vspace{12pt}
{\bfseries\color{neonyellow} \textbf{VISUAL} FX}\vspace{4pt}

\begin{codebox}
zpwr matrix --noir         \# monochrome green\\
zpwr matrix --vhs          \# VHS tracking distortion\\
zpwr matrix --crt          \# CRT scanlines\\
zpwr matrix --glitch       \# extra corruption\\
zpwr matrix --sound        \# terminal bells\\
zpwr matrix --no-fx        \# disable all FX\\
zpwr matrix -c 10          \# chaos level 1-10\\
\end{codebox}


\vspace{12pt}
{\bfseries\color{neonyellow} PLAYBACK CONTROL}\vspace{4pt}

\begin{codebox}
zpwr matrix -s 7           \# speed 1-9\\
zpwr matrix --turbo        \# skip all pauses\\
zpwr matrix -a 5           \# jump to act 5\\
zpwr matrix -r             \# skip story, rain only\\
\end{codebox}


\vspace{12pt}
{\bfseries\color{neonyellow} MULTIPLAYER}\vspace{4pt}

\begin{codebox}
zpwr matrix --battle       \# 2-player tmux typing battle\\
\end{codebox}

Player 1 (Null Pointer) vs Player 2 (Chad Electron). Type zpwr commands faster than your opponent. HP bars, taunts, and a winner-takes-all endgame. Requires tmux. Auto-splits into two panes.


\vspace{12pt}
{\bfseries\color{neonyellow} INFO}\vspace{4pt}

\begin{codebox}
zpwr matrix --stats        \# ecosystem stats\\
zpwr matrix --credits      \# end credits roll\\
\end{codebox}


\vspace{16pt}


\chapter{THE END}


\begin{center}
\includegraphics[width=0.85\textwidth]{/Users/wizard/.zpwr/docs/screenshots/art_hack_the_planet.png}
\end{center}

\section{The Epilogue: You Are the Wizard}


\vspace{10pt}
{\large\bfseries\color{neoncyan} THE EPILOGUE: YOU ARE THE WIZARD}\vspace{6pt}

The terminal bows to you.


\vspace{12pt}
{\bfseries\color{neonyellow} WHAT YOU'VE LEARNED}\vspace{4pt}

From keybindings to kernel. Every chapter, every verb. Autoloaded functions. Tmux pipelines. Fzf integration. Vim execution. Perl oneliners. Systemd services. 32-pane REPL layouts. Cross-platform installers. A 24-hour updater loop. A 60-file test suite. A colorization stack five tools deep. And a startup banner delivered by a rainbow pony.


\vspace{12pt}
{\bfseries\color{neonyellow} THE 5 COMMANDS TO MEMORIZE}\vspace{4pt}

zpwr verbs    list every zpwr subcommand zpwr doctor   run environment health check diagnostic zpwr bench    benchmark startup with percentiles zpwr wizard   interactive guided tour (you're here) zpwr study    split-pane learning: book left, shell right Five commands. That's your entire onboarding.


\vspace{12pt}
{\bfseries\color{neonyellow} HOW TO ADD YOUR OWN VERBS}\vspace{4pt}

1. Create an autoload function in \textbf{\$ZPWR/autoload/common/} 2. Add to \textbf{ZPWR\_VERBS} in \textbf{\$ZPWR\_ENV/zpwr.zsh} Format: \textbf{ZPWR\_VERBS}[myverb]='myFunc=description' 3. Reload: exec zsh. Your verb is live.


\vspace{12pt}
{\bfseries\color{neonyellow} HOW TO EXTEND THE WIZARD}\vspace{4pt}

Add a new page\_NNN.zsh to \textbf{\$ZPWR\_INSTALL/wizard\_pages/} Follow the format: PAGE\_TITLE, PAGE\_CHAPTER, PAGE\_CONTENT(). The wizard auto-discovers pages. No registration needed.


\vspace{12pt}
{\bfseries\color{neonyellow} HOW TO CONTRIBUTE}\vspace{4pt}

Fork github.com/MenkeTechnologies/zpwr Add your feature. Write tests. Submit a PR. The codebase is the documentation. Read it.


\vspace{12pt}
{\bfseries\color{neonyellow} THE \textbf{ZPWR} PHILOSOPHY}\vspace{4pt}

Verbs, not scripts.  Discoverability over memorization. Speed over bloat.    Fallbacks over hard dependencies. Convention over config.  Your shell should just work.


\vspace{12pt}
{\bfseries\color{neonyellow} ONE LAST THING}\vspace{4pt}

Run zpwr about to see the zpwr banner. Then open a new terminal. Watch the pony greet you. You are the wizard now. The terminal is yours.


\vspace{10pt}
{\large\bfseries\color{neoncyan} END OF TRANSMISSION}\vspace{6pt}

{\small\itshape\color{dimtext} // \textbf{ZPWR} WIZARD ENCYCLOPEDIA v1.0 // SIGNAL LOST //}

May your shell be fast, your pipes be clean, and your exit codes always zero.


\vspace{16pt}


\chapter{HEALTH \& DIAGNOSTICS}


\begin{center}
\includegraphics[width=0.85\textwidth]{/Users/wizard/.zpwr/docs/screenshots/zpwr_doctor_help.png}
\end{center}

\section{Expansion Stats Dashboard}


\vspace{10pt}
{\large\bfseries\color{neoncyan} EXPANSION STATS DASHBOARD}\vspace{6pt}

Your aliases are working for you — here's the proof.


\vspace{12pt}
{\bfseries\color{neonyellow} zpwrExpandStats}\vspace{4pt}

Every time zsh-expand fires — alias, global alias, suffix alias, self-ref escape, native glob/hist, or spelling correction — it logs the event. zpwrExpandStats reads that log and shows you a ranked dashboard.

\begin{codebox}
zpwrExpandStats            \# show full dashboard\\
zpwrExpandStats -h         \# help with all options\\
zpwrExpandStats -t 10      \# top 10 aliases\\
zpwrExpandStats -w 80      \# wider box (80 cols)\\
zpwrExpandStats -r         \# reset all stats\\
zpwrExpandStats -f FILE    \# custom stats file\\
\end{codebox}


\vspace{12pt}
{\bfseries\color{neonyellow} WHAT IT TRACKS}\vspace{4pt}

By trigger ......... space expansions vs history expansions By type ............ alias, global, suffix, escape, native Corrections ........ spelling fixes applied Total expansions ... how many times the engine fired Keystrokes saved ... chars you didn't have to type Top aliases ........ ranked bar chart with expansions shown


\vspace{12pt}
{\bfseries\color{neonyellow} HOW IT WORKS}\vspace{4pt}

The zsh-expand plugin hooks into the zle accept-line widget. Before each command runs, it expands aliases, globals, and suffixes inline. Each expansion is logged to \textbf{\$ZPWR\_LOCAL/zpwr-expand-stats.dat} with the trigger type, expansion type, and original alias name.


\vspace{12pt}
{\bfseries\color{neonyellow} CONFIG}\vspace{4pt}

\begin{codebox}
ZPWR\_EXPAND\_STATS\_FILE     \# override stats file path\\
ZPWR\_EXPAND\_STATS\_TOP      \# default top N (default: 20)\\
\end{codebox}

Every alias that fires is a keystroke investment paying dividends. This dashboard shows your ROI.


\vspace{16pt}

\end{document}
